\documentclass[a4paper]{report}

\usepackage[utf8]{inputenc}
\usepackage[T1]{fontenc}
\usepackage{mathpazo}
\usepackage[ngerman]{babel}
\usepackage[german=quotes]{csquotes}  % Proper German quotation marks with \enquote{}
\usepackage{altstyle}
\usepackage{amsmath}
\usepackage{amsthm}
\usepackage{dsfont}
\usepackage{amsfonts}
\usepackage{amssymb}
\usepackage{byname}
\usepackage[colorlinks=true, linkcolor=blue, citecolor=cyan]{hyperref}
\usepackage{smartref}
\usepackage{bookmark}
\usepackage{enumerate}
\usepackage{stmaryrd}
\usepackage{nicefrac}
\usepackage{ifsym}
\usepackage{mathtools}

% Double stroked N for natural numbers
\newcommand{\N}{\ensuremath{\mathbb{N}}}

% Double stroked Z for integers
\newcommand{\Z}{\ensuremath{\mathbb{Z}}}

% Double stroked Q for rational numbers
\newcommand{\Q}{\ensuremath{\mathbb{Q}}}

% Double stroked R for reals
\newcommand{\R}{\ensuremath{\mathbb{R}}}

% Double stroked K for a field (K as in "`Körper"'). Very common in Banach 
% space analysis. Same for F (as in field) for texts in english.
\newcommand{\K}{\ensuremath{\mathbb{K}}}
\newcommand{\F}{\ensuremath{\mathbb{F}}}

% Double stroked E, commonly used in discrete stochastics and asymptotic
% functional analysis to denote the expected value

\newcommand{\E}{\ensuremath{\mathbb{E}}}

% Double stroked P for the set of prime numbers
\newcommand{\Prim}{\ensuremath{\mathbb{P}}}

% Double stroked 1. Commonly used to denote indicator functions.
\newcommand{\Eins}{\ensuremath{\mathds{1}}}

% For flooring an expression without opening and closing any tags
\newcommand{\floor}[1]{\left\lfloor{#1}\right\rfloor}

% For ceiling an expression without opening and closing any tags
\newcommand{\ceil}[1]{\left\lceil{#1}\right\rceil}

% The following four commands create a quick and dirty way to type a sum,
% product, cap or cup in a $.$-math environment having the bounds above and 
% below the sign instead of its right. Use them economically as they
% destroy any flow.

\newcommand{\Sum}[2]{\stackrel{#2}{\underset{#1}{\sum}}}
\newcommand{\Prod}[2]{\stackrel{#2}{\underset{#1}{\prod}}}
\newcommand{\altCup}[2]{\overset{#2}{\underset{#1}{\bigcup}}}
\newcommand{\altCap}[2]{\overset{#2}{\underset{#1}{\bigcap}}}
\newcommand{\Sup}[1]{\underset{#1}\sup\,}
\newcommand{\Inf}[1]{\underset{#1}\inf\,}

% Simplified command for lower/upper integral bounds
\newcommand{\Int}[2]{\int\limits_{#1}^{#2}}

% Simplified command for an integral with just one lower bound, as in curvature
% integrals or integrals in measure theory
\newcommand{\Intmen}[1]{\underset{#1}{\int}}

% These three commands create brackets around the variable and add its 
% (single!) index (subscript) and its range. Handle with care ;).
% Requires the altstyle.sty to create the "Laue-Haken".

\newcommand{\finfol}[3]{(#1_{#2})_{#2\in\haken{#3}}}   %finite series
\newcommand{\finfolN}[3]{(#1_{#2})_{#2\in\haken{#3}_0}}%finite series with zero
\newcommand{\folge}[3]{(#1_{#2})_{#2 \in #3}}          %series

% Denotes the span of the argument in algebraic notation
\newcommand{\erz}[1]{\left\langle #1\right\rangle}

% Denotes the absolute value of the argument
\newcommand{\abs}[1]{\left|#1\right|}

% Use this to create "sets". The first argument defines the variable (x \in X
% for instance), the second one its defining property (e. g. x \neq y).

\newcommand{\menge}[2]{\left\{#1\,|\,#2\right\}}
\newcommand{\mengea}[2]{\left\{#1\,\vrule\,#2\right\}}

% Puts set brackets around the argument

\newcommand{\meng}[1]{\left\{#1\right\}}

% Same as \abs, but with "norm lines".
\newcommand{\norm}[1]{\lVert#1\rVert} %simple norm, always small delimiters
\newcommand{\norma}[1]{\left\|#1\right\|} %variable delimiters

% Creates the commonly used notation of the "operator norm". The first
% argument is the operator, the second its domain and the third its range.

\newcommand{\opnorm}[3]{\left\|#1\right\|_{\scriptscriptstyle{#2 \rightarrow #3}}}

% Puts [] around the argument
\newcommand{\klasse}[1]{\left[#1\right]}

% Simplifies the \lim command
\newcommand{\Limes}[2]{\underset{#1\to #2}{\lim}}

% limsup
\newcommand{\Lims}[2]{\underset{#1\to #2}{\overline{\lim}}}

% liminf
\newcommand{\Limi}[2]{\underset{#1\to #2}{\underline{\lim}}}

% lim from below
\newcommand{\Limup}[2]{\underset{#1\nearrow #2}{\lim}}

% lim from above
\newcommand{\Limdown}[2]{\underset{#1\searrow #2}{\lim}}

% creates a paragraph
\newcommand{\absatz}{\hfill \medskip \newline}

% creates a indented paragraph
\newcommand{\parag}{\hfill \medskip \par}

% denotes the "id" (as in identity)
\newcommand{\id}{\text{id}\,}

% Abbreviations
\newcommand{\Def}{\mathop{\rm Def}\nolimits}    % domain

\newcommand{\Kern}   {\mathop{\rm Kern}\nolimits}   % kernel
\newcommand{\Bild}   {\mathop{\rm Bild}\nolimits}   % range
\newcommand{\Spur}   {\mathop{\rm Spur}\nolimits}   % trace
\newcommand{\Inn}    {\mathop{\rm int}\nolimits}    % interior
\newcommand{\Syl}    {\mathop{\rm Syl}\nolimits}    % sylow groups
\newcommand{\minpol} {\mathop{\rm minPol}\nolimits} % minimal polynomial
\newcommand{\charpol}{\mathop{\rm charPol}\nolimits}% characteristic polynomial
\newcommand{\chark}  {\mathop{\rm char}\nolimits}   % characteristic (rings)
\newcommand{\grad}   {\mathop{\rm grad}\nolimits}   % degree
\newcommand{\sgn}    {\mathop{\rm sgn}\nolimits}    % sign
\newcommand{\Pot}    {\mathop{\rm Pot}\nolimits}    % power set
\newcommand{\End}    {\mathop{\rm End}\nolimits}    % endomorphisms
\newcommand{\Aut}    {\mathop{\rm Aut}\nolimits}    % automorphisms
\newcommand{\Hom}    {\mathop{\rm Hom}\nolimits}    % homomorphisms
\newcommand{\Epi}    {\mathop{\rm Epi}\nolimits}    % epimorphisms
\newcommand{\Mon}    {\mathop{\rm Mon}\nolimits}    % monomorphisms
\newcommand{\Iso}    {\mathop{\rm Iso}\nolimits}    % isomorphisms
\newcommand{\Fix}    {\mathop{\rm Fix}\nolimits}    % fix point set
\newcommand{\Inv}    {\mathop{\rm Inv}\nolimits}    % units
\newcommand{\ggt}    {\mathop{\rm ggT}\nolimits}    % greatest common divisor
\newcommand{\kgv}    {\mathop{\rm kgV}\nolimits}    % least common multiple
\newcommand{\inv}    {^{-1}}                        % to the power of -1
\newcommand{\ord}    {\mathop{\rm ord}\nolimits}    % degree of a group element
\newcommand{\Gal}    {\mathop{\rm Gal}\nolimits}    % galois group
\newcommand{\Mat}    {\mathop{\rm Mat}\nolimits}
\newcommand{\Umg}    {\mathop{\rm Umg}\nolimits}
\newcommand{\pr}     {\mathop{\rm pr}\nolimits}
\newcommand{\spann}  {\mathop{\rm span}\nolimits}
\newcommand{\conv}   {\mathop{\rm conv}\nolimits}
%induction stuff
\newcommand{\ia}{\textsl{Induktionsanker:}}
\newcommand{\iv}{\textsl{Induktionsannahme:}}
\newcommand{\is}{\textsl{Induktionsschritt:}}

% "`=>"' and "`<="'
\newcommand{\rueck}{"`$\Leftarrow$"':}
\newcommand{\hin}{"`$\Rightarrow$"':}

% same for subset and supset
\newcommand{\teilmenge}{"`$\subseteq$"':\\}
\newcommand{\obermenge}{"`$\supseteq$"':\\}

% some common distributions (not common notation)
\newcommand{\hyp}{\textsl{Hyp}} %hypergeometric
\newcommand{\ber}{\textsl{Ber}} %Bernoulli
\newcommand{\cau}{\textsl{Cau}} %Cauchy
\newcommand{\bin}{\textsl{Bin}} %binomial
\newcommand{\poi}{\textsl{Poi}} %Poisson
\newcommand{\Var}{\textsl{Var}} %variance

\newcommand{\erf}{\mathop{\rm erf}\nolimits}

% direct and orthogonal sum - if you find a single sign for this one, let
% me know ;)
\newcommand{\orth}{\text{\textcircled{$\bot$}}}

% litte shorter for single letters
\newcommand{\mca}{\mathcal{A}}
\newcommand{\mcb}{\mathcal{B}}
\newcommand{\mcc}{\mathcal{C}}
\newcommand{\mcd}{\mathcal{D}}
\newcommand{\mce}{\mathcal{E}}
\newcommand{\mcf}{\mathcal{F}}
\newcommand{\mcg}{\mathcal{G}}
\newcommand{\mch}{\mathcal{H}}
\newcommand{\mci}{\mathcal{I}}
\newcommand{\mcj}{\mathcal{J}}
\newcommand{\mck}{\mathcal{K}}
\newcommand{\mcl}{\mathcal{L}}
\newcommand{\mcm}{\mathcal{M}}
\newcommand{\mcn}{\mathcal{N}}
\newcommand{\mco}{\mathcal{O}}
\newcommand{\mcp}{\mathcal{P}}
\newcommand{\mcq}{\mathcal{Q}}
\newcommand{\mcr}{\mathcal{R}}
\newcommand{\mcs}{\mathcal{S}}
\newcommand{\mct}{\mathcal{T}}
\newcommand{\mcu}{\mathcal{U}}
\newcommand{\mcv}{\mathcal{V}}
\newcommand{\mcw}{\mathcal{W}}
\newcommand{\mcx}{\mathcal{X}}
\newcommand{\mcy}{\mathcal{Y}}
\newcommand{\mcz}{\mathcal{Z}}

\newcommand{\mfa}{\mathfrak{A}}
\newcommand{\mfb}{\mathfrak{B}}
\newcommand{\mfc}{\mathfrak{C}}
\newcommand{\mfd}{\mathfrak{D}}
\newcommand{\mfe}{\mathfrak{E}}
\newcommand{\mff}{\mathfrak{F}}
\newcommand{\mfg}{\mathfrak{G}}
\newcommand{\mfh}{\mathfrak{H}}
\newcommand{\mfi}{\mathfrak{I}}
\newcommand{\mfj}{\mathfrak{J}}
\newcommand{\mfk}{\mathfrak{K}}
\newcommand{\mfl}{\mathfrak{L}}
\newcommand{\mfm}{\mathfrak{M}}
\newcommand{\mfn}{\mathfrak{N}}
\newcommand{\mfo}{\mathfrak{O}}
\newcommand{\mfp}{\mathfrak{P}}
\newcommand{\mfq}{\mathfrak{Q}}
\newcommand{\mfr}{\mathfrak{R}}
\newcommand{\mfs}{\mathfrak{S}}
\newcommand{\mft}{\mathfrak{T}}
\newcommand{\mfu}{\mathfrak{U}}
\newcommand{\mfv}{\mathfrak{V}}
\newcommand{\mfw}{\mathfrak{W}}
\newcommand{\mfx}{\mathfrak{X}}
\newcommand{\mfy}{\mathfrak{Y}}
\newcommand{\mfz}{\mathfrak{Z}}

% common in probability theory and function theory
\newcommand{\spi}{\frac{1}{\sqrt{2\pi}}}

% \varepsilon looks great, but is way to long, same for \varphi
\newcommand{\eps}{\varepsilon}
\newcommand{\vi}{\varphi}

% tautology (quick and dirty)
\newcommand{\tteq}{\dashv \vdash}

% implication with extra space
\newcommand{\folgt}{\, \Rightarrow \,}

% creates two links to the same object, respecting the current value
% (ref). For instance "\shortref{fig:figure}{Fig.}"' will create a blue
% "`Fig. 1"', where both parts link to the figure.
\newcommand{\shortref}[2]{\hyperref[#1]{#2} \ref{#1}}

% inner produkt in "langle"-notation.
\newcommand{\sprod}[2]{\left\langle #1 |#2\right\rangle}

% used for the linear combinator (by me ;))
\newcommand{\bprod}[2]{\left\langle #1 ||#2\right\rangle}

% add the current section to the table of contents.

\newcommand{\tocsec}[2]
 {\addcontentsline{toc}{section}{\protect\numberline{#1} #2}}

% Some ball abbreviations.

% open ball in #1 with centre #2 and radius #3
\newcommand{\koff}[3]{K_{#1}\left(#2, #3\right)}

% closed ball in #1 with centre #2 and radius #3
\newcommand{\kabg}[3]{\overline K_{#1}\left(#2, #3\right)}

% the boundary of a closed ball in #1 with centre #2 and radius #3
\newcommand{\kran}[3]{\partial \kabg{#1}{#2}{#3}}

\newcommand{\koffk}[3]{K_{\scriptscriptstyle{#1}}
                        ^{\scriptscriptstyle{\left(#2, #3\right)}}}

% closed ball in #1 with centre #2 and radius #3
\newcommand{\kabgk}[3]{\overline K_{\scriptscriptstyle{#1}}
                        ^{\scriptscriptstyle{\left(#2, #3\right)}}}

% the boundary of a closed ball in #1 with centre #2 and radius #3
\newcommand{\krank}[3]{\partial \kabgk{#1}{#2}{#3}}

% overset an arrow with an index
\newcommand{\pfeil}[2]{ {\overrightarrow{#1}}^{{}_#2} }
\newcommand{\pfeilk}[2]{ {\scriptstyle{\overrightarrow{#1}}^{{}_#2}} }

\newcommand{\restrict}[2]{#1\vrule_{{}_{\scriptstyle #2}}}

\newcommand{\refclaim}[2]{\ref{#1}.(\ref{#2})}
\newcommand{\refcl}[2]{\ref{#1}.({\ref{#1#2}})}

\newcommand{\claim}[1]{\item \label{\thesection{#1}}}

%%%%%%%%%%%%%%%%%%%%%%%%%%%%%%%%%%%%%%%%%%%%%%%%%%%%%%%%%%%%%%%%%%%%%%%%%%%%%%%
% Theorem stuff, use "thm" for italic fonts inside theorem environments and   %
% "def" for normal fonts. These differ from the originals: newline is created %
% after the complete name (including optional parenthesis) of the theorem.    %
%%%%%%%%%%%%%%%%%%%%%%%%%%%%%%%%%%%%%%%%%%%%%%%%%%%%%%%%%%%%%%%%%%%%%%%%%%%%%%%

\newtheoremstyle{thm}   % name
{15pt}                  % Space above
{10pt}                  % Space below
{\itshape}              % Body font
{}                      % indent amount
{\bfseries}             % Theorem head font
{}                      % Punctuation after theorem head
{\newline}              % space after theorem head
{}                      % Theorem head spec

\newtheoremstyle{def}   % name
{15pt}                  % Space above
{10pt}                  % Space below
{}                      % Body font
{}                      % indent amount
{\bfseries}             % Theorem head font
{}                      % Punctuation after theorem head
{\newline}              % space after theorem head
{}                      % Theorem head spec

%%%%%%%%%%%%%%%%%%%%%%%%%%%%%%%%%%%%%%%%%%%%%%%%%%%%%%%%%%%%%%%%%%%%%%%%%%%%%%%
% Forces use of "Number.Subnumber Name" (Wrobel style ;) ). Enumeration       %
% follows the current chapter NOT the current section.                        %
% Do NOT use this, if you don't use chapters. Swap the comments below.        %
%%%%%%%%%%%%%%%%%%%%%%%%%%%%%%%%%%%%%%%%%%%%%%%%%%%%%%%%%%%%%%%%%%%%%%%%%%%%%%%

\swapnumbers
\theoremstyle{thm}
\newtheorem{sat}{Satz}[chapter]

\theoremstyle{def}
\newtheorem{bem}{Bemerkung}[section]

\theoremstyle{def}
\newtheorem{dfn}[sat]{Definition}

\theoremstyle{thm}
\newtheorem{kor}[sat]{Korollar}

\theoremstyle{thm}
\newtheorem{lem}[sat]{Lemma}

\theoremstyle{thm}
\newtheorem{defsat}[sat]{Definition/Satz}

\theoremstyle{thm}
\newtheorem{beisp}[sat]{Beispiel}

% These commands (definition through korollar) can be used to ensure the
% following:
% - you open and close a new "something", name it (#1) and label it (#2).
% - the label is set in the right place (i. e. hyperref will jump to the right
%   position; the label is "#2"
% - a tableofcontents-entry is created, which respects the same enumeration

\newenvironment{definition}[2] % 1=Name, 2=Label
{\begin{dfn}[#1]
\label{#2}
\addcontentsline{toc}{section}{\protect\numberline{\ref{#2}} {#1}} }
{\end{dfn}}

\newenvironment{lemma}[2] % 1=Name, 2=Label
{\begin{lem}[#1]
\label{#2}
\addcontentsline{toc}{section}{\protect\numberline{\ref{#2}} #1}}
{\end{lem}}

\newenvironment{satz}[2] % 1=Name, 2=Label
{\begin{sat}[#1]
\label{#2}
\addcontentsline{toc}{section}{\protect\numberline{\ref{#2}} #1}}
{\end{sat}}

\newenvironment{bemerkung}[2] % 1=Name, 2=Label
{\begin{bem}[#1]
\label{#2}
\addcontentsline{toc}{section}{\protect\numberline{\ref{#2}} #1}}
{\end{bem}}

\newenvironment{korollar}[2] % 1=Name, 2=Label
{\begin{kor}[#1]
\label{#2}
\addcontentsline{toc}{section}{\protect\numberline{\ref{#2}} #1}}
{\end{kor}}

\newenvironment{defsatz}[2] % 1=Name, 2=Label
{\begin{defsat}[#1]
\label{#2}
\addcontentsline{toc}{section}{\protect\numberline{\ref{#2}} #1}}
{\end{defsat}}

\newenvironment{beispiel}[2] % 1=Name, 2=Label
{\begin{beisp}[#1]
\label{#2}
\addcontentsline{toc}{section}{\protect\numberline{\ref{#2}} #1}}
{\end{beisp}}

\allowdisplaybreaks[1]
\numberwithin[\roman]{equation}{sat}

\begin{document}

\begin{titlepage}

\hypertarget{titlepage}{}

\begin{center}
{ \huge \textbf{Anwendungen des Gauß'schen Min--Max--Satzes}}\\[3cm]
\vspace*{3cm}
{\large
Diplomarbeit\\
von\\
{\textbf {Nikita Danilenko}} \\
}

\vspace*{\fill}

\begin{tabular}{rl}
Betreuung:& Priv.--Doz. Dr. Carsten Schütt\\
\end{tabular}
\end{center}
\vspace*{3cm}
\begin{center}
Mathematisches Seminar\\
Christian--Albrechts--Universität zu Kiel
\end{center}
\vspace*{3cm}
\begin{center}
Kiel, Juli 2010
\end{center}
\end{titlepage}

\tableofcontents

\chapter*{Einführung}

Ziel dieser Arbeit ist es, eine Reihe bekannter Sätze aus der Banachraumtheorie
unter Zuhilfenahme des Gauß'schen Min--Max--Satzes zu beweisen.
Um dieses zu bewerkstelligen, wird zunächst der besagte Satz auf eine bestimmte Klasse
von unendlichen Familien normalverteilter Zufallsgrößen verallgemeinert.

Unter Zuhilfenahme des verallgemeinerten Satzes wird es schließlich möglich
sein, die Sätze von Dvoretzky, Milman--Schechtman und Gluskin zu beweisen.
Außerdem folgen Beweise für das Johnson--Lindenstrauss--Lemma und ein
Satz über die \enquote{Restricted Isometry Property}.

In den ersten zwei Kapiteln werden funktionalanalytische und maßtheoretische
Grundlagen wiederholt, das dritte Kapitel verallgemeinert den Gauß'schen
Min--Max--Satz auf spezielle unendliche Familien und das vierte stellt
schließlich die Beweise der zuvor erwähnten Sätze vor.

Diese Arbeit beruht zu großen Teilen auf den Arbeiten von Yehoram Gordon.
Die im dritten Kapitel behandelten Gordon'schen Paare wurden von Gordon
in \cite{Gor} vorgestellt und in \cite{Gor2} in modifizierter Form verwendet,
die Hauptergebnisse über Gordon'sche Paare stammen ebenfalls aus diesen
Quellen.
Ebenso gehen die Folgerungen aus dem ursprünglichen Gauß'schen
Min--Max--Satz (besonders \ref{SatzVonFernique} und
\ref{MaximumZentrierterNormalverteilterZufallsgroessen}) weitestgehend auf
Gordon zurück.

Die Beweise der Resultate im vierten Kapitel beruhen zu wesentlichen Teilen auf
\cite{Gor2}, \cite{Gor}, \cite{Gued}.

\newpage
\setcounter{chapter}{-1}
\chapter{Bezeichnungen}
\label{ch:bezeichnungen}

Seien $n,m\in \N$, $p \in [1, \infty[$, $K$ ein Körper, $\K\in\meng{\R, \C}$.\\
Wenn wir von einem normierten Vektorraum reden, setzen wir implizit voraus,
dass dieser nicht der Nullvektorraum ist.

\begin{enumerate}[(1)]
 \item \label{Bez1}
       Wir schreiben verkürzend $\haken n := \N_{\le n}$ und 
       $\haken n _0:=\haken n \cup \meng 0$.
 \item \label{Bez2}
       Wir nennen eine Menge $M$ abzählbar, genau dann, wenn sie endlich oder
       abzählbar unendlich ist.
 \item \label{Bez3}
       Ist $A$ eine Menge und $a \in A^n$, so fassen wir $a$ als Abbildung 
       $a : \haken n \rightarrow A$, für die wir wiederum die
       Familienschreibweise nutzen. Spielt die Reihenfolge eine Rolle, so
       schreiben wir $a = (a_i)_{i=1}^n$ anstelle von $a=(a_i)_{i\in\haken n}$.
 \item \label{Bez13}
       Für eine Menge $M$ sei 
       $\Pot_{ne}(M):=\menge{A\in\Pot(M)}{A\neq\emptyset,\,A\text{ endlich}}$.
 \item \label{Bez4}
       Mit $(e^n_i)_{i =1}^n$ bezeichnen wir die kanonische Basis des
       $K^n$.
 \item \label{Bez5}
       Ist $(E, \norm{\cdot}_E)$ ein normierter $\K$--Vektorraum, so sei
  \[
   \norm{\cdot}_{p, n, E}: E^n \rightarrow [0, \infty[, 
   \finfol x i n\mapsto\bigg(\sum_{i=1}^n \norm{x_i}^p_E\bigg)^{\nicefrac 1p}  
  \]
  die $p$--Norm auf dem Vektorraum $E^n$.
 \item \label{Bez6}
       $\sprod{\cdot}{\cdot\cdot}_n$ sei das Standardskalarprodukt auf $K^n$.
 \item \label{Bez7}
       Ist $\Omega$ eine Menge und $f:\Omega \rightarrow \R^n$, so seien für
       alle $i \in \haken n$: $f_i := \text{pr}_i \circ f$, wobei $\text{pr}_i$
       die Projektion in die $i$--te Komponente sei. Wir werden dann häufig $f$
       mit $(f_i)_{i=1}^n$ identifizieren.
 \item \label{Bez8}
       Mit $K^{n \times m}$ bezeichnen wir den Raum der $n\times m$--Matrizen,
       wobei $n$ die Anzahl der Spalten und $m$ die der Zeilen ist.
       
       Für $A\in K^{n\times m}$ definieren wir $A_i := A e_i^n$ und 
       $A_{i,j}:=\sprod{A_i}{e_j^m}_m$ für alle $i\in\haken n$, 
       $j\in\haken m$, d. h. $A_i$ ist die $i$--te Spalte von $A$ und $A_{i,j}$
       ist der $j$--te Eintrag in der $i$--ten Spalten von $A$.\\
       Ist analog $\Omega$ eine Menge und 
       $\alpha:\Omega \rightarrow K^{m\times n}$, so bezeichnen wir für alle 
       $i \in \haken n$ und alle $j \in \haken m$
       \[\alpha_i: \Omega \rightarrow K^n, \omega \mapsto \alpha(\omega)_i
       \quad \text{und} \quad
       \alpha_{i,j}:\Omega\rightarrow K,\omega\mapsto\alpha(\omega)_{i,j}
       \text.\]
 \item \label{Bez9}
       Die $p$--Norm auf den $L_p$-- (bzw. $\mcl_p$--) Räumen bezeichnen wir
       mit $\norm\cdot_p$ und die $p$--Norm auf dem $\K^n$ speziell als 
       $\norm\cdot_{p, n}$.
 \item \label{Bez10}
       Sind $(X, \norm\cdot_X)$ und $(Y, \norm\cdot_Y)$ normierte
       $\K$--Vektorräume, so sei
       \[ ev : X \times L(X, Y) \rightarrow Y, (x, T) \mapsto Tx \]
       der Einsetzungshomomorphismus.\\
       Ferner sei
       \[\Inv(X, Y) := \menge{T \in L(X, Y)}{T \text{ invertierbar und }
       T\inv \in L(Y, X)}\]
 \item \label{Bez11}
       Ist $X$ ein $K$--Vektorraum, so seien $X^* := \Hom(X,\K)$ (algebraischer
       Dualraum) und $X':=L(X, \K)$ (stetiger Dualraum).
 \item \label{Bez12}
       Sei 
\[\Mat_{m,n}: \R^{mn} \rightarrow \R^{m \times n}, (x_i)_{i=1}^{mn} \mapsto
{\left(x_{m\cdot (i-1) + j}\right)_{i=1}^n}_{j=1}^m\text,\]
d. h. durch $\Mat_{m,n}$ können wir einen Vektor aus dem $\R^{mn}$ als Matrix
im $\R^{m \times n}$ auffassen. Offensichtlich ist $\Mat_{m,n}$ linear,
bijektiv und aufgrund der endlichen Dimension der beiden Räume auch stetig,
wenn man beide Räume mit Normen ausstattet.
 \item \label{Bez14} Ist $(M,d)$ ein metrischer Raum, $x \in M$ und $r \in
  \R_{>0}$, so seien
  \[\koff M x r := \menge{y \in X}{d(x,y) < r}\text{ und }
  \kabg M x r := \menge{y \in X}{d(x,y) \le r}\text.\]
  Ist $A \subseteq M$, so sei $\partial A := \overline A \setminus \Inn(A)$.
 \item \label{Bez15}
 Ist die Norm eines normierten Raumes $(X, \norm\cdot_X)$ von besonderer
 Bedeutung, so werden wir gelegentlich den Index bei Kugeln von $X$ zu
 $(X, \norm\cdot_X)$ abwandeln, um kenntlich zu machen, dass diese Norm
 gemeint ist. Analog werden wir verfahren, wenn $(Y, \norm\cdot_Y)$ ein
 weiterer normierter Raum ist und wir $L(X,Y)$ betrachten, d. h. wir werden
 statt dessen $L((X, \norm\cdot_X), (Y, \norm\cdot_Y))$ oder
 $L((X, \norm\cdot_X),Y)$ schreiben. Dies ist besonders dann wichtig, wenn es
 mehrere Normen auf $X$ (bzw. $Y$) gibt und aus dem Kontext nicht klar wird,
 welche gemeint ist.
 \item \label{Bez16}
 Sei $\ell_p^n := (\R^n, \norm\cdot_{p,n})$.
\end{enumerate}

\chapter{Funktionalanalytische Grundlagen}
\label{ch:kapitel1}

In diesem Kapitel werden wir einige Sätze und Lemmata im Zusammenhang mit
linearen Abbildungen und insbesondere Funktionalen beweisen.
Die meisten dieser Aussagen sind elementarer Natur, einige allerdings benutzen bekannte Sätze der
Funktionalanalysis wie den Satz von Hahn--Banach und den Satz von der offenen
Abbildung.

Ziel ist es, eine Reihe von nützlichen Werkzeugen zur Verfügung zu stellen, die
an späterer Stelle sehr häufig benötigt werden.

Wir beginnen mit zwei Begriffen, die in erster Linie dazu gut sind, aufwändige
Summendarstellungen zu vermeiden.

\begin{definition}{Erweiterung von Funktionen auf Tupel}
                  {ErweiterungVonFunktionenAufTupel}
Seien $A, B$ nichtleere Mengen, $\vi: A \rightarrow B$ und $n \in \N$.\\
Dann definieren wir 
\[\pfeil \vi n: A^n \rightarrow B^n, \quad
   (a_i)_{i=1}^n \mapsto (\vi(a_i))_{i=1}^n\]
Anders notiert: f�r $a \in A^n$ ist $\pfeil \vi n (a) = \vi \circ a$
(vgl. Bezeichnungen.(\ref{Bez3})).
\end{definition}
\parag
Wir fassen einige elementare Eigenschaften der Erweiterung im Falle linearer
Abbildungen zusammen.

\begin{lemma}{Eigenschaften der Erweiterung im Vektorraumfall}
             {EigenschaftenDerErweiterungImVektorraumfall}
Seien $K$ ein Körper, $V,W$ $K$--Vektorräume und $n \in \N$.\\
Dann gilt:
\begin{enumerate}[(1)]
	\item \label{EigenschaftenDerErweiterungImVektorraumfall1}
	 Für jedes $\alpha \in \Hom(V, W)$ ist auch 
    $\pfeil \alpha n:V^n \rightarrow W^n$ linear.
\item \label{EigenschaftenDerErweiterungImVektorraumfall2} Die Abbildung \[
\pfeil{\cdot}n: \Hom(V, W) \rightarrow \Hom(V^n, W^n), \varphi\mapsto\pfeil\varphi n\]
ist ebenfalls linear.
\end{enumerate}
\end{lemma}
\begin{proof}
Seien $\varphi, \tau \in \Hom(V, W)$, $k \in K$ und $v, w \in V^n$.\\
Es ist dann
\begin{align*}
\pfeil \varphi n(v+kw)
   & = \pfeil\varphi n \left((v_i+k w_i)_{i=1}^n\right)
     = (\varphi(v_i+kw_i))_{i=1}^n\\
   & = (\varphi(v_i) + k \varphi(w_i))_{i=1}^n &\quad \varphi \text{ linear}\\
   & = \left(\varphi(v_i)\right)_{i=1}^n + k\left(\varphi(v_i)\right)_{i=1}^n
     = \pfeil \varphi n (v) + k \pfeil \varphi n (w)
\end{align*}
Damit ist $\pfeil \varphi n$ linear.\\
Weiter gilt:
\begin{align*}
\pfeil {(\varphi + k \tau)} n (v)
  & = \left((\varphi + k \tau)(v_i)\right)_{i=1}^n
    = \left(\varphi(v_i) + k \tau(v_i)\right)_{i=1}^n\\
  & = \left(\varphi(v_i)\right)_{i=1}^n + k \left(\tau(v_i)\right)_{i=1}^n
    = \pfeil \varphi n (v) + k \pfeil \tau n (v)
\end{align*}
Also ist auch $\pfeil \cdot n$ linear.
\end{proof}
\parag
Im Falle normierter Vektorräume ergeben sich zusätzlich die nachfolgend
notierten Stetigkeitseigenschaften.

\begin{lemma}{Eigenschaften der Erweiterung im normierten Fall}
             {EigenschaftenDerErweiterungImNormiertenFall}
Seien $(V, \norm\cdot_V)$, $(W, \norm\cdot_W)$ normierte $\K$--Vektorräume,
$\varphi \in L(V, W)$, $p \in [1, \infty[$ und $n \in \N$. $V^n$ sei mit
$\norm\cdot_{p, n, V}$ und $W^n$ mit $\norm\cdot_{p, n, W}$ ausgestattet.\\
Dann gilt für alle $v \in V^n$:
    $\norm{\pfeil \varphi n(v)}_{p, n, W}
    \le \norm{\varphi}_{V \rightarrow W}\cdot \norm{v}_{p, n, V}$\\
Insbesondere ist $\pfeil \varphi n \in L(V^n, W^n)$
und $\opnorm{\pfeil \varphi n}{V^n}{W^n} = \opnorm{\varphi} V W$, also
\[\pfeil \cdot n: L(V, W) \rightarrow L(V^n, W^n), 
  \alpha \mapsto \pfeil \alpha n\]
eine lineare Isometrie (bezüglich der obigen Normwahl).
\end{lemma}
\begin{proof}
Nach~\ref{EigenschaftenDerErweiterungImVektorraumfall} ist $\pfeil\varphi n$
linear.\\
Wegen $p \ge 1$ sind $\norm \cdot_{p, n, V}$ und $\norm\cdot_{p,n,W}$ Normen.\\
Sei $v \in V^n$.
Dann gilt\[
       \norm{\pfeil \varphi n (v)}_{p, n, W}^p
   =   \sum_{i=1}^n \norm{\varphi(v_i)}_W^p
   \le \sum_{i=1}^n \opnorm\varphi V W^p \norm{v_i}_V^p
   =   \opnorm \varphi V W ^p \cdot \norm v _{p, n, V}^p\text.\]
Damit ist einerseits die erste Aussage gezeigt und andererseits ist 
$\pfeil \varphi n$ nach bekannten Charakterisierungen der Stetigkeit linearer
Abbildungen bereits stetig mit
$\opnorm {\pfeil \varphi n}{V^n}{W^n} \le \opnorm \varphi V W$.\\
Sei nun $\eps \in \R_{>0}$.\\
Dann existiert ein $v_\eps \in \overline K_V(0,1)$ mit
$\opnorm \varphi V W - \norm{\varphi(v_\eps)}_W< \eps$.\\
Sei
\[v' :\haken n \rightarrow V, i \mapsto \begin{cases}v_\eps & : i=1\\
0 &: \text{ sonst} \end{cases}\text.\]
Dann ist offenbar $\norm {v'}_{p, n, V} = \norm {v_\eps}_V \le 1$ und
$\norm {\pfeil \varphi n (v')}_{p, n, W} = \norm{\varphi(v_\eps)}_W$.\\
Also ist $
 \opnorm{\pfeil \varphi n}{V^n}{W^n} - \norm{\pfeil \varphi n (v')}
 \le \opnorm \varphi V W - \norm{\varphi(v_\eps)} < \eps$.\\
Damit ist $\opnorm \varphi V W
 = \sup\menge{\norm{\pfeil \varphi n (x)}_{p,n,W}}{x \in \overline K_{V^n}(0,1)}
 = \opnorm{\pfeil \varphi n}{V^n}{W^n}$.\\
Wegen $L(V, W) \le \Hom(V, W)$ ist $\pfeil \cdot n$ nach
\ref{EigenschaftenDerErweiterungImVektorraumfall} linear.
Die Isometrieeigenschaft ist genau das \enquote{Insbesondere}.
\end{proof}
\parag
Wir führen nun den Begriff des Linearkombinators ein, den wir dafür verwenden
werden, stets deutlich zu machen, dass ein Vektor Linearkombination bestimmter
Elemente ist, ohne die Summendarstellung zu verwenden.

\begin{definition}{Linearkombinator}
                  {Linearkombinator}
Seien $K$ ein Körper, $V$ ein $K$--Vektorraum und $n \in \N$.\\
Dann definieren wir \[
\bprod{\cdot}{\cdot\cdot}_{n, V}: K^n \times V^n \rightarrow V, 
    \left(s, v \right) \mapsto \sum_{i=1}^n s_iv_i
\]
und nennen diese Abbildung \textbf{Linearkombinator auf }$V$.\\
Wir werden dafür die Infixnotation verwenden.\\
Wenn aus dem Zusammenhang klar ist, auf welchem Vektorraum man sich befindet,
werden wir auf den zweiten Index $V$ verzichten und nur 
$\bprod{\cdot}{\cdot\cdot}_n$ schreiben.
\end{definition}
\parag
Auch hier ergeben sich einige einfache Eigenschaften.

\begin{lemma}{Eigenschaften des Linearkombinators}
             {EigenschaftenDesLinearkombinators}
Seien $K$ ein Körper, $V,W$ $K$--Vektorräume, $n\in\N$ und $\varphi\in\Hom(V,W)$.\\
Dann gilt:
\begin{enumerate}[(1)]
\item\label{EigenschaftenDesLinearkombinators1} 
$\bprod{\cdot}{\cdot\cdot}_{n, V}$ ist bilinear.
\item\label{EigenschaftenDesLinearkombinators2}
 $\forall\, t \in K^n\,\forall\, v \in V^n:
    \varphi(\bprod t v_{n, V}) = \bprod t {\pfeil \varphi n (v)}_{n, W}$.
\end{enumerate}
\end{lemma}
\begin{proof}
Die Bilinearität ist offensichtlich.\\
Seien $t \in K^n$ und $v \in V^n$.\\
Dann gilt \[
\bprod t {\pfeil \varphi n (v)} _{n, W} = \sum_{i=1}^n t_i \varphi(v_i)
                       = \varphi\left(\sum_{i=1}^n t_i v_i\right)
                       = \varphi(\bprod t v_{n, V})
\]
\end{proof}

\begin{lemma}{Stetigkeit des Linearkombinators}
             {StetigkeitDesLinearkombinators}
Seien $(E, \norm\cdot_E)$ ein normierter $\K$--Vektorraum und $n \in \N$.\\
Wir statten $\K^n$ mit $\norm\cdot_{2,n}$ und $E^n$ mit $\norm\cdot_{2,n,E}$
aus.\\
Dann gilt \[
\forall\, t \in \K^n \forall\, v \in E^n : 
   \norm{\bprod t v _n}_E \le \norm t_{2, n} \cdot \norm v _{2, n, E}
\]
Insbesondere ist $\bprod\cdot{\cdot\cdot}_n$ (mit genannter Normausstattung)
stetig.
\end{lemma}
\begin{proof}
Seien $t \in \K^n$ und $v \in E^n$.
Dann gilt
\begin{align*}
\norma{\bprod t v_n}_E 
    &= \norma{\sum_{i=1}^n t_iv_i}_E 
    \le\sum_{i=1}^n \abs{t_i}\cdot \norm{v_i}_E\\
    &= \sprod{\left(\pfeil {\abs\cdot}n(t)\right)}
            {\left(\pfeil {\norm\cdot_E}n(v)\right)}_n\\
    &\le\norma{\pfeil {\abs\cdot}n(t)}_{2,n}
        \cdot
        \norma{\pfeil{\norm\cdot_E}n(v)}_{2, n} & \quad\text{Cauchy--Schwarz}\\
    &= \norm t_{2, n}\cdot \norm v_{2, n, E}
\end{align*}
Da $\bprod\cdot{\cdot\cdot}_n$ nach \ref{EigenschaftenDesLinearkombinators}
bilinear ist, folgt mit der bekannten Charakterisierung der Stetigkeit
bilinearer Abbildungen die Stetigkeit für den Linearkombinator.
\end{proof}
\parag
Das folgende Lemma ist eine Zusammenfassung bekannter Ergebnisse aus dem
Grundstudium und der Funktionalanalysis und wird üblicherweise als ein
Korollar aus dem Satz von Hahn--Banach vorgestellt. Die hier aufgeführte
Version fasst also nur Bekanntes zusammen.

\begin{lemma}{Norm und duale Norm}{NormUndDualeNorm}
Seien $(E, \norm\cdot_E)$ ein normierter $\K$--Vektorraum und $x \in E$.\\
Dann existiert ein $\vi \in \kran{E'}01$ 
mit $\norm x_E = \vi(x)$.\\
Insbesondere ist damit $\opnorm {ev(x, \cdot)} {E'} \K  = \norm x_E$.\\
Für alle $K \in \meng{\kabg{E'}01, \kran{E'}01}$
gilt daher $\norm x _ E = \sup\menge{\abs{\vi(x)}}{\vi \in K}$.\\
Für $\K = \R$ gilt außerdem $\norm x _ E = \sup\menge{\vi(x)}{\vi \in K}$.
\end{lemma}
\begin{proof}
Wir verwenden den folgenden Spezialfall des Satzes von Hahn--Banach:
\begin{quote}
Seien $(X, \norm\cdot)$ ein normierter $\K$--Vektorraum, $Y \le X$ und
$l \in Y'$.\\
Dann existiert ein $L \in X'$ mit $L\vert _Y = l$ und $\opnorm L X \K = \opnorm l Y \K$.
\end{quote}
Falls $x = 0$, so ist nichts zu tun, denn es gibt stetige Funktionale außer dem
Nullfunktional (wegen $E \neq 0$ kann man $y \in E\setminus \meng 0$ wählen
und damit wie unten beschrieben verfahren).\\
Sei also zusätzlich $x \neq 0$ und $x' \coloneqq  \frac 1 {\norm x _E} x$.\\
Seien $F\coloneqq \erz {x'}_\K$ und $\iota : \K \rightarrow F, k \mapsto k\cdot x'$.\\
Dann ist $\iota$ wegen $\norm{x'}_E=1$ eine bijektive, lineare Isometrie.\\
Sei nun $l \coloneqq  \iota\inv$. Dann ist $l$ ebenfalls eine lineare Isometrie, also
$\opnorm l F \K = 1$.\\
Also existiert nach obigem Satz ein $L \in E'$ mit $L \vert_ F = l$ und
$\opnorm L E \K = 1$.\\
Wegen $x \in F$ gilt weiter $L(x) = L(\norm x_E \cdot x') = \norm x_E$, womit
der erste Teil der Behauptung gezeigt ist.\\
Seien $K \in \meng{\kabg{E'}01, \kran{E'}01}$ und $\vi \in K$.\\
Dann gilt $\abs{\vi(x)} \le \opnorm \vi E \K \cdot \norm x_E \le \norm x_E$.\\
Damit ist bereits $\opnorm {ev(x, \cdot)} {E'} \K \le \norm x _E$.\\
Da aber weiter $L \in K$ ist, gilt auch Gleichheit, denn nach Definition und
einer aus dem Grundstudium bekannten Rechnung ist
\[\opnorm {ev(x, \cdot)} {E'} \K = \sup\menge{\abs{ev(x,\psi)}}{\psi \in K} =
                          \sup\menge{\abs{\psi(x)}}{\psi \in K}\text.\]
Im Falle $\K = \R$ kann man offensichtlich auf den Betrag verzichten.
\end{proof}
\parag
Da man im Hilbertraum jedes stetige Funktional explizit ausschreiben kann,
ergibt sich das folgende Korollar.

\begin{korollar}{Norm und duale Norm in Hilbertr�umen}
                {NormUndDualeNormImHilbertraum}
Seien $(H, \sprod\cdot{\cdot\cdot})$ ein $\K$--Hilbertraum, $x \in H$ und
$K \in \meng{ \kran H 0 1,\kabg H 0 1}$.\\
Dann gilt $\norm x_H = \sup\menge{\abs{\sprod x y}}{y \in K}$.\\
Im Falle $\K = \R$ gilt au�erdem $\norm x_H = \sup\menge{\sprod x y}{y \in K}$.
\end{korollar}
\begin{proof}
Nach dem aus der Funktionalanalysis bekannten Darstellungssatz von
Fr�chet--Riesz ist $\iota : H \rightarrow H', x \mapsto \sprod\cdot x$
eine isometrische, antilineare Bijektion.\\
Daher ist $\menge{\abs{\sprod x y}}{y \in K} 
= \menge{\abs{\vi(x)}}{\vi \in \iota(K)}$.\\
Da $\iota$ eine bijektive Isometrie ist, gilt 
$\iota(K) \in \meng{\kabg{H'}01, \kran{H'}01}$.\\
Somit folgt der Rest der Behauptung direkt aus \ref{NormUndDualeNorm}.
\end{proof}
\parag
Das folgende Korollar ist ein Spezialfall des vorigen im
$(\R^n, \sprod \cdot {\cdot\cdot}_n)$.

\begin{korollar}{Iterierte euklidische Norm}{IterierteEuklidischeNorm}
Seien $n \in \N$, $\Omega$ eine nichtleere Menge und
$f: \Omega \rightarrow \R^n$.\\
Dann sind
\[\inf_{t \in S^{n-1}}\bprod t f_n
\quad \text{ und }\quad \sup_{t \in S^{n-1}}\bprod t f_n\]
wohldefiniert und es gilt
\[\inf_{t \in S^{n-1}}\bprod t f_n = - \norm\cdot_{2, n} \circ f \quad
\text{und}\quad
\sup_{t \in S^{n-1}}\bprod t f_n = \norm\cdot_{2, n} \circ f\text.\]
\end{korollar}
\begin{proof}
Da $\R^n$ mit dem Standardskalarprodukt ein Hilbertraum ist, dessen
Norm genau $\norm\cdot_{2, n}$ ist, gilt $S^{n-1} = \kran{\ell_2^n}01$.\\
Sei $\omega \in \Omega$. Dann gilt nach \ref{NormUndDualeNormImHilbertraum}:
\[\norm {f(\omega)}_{2, n} = \sup_{t \in S^{n-1}} \sprod {f(\omega)}t_n
= \sup_{t \in S^{n-1}} \sprod t {f(\omega)}_n =
\left(\sup_{t \in S^{n-1}} \bprod t f _n\right)(\omega)\text.\]
Au�erdem gilt analog wegen $S^{n-1}=-S^{n-1}$:
\begin{align*}
-\norm{f(\omega)}_{2, n}& = -\sup_{t \in S^{n-1}} \sprod t {f(\omega)}_n
= \inf_{t \in S^{n-1}} \sprod {-t}{f(\omega)}_n
= \inf_{t \in - S^{n-1}}\sprod {t}{f(\omega)}_n\\
&= \inf_{t \in S^{n-1}}\sprod {t}{f(\omega)}_n
= \left(\inf_{t \in S^{n-1}} \bprod t f _n\right)(\omega)\text.
\end{align*}
\end{proof}
\parag
Wir werden nun Vorbereitungen für zwei nützliche Aussagen in Bezug auf den
Banach--Mazur--Abstand zweier Vektorräume treffen.
Zwar werden wir diese Aussagen nur im Falle endlichdimensionaler Räume anwenden, allerdings ist der
Mehraufwand für die Beweise im Unendlichdimensionalen durchaus vertretbar.\\
Wir beginnen mit einem ebenso trivialen wie nützlichen Lemma.

\begin{lemma}{Injektivit�t linearer Abbildungen zwischen norm. R�umen}
             {Injektivit�tLinearerAbbildungenZwischenNormiertenR�umen}
Seien $(E, \norm\cdot_E), (F, \norm\cdot_F)$ normierte $\K$--Vektorr�ume und
$T \in \Hom(E, F)$ mit der Eigenschaft
$\inf\menge{\norm{Tx}_F}{x \in \kran E 0 1} \neq 0$.\\
Dann ist $T$ injektiv.
\end{lemma}
\begin{proof}
Sei $T$ nicht injektiv. Dann gibt es ein $x\in E\setminus\meng 0$ mit $Tx=0$.\\
Da $\norm x _E > 0$ ist $x' := \frac 1 {\norm x_E}x$ wohldefiniert.\\
Es ist also $\norm {x'}_E = 1$ und 
$\norm{Tx'}_F = \frac 1 {\norm x _E}\norm{Tx}_F = 0$, da $Tx = 0$.\\
Wegen 
$\menge{\norm{Ta}_F}{a \in \kran E 0 1} \subseteq [0,\infty[$,
ist 0 eine untere Schranke dieser Menge und nach obigem auch darin
enthalten, also gleich ihrem Infimum.
\end{proof}

Es stellt sich heraus, dass das in diesem Lemma beschriebene Infimum eine
weitere nützliche Rolle spielt.

Die Umkehrung dieses Lemmas ist leider falsch.
Als Beispiel betrachte man die Räume
$(E,\norm\cdot_E)\coloneqq \ell_2(\N,\K)$, $(F,\norm\cdot_F)\coloneqq (E,\norm\cdot_E)$ und
\[T:E\rightarrow F, \folge x n \N \mapsto \left(\frac{x_n}n\right)_{n\in \N}
\text.\]
Dann ist $T$ offenbar linear, stetig und injektiv.\\
Bezeichne $(e_n)_{n\in \N}$ die Standardeinheitsbasis von $\ell_2(\N, \K)$.\\
Dann ist
\begin{align*}\inf\menge{\norm {T x}_2}{x \in \kran{\ell_2(\N, \K)}01}
\le\inf_{n \in \N}\norm{Te_n}_2=\inf_{n \in \N}\nicefrac 1n=0\text.
\end{align*}
In der Tat liegen in diesem Beispiel schärfere Voraussetzungen vor, denn $E$
und $F$ sind $\K$--Banachräume und $T$ ist stetig, was im Lemma nicht
vorausgesetzt wurde. Allerdings ist $T$ offensichtlich weder surjektiv
(z.B. hat $(\frac 1 n)_{n \in \N} \in F$ kein Urbild), noch hat $T$ ein
abgeschlossenes Bild in $F$.\\
Nimmt man jedoch die erwähnten Voraussetzungen hinzu, so erhält man den
folgenden Sachverhalt.

\begin{satz}{Injektivität linearer Operatoren und Operatornorm der Inversen}
            {InjektivitaetLinearerOperatorenUndOperatornormDerInversen}
Seien $(E,\norm\cdot_E),(F,\norm\cdot_F)$ $\K$--Banachräume, $T\in L(E,F)$
injektiv und $T(E)$ sei abgeschlossen in $F$.\\
Dann gilt \[\inf\menge{\norm{Tx}_F}{x \in\kran E 0 1}>0\text.\]
Weiter ist $T\inv$ eine Abbildung, $T\inv \in L(T(E), E)$
und es gilt \[
\opnorm{T\inv} {T(E)} E 
   = \frac 1 {\inf\menge{\norm{Tx}_F}{x \in\kran E 0 1}}\text.
\]
\end{satz}
\begin{proof}
Sei $G\coloneqq T(E)$.\\
Dann ist $G$ wegen der Linearität von $T$ ein Vektorraum.\\
Offenbar ist $\norm\cdot_G\coloneqq \restrict{\norm\cdot_F}G$ eine Norm auf $G$.\\
Da $G$ nach Voraussetzung abgeschlossen in $F$ ist, ist $G$ als abgeschlossene
Menge eines vollständigen metrischen Raums selbst vollständig und damit
insgesamt ein Banachraum.\\
Sei $S : E \rightarrow G, x \mapsto Tx$.\\
Dann ist $S$ bijektiv (weil $T$ injektiv ist), insbesondere surjektiv.\\
Offenbar ist auch $S$ in dieser Normwahl stetig.\\
Nach dem Satz von der offenen Abbildung ist damit $S$ offen.\\
Nun ist $S$ aber auch invertierbar, sodass damit auch $S\inv$ stetig ist.\\
Da $T$ als Relation links-- und rechtseindeutig ist, ist es auch $T\inv$.\\
Damit ist $T\inv$ eine injektive Abbildung von $T(E)$ nach $E$ und stimmt als
solche mit $S\inv$ überein, ist also inbesondere linear und stetig.\\
Seien $c\coloneqq  \opnorm {S\inv} G E$ und $x \in \kran E 0 1$.\\
Dann ist $c > 0$ und es gilt:
\[1 =\norm x _E = \norm{S\inv(Sx)}_E \le \opnorm{S\inv} G E \norm{Tx}_G 
             = c \norm{Tx}_G\text.\]
Also ist $\frac 1 c$ eine untere Schranke der Menge
$\menge{\norm{Tz}_F}{z\in\kran E01}$, sodass auch
$\inf\menge{\norm{Tz}_F}{z\in\kran E01} \ge \frac 1 c > 0$ gilt.\\
Offenbar ist 
$\menge{\frac{Tz}{\norm{Tz}_F}}{z \in \kran E 0 1} \subseteq \kran G01$,
während für $w \in \kran G01$ mit der Setzung 
$v\coloneqq \frac{S\inv w}{\norm{S\inv w}_E}$ gilt:
$\norm v_E = 1$ und
\[\frac{Tv}{\norm{Tv}_F} = \frac{T\left(\frac{S\inv w}{\norm{S\inv w}_E} \right)}{\norma{T\left(\frac{S\inv w}{\norm{S\inv w}_E}\right)}_F}
= \frac{\frac 1 {\norm{S\inv w}_E} w}{\frac{\norm w_F}{\norm{S\inv w}_E}}
\overset{\norm w _F=1} = w\text.\]
Somit ist
$\mengea{\frac{Tz}{\norm{Tz}_F}}{z \in \kran E 0 1} = \kran G01$, also:
\begin{align*}
\opnorm {T\inv}{T(E)}E & = \opnorm {S\inv} G E 
        = \sup\menge{\norm{S\inv z}_E}{z \in \kran G 0 1}\\
        &= \sup\mengea{\norma{S\inv \left(
                  \frac{Ty}{\norm{Ty}_F}
                  \right)}_E}{y \in \kran E 0 1}\\
        &= \sup\mengea{\frac 1 {\norm {Ty}_F}}{y \in \kran E 0 1}\\
        &= \frac 1 {\inf\menge{\norm {Ty}_F}{y \in \kran E 0 1}}\text.
\end{align*}
\end{proof}
\parag
Das zuvor beschriebene Infimum kann also benutzt werden, um die Operatornorm
der inversen Abbildung zu bestimmen. Wir notieren das.

\begin{korollar}{Operatornorm der Inversen}
                {OperatornormDerInversen}
Seien $(E,\norm\cdot_E),(F,\norm\cdot_F)$ $\K$--Banachräume und $T\in L(E,F)$
bijektiv.\\
Dann ist $T\inv \in L(F, E)$, es gilt
$\inf\menge{\norm{Tx}_F}{x \in\kran E 0 1}>0$ und\[
\opnorm{T\inv} FE = \frac 1 {\inf\menge{\norm{Tx}_F}{x \in\kran E 0 1}}\text.\]
\end{korollar}
\begin{proof}
Da $T$ bijektiv ist, ist es insbesondere injektiv.\\
Wegen $T(E) = F$ ist $T(E)$ abgeschlossen in $F$.
Somit folgen alle Aussagen aus
~\ref{InjektivitaetLinearerOperatorenUndOperatornormDerInversen},
indem man dort überall $T(E)$ durch $F$ ersetzt.
\end{proof}
\parag
Unter Zuhilfenahme der letzten zwei Resultate wird es möglich sein, gewisse
Normen elegant gegeneinander abzuschätzen. Dazu führen wir einen neuen Begriff
ein, der diesen Sachverhalt zusammenfasst.

\begin{definition}{Bildhalbnorm}{Bildhalbnorm}

Seien $E,F$ $\K$--Vektorr�ume $\norm\cdot_F$ eine Norm auf $F$ und
$T \in \Hom(E, F)$.\\
Dann definieren wir
\[\norm\cdot_T: E \rightarrow [0, \infty[, x \mapsto \norm{Tx}_F\] 
und nennen $\norm\cdot_T$ die \textbf{Bildhalbnorm auf $E$ bez�glich $T$}.
\end{definition}

Natürlich ist die Namensgebung nicht willkürlich, denn die Bildhalbnorm ist
in der Tat eine Halbnorm und erfüllt noch eine Reihe weiterer Eigenschaften,
die wir in dem folgenden Lemma festhalten.

\begin{lemma}{Eigenschaften der Bildhalbnorm}{EigenschaftenDerBildhalbnorm}
Seien $(E, \norm\cdot_E), (F, \norm\cdot_F)$ normierte $\K$-Vektorräume,
$T \in \Hom(E, F)$.\\
Dann gilt:
\begin{enumerate}[(1)]
 \item \label{EigenschaftenDerBildhalbnorm1}
 $\norm\cdot_T$ ist eine Halbnorm auf $E$.
 \item\label{EigenschaftenDerBildhalbnorm2}
 Genau dann ist $\norm \cdot_T$ eine Norm auf $E$, wenn $T$ injektiv ist.
 \item\label{EigenschaftenDerBildhalbnorm3}
 Genau dann gibt es ein $C \in \R_{>0}$
    mit $\norm\cdot_T \le C\cdot\norm\cdot_E$, wenn $T$ stetig ist.
 \item\label{EigenschaftenDerBildhalbnorm4}
 Sind $E, F$ Banachräume, $T$ stetig, injektiv und $T(E)$ abgeschlossen, so
    ist $\norm \cdot_T$ eine zu $\norm\cdot_E$ äquivalente Norm und es gilt:
    \[\inf\menge{ \norm{Tx}_F }{x \in \kran E 0 1} \cdot \norm\cdot_E
    \le \norm\cdot_T \le \opnorm T E F \cdot \norm\cdot _E\text.\]
\end{enumerate}
\end{lemma}
\begin{proof}
Die ersten drei Aussagen sind trivial.\\
Für die vierte sei $c \coloneqq  \inf\menge{ \norm{Tx}_F }{x \in \kran E 0 1}$.
Die Voraussetzung in $(4)$ ist genau jene von~\ref{InjektivitätLinearerOperatorenUndOperatornormDerInversen},
sodass also $c > 0$ gilt.\\
Sei $x \in E\setminus\meng 0$.
Dann gilt wegen $\norma{\frac 1 {\norm{x}_E} \cdot x}_E = 1$
\begin{align*} \norm x _T = \norm {Tx}_F 
  = \norma{\norm x _E\cdot T\left(\frac 1 {\norm x _E}x \right)}_F
  = \norm x _E \cdot \norma{T\left(\frac 1 {\norm x _E}x \right)}_F
  \ge \norm x _E \cdot c\text.\end{align*}
Für $0$ ist diese Abschätzung trivial erfüllt.\\
Die zweite Abschätzung folgt direkt aus der Stetigkeit von $T$ und der
Definition der Bildhalbnorm.
\end{proof}
\parag
Als letztes vorbereitendes Mittel beschäftigen wir uns mit der Frage, wann der
Grenzwert einer Folge invertierbarer linearer Operatoren ebenfalls invertierbar
ist. Die Antwort ist sehr handlich, es ist nämlich genau dann der Fall, wenn
die Folge der Inversen beschränkt ist.

\begin{lemma}{Invertierbarkeit des Grenzwerts invertierbarer Operatoren}
             {InvertierbarkeitDesGrenzwertsInvertierbarerOperatoren}
Seien $(E, \norm\cdot_E)$ ein $\K$--Banachraum, $(F, \norm\cdot_F)$ ein
normierter $\K$--Vektorraum, $\folge T n \N \in \Inv(E, F)^\N$ konvergent
und $T \in L(E, F)$ sei der Grenzwert dieser Folge. 
Dann gilt
\begin{enumerate}[(1)]
 \item\label{InvertierbarkeitDesGrenzwertsInvertierbarerOperatoren1}
 $\left(T_n\inv\right)_{n \in \N}$ ist konvergent $\iff$
 $\left(T_n\inv\right)_{n \in \N}$ ist beschränkt.
 \item\label{InvertierbarkeitDesGrenzwertsInvertierbarerOperatoren2}
 Ist eine der äquivalenten Bedingungen aus (1) erfüllt, so ist
 $T$ invertierbar und $T\inv = \Limes n \infty T_n\inv$.
\end{enumerate}
\end{lemma}
\begin{proof} (1) Da konvergente Folgen stets beschränkt sind, ist lediglich
die Umkehrung zu zeigen.\\
Sei also $\left(T_n\inv\right)_{n \in \N}$ beschränkt.\\
Dann gibt es ein $c \in \R_{>0}$ mit
$\forall\, n \in \N: \opnorm{T_n\inv}F E \le c$.\\
Da $(E, \norm\cdot_E)$ vollständig ist, ist $(L(F, E), \opnorm \cdot F E)$ 
ein Banachraum, sodass es ausreicht zu zeigen, dass $\folge {T\inv} n \N$ eine
Cauchyfolge ist.\\
Sei also $\eps \in \R_{>0}$.\\
Da $\folge T n \N$ konvergiert, ist dies eine Cauchyfolge. Somit existiert also
ein $n_0 \in \N$ mit
$\forall\, n, m \in \N_{\ge n_0}: \opnorm{T_n - T_m} E F < \frac \eps {c^2}$.\\
Seien $n, m \in \N_{\ge n_0}$. Dann gilt:
\begin{align*}
\opnorm{T_n\inv - T_m\inv}FE 
 & = \opnorm{\left(T_n\inv T_n \right)(T_n\inv - T_m\inv)}FE\\
   &\le \opnorm{T_n\inv}FE\opnorm{T_nT_n\inv-T_nT_m\inv}F F\\
   &=c\opnorm{T_mT_m\inv -T_nT_m\inv}FF & \quad T_nT_n\inv = T_mT_m\inv\\
   &=c\opnorm{\left(T_m -T_n\right)T_m\inv}FF\\
   & \le c \opnorm{T_m\inv} F E \opnorm{T_m - T_n} E F\\
   & \le c^2 \opnorm{T_m - T_n}EF < \eps
\end{align*}
Damit ist die erste Aussage gezeigt.\absatz
(2) Sei also $\left(T_n\inv\right)_{n \in \N}$ beschränkt. Nach (1) ist diese
Folge dann konvergent gegen ein $T_0 \in L(X, Y)$.\\
Sei $L(F,E) \times L(E, F)$ mit der Produktnorm ausgestattet und
\[\alpha : L(F, E) \times L(E, F) \rightarrow L(E, E), (A,B) \mapsto AB\text,\]
\[\beta  : L(F, E) \times L(E, F) \rightarrow L(F, F), (A,B) \mapsto BA\text.\]
Diese beiden Abbildungen sind offensichtlich bilinear und wegen der
Submultiplikativität der Operatornorm stetig.\\
Dann gilt mit $(1)$: $\left( \alpha(T_n\inv, T_n)\right)_{n \in \N}$ ist
konvergent und
\begin{align*}
\id_E &= \Limes n \infty T_n\inv T_n = \Limes n \infty \alpha (T_n\inv, T_n)
      = \alpha\left( \Limes n \infty T_n\inv, \Limes n \infty T_n\right)
      = \alpha\left(T_0, T\right)\text.
\end{align*}
Damit ist also $T_0$ eine Linksinverse von $T$.\\
Analog:
\begin{align*}
\id_F & = \Limes n  \infty T_n T_n\inv  = \Limes n \infty \beta (T_n\inv, T_n)
      = \beta\left( \Limes n \infty T_n\inv, \Limes n \infty T_n\right)
      = \beta\left(T_0, T\right)\text.
\end{align*}
Daher ist $T_0$ auch eine Rechtsinverse von $T$.\\
Also ist $T$ invertierbar mit $T\inv = T_0$.
\end{proof}
\parag
Wir definieren nun den Banach--Mazur--Abstand zweier normierter Räume. Dieser
kann dazu verwendet werden, optimale untere Normäquivalenzkonstanten zu
erhalten, sofern die obere bereits 1 ist.\\
In der Definition wird ein Infimum einer Menge verwendet, die zunächst leer
sein könnte -- in diesem Fall vereinbaren wir, dass $\inf \emptyset = \infty$
gelten möge, um keine weitere Forderung stellen zu müssen.
In sämtlichen späteren Anwendungen (vor allem im endlichdimensionalen Fall) wird die erwähnte Menge stets nichtleer sein.

\begin{definition}{Banach--Mazur--Abstand, Banach--Mazur--Abbildung}
                  {BanachMazurAbstand}
Seien $(E, \norm\cdot_E), (F, \norm\cdot_F)$ normierte $\K$--Vektorr�ume.\\
Dann hei�t
\[d_{BM}\left( (E, \norm\cdot_E), (F, \norm\cdot_F) \right) := 
\inf\menge{\opnorm S E F \cdot \opnorm {S\inv} F E}{S \in \Inv(E, F)}\]
\textbf{Banach--Mazur--Abstand zwischen $(E, \norm\cdot_E)$ und 
$(F, \norm\cdot_F)$}.\\
Wenn aus dem Zusammenhang klar ist, welche Norm auf $E$ bzw. $F$ vorliegt,
werden wir auch $d_{BM}(E,F)$ schreiben.\\
Durch oben erw�hnte Festsetzung $\inf \emptyset = \infty$ ist $d_{BM}$ f�r je
zwei normierte R�ume sinnvoll erkl�rt.\\
Ist $\Inv(E, F) \neq \emptyset$, so hei�t ein $T \in \Inv(E, F)$ \textbf{Banach--Mazur--Abbildung zwischen $(E, \norm\cdot_E)$ und 
$(F, \norm\cdot_F)$}
\[:\iff \opnorm T E F = 1 \quad \wedge \quad 
\opnorm {T\inv} F E 
 = d_{BM}\left( (E, \norm\cdot_E), (F, \norm\cdot_F) \right)\text.\]
Insbesondere gilt also f�r eine Banach--Mazur--Abbildung $T$:
\[\opnorm T E F \cdot \opnorm {T\inv} F E 
= d_{BM}\left( (E, \norm\cdot_E), (F, \norm\cdot_F) \right)\text.\]
\end{definition}
\parag
Die Normiertheitsbedingung einer Banach--Mazur--Abbildung dient lediglich der
Bequemlichkeit.
Als Nächstes notieren wir einige Eigenschaften des
Banach--Mazur--Abstandes.

\begin{lemma}{Eigenschaften des Banach--Mazur--Abstands}
             {EigenschaftenDesBanachMazurAbstands}
Seien $(E, \norm\cdot_E), (F, \norm\cdot_F), (G, \norm\cdot_G)$
normierte $\K$--Vektorräume.\\
Dann gilt:
\begin{enumerate}[(1)]
	\item \label{EigenschaftenDesBanachMazurAbstands1}
	Sind $\Inv(E, F) \neq \emptyset$ und $\Inv(F, G) \neq \emptyset$,
	so auch $\Inv(E, G) \neq \emptyset$ und es gilt
	\[d_{BM}(E,G)\le d_{BM}(E,F)d_{BM}(F,G)\text.\]
	\item \label{EigenschaftenDesBanachMazurAbstands2}
	Sind $(E, \norm\cdot_E)$ und $(F, \norm\cdot_F)$ isometrisch isomorph, so ist
	$d_{BM}(E,F) = 1$.
	\item \label{EigenschaftenDesBanachMazurAbstands3}
	Sind $\norm\cdot$ eine weitere Norm auf $E$ und $c, d \in \R_{>0}$
	mit $c \norm\cdot \le \norm\cdot_E \le d \norm\cdot$, so ist
	\[d_{BM}\bigg((E, \norm\cdot_E), (E, \norm\cdot) \bigg) \le \frac d c
	\text.\]
	\item\label{EigenschaftenDesBanachMazurAbstands4}
	$d_{BM}(E,F)= d_{BM}(F,E)$.
\end{enumerate}
\end{lemma}
\begin{proof}
Es gilt für alle $T \in \Inv(E, F)$:
\[ 1 = \opnorm{\id_E}EE = \opnorm{T\inv \circ T} E E
\le \opnorm T EF \opnorm{T\inv}FE\]
und daher auch stets
\begin{align}d_{BM}\bigg((E, \norm\cdot_E), (F, \norm\cdot_F)\bigg) \ge 1\text.
\tag{$\ast$}\label{edbmaeq:1}
\end{align}
(1) Seien $R \in \Inv(E, F)$ und $S \in \Inv(F, G)$.\\
Dann ist $S \circ R$ linear, invertierbar und stetig. Wegen $(S \circ R)\inv
= R\inv \circ S\inv$ ist aber auch $(S \circ R)\inv$ stetig und linear und
damit $S \circ R \in \Inv(E, G)$.\\
Daher gilt auch
\begin{align*}d_{BM}\left( (E, \norm\cdot_E), (G, \norm\cdot_G)\right)
&\le \opnorm{S\circ R}EG \opnorm{(S \circ R)\inv}G E\\
&\le \opnorm SFG \opnorm{S\inv}GF \opnorm REF \opnorm{R\inv}FE\text.
\end{align*}
Wegen $\opnorm REF \opnorm{R\inv}FE > 0$ ist
$\frac{d_{BM}\left( (E, \norm\cdot_E), (G, \norm\cdot_G)\right)}
{\opnorm REF \opnorm{R\inv}FE}$ eine untere Schranke von
$\menge{\opnorm A F G \opnorm{A\inv}GF}{A \in \Inv(F,G)}$ und damit
\[\frac{d_{BM}\left( (E, \norm\cdot_E), (G, \norm\cdot_G)\right)}
{\opnorm REF \opnorm{R\inv}FE} \le 
d_{BM}\left( (F, \norm\cdot_F), (G, \norm\cdot_G)\right)\text.\]
Analog ist nun (wegen \eqref{edbmaeq:1} wohldef.) $\frac{d_{BM}\left((E, \norm\cdot_E), 
(G, \norm\cdot_G)\right)}{d_{BM}\left( (F, \norm\cdot_F), 
(G, \norm\cdot_G)\right)}$ eine untere Schranke von
\[\menge{\opnorm B E F \opnorm{B\inv}FE}{B \in \Inv(E,F)}\] und damit
\[\frac{d_{BM}\left( (E, \norm\cdot_E), 
(G, \norm\cdot_G)\right)}{d_{BM}\left( (F, \norm\cdot_F), 
(G, \norm\cdot_G)\right)} \le d_{BM}\left( (E, \norm\cdot_E), 
(F, \norm\cdot_F)\right)\text.\]
(2) Seien $(E, \norm\cdot_E)$ und $(F, \norm\cdot_F)$ isometrisch isomorph.\\
Dann existiert $\varphi : E \rightarrow F$ linear, stetig und bijektiv mit
$\norm \cdot _F \circ \varphi = \norm\cdot _E$.\\
Da $\varphi$ insbesondere eine bijektive Isometrie ist, ist auch $\varphi\inv$ stetig
und es gilt $\norm\cdot_F = \norm\cdot_E \circ \varphi \inv$.\\
Somit ist $\varphi \in \Inv(E, F)$ und es ist $\opnorm \varphi E F = 1$, sowie
$\opnorm {\varphi\inv} F E = 1$.
Also
\[d_{BM}\bigg((E, \norm\cdot_E), (F, \norm\cdot_F)\bigg)
\le \opnorm \varphi E F \opnorm {\varphi\inv} F E = 1\text.\]
Wegen \eqref{edbmaeq:1} gilt aber auch die umgekehrte Ungleichung, also
Gleichheit.\\
(3) Seien $\norm\cdot$ eine Norm auf $E$ und $c, d \in\!\ ]0,\infty[$
mit $c \norm\cdot \le \norm\cdot_E \le d \norm\cdot$.\\
Sei $I : (E, \norm\cdot) \rightarrow (E, \norm\cdot_E), x \mapsto x$.\\
Offenbar ist $I$ linear und invertierbar.\\
Wegen $\norm\cdot_E \le d \norm\cdot$ ist $I$ stetig und es gilt
$\opnorm I {(E, \norm\cdot)}{(E, \norm\cdot_E)} \le d$.\\
Analog ist wegen $\norm\cdot \le \frac 1 c \norm \cdot_E$ auch $I\inv$ stetig
und es gilt $\opnorm {I\inv}{(E, \norm\cdot_E)}{(E, \norm\cdot)}\le\frac 1c$.\\
Also ist
\[d_{BM}\bigg((E, \norm\cdot_E), (E, \norm\cdot) \bigg) \le
\opnorm I {(E, \norm\cdot)}{(E, \norm\cdot_E)}
\opnorm {I\inv}{(E, \norm\cdot_E)}{(E, \norm\cdot)} \le \frac d c\text.\]
(4) Offenbar gilt für $T \in L(E, F)$
\[T \in \Inv(E, F) \iff T\inv \in \Inv(F, E)\]
und $(T\inv)\inv = T$, sodass gilt:
\[\mengea{\opnorm S E F \opnorm{S\inv}FE}{S \in \Inv(E, F)}
= \mengea{\opnorm R F E \opnorm{R\inv}FE}{R \in \Inv(F, E)}\text.\]
Daher stimmen die Infima dieser Mengen überein.
\end{proof}

Wir stellen uns nun die Frage, ob es stets eine Banach--Mazur--Abbildung
zwischen zwei isomorphen Räumen gibt.
Im Unendlichdimensionalen braucht es
diese Abbildungen nicht zu geben, im Endlichdimensionalen hingegen existieren
stets Banach--Mazur--Abbildungen.

\begin{satz}{Existenz einer Banach--Mazur--Abbildung}
            {ExistenzEinerBanachMazurAbbildung}
Seien $(E, \norm\cdot_E)$, $(F, \norm\cdot_F)$
endlichdimensionale normierte $\K$--Vektorräume mit $\dim_\K E = \dim_\K F$.\\
Dann existiert eine Banach--Mazur--Abbildung zwischen $(E, \norm\cdot_E)$ und
$(F, \norm\cdot_F)$.
\end{satz}
\begin{proof}
Da die Vektorräume $E, F$ endlichdimensional sind, gibt es bijektive lineare
Abbildungen von $E$ nach $F$ und diese, sowie deren Umkehrabbildungen, sind
stetig.
Damit ist $\Inv(E, F) \neq \emptyset$.\\
Sei $d\coloneqq d_{BM}\left( (E, \norm\cdot_E), (F, \norm\cdot_F) \right)$.\\
Aufgrund der Infimumseigenschaft existiert zu jedem $n \in \N$ ein
$S_n \in \Inv(E, F)$ mit 
$d + \frac 1 n > \opnorm {S_n} E F \opnorm {S_n\inv}FE$.\\
Sei $\folge S n \N$ eine Auswahl solcher Elemente und
$T_n \coloneqq  \frac 1 {\opnorm {S_n} E F} S_n$ für alle $n \in \N$.\\
Dann gilt offenbar: 
$T_n \in \Inv(E, F)$ und $T_n\inv = \opnorm{S_n} F E \cdot S_n\inv$
für alle $n \in \N$.\\
Somit gilt für alle $n \in \N$:\[
\opnorm{T_n\inv}FE=\opnorm{S_n}EF\cdot\opnorm{S_n\inv}FE<d+\frac1n\le d+1
\text.\]
Nun ist auch $\folge T n \N \in \kran {L(E, F)} 0 1 ^\N$.\\
Da $E$ und $F$ endlichdimensional sind, ist auch $L(E, F)$ endlichdimensional
und somit ist $\kran{L(E,F)}01$ kompakt,
also existiert eine
Teilfolge $\left(T_{n_j}\right)_{j \in \N}$ von $\folge T n \N$, die gegen ein
$T_0 \in \kran{L(E, F)}01$ konvergiert.\\
Die Folge $(T_{n_j}\inv)_{j\in \N}$ ist offenbar ebenfalls
beschränkt.\\
Da $(E, \norm\cdot_E)$ als endlichdimensionaler normierter Raum vollständig
ist, erhalten wir mit~\ref{InvertierbarkeitDesGrenzwertsInvertierbarerOperatoren} die Konvergenz von
$\left(T_{n_j}\inv \right)_{j\in \N}$ und es gilt
\[T_0\text{ ist invertierbar mit }T_0\inv=\Limes j \infty T_{n_j}\inv\text. \]
Nach Konstruktion gilt für alle $n \in \N$\[ 
\abs{\opnorm {T_n} E F \opnorm{T_n\inv} F E - d} 
=\opnorm{T_n}EF\opnorm{T_n\inv}FE-d=\opnorm{T_n\inv}FE-d<\frac 1n\text.\]
Daher gilt\[
\opnorm {T_0\inv} F E = \Limes j \infty \opnorm {T_{n_j}\inv} F E = d\text.\]
Wegen $T_0 \in \kran{L(E, F)}01$ ist $T_0$ also eine Banach--Mazur--Abbildung.
\end{proof}
\parag
Im Beweis ist eingegangen, dass $(E, \norm\cdot_E)$
ein Banachraum ist, $\Inv(E, F)\neq \emptyset$ gilt und $\kran{L(E,F)}01$
kompakt ist.
Die beiden letztgenannten Voraussetzungen implizieren aber die
endliche Dimension von $E$ und $F$, wie das nachfolgende Lemma zeigt, sodass
dieser Satz mit dem obigen Beweis nicht weiter verallgemeinert werden kann.

\begin{lemma}{Kompaktheit der Sph�re}{PraekompaktheitDerSphaere}
Sei $(E, \norm\cdot_E)$ ein normierter $\K$--Vektorraum.\\
Dann gilt:
\begin{enumerate}[(1)]
	\item Ist $\kran E01$ kompakt, so auch $\kabg E01$.
	\item Ist $(F, \norm\cdot_F)$ ein weiterer normierter $\K$--Vektorraum,
	      $\kran {L(E,F)}01$ kompakt und $\Inv(E,F) \neq \emptyset$, so sind
	      $E$ und $F$ endlichdimensional.
\end{enumerate}
\end{lemma}
\begin{proof}
(1) Die Aussage ist trivial, wenn $E = \meng 0$.\\
Sei also zus�tzlich $E \neq \meng 0$ angenommen.\\
Dann ist $\kran E 01 \neq \emptyset$ und offensichtlich gilt
\begin{align}\kabg E 01 = \menge{sy}{s \in [0,1], y \in \kran E 01}\text. 
\tag {$\ast$}\label{pdseq:1}\end{align}
Sei nun $K:=[0,1] \times \kran E01$ mit der Produktmetrik ausgestattet.\\
Ist $\kran E01$ kompakt, so ist auch $K$ kompakt.\\
Seien $f: K \rightarrow \kabg E 01, (s,y) \mapsto sy$. Dann ist $f$ stetig.\\
Wegen $\eqref{pdseq:1}$ ist $f(K)=\kabg E01$ und damit ist $\kabg E01$ als
stetiges Bild einer kompakten Menge selbst kompakt.\\
(3) Sei $(F, \norm\cdot_F)$ ein normierter $\K$--Vektorraum und es gelte
$\kran {L(E,F)}01$ ist kompakt und $\Inv(E, F) \neq \emptyset$.\\
Nach (2) ist dann $\kabg{L(E,F)}01$ kompakt, also $L(E,F)$
endlichdimensional.\\
Ist $F= \meng 0$, so ist wegen $\Inv(E,F) \neq \emptyset$ auch $E = \meng 0$.\\
Ist $F\neq0$, so existiert nach Hahn--Banach ein $\vi\in F'\setminus\meng 0$.\\
Sei $\Phi: F \rightarrow L(E,F), y \mapsto (x \mapsto \vi(x)y)$.\\
Offensichtlich ist $\Phi$ linear und wegen $\vi \neq 0$ auch injektiv.\\
Damit ist $F$ isomorph zu einem Unterraum des endlichdimensionalen Raums
$L(E,F)$, also endlichdimensional und
wegen $\Inv(E,F) \neq \emptyset$, ist es auch $E$.
\end{proof}
\parag
Eine nützliche Folgerung aus diesem Satz \ref{ExistenzEinerBanachMazurAbbildung}
ist die Existenz einer Basis, welche die Äquivalenz der Bildhalbnorm auf dem
$\K^n$ zu jeder anderen Norm nur mit Hilfe des Banach--Mazur--Abstands
ausdrückt.

\begin{korollar}{Normäquivalenzbasis im Endlichdimensionalen}
                {NormaequivalenzbasisImEndlichdimensionalen}
Seien $n \in \N$, $(E, \norm\cdot_E)$ ein $n$--dimensionaler normierter
$\K$--Vektorraum, $\norm\cdot$ eine Norm auf dem $\K^n$ und 
$d \coloneqq  d_{BM}((\K^n, \norm\cdot), (E, \norm\cdot_E))$.\\
Dann existiert eine Basis $x \in E^n$ mit der Eigenschaft
\[\frac 1 d \norm\cdot \le\norm\cdot_E \circ \bprod\cdot x_{n, E} \le \norm\cdot\text.\]
\end{korollar}
\begin{proof}
Da $E$ und $\K^n$ endlichdimensional sind und die gleiche Dimension haben,
existiert nach~\ref{ExistenzEinerBanachMazurAbbildung} eine
Banach--Mazur--Abbildung 
$T \in \Inv(\K^n, E)$ zwischen $(\K^n, \norm\cdot)$ und $(E, \norm\cdot_E)$,
also sind 
\begin{align}\opnorm T{\K^n}E=1 \quad \text{ und }\quad 
\opnorm{T \inv}E{\K^n} = d\text. \label{nieeq:1}\end{align}
Seien $e \coloneqq  (e_i^n)_{i=1}^n$ und $x \coloneqq  \pfeil T n (e)$. Offenbar ist $x$ eine Basis von $E$.\\
Sei $t \in \K^n$. Dann gilt:
\[ \bprod t x _{n, E}= \bprod t {\pfeil T n (e)} _{n, E}
\overset{\refclaim{EigenschaftenDesLinearkombinators}
{EigenschaftenDesLinearkombinators2}}{=} 
T\left(\bprod t e _{n, \K^n}\right) = T(t)\text.\]
Somit ist also $T = \bprod \cdot x_{n, E}$ und daher 
\begin{align}\norm\cdot_T = \norma\cdot_E \circ \bprod \cdot x _{n, E}\text.
\label{nieeq:2}\end{align}
Wegen der endlichen Dimension sind $(E, \norm\cdot_E)$ und $(\K^n, \norm\cdot)$
vollständig, $T$ ist injektiv und hat abgeschlossenes Bild.\\
Daher liefert~\ref{EigenschaftenDerBildhalbnorm}:
\begin{align}\inf\menge{ \norm{Ta}_E }{a \in \kran {\K^n} 0 1} \cdot \norm\cdot
    \le \norm\cdot_T \le \opnorm TE{\K^n}\cdot \norm\cdot\text. \label{nieeq:3}
    \end{align}
Auch die Voraussetzungen von~\ref{OperatornormDerInversen} sind erfüllt,
sodass wir weiterhin erhalten:
\begin{align}\label{nieeq:4}\inf\menge{ \norm{Ta}_E }{a \in \kran {\K^n} 0 1} 
 = \frac 1 {\opnorm{T\inv}E {\K^n}} \overset{\eqref{nieeq:1}}= \frac 1 d\text. \end{align}
Insgesamt folgt:
\begin{align*}
\frac 1 d \norm\cdot &\overset{\eqref{nieeq:4}}=
 \left(\inf\menge{\norm{Ta}_E }{a\in\kran{\K^n}01}\right) \cdot \norm\cdot \\
 &\overset{\eqref{nieeq:3}}\le \norm\cdot_T \overset{\eqref{nieeq:2}}= 
 \norm\cdot_E \circ \bprod \cdot x _{n, E}
 \overset{\substack{\eqref{nieeq:2},\\\eqref{nieeq:3}}}\le 
 \opnorm T {\K^n} E \norm\cdot \overset{\eqref{nieeq:1}}
 = \norm \cdot \text.
\end{align*}
\end{proof}
\parag
Insbesondere gibt es eine Basis $x \in E^n$ mit der Eigenschaft
\[\frac 1 d \norm \cdot _{\infty, n} \le \frac 1 d \norm\cdot_{2,n}
\le \norm\cdot_E \circ \bprod \cdot x_n \le \norm\cdot_{2,n}\text.\]
Interessiert man sich für Abschätzungen der Art
$c\norm\cdot_{\infty,n}\le\norm\cdot_E\circ\bprod \cdot y_n\le\norm\cdot_{2,n}$
für ein $c \in \R_{>0}$ und $y \in E^n$ (nicht notwendigerweise Basis), so
liefert das Lemma von Dvoretzky--Rogers eine bessere Aussage.

\begin{lemma}{Dvoretzky--Rogers--Lemma}{DvoretzkyRogersLemma}
Seien $n, m \in \N$ mit $n \ge 4m^2$, sowie $(E, \norm\cdot_E)$ ein normierter
$\R$--Vektorraum mit $\dim E = n$.\\
Dann gibt es $x \in E^m$ mit
\[\frac 1 2 \norm\cdot_{\infty, m} \le \norm\cdot_E \circ \bprod \cdot x_m
\le \norm\cdot_{2,m}\]
\end{lemma}
\begin{proof}
Ein Beweis dieser Aussage befindet sich in \cite{Kad}, dort Lemma 6.2.1.
Leider enth�lt die dortige Absch�tzung in der letzten Zeile einen kleinen
Fehler, der aber wie folgt (gem�� dem urspr�nglichen Beweis von Dvoretzky und
Rogers, vgl. \cite{DR}) behoben werden kann (die Notation lehnt sich an jene
aus \cite{Kad} an):
\begin{quote}
Sei $i \in \haken m$ mit $i+1 \in \haken m$ und die Behauptung gelte f�r alle
$k \le i$.\\
Wegen $i < i+1 \le m \le \frac {\sqrt n} 2$ gilt:
\[i\sqrt n + i^2 + i = i\left(\sqrt n+i+1\right)
\le\frac{\sqrt n}2\left(\sqrt n+\frac{\sqrt n}2\right)
= \frac 3 4 n\text.\]
Daher ist
\[4\frac i {\sqrt n} + 4\frac{i^2}n+ 4 \frac i n \le 3,\]
also
\begin{align}\left(1+\frac {2i}{\sqrt n} \right)^2
&= 1 + 4 \frac i{\sqrt n} +  4\frac {i^2}n \le 4 - 4\frac i n 
= 4 \left(1- \frac i n \right)\text. \tag{$\ast$}\label{drleq:1}
\end{align}
Damit gilt wie im Beweis angegeben:
\begin{align*}
\abs{t_{i+1}}^2 &\le \frac n {n-i}\left(1+ 2\sqrt{i}\sqrt{\frac i n} \right)^2
=\frac 1{1-\frac in}\left(1+\frac{2i}{\sqrt n}\right)^2
\overset{ \eqref{drleq:1}}\le 4\text.
\end{align*}
\end{quote}
Dies liefert genau die f�r den Beweis ben�tigte Aussage.
\end{proof}
\parag
Wir halten nun kanonische, aber scharfe Äquivalenzkonstanten für
$p$--Normen auf dem $\K^n$ fest. Ein Teil dieser Abschätzungen (der nicht von
$n$ abhängt) ist auch für $\ell_p(\N, \K)$ mit dem jeweils gleichen Beweis
gültig. In der gewählten Formulierung sieht man jedoch die Äquivalenzkonstanten
direkt in der Formulierung.

\begin{lemma}{Äquivalenzkonstanten für p--Normen}
             {AequivalenzkonstantenFuerPNormen}
Seien $n \in \N$ und $p \in [1, \infty[$.\\
Dann gilt:
\begin{enumerate}[(1)]
\item\label{AequivalenzkonstantenFuerPNormen2}
 $\norm\cdot_{\infty, n} \le \norm\cdot_{p, n} 
     \le n^{\frac 1 p}\norm\cdot_{\infty, n}$.
\item\label{AequivalenzkonstantenFuerPNormen3}
 Ist zusätzlich $q \in [1, \infty[$ mit $p \le q$, so gilt
$\norm\cdot_{q,n}\le\norm\cdot_{p,n}\le n^{\frac 1p-\frac 1q}\norm\cdot_{q,n}$.
\item\label{AequivalenzkonstantenFuerPNormen1}
 $\norm\cdot_{p,n}\le\norm\cdot_{1, n}\le n^{\frac{p-1}p}\norm\cdot_{p,n}$.
\end{enumerate}
Ferner sind alle Ungleichungen optimal.
\end{lemma}
\begin{proof}
Offensichtlich gelten alle Aussagen für den Nullvektor.\\
Sei $x \in \K^n\setminus\meng 0$.\\
(1)
Da $f: \R_{>0} \rightarrow \R, t \mapsto t^{\frac 1 p}$ monoton wachsend
ist, gilt für alle $j \in \haken n$:
\[\abs{x_j} = \left(\abs{x_j}^p \right)^{\frac 1p}
\le \left(\sum_{i=1}^n \abs{x_i}^p\right)^{\frac 1 p}\]
und damit $\norm x _{\infty, n} \le \norm x _{p,n}$.
Außerdem gilt\[
\norm x _{p, n}^p = \sum_{i=1}^n \abs{x_i}^p \le 
\sum_{i=1}^n \max\menge{\abs{x_i}}{i\in\haken n}^p=n \norm x _{\infty, n}^p
\text.\]
Damit ist auch $\norm \cdot_{p, n} \le n^{\frac 1 p}\norm\cdot_{\infty, n}$
gezeigt.\\
(2) Sei $q \in [1, \infty[$ mit $p \le q$.\\
Falls $p=q$ ist nichts zu zeigen, sodass wir $p < q$ annehmen können.\\
Dann gilt:
\begin{align*}
\norm x _{q,n}^q&= \sum_{i=1}^n \abs{x_i}^q
=\sum_{i=1}^n \abs{x_i}^p\abs{x_i}^{q-p}
\le \left(\sum_{i=1}^n \abs{x_i}^p \right)\cdot 
  \max_{i \in \haken n}\left(\abs{x_i}^{q-p}\right)\\
 &= \norm x_{p,n}^p \max_{i \in \haken n}\left(\abs{x_i}^{q-p}\right)
 \le \norm x_{p,n}^p \left(\max_{i \in \haken n}\abs{x_i}\right)^{q-p}
 = \norm x _{p, n}^p \norm x _{\infty, n}^{q-p}\\
 &\overset{(1)}\le \norm x _{p, n}^p \norm x _{p, n}^{q-p} = \norm x _{p,n}^q
 \text.\end{align*}
Also auch $\norm x _{q, n} \le \norm x _{p, n}$.\\
Seien nun $p'\coloneqq\frac qp$ und $q'\coloneqq  \frac q{q-p}$ (wohldefiniert wegen $p<q$).\\
Dann gilt
\[\frac 1 {p'} + \frac 1 {q'} = \frac p q + \frac{q-p}q = 1\text{.}\]
Also liefert die diskrete Hölderungleichung
\begin{align*}
\norm x _{p,n}^p &= \sum_{i=1}^n \abs{x_i}^p \cdot 1 \le
\left(\sum_{i=1}^n \left(\abs{x_i}^p\right)^{p'} \right)^{\frac 1 {p'}}
\left(\sum_{i=1}^n 1 \right)^{\frac 1 {q'}}\\
&= \left(\sum_{i=1}^n  \left(\abs{x_i}^p\right)^{\frac q p}\right)^{\frac pq}
 n^{\frac {q-p}q}
 = \left(\sum_{i=1}^n \abs{x_i}^q \right)^{\frac p q} n^{\frac {q-p}q}
 = \norm x _{q, n}^p n ^{\frac {q-p}q}
\end{align*}
und damit
\[\norm x_{p, n} \le n ^\frac{q-p}{qp} \norm x _{q, n} 
= n ^{\frac 1 p - \frac 1 q}\norm x _{q, n}\text.\]
(3) Wegen $p \ge 1$ gilt nach (2)
\[\norm\cdot_{p,n} \le \norm\cdot_{1, n} 
\le n^{1 - \frac 1 p}\norm \cdot_{p,n}
= n ^{\frac {p-1}p}\norm \cdot_{p,n}\text.\]
Die Optimalität der ersten Ungleichung in allen drei Punkten liefert $e_1^n$,
jede der zweiten $\sum_{i=1}^n e_i^n$.
\end{proof}
\parag
Die beiden folgenden Lemmata werden wir dafür verwenden, das Supremum bzw.
Infimum einer stetigen Funktion über einer dichten Teilmenge darzustellen
anstelle der Menge selbst. Außerdem ist es möglich, Häufungspunkte aus
einer dichten Teilmenge zu entfernen und weiterhin Dichtheit zu haben.

\begin{lemma}{Stetige Funktionen und dichte Teilmengen}
             {StetigeFunktionenUndDichteTeilmengen}
Seien $(M, d)$ ein metrischer Raum, $f: M \rightarrow \R$ stetig.\\
Seien weiter $\emptyset \neq A \subseteq M$ und $A'\subseteq A$ dicht in $A$.\\
Wegen $A \neq \emptyset$ ist dann $A' \neq \emptyset$, sodass auch
$f(A)\neq\emptyset\neq f(A')$ gilt.\\
Dann gilt:
\begin{enumerate}[(1)]
  \item \label{StetigeFunktionenUndDichteTeilmengen1}
	Ist $f(A)$ nach oben beschr�nkt, so auch $f(A')$ mit
	 $\sup f(A) = \sup f(A')$.
	\item \label{StetigeFunktionenUndDichteTeilmengen2}
	Ist $f(A)$ nach unten beschr�nkt, so auch $f(A')$ mit 
	$\inf f(A) = \inf f(A')$.
\end{enumerate}
\end{lemma}
\begin{proof}
(1) Sei also $f(A)$ nach oben beschr�nkt.\\
Wegen $f(A) \neq \emptyset$ besitzt $f(A)$ also ein Supremum.\\
Da $f(A') \subseteq f(A)$ und $f(A') \neq \emptyset$, besitzt
auch $f(A')$ ein Supremum und es gilt $\sup f (A') \le \sup f(A)$.\\
Sei nun $\eps \in \R_{>0}$.
Dann gibt es ein $a \in A$ mit $\sup f(A) - f(a) < \frac \eps 2$.\\
Weil $f$ stetig ist, gibt es $\delta \in \R_{>0}$, sodass gilt:
$\forall\, b \in K_M(a,\delta): \abs{f(a)-f(b)} < \frac \eps 2$.\\
Da $A'$ dicht in $A$ ist, gibt es ein $a' \in A'$ mit $d(a, a') < \delta$,
also $\abs{f(a)-f(a')}< \frac \eps 2$.\\
Somit gilt:
\begin{align*}
\sup f(A) - f(a') & = \sup f(A) - f(a) + f(a) -f(a')\\
&\le(\sup f(A)-f(a))+\abs{f(a)-f(a')}<\frac\eps 2+\frac\eps 2 = \eps\text.
\end{align*}
Also ist $\sup f(A) = \sup f(A')$.\\
(2) Sei also $f(A)$ nach unten beschr�nkt. Dann ist $-f(A)$ nach oben
beschr�nkt und daher auch $-f(A')$ nach oben beschr�nkt und nichtleer.\\
Es ist auch $-f$ stetig und $-(f(A)) = (-f)(A)$, sodass gilt \[
\inf f(A) = - \sup(-f(A)) \overset{(1)}= -\sup(-f(A'))= \inf f(A')\text. \]
\end{proof}

\begin{lemma}{Dichtheit und H�ufungspunkte}
             {DichtheitUndHaeufungspunkte}
Seien $(M, d)$ ein metrischer Raum, $x_0\in M$
ein H�ufungspunkt von $M$, sowie $D \subseteq M$ dicht in $M$.\\
Dann gilt:
\begin{enumerate}[(1)]
	\item \label{DichtheitUndHaeufungspunkte1}
	Es gibt $A \subseteq M\setminus\meng{x_0}$ abz�hlbar, sodass
	$(D\setminus\meng{x_0}) \cup A$ dicht in $M$ ist.
	\item \label{DichtheitUndHaeufungspunkte2}
	Ist $D$ abz�hlbar, so gibt es $D'' \subseteq M$ abz�hlbar und dicht in
	$M$ mit $x_0 \notin D''$.
\end{enumerate}
\end{lemma}
\begin{proof}
Da $x_0$ ein H�ufungspunkt von $M$ ist, gibt es eine gegen $x_0$ konvergente
Folge $\folge xn\N\in(M\setminus\meng{x_0})^\N$.\\Dann ist
$A:=\menge{x_n}{n\in \N} \subseteq M\setminus\meng{x_0}$ und abz�hlbar.\\
Sei $D':= \left(D \setminus \meng{x_0} \right) \cup A$.\\
Dann ist $x_0\notin D'$ und falls $D$ abz�hlbar ist, ist auch $D'$ als endliche
Vereinigung abz�hlbarer Mengen abz�hlbar.\\
Da $\folge x n \N$ gegen $x_0$ konvergiert, ist
$\menge{x_n}{n \in \N} \cup \meng{x_0}$ kompakt und damit insbesondere
abgeschlossen, also 
$\overline{\menge{x_n}{n\in\N}}\subseteq\menge{x_n}{n\in\N}\cup\meng{x_0}$.\\
Offenbar ist $x_0$ als Grenzwert von $\folge x n \N$ wegen $x_n\neq x_0$ f�r
alle $n \in \N$ ein H�ufungspunkt von $\menge{x_n}{n \in \N}$ und damit 
$\overline{\menge{x_n}{n\in\N}}\supseteq\menge{x_n}{n\in\N}\cup\meng{x_0}$.\\
Also $\overline{\menge{x_n}{n\in\N}}=\menge{x_n}{n\in\N}\cup\meng{x_0}$.\\
Es gilt weiter:
\begin{align*}\overline{D'} 
&=\overline{\left(D \setminus \meng{x_0} \right) \cup \menge{x_n}{n \in \N}}
=\overline{D \setminus \meng{x_0}} \cup \overline {\menge{x_n}{n \in \N}}\\
&= \overline{D \setminus \meng{x_0}} \cup \menge{x_n}{n \in \N}\cup\meng{x_0}\\
&\supseteq D \setminus \meng{x_0} \cup \menge{x_n}{n \in \N}\cup\meng{x_0}
= D \cup\menge{x_n}{n \in \N}\text.
\end{align*}
Da $\overline {D'}$ abgeschlossen ist, enth�lt es auch den Abschluss von
$D \cup\menge{x_n}{n \in \N}$. Somit gilt wegen der Dichtheit von $D$ in M:
\[\overline{D'} \supseteq \overline{D \cup\menge{x_n}{n \in \N}} \supseteq
\overline D = M\text.\]
Die umgekehrte Inklusion ist trivial.
\end{proof}

Ebenso wie der punktweise Limes gleichmäßig lipschitz--stetiger Funktionen
bereits selbst lipschitz--stetig ist, ist es auch bei dem Supremum der Fall.

\begin{lemma}{Gleichmäßige Lipschitzstetigkeit}
             {GleichmaessigeLipschitzstetigkeit}
Seien $(M,d)$ ein metrischer Raum und $I$ eine nichtleere Menge.\\
Sei $(f_i)_{i \in I} \in \mcf(M, \R)^I$ mit den Eigenschaften:
\begin{enumerate}[(a)]
\item $\forall\, m \in M: \menge{f_i(m)}{i \in I}$ ist nach oben beschr�nkt,
\item $\exists\, c \in \R_{>0}\,\forall\, i \in I \, \forall\, x, y \in M:
\abs{f_i(x)-f_i(y)}\le c\cdot d(x, y)$.
\end{enumerate}
Dann ist die Abbildung
\[f: M \rightarrow \R, m \mapsto \sup\menge{f_i(m)}{i \in I}\]
wohldefiniert und lipschitzstetig mit Lipschitzkonstante $c$.
\end{lemma}
\begin{proof}
Die Wohldefiniertheit ist offensichtlich, denn wegen $I \neq \emptyset$ ist
für jedes $m \in M$ die Menge $\menge{f_i(m)}{i \in I}$ nichtleer und wegen
$(a)$ nach oben beschränkt, besitzt also ein Supremum (und auch genau eins).\\
Seien nun $x, y \in M$. O. B. d. A.: $f(x) \ge f(y)$.\\
Dann gilt:
\begin{align*}
\abs{f(x)-f(y)} &= f(x) - f(y) = \left(\sup_{i \in I}f_i(x)\right) -
 \left(\sup_{i \in I}f_i(y)\right)\\
 & = \left(\sup_{i \in I}(f_i(x)-f_i(y) + f_i(y))\right) -
 \left(\sup_{i \in I}f_i(y)\right)\\
 &\le \left(\sup_{i \in I}(f_i(x) -f_i(y)) \right)
 + \left(\sup_{i\in I}f_i(y)\right) - \left(\sup_{i\in I}f_i(y)\right)\\
 & = \sup_{i \in I}\underbrace{(f_i(x) -f_i(y))}_{\le c\cdot d(x, y)}
 \le c \cdot d(x, y)\text.
\end{align*}
Also ist $f$ lipschitzstetig mit Lipschitzkonstante $c$.
\end{proof}

\begin{lemma}{Gerade und ungerade Funktionen}{GeradeUndUngeradeFunktionen}
Seien $a \in\!\ ]0,\infty]$, $f:\!\  ]-a,a[ \rightarrow \R$ ungerade und
$g:\!\ ]-a,a[ \rightarrow \R$ gerade.\\
Dann gilt:
\begin{enumerate}[(1)]
	\item \label{GeradeUndUngeradeFunktionen1}
	Ist $f$ stetig in $0$, so ist $f(0) = 0$.
	\item \label{GeradeUndUngeradeFunktionen2}
	Ist $f$ differenzierbar, so ist $f'$ gerade.
	\item \label{GeradeUndUngeradeFunktionen3}
	Ist $g$ differenzierbar, so ist $g'$ ungerade.
	\item \label{GeradeUndUngeradeFunktionen4}
	Ist $f$ injektiv und stetig, so gibt es $b \in \!\ ]0,\infty]$ mit 
	$\Bild f=\!\  ]-b,b[$
	und $f\inv: \Bild f \rightarrow \R$ ist ungerade.
	\item \label{GeradeUndUngeradeFunktionen5}
	Ist $f \in \mcc^2(]-a,a[, \R)$, $f'(0) \ge 0$ und gibt es $c \in \R_{>0}$ mit
	$f'' = c\cdot f\cdot(f')^2$, so ist bereits $f \in \mcc^\infty(]-a,a[,\R)$
	und es gilt \[\forall\, n \in \N: f^{(2n+1)}(0) \ge 0\] 
\end{enumerate}
\end{lemma}
\begin{proof}
(1) Sei also $f$ stetig in $0$.\\
Dann ist $\frac a {2n} \in\!\  ]-a,a[$ für alle $n \in \N$ und es gilt
\begin{align*}f(0) &\overset{f \text{ stet.}}
= \Limes n \infty f\left(\frac a {2n}\right) 
= \Limes n \infty f\left(-\left(-\frac a {2n}\right)\right)
\overset{f \text{ ung.}}= \Limes n \infty - f\left(-\frac a {2n}\right)\\ 
&= - \Limes n \infty f\left(-\frac a {2n}\right)
\overset{f \text{ stet.}}= - f(0)\text,
\end{align*}
also $f(0)=0$.\\
(2)/(3) Seien $f,g$ differenzierbar und $\mu:\R\to\R, x\mapsto-x$.\\
Es gilt dann für alle $x \in \!\ ]-a,a[$:
\begin{align}
\label{guufeq:1}f(x) &= f(-(-x)) = - f(-x) = - (f\circ \mu)(x)\text, \\
\label{guufeq:2}g(x) &= g(-(-x)) = g(-x) = (g \circ \mu)(x)\text. 
\end{align}
Also ist für alle $x \in \!\ ]-a,a[$:
\[f'(-x) \overset{\eqref{guufeq:1}}= 
(-f \circ \mu)'(-x) = - (-1)\cdot f'(\mu(-x))= f'(x)\]
und analog
\[g'(-x) \overset{\eqref{guufeq:2}}= 
(g \circ \mu)'(-x) = - g'(\mu(-x)) = -g'(x)\text.\]
Also ist $f'$ gerade und $g'$ ungerade.\\
(4) Sei $f$ injektiv und stetig. Nach dem Zwischenwertsatz ist $\Bild f$
ein Intervall und es gilt:
\begin{align*}
-\Bild f &= \menge{-f(x)}{x \in \!\ ]-a,a[}=\menge{f(-x)}{x \in \!\ ]-a,a[}\\
&= \menge{f(x)}{x \in \!\ ]-a,a[}= \Bild f\text.\end{align*}
Somit ist $\Bild f$ symmetrisch, woraus die Behauptung über das Bild folgt.\\
Sei nun $x \in \Bild f$. Dann gilt
\[f\inv(-x) = f\inv(-f(f\inv(x))) \overset{f \text{ ung.}} = 
f\inv(f(-f\inv(x))) = -f\inv(x)\text,\]
also $f\inv$ ungerade.\\
(5) Seien $f$ zweimal stetig differenzierbar und $c \in \R_{>0}$ mit
$f'' = cf(f')^2$.\\
Da $cf(f')^2 = cff'f'$ als Produkt differenzierbarer Funktionen differenzierbar
ist, ist $f$ bereits dreimal stetig differenzierbar.\\
Ist nun $f$ schon $2+n$--mal stetig differenzierbar für ein $n \in\N$, so ist
$f'$ noch $1+n$--mal stetig differenzierbar und damit $f''= cff'f'$ als Produkt
$1+n$--mal stetig differenzierbarer Funktionen ebenfalls noch $1+n$--mal stetig
differenzierbar, also $f$ bereits $3+n$--mal stetig differenzierbar.\\
Damit ist $f$ unendlich oft stetig differenzierbar.\\
Nach Voraussetzung ist $f'(0) \ge 0$.\\
Sei nun $N \in \N$ und es gelte $\forall\, m \in \haken N: f^{(2m+1)}\ge 0$.\\
Dann gilt
\begin{align}
&\,f^{(2(N+1)+1)}(0) = f^{(2N+3)}(0) = 
\left(\left(\frac d {dx} \right)^{2N+1}f'' \right)(0)\notag\\
=&\, \left(\left(\frac d {dx} \right)^{2N+1} c f (f')^2\right) (0)\notag\\
\overset{\text{Leibniz}}=&\, c\sum_{k=0}^{2N+1}\binom{2N+1}k f^{(2N+1-k)}(0) 
\left(\left(\frac d {dx} \right)^{k}f'f' \right)(0)\notag\\
\overset{\text{Leibniz}}=&\,c\sum_{k=0}^{2N+1}\binom{2N+1}k f^{(2N+1-k)}(0)
\sum_{j=0}^k\binom k j f^{(k+1-j)}(0)f^{(j+1)}(0) \label{guufeq:3}
\end{align}
Sei nun $i \in \haken {2N+2}_0$.\\
Falls $i$ gerade ist, so ist $f^{(i)}$ nach $(2),(3)$ ungerade, also nach
$(1)$: $f^{(i)}(0) = 0$.\\
Falls $i$ ungerade ist, so ist $i \in \haken{2N+1}$ und damit nach
Induktionsvoraussetzung $f^{(i)}(0) \ge 0$.\\
Gemäß $\eqref{guufeq:3}$ ist dann $f^{2N+3}(0)$ als Summe von Produkten
nichtnegativer Zahlen selbst nichtnegativ.
\end{proof}
\parag
Offensichtlich kann man jede endliche Menge reeller Zahlen der Größe nach
anordnen. Wir sind daran interessiert, diese Anordnung mengentheoretisch
deskriptiv zu halten, um die Stetigkeit dieser Anordnung einfach einsehen zu
können.

\begin{lemma}{Sortierungsmengen}{Sortierungsmengen}
Seien $n \in \N$ und $x \in \R^n$.\\
Für jedes $r \in \haken n$ sei
$ M_{r,x}:=\menge{x_i}{i\in\haken n,\abs{\menge{j\in\haken n}{x_j>x_i}}<r}$.\\
Dann gilt:
\begin{enumerate}[(1)]
	\item \label{Sortierungsmengen1}
	$\forall\, k \in \haken {n-1}: M_{k, x} \subseteq M_{k+1,x}$.
	\item \label{Sortierungsmengen2}
	$\forall\, k \in \haken n: M_{k, x} \neq \emptyset$.
\end{enumerate}
\end{lemma}
\begin{proof}
(1) Sei $k \in \haken n$ und $y \in M_{k, x}$.\\
Dann gibt es ein $i \in \haken n$ mit $y = x_i$ und
$\abs{\menge{j\in\haken n}{x_j>x_i}}<k < k+1$.\\
Also ist $y=x_i \in M_{k+1,x}$.\\
(2) Wegen (1) genügt es zu zeigen, dass $M_{1, x} \neq \emptyset$ gilt.\\
Es gilt:
\begin{align*}
M_{1,x} &= \menge{x_i}{i\in\haken n,\abs{\menge{j\in\haken n}{x_j>x_i}}<1}
= \menge{x_i}{i\in\haken n, \menge{j\in\haken n}{x_j>x_i} = \emptyset}\\
&= \menge{x_i}{i\in\haken n, \forall\, j \in \haken n: x_i \ge x_j}
= \menge{x_i}{i \in \haken n, x_i = \max\menge{x_j}{j\in \haken n}}\text.
\end{align*}
Da $\menge{x_j}{j \in \haken n}$ endlich und nichtleer ist, gibt es
$l \in \haken n$ mit $x_l = \max\menge{x_j}{j \in \haken n}$ und damit ist
$M_{1, x}$ wegen obiger Umformung nichtleer.
\end{proof}
\parag
Unter Zuhilfenahme der Sortierungsmengen kann nun das $k$--te Maximum definiert
werden, welches sich als das Minimum der jeweiligen Sortierungsmenge ergibt.

\begin{definition}{k--tes Maximum}{ktesMaximum}
Seien $n \in \N$ und $k \in \haken n$.\\
Dann definieren wir:
\[\mu_{k, n}:\R^n \rightarrow \R, 
x\mapsto \min\menge{x_i}{i\in\haken n,\abs{\menge{j\in\haken n}{x_j>x_i}}<k}
\text.\]
Nach \refclaim{Sortierungsmengen}{Sortierungsmengen2} ist $\mu_{k,n}$
wohldefiniert und wir nennen es \textbf{das $k$--te Maximum}.
\end{definition}
\parag
Einige einfache Eigenschaften des $k$--ten Maximums halten wir in dem folgenden
Lemma fest.

\begin{lemma}{Eigenschaften des k--ten Maximums}
             {EigenschaftenDesKtenMaximums}
Sei $n \in \N$.\\
Dann gilt:
\begin{enumerate}[(1)]
	\item \label{EigenschaftenDesKtenMaximums1}
	$\forall\, k \in \haken {n-1}: \mu_{k+1,n} \le \mu_{k,n}$.
	\item \label{EigenschaftenDesKtenMaximums2}
	$\forall\, k \in \haken n\, \forall\, x \in \R^n\, \forall\,a \in \R:
	\mu_{k,n}(x) \ge a \iff \abs{\menge{i \in \haken n}{x_i \ge a}} \ge k$.
	\item\label{EigenschaftenDesKtenMaximums2a}
	$\forall\, k \in \haken n\, \forall\, x \in \R^n\, \forall\,a \in \R:
	\mu_{k,n}(x) < a \iff \abs{\menge{i \in \haken n}{x_i < a}} > n - k$.
	\item\label{EigenschaftenDesKtenMaximums2b}
	$\forall\, k \in \haken n\, \forall\, x \in \R^n\, \forall\,a \in \R:
	\mu_{k,n}(x) > a \iff \abs{\menge{i \in \haken n}{x_i > a}} \ge k$.
	\item\label{EigenschaftenDesKtenMaximums2c}
	$\forall\, k \in \haken n: \mu_{k, n}$ stetig.
	\item \label{EigenschaftenDesKtenMaximums3}
	$\forall\, k \in \haken n: \mu_{k, n}$ nichtnegativ--homogen.
\end{enumerate}
\end{lemma}
\begin{proof}
(1) Seien $k \in \haken {n-1}$ und $x \in \R^n$.\\
Seien
\[M_{k, x}\coloneqq \menge{x_i}{i\in\haken n,\abs{\menge{j\in\haken n}{x_j>x_i}}<k}\]
und
\[M_{k+1,x}\coloneqq
\menge{x_i}{i\in\haken n,\abs{\menge{j\in\haken n}{x_j>x_i}}<k+1}.\]
Nach~\refclaim{Sortierungsmengen}{Sortierungsmengen1} ist dann
$M_{k,x} \subseteq M_{k+1,x}$ und daher:
\[\mu_{k+1,n}(x) = \min M_{k+1,x} \le \min M_{k,x} = \mu_{k,n}(x).\]
(2) Seien $k \in \haken n$, $x \in \R^n$ und $a \in \R$.\\
\enquote{$\Rightarrow$}: Es gelte $\mu_{k,n}(x) \ge a$.\\
Nach (1) gilt: $\forall\,j\in\haken k: \mu_{j, n}(x) \ge \mu_{k, n}(x)\ge a$.\\
Es gibt $\iota: \haken k \rightarrow \haken n$ injektiv mit
$\forall\, j \in \haken k: x_{\iota(j)}= \mu_{j,n}(x)$.\\
Ferner gilt $\iota(\haken k) \subseteq \menge{i \in \haken n}{x_i \ge a}$ und
damit wegen der Injektivität von $\iota$
\[k = \abs{\iota(\haken k)} \le \abs{\menge{i \in \haken n}{x_i \ge a} }.\]
\enquote{$\Leftarrow$:} Es gelte $\mu_{k,n}(x) < a$.\\
Es gibt ein $j \in \haken n$ mit $x_j = \mu_{k,n}(x) < a$.\\
Nach Definition von $\mu_{k,n}$ gilt weiter: 
$\abs{\menge{i \in \haken n}{x_i> x_j}}<k$.\\
Somit gilt 
$\menge{i \in \haken n}{x_i\ge a} \subseteq \menge{i \in \haken n}{x_i > x_j}$
und daher
\[\abs{\menge{i\in\haken n}{x_i\ge a}}
\le\abs{\menge{i\in\haken n}{x_i>x_j}}<k.\]
(3) Seien $k \in \haken n$, $x \in \R^n$ und $a \in \R$. Dann gilt:
\begin{align*}
\mu_{k, n}(x) < a & \overset{(2)}\iff \abs{\menge{i \in \haken n}{x_i \ge a}}<k
\text.
\end{align*}
Offensichtlich gilt auch:
\[\abs{\menge{i \in \haken n}{x_i \ge a}}<k \iff
  \abs{\menge{i \in \haken n}{x_i < a}} > n-k\text.\]
Damit ist (3) gezeigt.\\
(4) Analog zu (2).\\
(5) Es genügt zu zeigen, dass Urbilder offener Intervalle offen sind.\\
Seien also $a, b \in \R$ mit $a < b$ und $k \in \haken n$. Dann gilt:
\begin{align*}
&\menge{x \in \R^n}{\mu_{k, n}(x)> a}
\overset{(4)}= \menge{x \in \R^n}{\exists A \subseteq \haken n: \abs A
\ge k \wedge (\forall\, i \in A: x_i > a)}\\
=\,&\bigcup_{\substack{A \subseteq \haken n,\\ \abs A \ge k}}
\bigcap_{i \in \haken A}\menge{x \in \R^n}{x_i> a}
= \bigcup_{\substack{A \subseteq \haken n,\\ \abs A \ge k}}
\bigcap_{i \in \haken A} \pr_i\inv(]a, \infty[) \eqqcolon G\text.
\end{align*}
Da $\pr_i$ für alle $i \in \haken n$ unabhängig von der Norm stetig sind, ist
$G$ als Vereinigung endlicher Schnitte offener Mengen offen.\\
Analog gilt:
\begin{align*}
&\menge{x \in \R^n}{\mu_{k, n}(x)< b}
\overset{(3)}= \menge{x \in \R^n}{\exists B \subseteq \haken n: \abs B
>n- k \wedge (\forall i \in B: x_i < b)}\\
=&\bigcup_{\substack{B \subseteq \haken n,\\ \abs B > n - k}}
\bigcap_{i \in \haken B}\menge{x \in \R^n}{x_i< b}
= \bigcup_{\substack{B \subseteq \haken n,\\ \abs B > n - k}}
\bigcap_{i \in \haken B} \pr_i\inv(]-\infty, b[) \eqqcolon K\text.
\end{align*}
Auch $K$ ist (mit dem gleichen Argument wie $G$) offen.
Außerdem ist
\[\mu_{k,n}\inv(]a,b[) = \menge{x \in \R^n}{\mu_{k,n}(x) > a \wedge
\mu_{k,n}(x) < b} = G \cap K\]
und damit offen, weil $G$ und $K$ es sind.\\
(6) Seien $k \in \haken n$, $x \in \R^n$ und $t \in \R_{>0}$.\\
Dann gilt:
\begin{align*}
&\menge{(tx)_i}{i\in\haken n,\abs{\menge{j\in\haken n}{(tx)_j>(tx)_i}}<k}\\
=&\,\menge{tx_i}{i\in\haken n,\abs{\menge{j\in\haken n}{tx_j>tx_i}}<k}\\
\overset{t>0}= &\,
\menge{tx_i}{i\in\haken n,\abs{\menge{j\in\haken n}{x_j>x_i}}<k}\\
=&\, t\menge{x_i}{i\in\haken n,\abs{\menge{j\in\haken n}{x_j>x_i}}<k}\text.
\end{align*}
Also nach Definition:
\begin{align*}
\mu_{k, n}(tx) &=
\min\menge{(tx)_i}{i\in\haken n,\abs{\menge{j\in\haken n}{(tx)_j>(tx)_i}}<k}\\
&=\min
\left(t\menge{x_i}{i\in\haken n,\abs{\menge{j\in\haken n}{x_j>x_i}}<k}\right)
\overset{t>0}=t\mu_{k, n}(x)\text.
\end{align*}
Zudem gilt offenbar $\mu_{k,n}(0) = 0$, sodass $\mu_{k,n}$
nichtnegativ--homogen ist.
\end{proof}
\parag
Wir definieren nun eine spezielle Norm auf dem $\R^n$, welche an anderer Stelle
von Interesse sein wird.

\begin{lemma}{Sortierungsnorm}{Sortierungsnorm}
Seien $n \in \N$, $m \in \haken n$, 
\[\norm\cdot_{(m)}: \R^n \rightarrow [0, \infty[,
 v\mapsto\frac 1m\sum_{i=1}^m \mu_{i, n}\left(\pfeil {\abs\cdot}n(v) \right)\]
und
\[V\coloneqq  \mengea{z \in \R^n\setminus\meng 0}
    {\left(\forall\, i \in \haken n: \abs{z_i} \in\meng{0,\nicefrac 1m} \right)
    \wedge \left(\abs{\menge{i \in \haken n}{\abs{z_i} \neq 0}}\le m \right)}
    \text.\]
Dann ist $V$ endlich, nichtleer und es gilt:
\begin{enumerate}[(1)]
	\item \label{Sortierungsnorm1}
	$\sup \menge{\norm v_{2, n}}{v \in V} = \frac 1 {\sqrt m}$.
	\item \label{Sortierungsnorm2}
	$\forall\, w \in \R^n: \sup\menge{\sprod w v_n}{v \in V} = \norm w _{(m)}$.
	\item \label{Sortierungsnorm3}
	$\norm\cdot_{(m)}$ ist eine Norm auf dem $\R^n$ mit $\norm\cdot_{(m)} \le
	\norm\cdot_{\infty, n}$.
\end{enumerate}
\end{lemma}
\begin{proof}
Wegen $\frac 1 m e_1^n \in V \subseteq \meng{-\frac 1 m, 0, \frac 1 m}^n$ ist
$V$ nichtleer und endlich.\\
(1) Sei $v \in V$. Dann gilt:
\begin{align*}
\norm v_{2, n}^2 &= \sum_{i \in \haken n}\abs{v_i}^2
= \sum_{\substack{i\in \haken n\\ \abs{v_i} \neq 0}} \abs{v_i}^2
= \frac 1 {m^2} \sum_{\substack{i\in \haken n\\ \abs{v_i} \neq 0}} 1
\le \frac 1 {m^2} m = \frac 1 m\text.
\end{align*}
Also ist $\norm v _{2, n} \le\frac 1 {\sqrt m}$ und es gilt
\[v_0 \coloneqq  \sum_{i=1}^m \frac 1 m e_i^n \in V \quad \text{ und } \quad
\norm{v_0}_{2, n} = \frac 1 {\sqrt m}\text.\]
Somit gilt $\sup\menge{\norm w_{2, n}}{w \in V} = \frac 1 {\sqrt m}$.\\
(2) Sei $w \in \R^n$. Falls $w = 0$, so ist die Aussage trivial.\\
Sei also $w \neq 0$. Dann gilt für alle $v\in V$:
\begin{align*}\sprod w v_n& \le \sum_{i=1}^n \abs{w_i}\abs{v_i}
= \sum_{\substack{i \in \haken n\\ \abs{v_i} \neq 0} }\abs{w_i}\abs{v_i}
= \sum_{\substack{i \in \haken n\\ \abs{v_i} \neq 0} }\abs{w_i}\frac 1 m
\le \frac 1 m \sum_{i=1}^m \mu_{i, n}\left(\pfeil {\abs\cdot}n (w) \right)\\
&= \norm w_{(m)}\text.
\end{align*}
Ferner gibt es $\iota: \haken m \rightarrow \haken n$ injektiv mit
$\forall\, i \in \haken m: \abs{w_{\iota(i)}}
= \mu_{i, n}\left(\pfeil {\abs\cdot}n (w) \right)$.\\
Sei
\[v_w : \haken n \rightarrow \R,\, j \mapsto \begin{cases}
\frac 1 m \sgn(w_{j})&: j \in \Bild(\iota)\\
0 & : \text{ sonst}
\end{cases}.\]
Dann ist $v_w \in V$ und es gilt:
\begin{align*}
\sprod w {v_w}_n &= \sum_{i \in \haken n}w_i (v_w)_i
= \sum_{\substack{i \in \haken n\\ (v_w)_i \neq 0}}w_i (v_w)_i
= \sum_{i \in \Bild(\iota)}w_i \frac 1 m \sgn(w_i)\\
&= \frac 1 m\sum_{i \in \Bild(\iota)} \abs{w_i}
= \frac 1 m \sum_{i \in \haken m} \abs{w_{\iota(i)}}
= \frac 1 m \sum_{i=1}^m\mu_{i, n}\left(\pfeil {\abs\cdot}n (w) \right)
= \norm w _{(m)}\text.
\end{align*}
Somit ist auch (2) (analog zu (1)) gezeigt.\\
(3) Offenbar ist $\norm\cdot_{(m)}$ nichtnegativ--definit.\\
Seien $x, y \in \R^n$ und $t \in \R$. Dann gilt wegen $\abs t \ge 0$:
\begin{align*}
\norm{tx}_{(m)}&= 
\frac 1 m \sum_{i=1}^n\mu_{i, n}\left(\pfeil{\abs\cdot}n(tx) \right)
= \frac 1 m \sum_{i=1}^n\mu_{i, n}\left(\abs{t} \pfeil{\abs\cdot}n(x) \right)\\
&\overset{\refclaim{EigenschaftenDesKtenMaximums}
{EigenschaftenDesKtenMaximums3}}=
\frac 1 m\sum_{i=1}^n\abs{t}\mu_{i, n}\left(\pfeil{\abs\cdot}n(x) \right)
= \abs{t}\frac 1 m\sum_{i=1}^n\mu_{i, n}\left(\pfeil{\abs\cdot}n(x) \right)
= \abs{t}\norm{x}_{(m)}\text.
\end{align*}
Schließlich gilt:
\begin{align*}
\norm{x+y}_{(m)} &\overset{(2)}=
\sup_{v \in V} \sprod {x+y}v_n 
= \sup_{v \in V}\left(\sprod x v_n + \sprod y v_n \right)\\
&\le \left(\sup_{v \in V}\sprod x v_n \right) 
+ \left(\sup_{v \in V}\sprod y v_n \right)
\overset{(2)}= \norm x_{(m)} + \norm y _{(m)}\text.
\end{align*}
Also ist $\norm\cdot_{(m)}$ in der Tat eine Norm und es gilt
\[\norm\cdot _{(m)}
 = \frac 1 m \sum_{i=1}^m \mu_{i, n} \circ \pfeil{\abs\cdot}n
 \overset{\refclaim{EigenschaftenDesKtenMaximums}
 {EigenschaftenDesKtenMaximums1}}\le
 \frac 1 m \sum_{i=1}^m \mu_{1, n} \circ \pfeil{\abs\cdot}n
 = \mu_{1, n} \circ \pfeil{\abs\cdot}n = \norm\cdot_{\infty, n}\text.\]
\end{proof}
\parag
Zum Schluss dieses Kapitels beschäftigen wir uns kurz mit einigen speziellen
Funktionen und gewissen Berechnungen.

\begin{lemma}{Elementare Eigenschaften der Gamma-- und Betafunktion}
             {ElementareEigenschaftenDerGammafunktion}
Seien $\Gamma:\R_{>0}\rightarrow\R,\,
t\mapsto\int_0^\infty x^{t-1}\exp(-x)dx$ und 
\[B: (\R_{>0})^2 \rightarrow \R,
(x,y)\mapsto\int_0^1 (1-s)^{x-1}s^{y-1}ds\text.\]
Dann gilt:
\begin{enumerate}[(1)]
	\item \label{ElementareEigenschaftenDerGammafunktion1}
	$\forall\, t \in \R_{>0}: \Gamma(t) > 0$.
	\item \label{ElementareEigenschaftenDerGammafunktion2}
	$\forall\, t \in \R_{>0}: \Gamma(t+1) = t \Gamma(t)$.
	\item\label{ElementareEigenschaftenDerGammafunktion3}
	$\forall\, n \in \N: \Gamma(n) = (n-1)!$
	\item\label{ElementareEigenschaftenDerGammafunktion4}
	$\forall\,x,y\in\R_{>0}: B(x,y) = \frac{\Gamma(x)\Gamma(y)}{\Gamma(x+y)}$.
\end{enumerate}
\end{lemma}
\begin{proof}
(1) Sei $t \in \R_{>0}$.\\
Sei $f_t: \R_{>0} \rightarrow \R, x \mapsto x^{t-1}\exp(-x)$.\\
Dann gilt f�r alle
$x \in \R_{>0}: f_t(x) = \exp((t-1)\log(x))\exp(-x) > 0$, also $f_t > 0$.\\
Da $f_t$ offensichtlich stetig ist, ist auch
$\Gamma(t)=\int_0^\infty f_t(y)dy > 0$.\\
(2)/(3) siehe \cite{K1} dort 18.1.(5) und 18.1.(4).\\
(4) siehe \cite{K2}, dort 8.4.(8).
\end{proof}

\begin{lemma}{Summen von Logarithmen}{SummenVonLogarithmen}
Seien $p \in \!\ ]0,1]$, $k, n \in \N$ mit $k \le \frac n {\exp(1)}$.\\
Dann gilt
\[k\left(\log\left(\frac n k \right) \right)^p
\le \sum_{j=1}^k \left(\log\left(\frac n j \right) \right)^p
\le 2^p \cdot k \cdot \left(\log\left(\frac n k \right) \right)^p\text.\]
\end{lemma}
\begin{proof}
Die erste Ungleichung ist trivial.\\
Wegen $\exp(k)\ge\frac{k^k}{k!}$ ist
\begin{align}
k!\ge\left(\frac k{\exp(1)}\right)^k\text.\label{svleq:1}
\end{align}
Auflerdem ist mit $\frac n k \ge \exp(1)$ auch
\begin{align}
\log\left(\frac n k \right) \ge 1\label{svleq:2}
\end{align}
Somit gilt:
\begin{align}
\sum_{j=1}^k \log\left(\nicefrac n j \right) &=
\log\left(\frac{n^k}{k!} \right) \overset{\eqref{svleq:1}}\le
\log\left(\frac{n^k}{\left(\nicefrac{k}{\exp(1)} \right)^k} \right)
= k \log\left(\frac n k \exp(1) \right)\notag\\
&= k\left(\log\left(\nicefrac n k\right) + 1 \right)
\overset{\eqref{svleq:2}}\le 2k\log\left(\nicefrac n k \right)\text.
\label{svleq:3}
\end{align}
Da $p \in \!\ ]0, 1]$, ist $\frac 1 p \in [1, \infty[$ und damit ist
$\norm\cdot_{\frac 1 p, k}$ eine Norm auf dem $\R^k$.\\
Daher gilt weiter
\begin{align*}
&\sum_{j=1}^k \left(\log\left(\nicefrac n j \right) \right)^p
=\norma{\left(\left(\log\left(\nicefrac nj\right)\right)^p\right)_
{j\in \haken k}}_{1,k}\\
\overset{\refclaim{AequivalenzkonstantenFuerPNormen}
{AequivalenzkonstantenFuerPNormen1}}\le &\,
k^{\frac{\frac 1 p -1}{\frac 1 p}}
\norma{\left(\left(\log\left(\nicefrac nj\right)\right)^p\right)
_{j\in\haken k}}_{\frac 1p,k}
= k^{1-p}\left(\sum_{j=1}^k \left( \left(
\log\left(\nicefrac nj\right)\right)^p\right)^{\frac 1 p} \right)^p\\
=&\, k^{1-p}\left(\sum_{j=1}^k \log\left(\nicefrac n j \right) \right)^p
\overset{\eqref{svleq:3}}\le
k^{1-p}\cdot 2^p \cdot k^p \cdot 
\left(\log\left(\nicefrac nk \right) \right)^p\\
=&\, 2^p k \left(\log\left(\nicefrac n k \right) \right)^p\text.
\end{align*}
\end{proof}
\parag
In das folgende Lemma lagern wir eine Reihe nur bedingt zusammengehöriger
Rechnungen aus, welche an späteren Stellen benötigt werden. Die Aussagen (2)
und (3) des Lemmas sind zu unhandlich, um in einem Beweis untergebracht zu
werden, der sie benötigt. Die Aussagen (1) und (4) hingegen können auch
unabhängig von der Situation von Interesse sein, denn die erste gibt eine
Reihenentwicklung des Logarithmus' mit positiven Koeffizienten an, während die
vierte konkret lineares und logarithmisches Wachstum vergleicht.

\begin{lemma}{Numeriklemma}{Numeriklemma}
Seien 
\[h_1: [1, \infty[ \rightarrow \R, 
t\mapsto\frac{\sqrt t}{t+1}
\bigg((\pi-1)\sqrt{\log(t)}+\sqrt{\log(t)+2\pi}\bigg)\text,\]
\[h_2: \R_{>0} \rightarrow \R,
t \mapsto \frac{\sqrt{t+1}}t
\bigg( (\pi-1)\sqrt{\log(1+t)} + \sqrt{\log(1+t) + 2\pi} \bigg)\text.\]
Dann gilt:
\begin{enumerate}[(1)]
	\item \label{Numeriklemma0}
	$\forall\, t \in \R_{>0}: \log(t) = 2\Sum{k=0}\infty \frac 1 {2k+1}
	\left(\frac{t-1}{t+1} \right)^{2k+1}$.
	\item \label{Numeriklemma1}
	$h_1 < \sqrt{2 \pi}$.
	\item \label{Numeriklemma2}
	$h_2$ ist monoton fallend und für alle $t \in [\frac{13}2, \infty[:
	h_2(t) < \sqrt{2 \pi}$.
	\item\label{Numeriklemma4}
	$\forall\, \eta\in\!\  ]0,1[ 
	 \forall \,t\in 
	 \bigg[\frac 4\eta \log\left(\frac 2 \eta\right), \infty\bigg[
	 : \log(2t)< \eta t$.
\end{enumerate}
\end{lemma}
\begin{proof}
(1) Sei $t \in \R_{>0}$.\\
Die Taylorentwicklung des Logarithmus' liefert
\begin{align}\forall\, x \in\!\ ]-1, 1[: \log(1+x) = \sum_{k=0}^\infty (-1)^k 
\frac {x^{k+1}}{k+1}\text. \label{nleq:1}\end{align}
Sei $x_t \coloneqq  \frac{t-1}{t+1}$. Dann gilt:
\[ \frac{1+x_t}{1-x_t} = \frac{1+\frac{t-1}{t+1}}{1-\frac{t-1}{t+1}}
= \frac{\frac{t+1+t-1} {t+1}}{\frac{t+1-t+1}{t+1}}
= \frac{\frac{2t}{t+1}}{\frac 2 {t+1}}=t \text.\]
Ferner gilt: $-t < t$, sodass $-(t+1)=-t -1< t-1$, also $x_t > -1$,
sowie $t-1 < t+1$, sodass $x_t < 1$. Also $x_t, -x_t \in\!\  ]-1,1[$.\\
Erhalte:
\begin{align*}
\log(t) &= \log\left(\frac{1+x_t}{1-x_t} \right) = \log(1+x_t)-\log(1-x_t)\\
&\overset{\eqref{nleq:1}}=
\left(\sum_{k=0}^\infty (-1)^k\frac{x_t^{k+1}}{k+1} \right)
-\left(\sum_{k=0}^\infty (-1)^k(-1)^{k+1}\frac{x_t^{k+1}}{k+1} \right)\\
&=\sum_{k=0}^\infty\left((-1)^k-(-1)^k(-1)^{k+1} \right)\frac{x_t^{k+1}}{k+1}\\
&=2\sum_{\substack{k=0,\\ k \text{ gerade}}}^\infty\frac{x_t^{k+1}}{k+1}
= 2\sum_{k=0}^\infty\frac{x_t^{2k+1}}{2k+1} 
= 2\sum_{k=0}^\infty\frac{1}{2k+1}\left(\frac{t-1}{t+1}\right)^{2k+1}\text.
\end{align*}
(2) Sei zunächst 
$\sigma:\!\  ]1, \infty[ \rightarrow \R, t \mapsto \frac{\sqrt t}{\log(t)}$.\\
Dann ist $\sigma$ offensichtlich stetig differenzierbar und es gilt für alle
$t \in\!\  ]1, \infty[$:
\[\sigma'(t)=\frac {\frac 1{2\sqrt t}\log(t)-\sqrt t\frac 1 t}{(\log(t))^2}
= \frac{\frac 1 {2\sqrt t}(\log(t) - 2)}{(\log(t))^2}\text.\]
Daher hat $\sigma'$ genau eine Nullstelle in $\exp(2)$, $\sigma$ ist auf 
$]1, \exp(2)[$ streng mototon
fallend, auf $]\exp(2), \infty[$ streng monoton wachsend und hat somit ein
Minimum in $\exp(2)$, sodass also
$\sigma(z)\ge\sigma(\exp(2))=\frac{\exp(1)}2$ für alle $z\in\!\  ]1,\infty[$
gilt.\\
Ferner gilt: $\exp(1)\ge 1 + 1 + \frac 1 2 + \frac 1 6 = \frac 8 3$ und somit
$\sigma(t) \ge \frac 1 2 \cdot\frac 8 3 = \frac 4 3$.\\
Somit gilt (die folgende Aussage ist für $x=1$ offensichtlich):
\begin{align}
\forall\,x\in [1,\infty[: \log(x) \le \frac 3 4 \sqrt x \text.
\label{nleq:2}
\end{align}
Sei weiter
$\alpha:\R_{>0} \rightarrow \R, s \mapsto \frac{s^{\frac 3 4}}{s+1}$.\\
Auch $\alpha$ ist stetig differenzierbar und es gilt für alle $s\in\R_{>0}$
\[\alpha'(s) = \frac{\frac 3 4 s^{-\frac 1 4}(s+1) - s^{\frac 3 4}}{(s+1)^2}
= \frac{\frac 34(s+1)-s}{s^{\frac 14}(s+1)^2}
=\frac{3-s}{4s^{\frac 1 4}(s+1)^2}\text.\]
Daher hat $\alpha'$ genau eine Nullstelle (in 3), $\alpha$ ist auf $]0, 3[$
streng monoton wachsend, auf $]3, \infty[$
streng monoton fallend.\\
Also gilt:
\begin{align}
 \alpha \text{ hat ein globales Maximum in }3\text.\label{nleq:3}
\end{align}
Ferner ist wegen
$3^5=243< \frac{24343800625}{100000000}=\left(\frac{395}{100}\right)^4$
auch 
\begin{align}3^{\frac 5 4} < \frac {395}{100}
\label{nleq:4}
\end{align}
Sei schließlich $\beta: \R_{>0} \rightarrow \R,
r \mapsto \frac 1 {r+1}\sqrt{\frac 3 4 r^2 + 2\pi r}$.\\
Dann ist $\beta$ stetig differenzierbar und es gilt für alle
$r \in \R_{>0}$
\begin{align*}
\beta'(r)&=\frac{\frac{\frac 64r+2\pi}{2\sqrt{\frac 34r^2+2\pi r}}(r+1)
 - \sqrt{\frac 3 4 r^2 + 2 \pi r}}
 {(r+1)^2}
 = \frac{\frac{\left(\frac 3 4 r + \pi\right)(r+1)- \frac 3 4 r^2 - 2\pi r}
              {\sqrt{\frac 3 4 r^2 + 2 \pi r}}}{(r+1)^2}\\
 &= \frac{\frac 3 4 r^2 + \pi r + \frac 3 4 r + \pi - \frac 3 4 r^2 - 2 \pi r}
 {(r+1)^2\sqrt{\frac 3 4 r^2 + 2 \pi r}}
= \frac{\left(\frac 34-\pi\right)r+\pi}{(r+1)^2\sqrt{\frac 34r^2+2\pi r}}\text.
\end{align*}
Sei $r_0 \coloneqq  \frac{\pi}{\pi -  \frac 3 4}$. Dann ist $\beta$ auf $]0,r_0[$
streng monoton wachsend, auf $]r_0, \infty[$ streng monoton fallend
und $\beta'$ hat genau eine Nullstelle (nämlich $r_0$).\\
Damit gilt:
\begin{align}
\beta \text{ hat ein globales Maximum in } r_0\text. \label{nleq:5}
\end{align}
Zuletzt gilt nach Archimedes:
\begin{align}\frac {223}{71} < \pi < \frac{22}7\text.
\label{nleq:6}
\end{align}
Damit ist auch wegen $6248\cdot 100 =624800<627396=4753\cdot 132$
\begin{align}\frac{6244}{4757} &=\frac{28\cdot 223}{67\cdot 71}
= \frac{\frac{223}{71}}{\frac {22} 7 - \frac 3 4}\notag\\ 
&< \frac{\pi}{\pi - \frac 3 4} < \frac{\frac {22}7}{\frac{223}{71}-\frac 3 4}
 = \frac{\frac {22}7}{\frac{679}{284}}= \frac{6248}{4753} < 
 \frac{132}{100}\text. \label{nleq:7} \end{align}
Sei nun $t \in [1, \infty[$. Dann gilt:
\begin{align*}
h_1(t)& = \frac{\sqrt t}{t+1}
\bigg((\pi-1)\sqrt{\log(t)}+\sqrt{\log(t)+2\pi}\bigg)\\
&\overset{\eqref{nleq:2}}\le
\frac{\sqrt t}{t+1}
\left((\pi-1)\sqrt{ \frac 3 4 \sqrt t }+\sqrt{\frac 3 4 \sqrt t +2\pi}\right)\\
&= (\pi-1)\sqrt{\frac 3 4}\frac {t^{\frac 3 4}}{t+1}
 + \frac 1 {t+1}\sqrt{\frac 3 4 t \sqrt t +  2\pi t}\\
&\overset{t \ge 1}\le(\pi-1)\sqrt{\frac 3 4}\frac {t^{\frac 3 4}}{t+1} + 
\frac 1 {t+1}\sqrt{\frac 3 4 t^2 +  2\pi t} 
= (\pi-1)\sqrt{\frac 3 4}\alpha(t) + \beta(t)\\
&\overset{\substack{\eqref{nleq:3},\\\eqref{nleq:5}}}\le 
(\pi-1)\sqrt{\frac 3 4}\alpha(3)
 + \beta\left(\frac{\pi}{\pi -  \frac 3 4} \right)\\
 &= (\pi-1)\frac 1 2 \cdot \frac{3^{\frac 5 4}}4 + 
 \frac 1 {\frac{\pi}{\pi -  \frac 3 4} + 1}
 \sqrt{\frac 3 4 \left(\frac{\pi} {\pi -  \frac 3 4} \right)^2 
 + 2\pi\frac{\pi}{\pi -  \frac 3 4}}\\
 &\overset{\substack{\eqref{nleq:6},\\\eqref{nleq:7}}}< 
 \frac{15}{56}3^{\frac 5 4} +
 \frac 1 {\frac{6244}{4757}+1}
 \sqrt{\frac 3 4 \cdot
  \frac{17424}{10000}+\frac{44}{7}\cdot\frac{132}{100}}\\
 &\overset{\eqref{nleq:4}}< \frac{5925}{5600} + \frac{4757}{11001}\sqrt{ 
  \frac{13068}{10000}+\frac{5808}{700}}
  < \frac{5925}{5600} + \frac{4757}{11001}
  \sqrt{\frac{13068}{10000}+\frac{58080}{7000}} \\
  &<\frac{237}{224} + \frac{4757}{11001}
  \sqrt{\frac{13068}{10000}+\frac{58086}{7000}}
  =\frac{237}{224} + \frac{4757}{11001}
  \sqrt{\frac{13068}{10000}+\frac{8298}{1000}}\\
  &=\frac{237}{224} + \frac{4757}{11001}
  \sqrt{\frac{96048}{10000}}
  < \frac{237}{224} + \frac{4757}{11001}
  \sqrt{\frac{96100}{10000}}
  = \frac{237}{224} + \frac{4757}{11001} \cdot \frac{31}{10}\\
  &= \frac{237}{224} + \frac{147467}{110010}
  < \frac{237}{224} + \frac{148}{110}
  = \frac{237}{224} + \frac{74}{55}
  =\frac{29611}{12320} < \frac{300}{123}\\
  &= \sqrt{\frac{90000}{15129}}\sqrt{\frac{71}{446}}\sqrt{\frac{446}{71}}
  = \sqrt{\frac{6390000}{6747534}}\sqrt{2\cdot \frac{223}{71}}
  < \sqrt{2\cdot \frac{223}{71}} \overset{\eqref{nleq:6}}< \sqrt{2\pi}\text.
\end{align*}
(3) Offenbar ist $h_2$ auf $\R_{>0}$ stetig differenzierbar und es gilt
für alle $t \in \R_{>0}$:
\begin{align}
&\,h_2'(t)\notag\\
 =&\, \frac{\frac t {2\sqrt{t+1}} - \sqrt{t+1}}{t^2}
\left((\pi-1)\sqrt{\log(1+t)} + \sqrt{\log(1+t)+2\pi} \right)\notag\\
&+ \frac{\sqrt{t+1}}t
\left(\frac{\pi-1}{2(t+1)\sqrt{\log(t+1)}} 
+ \frac 1 {2(t+1)\sqrt{\log(1+t) + 2\pi}} \right)\notag\\
=& \frac{-t-2}{2t^2\sqrt{t+1}}
  \left((\pi-1)\sqrt{\log(1+t)} + \sqrt{\log(1+t)+2\pi} \right)\notag\\
 & + \frac 1 {2t\sqrt{t+1}}\left(\frac{\pi-1}{\sqrt{\log(t+1)}} 
+ \frac 1 {\sqrt{\log(1+t) + 2\pi}}  \right)\notag\\
=&\,\scriptstyle{ \frac 1 {2t\sqrt{t+1}}
\left(- \frac {t+2}t\left((\pi-1)\sqrt{\log(1+t)} 
          + \sqrt{\log(1+t)+2\pi} \right)
          + \frac{\pi-1}{\sqrt{\log(t+1)}} 
+ \frac 1 {\sqrt{\log(1+t) + 2\pi}}  \right)}\text.
\label{nleq:8}
\end{align}
Nach (1) gilt für alle $x \in \R_{>0}: \log(x+1)\ge 2\frac{x}{x+2}$.\\
Sei $t \in \R_{>0}$. Dann gilt:
\[ \sqrt{\log(1+t)}\sqrt{\log(1+t)}=\log(1+t) \ge 2 \frac t{t+2} > 
\frac t {t+2}\text.\]
Also auch
\begin{align}
 \frac{1}{\sqrt{\log(1+t)}} < \frac{t+2}t\sqrt{\log(1+t)}\text.
 \label{nleq:9}
 \end{align}
Ferner gilt für alle $y \in [1, \infty[: -y + \frac 1 y \le 0$.\\
Es ist $y_t\coloneqq \sqrt{\log(1+t) + 2\pi}\ge 1$, also:
\begin{align} -\frac{t+2}t y_t + \frac 1 {y_t}
< - y_t + \frac 1 {y_t} < 0 \text.
\label{nleq:10}\end{align}
Somit gilt:
\begin{align*}
&2t\sqrt{t+1}\cdot h_2'(t)\\
& \overset{\eqref{nleq:8}}\le 
\scriptstyle{\left(- \frac {t+2}t\left((\pi-1)\sqrt{\log(1+t)} 
          + \sqrt{\log(1+t)+2\pi} \right)
          + \frac{\pi-1}{\sqrt{\log(t+1)}} 
+ \frac 1 {\sqrt{\log(1+t) + 2\pi}}  \right)}\\
& =(\pi-1)\underbrace{\left(\frac 1 {\sqrt{\log(1+t)}} 
- \frac{t+2}t\sqrt{\log(1+t)} \right)}_{\overset{\eqref{nleq:9}}<0}
+ \underbrace{\frac 1 {y_t} - \frac{t+2}ty_t}
             _{\overset{\eqref{nleq:10}}<0} < 0\text.
\end{align*}
Wegen $t > 0$ ist $h_2$ damit monoton fallend.\\
Es gilt:
\begin{align*}
 \exp\left(\frac{202}{100} \right) 
&\ge \sum_{k=0}^6\frac 1 {k!}\left(\frac{202}{100} \right)^k\\
&=\scriptstyle{1+\frac{202}{100} + \frac{40804}{20000}
+\frac{8242408}{6000000} + \frac{1664966416}{2400000000}
+ \frac{336323216032}{1200000000000}
+\frac{67937289638464}{720000000000000}}\\
&>\scriptstyle{1+\frac{202}{100} + \frac{204}{100}
+\frac{8242404}{6000000} + \frac{1664966400}{2400000000}
+ \frac{336323215920}{1200000000000}
+\frac{67937289638400}{720000000000000}}\\
&= \scriptstyle{1+\frac{202}{100} + \frac{204}{100}
+\frac{1373734}{1000000} + \frac{69373600}{100000000}
+ \frac{2802693466}{10000000000}
+\frac{94357346720}{1000000000000}}\\
&> 1+\frac{202}{100} + \frac{204}{100}
+\frac{1373}{1000} + \frac{693}{1000}
+ \frac{280}{1000}
+\frac{94}{1000} = \frac{7500}{1000}= \frac{15}2\text.
\end{align*}
Somit ist also 
\begin{align}
\log\left(\frac{15}{2}\right) < \frac{202}{100}\text.\label{nleq:11}
\end{align}
Weiter ist wegen $\frac {15} 2 < \frac{7502121}{1000000} =
\left(\frac{2739}{1000}\right)^2$ auch
\begin{align}
\sqrt{\frac{15}2}< \frac{2739}{1000}\text. \label{nleq:12}
\end{align}
Ebenso ist wegen $\frac{202}{100} < \frac{2022084}{1000000} =
\left(\frac{1422}{1000} \right)^2$ auch
$\sqrt{\frac{202}{100}}<\frac{1422}{1000}$
und analog ist auch wegen
$8305924\cdot 700 = 5814146800 > 5814000000 = 5814\cdot 1000000$
zunächst $\frac{5814}{700} < \frac{8305924}{1000000}=\left(\frac{2882}{1000}
\right)^2$ und damit
\begin{align}
\sqrt{\frac{5814}{700}}< \frac{2882}{1000}\text.\label{nleq:13}
\end{align}
Somit gilt wegen des monotonen Fallens von $h_2$ für alle 
$s \in [\frac{13}2, \infty[$:
\begin{align*}
h_2(s)&\le h_2\left(\frac{13}2\right) = \frac{\sqrt{\frac{15}2}}{\frac{13}{2}}
\left((\pi-1)\sqrt{\log\left(\frac{15}2 \right)} 
 + \sqrt{\log\left(\frac{15}2 \right) + 2\pi} \right)\\
 &\overset{\eqref{nleq:6}}<\frac{\sqrt{\frac{15}2}}{\frac{13}{2}}
\left(\frac{15}7\sqrt{\log\left(\frac{15}2 \right)} 
 + \sqrt{\log\left(\frac{15}2 \right) + \frac{44}7} \right)\\
& \overset{\substack{\eqref{nleq:11},\\\eqref{nleq:12}}}<\frac{5478}{13000}
 \left(\frac{15}7\sqrt{\frac{202}{100}} 
 + \sqrt{\frac{202}{100} + \frac{44}7} \right)\\
 &= \frac{5478}{13000}
 \left(\frac{15}7\sqrt{\frac{202}{100}}
 + \sqrt{\frac{5814}{700}}\right)
 \overset{\eqref{nleq:13}}< \frac{5478}{13000}
 \left(\frac{15}7 \cdot \frac{1422}{1000}
 + \frac{2882}{1000}\right)\\&= \frac{5478}{13000}\cdot\frac{41504}{7000}
 = \frac{227358912}{91000000}<\frac 5 2 
 =\sqrt{\frac{25}4 \cdot \frac{71}{446}}\sqrt{2 \frac{223}{71}}\\
 &= \sqrt{\frac{1775}{1784}}\sqrt{2 \frac{223}{71}}\overset{\eqref{nleq:6}}<
  \sqrt{2\pi}
 \text.
\end{align*}
(4) Seien $\eta \in\!\  ]0, 1[$ und $f_\eta:\!\  ]1, \infty[ \rightarrow \R,
t \mapsto \log(2t)-\eta t$.\\
Offenbar ist $f_\eta$ stetig differenzierbar und es gilt
$f_\eta'=\frac 1\cdot- \eta$.\\
Daher ist $\forall\, t \in [\frac 1 \eta, \infty[: f_\eta'(t) \le 0$.\\
Somit ist $\restrict{f_\eta}{[\nicefrac 1 \eta, \infty[}$ monoton fallend.\\
Wegen
\[4\log\left(\frac 2 \eta \right) = \log\left(\frac{16}{\eta^4} \right)
> \log(16) > 1\]
und damit $\frac 4 \eta \log\left(\frac 2 \eta\right) > \frac 1 \eta$ genügt es
also zu zeigen, dass
$f_\eta\left(\frac 4\eta\log\left(\frac 2 \eta\right)\right)<0$
gilt.\\
Es gilt allerdings $\forall\, x \in \R: \exp(x) \ge 1+x > x$, also
$\forall\, x \in \R: \log(x) < x$ und damit:
\begin{align*}
f_\eta\left(\frac 4 \eta \log\left(\frac 2 \eta\right) \right) & =
\log\left(\frac 8 \eta \log\left(\frac 2 \eta\right) \right)
 -4\log\left(\frac 2 \eta\right)\\
&= \log\left(\frac 8 \eta \log\left(\frac 2 \eta\right) \right)
  -\log\left(\frac{16}{\eta^4} \right)
  = \log\left(\frac{\frac{8 \log\left(\frac 2 \eta\right)}{\eta}}
    {\frac{16}{\eta^4}} \right)\\
&= \log\left(\frac{8\eta^4\log\left(\frac 2 \eta \right)}{16\eta} \right)
= \log\left(\frac 1 2 \eta^3 \log\left(\frac{2}{\eta} \right) \right)
< \log\left(\frac 1 2 \eta ^3 \frac 2 \eta \right)\\
& = \log(\eta^2) < 0\text.
\end{align*}
Somit gilt für
$t \in \bigg[\frac 4 \eta \log\left(\frac 2 \eta\right),\infty  \bigg[$ mit
$f_\eta(t)< 0$ auch $\log(2t)< \eta t$.
\end{proof}

\chapter{Wahrscheinlichkeitstheoretische Grundlagen}
\label{ch:kapitel2}

Dieses Kapitel dient dazu, einige für uns wichtige Ergebnisse der Maß-- bzw.
Wahrscheinlichkeitstheorie vorzustellen.

Sei dazu für dieses Kapitel $(\Omega, \mfa, \Prim)$
ein Wahrscheinlichkeitsraum. Für alle $n \in \N$ setzen wir auf dem $\R^n$
die Borel'sche $\sigma$--Algebra $\mfb_n$ voraus, wobei wir speziell
$\mfb \coloneqq  \mfb_1$ schreiben.\\
Weiter setzen wir für alle $n \in \N$:
\[\mcm_n(\Omega, \mfa, \R^n)
 \coloneqq \menge{f\in\mcf(\Omega,\R^n)}{f\quad \mfa - \mfb_n - \text{messbar}}\text.\]
Schließlich sei für alle $n \in \N$:
\[\mct_n(\Omega, \mfa, \R^n)
\coloneqq  \menge{f \in \mcm_n(\Omega, \mfa, \R^n)}{\Bild f \text{ endlich}}\text.\]
Auch hier treffen wir die Vereinbarung, dass für $n = 1$ der Index weggelassen
werden möge.

\begin{definition}{(vollst\"andig) unabhängig/identisch verteilt/zentriert}
                  {UnabhaengigIdentischVerteilt}
Seien $n, m \in \N$, $\gamma \in \mcm_n(\Omega, \mfa, \R^n)$,
$\delta \in \mcm_m(\Omega, \mfa, \R^m)$ und $F: \R \rightarrow [0, 1]$ eine
Verteilungsfunktion.\\
$\gamma$ heißt \textbf{unkorreliert}
\[: \iff \gamma \in L_2(\Omega, \mfa, \Prim, \R^n) \; \wedge \;
    \forall\, i, j \in \haken n: i \neq j \folgt \E(\gamma_i\gamma_j)
  = \E(\gamma_i)\E(\gamma_j)\text.\]
$\gamma$ und $\delta$ heißen \textbf{vollständig unkorreliert}
\[:\iff \gamma, \delta \text{ unkorreliert } \; \wedge \;
        \forall\, i \in \haken n\, \forall\, j \in \haken m:
        \E(\gamma_i\delta_j)= \E(\gamma_i)\E(\delta_j)\text.\]
$\gamma$ heißt \textbf{unabhängig}
\[:\iff \forall A \in \mathfrak B^n: 
\Prim\left(\bigcap_{i=1}^n \meng{\gamma_i\in A_i}\right)
=\prod_{i=1}^n \Prim(\gamma_i\in A_i)\text.\]
$\gamma$ und $\delta$ heißen \textbf{vollständig unabhängig}
\begin{align*}
&:\iff \,
\forall\, A\in \mfb^n\, \forall\, B\in \mfb^m:\\
& \Prim\left(\bigcap_{i=1}^n \meng{\gamma_i \in A_i} \cap 
\bigcap_{j=1}^m \meng{\delta_j \in B_j} \right)
= \left(\prod_{i=1}^n \Prim(\gamma_i\in A_i)\right)
\left(\prod_{j=1}^m \Prim(\delta_j\in B_i)\right)\text.
\end{align*}
$\gamma$ heißt \textbf{identisch verteilt}
$: \iff \forall \, i, j \in \haken n\, \forall \, t \in \R:
    \Prim(\gamma_i \le t) = \Prim(\gamma_j \le t)$.\\
$\gamma$ heißt \textbf{$F$--verteilt}
$:\iff \gamma \text{ identisch verteilt } \; \wedge \;
 \forall \, t \in \R: \Prim(\gamma_1 \le t) = F(t)$.\\
$\gamma$ heißt \textbf{zentriert}
$:\iff \forall\, i \in \haken n: \E(\gamma_i) = 0$.\\
$\gamma$ heißt \textbf{$L_2$--normiert}
$:\iff \gamma \in L_2(\Omega, \mfa, \Prim, \R^n) \; \wedge \;
   \forall i \in \haken n: \norm{f_i}_2 = 1$.
\end{definition}

Für Zufallsgrößen ist bekannt, dass die Unabhängigkeit die
Unkorreliertheit impliziert (siehe z. B. \cite{Irle}, 10.16).
Dies wird später von Nutzen sein, weil wir in
diesem Kapitel weitgehend mit Unkorreliertheit auskommen werden und später
Unabhängigkeit voraussetzen werden. Da die Unkorreliertheit einfacher zu
prüfen ist, werden wir somit viel Arbeit einsparen können.

Wir beobachten, dass für zentrierte Zufallsgrößen, die quadratintegrierbar
sind, die Varianz genau das Quadrat der $L_2$--Norm ist. Außerdem wissen wir
aus der Funktionalanalysis, dass $L_2(\Omega, \mfa, \Prim, \R)$ ein
Hilbertraum ist, dessen Skalarprodukt von $\E$ (bzw. $\int_\Omega$) stammt.

\begin{lemma}{Elementare Eigenschaften quadratintegrierbarer Abbildungen}
             {ElementareEigenschaftenQuadratintegrierbarerAbbildungen}

Seien $n \in \N$ und $\gamma \in L_2(\Omega, \mfa, \Prim, \R^n)$ zentriert,
unkorreliert und $L_2$--normiert.\\
Dann gilt:
\begin{enumerate}[(1)]
	\item \label{ElementareEigenschaftenQuadratintegrierbarerAbbildungen1}
 $\forall\, i, j \in \haken n : i\neq j\folgt\sprod{\gamma_i}{\gamma_j}=0$.
\item\label{ElementareEigenschaftenQuadratintegrierbarerAbbildungen2}
 $\bigg\|\Sum{i=1}n \gamma_i\bigg\|_2^2=\Sum{i=1}n \norm{\gamma_i}_2^2= n$.
\item\label{ElementareEigenschaftenQuadratintegrierbarerAbbildungen3}
 $\forall\, c \in \R^n: \norma{\bprod c \gamma_{n, L_2}}_2=\norm c_{2, n}$.
\item\label{ElementareEigenschaftenQuadratintegrierbarerAbbildungen4}
 $\norm\cdot_{2, n} \circ\gamma$ ist $\mfa-\mfb$--messbar und
    $\E(\norm\cdot_{2, n} \circ \gamma) \le \sqrt n$.
\end{enumerate}
\end{lemma}
\begin{proof}
(1) Seien $i \neq j \in \haken n$.\\
Dann ist $\sprod {\gamma_i} {\gamma_j}
= \E(\gamma_i\gamma_j) = \E(\gamma_i)\E(\gamma_j) = 0$, da $\gamma$ zentriert
ist.\\
(2) Die erste Gleichheit ist Pythagoras im Hilbertraum 
$L_2(\Omega, \mfa, \Prim, \R)$, die zweite folgt aus der Forderung, dass f�r
alle $i \in \haken n$ wegen $\gamma$ $L_2$--normiert $\norm{\gamma_i}_2 = 1$
gilt.\\
(3) Sei $c \in \R^n$. Dann gilt mit Pythagoras:
\[\norma{\bprod c \gamma _{n, L_2}}_2^2 = \norma{\sum_{i=1}^nc_i\gamma_i}_2^2
= \sum_{i=1}^n c_i^2 \norm{\gamma_i}_2^2 = \norm c _{2, n}^2\text.\]
(4) Sei $Y := \norm\cdot_{2,n}^2 \circ \gamma$.\\
Da $\gamma$ messbar,
$\norm\cdot_{2, n}$ stetig bez�glich sich selbst, $q: \R \rightarrow \R,
t\mapsto t^2$ stetig sind und auf $\R$ die Borel'sche $\sigma-$Algebra
vorliegt, sind $Y$ und $\sqrt\cdot \circ Y$ als Hintereinanderausf�hrung
messbarer und stetiger Funktionen selbst messbar.\\
Da $\norm\cdot_{2, n} \circ \gamma \ge 0$, gilt
$\norm\cdot_{2, n} \circ \gamma = \sqrt\cdot \circ Y$.\\
Die Wurzelabbildung ist konkav, also gilt mit der Jensenschen Ungleichung und
der Normiertheit der Komponentenfunktionen von $\gamma$:
\[\E\left(\sqrt\cdot \circ Y\right) \le \sqrt{\E(Y)} 
= \sqrt{\sum_{i=1}^n\E(\gamma_i^2)} = \sqrt n\text.\]
\end{proof}
\parag
Linearkombinationen unabhängiger Zufallsgrößen werden uns an späterer Stelle
oft interessieren. Insofern ist das folgende Lemma von Interesse, welches
besagt, dass der Linearkombinator vollständig unkorrelierte Abbildungen auf
unkorrelierte abbildet.

\begin{lemma}{Unkorreliertheit im Linearkombinator}
             {UnkorreliertheitImLinearkombinator}
Seien $n, m \in \N$, $\gamma\in \mcm_n(\Omega, \mfa, \R^n)$,
$\delta\in\mcm_m(\Omega, \mfa, \R^m)$, $x \in \R^n$ und $y \in \R^m$.\\
$\gamma$ und $\delta$ seien vollst�ndig unkorreliert.\\
Dann sind
$\bprod x \gamma _{n}, \bprod y \delta _m \in L_2(\Omega, \mfa, \Prim, \R)$
und es gilt:
\begin{align*}\E\bigg(\bprod x \gamma _{n} \cdot \bprod y \delta _m \bigg)
  = \E\left(\bprod x \gamma_{n}\right)\cdot\E\left(\bprod y\delta_{n} \right)
= \sprod x {\pfeil \E n (\gamma)} _n \cdot \sprod y {\pfeil \E m (\delta)} _m
\text.
\end{align*}
\end{lemma}
\begin{proof}
Die Quadratintegrierbarkeit der beiden Funktionen folgt einfach aus der
Tatsache, dass die Komponentenfunktionen von $\gamma$ und $\delta$ selbst
quadratintegrierbar sind. Da $L_2(\Omega, \mfa, \Prim, \R)$ ein Vektorraum ist,
liegen darin also auch Linearkombinationen quadratintegrierbarer Funktionen.\\
Da $\gamma, \delta$ vollst�ndig unkorreliert sind, gilt:
\[\forall\, i \in \haken n\, \forall\, j \in \haken m: \E(\gamma_i\delta_j)
= \E(\gamma_i)\E(\delta_j)\text.\]
Daher erhalten wir:
\begin{align*}
& \E\left(\bprod x \gamma _{n} \cdot \bprod y \delta _m \right)
= \E\left( \left(\sum_{i=1}^n x_i \gamma_i\right) 
           \left(\sum_{j=1}^m y_i \delta_j \right) \right)\\
= & \E\left(\sum_{i=1}^n \sum_{j=1}^m x_iy_j\gamma_i\delta_j \right)
 = \sum_{i=1}^n \sum_{j=1}^m x_iy_j\E(\gamma_i\delta_j)\\
 =& \sum_{i=1}^n \sum_{j=1}^m x_iy_j\E(\gamma_i)\E(\delta_j)
 = \sum_{i=1}^n \sum_{j=1}^m \E(x_i\gamma_i)\E(y_j\delta_j)\\
 =& \left(\sum_{i=1}^n\E(x_i\gamma_i) \right) 
    \left(\sum_{j=1}^m\E(y_j\delta_j) \right)
 = \E\left(\sum_{i=1}^n x_i\gamma_i \right)
   \E\left(\sum_{j=1}^m y_j\delta_j \right)\\
 =& \E(\bprod x \gamma _n) \E(\bprod y \delta _m)\text.
\end{align*}
Die zweite Gleichheit folgt aus 
\refclaim{EigenschaftenDesLinearkombinators}
{EigenschaftenDesLinearkombinators2}, denn $\E$ ist linear.
\end{proof}
\parag
In diesem Lemma geht entscheidend ein, dass $\gamma$ und $\delta$ vollständig
unkorreliert sind, sodass man es nicht für die Berechnung des obigen Ausdrucks
im Falle $\gamma = \delta$ verwenden kann.\\
Fordert man zusätzlich, dass $\gamma$ zentriert und $L_2$--normiert ist, so
erhält man allerdings eine ähnliche Aussage.

\begin{lemma}{Unkorreliertheit im Linearkombinator (eine Funktion)}
             {UnkorreliertheitImLinearkombinatorEinfach}
Seien $n \in \N$, $\gamma \in\mcm_n(\Omega, \mfa, \R^n)$ zentriert,
unkorreliert und $L_2$--normiert.\\
Dann sind $\bprod x \gamma _n, \bprod y \gamma_n\in L_2(\Omega,\mfa,\Prim, R)$
und es gilt:
\[\forall\, x, y \in \R^n : \E\bigg(\bprod x \gamma _n\bprod y \gamma_n\bigg)
 = \sprod x y _n\text.\]
\end{lemma}
\begin{proof}
Die erste Aussage folgt mit dem gleichen Argument wie in~\ref{UnkorreliertheitImLinearkombinator}.\\
Seien $x, y \in \R^n$.\\
Da $\gamma$ unkorreliert und zentriert ist, gilt 
$\forall i\neq j\in\haken n:\E(\gamma_i\gamma_j)=\E(\gamma_i)\E(\gamma_j)=0$.\\
Daher gilt:
\begin{align*}
 \E\bigg(\bprod x \gamma _n\bprod y \gamma_n\bigg)
& = \E\left(\sum_{i=1}^n\sum_{j=1}^n  x_iy_j\gamma_i\gamma_j\right)
= \sum_{i=1}^n\sum_{j=1}^n x_iy_j\E(\gamma_i\gamma_j)\\
&= \sum_{i=1}^n x_iy_i \underbrace{\E(\gamma_i^2)}_{=\norm{\gamma}_2^2 = 1}
= \sprod x y_n\text.
\end{align*}
\end{proof}
\parag
Will man die Integrierbarkeit einer messbaren Abbildung nachweisen, so ist es
oft einfach, dieses dadurch zu erreichen, dass man eine integrierbare Majorante
und Minorante angibt. Wir halten diesen Umstand in dem folgenden Lemma fest.

\begin{lemma}{Integrierbarkeitskriterien}
             {Integrierbarkeitskriterien}
Seien $f, g, h \in \mcm(\Omega, \mfa, \R)$.\\
Dann gilt:
\begin{enumerate}[(1)]
	\item \label{Integrierbarkeitskriterien1}
 Sind $f, g \ge 0$, $f \le g$ und $g$ integrierbar, so auch $f$ mit
$0 \le \E(f) \le \E(g)$.
\item \label{Integrierbarkeitskriterien2}
 Sind $h \le f \le g$ und $g, h$ integrierbar, so auch $f$ und
$\E(h) \le \E(f) \le \E(g)$.
\end{enumerate}
\end{lemma}
\begin{proof}
Ist die Integrierbarkeit von $f$ in den beiden F�llen gekl�rt, so folgt die
jeweils zweite Aussage einfach �ber die Monotonie des Erwartungswerts.\\
(1) Seien $f, g \ge 0$, $f\le g$ und $g$ integrierbar.\\
Sei $\vi \in \mct(\Omega, \mfa, \R)$ mit $0 \le \vi \le f$.\\
Dann ist auch $0 \le \vi \le f \le g$ und daher
\[\menge{\E(\psi)}{\psi \in \mct(\Omega,\mfa,\R), 0 \le \psi\le f}
\subseteq
 \menge{\E(\psi)}{\psi \in \mct(\Omega,\mfa,\R), 0 \le \psi\le g}\text.\]
Die gr��ere Menge ist allerdings beschr�nkt, weil $g$ integrierbar ist.\\
Also ist $\menge{\E(\psi)}{\psi \in \mct_1(\Omega,\mfa,\R), 0 \le \psi\le f}$
nichtleer und nach oben beschr�nkt, besitzt also ein Supremum in $\R$.\\
Daher
\begin{align*}
\E(f)&=\sup\menge{\E(\psi)}{\psi \in \mct(\Omega,\mfa,\R), 0 \le \psi\le f}\\
   &\le\sup\menge{\E(\psi)}{\psi \in \mct(\Omega,\mfa,\R), 0 \le \psi\le g}
   = \E(g) < \infty\text.
\end{align*}
Somit ist $f$ integrierbar.\\
(2) Seien $h \le f \le g$ und $g, h$ integrierbar.\\
Es gilt also $0 \le f-h \le g -h$.\\
Da $g, h$ integrierbar sind, ist es auch $g-h$ und damit nach (1) auch $f-h$.\\
Analog ist also $(f-h)+ h = f$ integrierbar.
\end{proof}
\parag
Ähnlich wie in dem vorigen Lemma bei dem Vorliegen einer integrierbaren
Minorante auf Nichtnegativität verzichtet werden konnte, kann man den Satz von
der monotonen Konvergenz auf den Spezialfall einer integrierbaren Folge
anpassen.

\begin{lemma}{monotone Konvergenz im integrierbaren Fall}
             {MonotoneKonvergenzImIntegrierbarenFall}
Seien $\folge f n \N \in L_1(\Omega, \mfa, \Prim, \R)^\N$ und 
$f \in L_1(\Omega, \mfa, \Prim, \R)$.\\
Die Folge $\folge f n \N$ sei monoton wachsend und punktweise konvergent gegen
$f$.\\
Dann ist $\left(\E(f_n) \right)_{n \in \N}$ konvergent und es gilt:
$ \Limes n \infty \E(f_n) = \E(f)$.
\end{lemma}
\begin{proof}
Seien $g_n := f_n - f_1$ f�r alle $n \in \N$.\\
Dann ist $g_n \in L_1(\Omega, \mfa, \R)$ und $g_n \ge 0$ f�r alle $n \in \N$.\\
Wegen der Monotonie von $\folge fn\N$ ist auch $\folge gn\N$ monoton
wachsend.\\
Au�erdem ist $\folge g n \N$ punktweise konvergent gegen $f - f_1$.\\
Nach dem Satz von der monotonen Konvergenz ist also 
$\left(\E(g_n) \right)_{n \in \N}$ konvergent und es gilt wegen der 
Integrierbarkeit von $f$ und $f_1$:
\[\Limes n \infty \E(g_n) = \E(f - f_1) = \E(f) - \E(f_1)\text.\]
Also auch
\begin{align*}
\E(f) & = \left(\Limes n \infty \E(g_n) \right) + \E(f_1)
        = \Limes n \infty \bigg(\E(g_n) + \E(f_1) \bigg)
       = \Limes n \infty \E(g_n + f_1)\\
        &= \Limes n \infty \E(f_n)\text.
\end{align*}
\end{proof}
\parag
Aus der Stochastik ist bekannt (siehe z. B. \cite{Irle}, dort 8.12), dass $\E$
monoton ist. Es gilt allerdings auch, dass $\E$ streng monoton ist. In der Tat
genügt eine hinreichend große Menge (bezüglich $\Prim$), auf der ein $<$
vorliegt, um selbiges im Erwartungswert zu erreichen.

\begin{lemma}{Strenge Monotonie des Erwartungswerts}
             {StrengeMonotonieDesErwartungswerts}
Seien $f,g\in L_1(\Omega,\mfa,\Prim, \R)$ mit $f \le g$.\\
Dann ist $\meng{f < g}\in \mfa$ und es gilt:
\[\Prim(f < g)>0 \folgt \E(f) < \E(g)\]
\end{lemma}
\begin{proof}
Die Messbarkeit von $\meng{f < g}$ ist bekannt.\\
Sei also $\E(f) \ge \E(g)$.\\
Wegen $f \le g$ ist auch $\E(f) \le \E(g)$ und somit $\E(f) = \E(g)$.\\
Es ist $g-f\ge 0$ und $\E(g-f)=\E(g) -\E(f) = 0$, da $f,g$ integrierbar sind.\\
Daher ist also $0 = \Prim(g-f > 0) = \Prim(f < g)$.
\end{proof}
\parag
Sind $f,g$ nichtnegative, integrierbare Funktionen mit $0 < \E(f) \le \E(g)$,
so heißt es i. A. nicht, dass die gleiche Abschätzung schon für die Funktionen
gilt. Allerdings hat die Menge, auf welcher eine vergleichbare Aussage schon
für $f, g$ gilt, positives Maß.
Dieser Sachverhalt ist Inhalt des folgenden Lemmas.

\begin{lemma}{Erwartungswertlemma}{Erwartungswertlemma}
Seien $f,g\in L_1(\Omega,\mfa,\Prim,\R)$ mit $f,g\ge 0$ und $\E(f)\le\E(g)$.\\
Seien $a, b \in \R_{>0}$ mit $a \le \E(f) \le \E(g) \le b$.\\
Dann ist $\meng{f > 0}\cap \meng{ag \le bf} \in \mfa$ und es gilt:
$\Prim\big(\meng{f > 0}\cap \meng{ag \le bf}\big) > 0\text.$\\
Insbesondere existiert also ein $\omega \in \Omega$ mit $f(\omega) > 0$ und
$\frac {g(\omega)}{f(\omega)} \le \frac b a$.
\end{lemma}
\begin{proof}
Da $f$ und $g$ integrierbar sind, sind auch $ag$ und $bf$ integrierbar. Damit
sind sie messbar und daher sind $\meng{f >0}$,
$\meng{ag \le bf} \in \mfa$. Also auch deren Schnitt.\\
Annahme: $\Prim(\meng{f > 0}\cap \meng{ag \le bf}) = 0$.\\
Dann ist $\Prim(f = 0 \vee ag > bf) = 1$.\\
Es ist 
\begin{align*}\meng{f = 0 \vee ag > bf} &= \meng{f=0} \cup \meng{ag > bf} \\
&= \meng{f=0} \overset{\cdot}\cup \bigg(\meng{ag>bf}\setminus\meng{f=0}
\bigg)\text.\end{align*}
Da $\E(f) > 0$ und $f \ge 0$ ist $\Prim(f = 0) < 1$.\\
Wegen der Disjunktheit von $\meng{ag >bf}\setminus\meng{f=0}$ und $\meng{f=0}$
gilt daher
\[\Prim\left(ag > bf \wedge f > 0 \right) = 1 - \Prim(\meng{f= 0}) > 0\text.\]
Sei $A \coloneqq  \meng{ag>bf}$.\\
Dann ist $\meng{ag>bf}\setminus\meng{f=0}\subseteq A$ und daher $\Prim(A)>0$.\\
Wegen $b>0$ ist $A = \meng{f < \frac a b g}$.\\
Seien $X \coloneqq  f \Eins_A$ und $Y\coloneqq  \left(\frac a b g\right) \Eins_A$.\\
Dann sind $0 \le X \le f$ und $0 \le Y \le \frac a b g$, also nach
\ref{Integrierbarkeitskriterien}: $X, Y$ integrierbar.\\
Weiter sind $X \le Y$ und $\meng{X < Y} = A$.\\
Nach~\ref{StrengeMonotonieDesErwartungswerts} gilt somit: $\E(X) < \E(Y)$.\\
Erhalte wegen der Integrierbarkeit von $g$:
\begin{align*}
& \E(f) = \E\left(f\Eins_{\meng{f = 0 \vee ag > bf}}\right) 
= \E\left(f\Eins_{\meng{f=0}} +f\Eins_{\meng{ag > bf}\setminus\meng{f=0}}\right)\\
=&\, \E\left(f\Eins_{\meng{ag > bf}\setminus\meng{f=0}} \right) 
\le \E\left(f\Eins_A \right)= \E(X) < \E(Y) = \E\left(\frac abg \right)
=\frac a b \E\left(g\right) \le a
\end{align*}
Was aber wegen $\E(f) \ge a$ nicht sein kann. $\lightning$
\end{proof}
\parag
Das nachfolgende Lemma wird für uns von großer Wichtigkeit sein, denn
einerseits sichert es die Existenz beliebig vieler unabhängiger, identisch
verteilter Zufallsgrößen zu und garantiert dabei noch andererseits, dass all
diese überall endlich sind.

\begin{lemma}{Existenz (unabh\"angiger) F--verteilter Zufallsgrößen}
             {ExistenzFVerteilterZufallsgroessen}
Seien $F: \R \rightarrow [0,1]$ eine Verteilungsfunktion und $n \in \N$.
Dann gilt:
\begin{enumerate}[(1)]
	\item \label{ExistenzFVerteilterZufallsgroessen1}
	Es gibt einen Wahrscheinlichkeitsraum $(\Omega_1, \mfc, \Prim_1)$
	und $\xi \in \mcm(\Omega_1, \mfc, \R)$, sodass $\xi$ $F$--verteilt
	 ist.\\
	Insbesondere ist $\xi$ sicher endlich.
	\item \label{ExistenzFVerteilterZufallsgroessen2}
	Es gibt einen Wahrscheinlichkeitsraum $(\Omega_2, \mfd, \Prim_2)$ und
	$\eta \in \mcm_n(\Omega_2, \mfd, \R^n)$ derart, dass $\eta$ $F$--verteilt  
	und unabhängig ist (im Sinne von~\ref{UnabhaengigIdentischVerteilt}).\\
	Insbesondere ist $\eta_i$ sicher endlich für alle $i \in \haken n$.
\end{enumerate}
\end{lemma}
\begin{proof}
Als Verteilungsfunktion ist 
$F$ monoton wachsend, rechtsseitig stetig und es gilt
\[\Limes t {-\infty}F(t) = 0\text{, sowie }\Limes t \infty F(t) = 1\text.\]
(1) Aus~\cite{Irle}, dort Beweis von 6.6, geht hervor, dass gilt:
\[G:\!\  ]0,1[ \rightarrow \R, p \mapsto \inf\menge{x \in \R}{F(x) \ge p}\]
ist wohldefiniert und genügt der Äquivalenz
\begin{align}
\forall\,t\in\R\,\forall\,s\in\!\ ]0,1[:
G(s)\le t\iff s\le F(t)\text.\label{efvzeq:1}
\end{align}
Seien $\Omega_1\coloneqq \!\ ]0,1[$, $\mfc \coloneqq  \Pot(]0,1[)\: \cap\: \mfb$ und
$\Prim_1\coloneqq \restrict{\lambda}{\mfc}$, wobei $\lambda$ das eindimensionale Lebesguemaß
bezeichne.\\
Sei $t \in \R$.\\
Dann ist
\[\meng{G \le t} = \menge{x \in\!\  ]0,1[}{G(x) \le t}
\overset{\eqref{efvzeq:1}}= 
\menge{x \in\!\  ]0,1[}{x \le F(t)} = \!\ ]0, F(t)] \in \mfc\]
und daher $\Prim_1(G \le t) = \lambda (]0, F(t)]) = F(t)$, sodass
$G$ $F$--verteilt ist.\\
Damit wurde ein Wahrscheinlichkeitsraum und eine Abbildung wie behauptet
gefunden.\absatz
(2) Seien $\Omega_2\coloneqq \!\ ]0,1[^n$, $\mfd\coloneqq \Pot\left(]0,1[^n\right) \cap \mfb_n$
und $\Prim_2\coloneqq  \restrict{\lambda_n}{\mfd}$, wobei $\lambda_n$ das $n$--dimensionale
Lebesguemaß bezeichne.\\
Nach (1) existiert ein $F$--verteiltes $\xi:\!\ ]0,1[\rightarrow\R$.\\
Sei $\eta: \Omega_2 \rightarrow \R^n, \omega \mapsto \pfeil \xi n (\omega)$.\\
Offenbar gilt dann für alle $j \in \haken n$ und alle $\omega \in \Omega_2$:
$\eta_j(\omega) = \xi(\omega_j)$.\\
Wir können o. B. d. A. annehmen, dass $n > 1$ ist, denn für $n = 1$ sagt (2)
dasselbe aus, wie (1).\\
Seien $i \in \haken n$ und $t \in \R$.\\
Falls $1<i<n$, so gilt:
\begin{align*}
\meng{\eta_i \le t} & = \menge{\omega \in \Omega_2}{\eta_i(\omega) \le t}
 = \menge{\omega \in \Omega_2}{\xi(\omega_i) \le t}\\
 &=\!\  ]0,1[^{i-1}\times\menge{x \in\!\  ]0, 1[}{\xi(x)\le t}\times\!\ 
 ]0,1[^{n-i}
 \in \mfd
\end{align*}
falls $i = 1$:
\begin{align*}
\meng{\eta_i \le t} & = \menge{\omega \in \Omega_2}{\eta_1(\omega) \le t}
 = \menge{\omega \in \Omega_2}{\xi(\omega_1) \le t}\\
 &= \menge{x \in\!\  ]0,1[}{\xi(x) \le t}\times\!\  ]0,1[^{n-1} \in \mfd
\end{align*}
und falls $i = n$:
\begin{align*}
\meng{\eta_i \le t} & = \menge{\omega \in \Omega_2}{\eta_n(\omega) \le t}
 = \menge{\omega \in \Omega_2}{\xi(\omega_n) \le t}\\
 &=\!\  ]0,1[^{n-1} \times \menge{x \in \!\ ]0,1[}{\xi(x) \le t}  \in \mfd
\end{align*}
Außerdem gilt
\begin{align}
\Prim_2(\eta_i \le t) & = \lambda_n(\eta_i \le t) \notag\\
&= \begin{cases} 
\lambda_n\bigg(
]0,1[^{i-1}\times\menge{x\in\!\ ]0,1[}{\xi(x)\le t}\times\ \!]0,1[^{n-i} \bigg)
 & : \, 1 < i < n\notag\\
\lambda_n\bigg(
\menge{x \in\!\  ]0,1[}{\xi(x) \le t}\times\!\ ]0,1[^{n-1} \bigg) & : \,i = 1
\notag\\
\lambda_n\bigg(]0,1[^{n-1}\times\menge{x\in\!\  ]0,1[}{\xi(x) \le t} \bigg)&:\,
i=n\end{cases}\notag\\
&= \lambda\left(\menge{x \in\!\  ]0,1[}{\xi(x) \le t}\right) \cdot
 \lambda_{n-1}\left(]0,1[^{n-1} \right)\notag\\
 &= \lambda\left(\menge{x \in\!\  ]0,1[}{\xi(x) \le t}\right)
  \overset{(1)}= F(t) \label{efvzeq:2}
\end{align}
Damit ist $\eta$ $F$--verteilt.\\
Sei nun $s \in \R^n$. Dann gilt:
\begin{align*}
& \Prim_2\left(\bigcap_{k=1}^n \meng{\eta_k \le s_k}\right)
= \Prim_2\left(\forall\, k \in \haken n: \eta_k \le s_k\right)\\
 =&\, \lambda_n\left( \menge{\omega \in \Omega_2}
  {\forall\, k \in \haken n: \eta_k(\omega) \le s_k} \right)
= \lambda_n\left(\menge{\omega \in \Omega_2}
  {\forall\, k \in \haken n: \xi(\omega_k) \le s_k} \right) \\
= &\,\lambda_n \left(\prod_{k = 1}^n \menge{x \in \!\ ]0,1[}
  {\xi(x) \le s_k}\right)
= \prod_{k=1}^n \lambda\left(\menge{x \in\!\  ]0,1[}
  {\xi(x) \le s_k} \right)\\
\overset{(1)}=&\, \prod_{k=1}^n F(s_k) \overset{\eqref{efvzeq:2}}= 
\prod_{k=1}^n \Prim_2(\eta_k \le s_k)
\end{align*}
Damit ist $\eta$ auch unabhängig und die Behauptung gezeigt.
\end{proof}

\begin{definition}{Unabh�ngiges (\textsl N(0,1),
\textsl n)--Tupel}{NormalNTupel}
Seien $(\Omega', \mfa', \Prim')$ ein Wahrscheinlichkeitsraum,
$n \in \N$, $\alpha \in \N^n$ und f�r alle $i \in \haken n$ sei
$\gamma_i :\Omega' \rightarrow \R^{\alpha_i}$.\\
Wir nennen $\bigg((\Omega', \mfa', \Prim'), (\gamma_i)_{i=1}^n,
(\alpha_i)_{i=1}^n \bigg)$ \textbf{unabh�ngiges $(N(0,1), n)$--Tupel
zu $(n, \alpha)$}
\begin{align*}
:\iff &(1)\quad\forall\,i\in\haken n:\gamma_i \text{ ist }
N(0,1)\text{--verteilt } \wedge\\
&(2)\quad \forall\,i, j \in \haken n: i \neq j \folgt \gamma_i, \gamma_j 
     \text{ vollst�ndig unabh�ngig.}
\end{align*}
\end{definition}

\begin{korollar}{Existenz unabhängiger (\textsl N(0,1),
\textsl n)--Tupel}{ExistenzNormalNTupel}
Seien $n \in \N$ und $\alpha \in \N^n$.\\
Dann existiert ein unabhängiges $(N(0,1),n)$--Tupel zu $(n, \alpha)$.
\end{korollar}
\begin{proof}
Sei $N\coloneqq  \Sum{i=1}n \alpha_i$.\\
Nach~\refclaim{ExistenzFVerteilterZufallsgrössen}
{ExistenzFVerteilterZufallsgrössen2} gibt es einen Wahrscheinlichkeitsraum
$(\Omega', \mfa', \Prim')$ und $\eta: \Omega' \rightarrow \R^N$ unabhängig und
$N(0,1)$--verteilt.\\
Definiere für alle $i \in \haken n$:
\[s_i \coloneqq  \left(\sum_{k=1}^{i-1}\alpha_i\right)+1\text{, }
  t_i \coloneqq  \sum_{k=1}^{i}\alpha_i \text{ und }
 \gamma_i\coloneqq \left(\eta_j \right)_{j = s_i}^{t_i}\text.\]
Seien nun $k, l \in \haken n$ mit $k \neq l$. Dann sind $\gamma_k$ und
$\gamma_l$ vollständig unabhängig als Teilfamilien einer unabhängigen Familie
und damit ist \[\bigg((\Omega',\mfa',\Prim'),(\gamma_i)_{i=1}^n, (\alpha_i)_{i=1}^n\bigg)\]
ein $(N(0,1),n)$--Tupel zu $(n, \alpha)$.
\end{proof}
\parag
Zu Beginn des Kapitels
(\ref{ElementareEigenschaftenQuadratintegrierbarerAbbildungen})
hatten wir schon einmal mit dem Ausdruck  $\E(\norm\cdot_{2,n}\circ\gamma)$
für $n \in \N$ und $\gamma \in \mcm(\Omega, \mfa, \R^n)$ zu tun. An dieser
Stelle wollen wir diesen Ausdruck für bestimmte $\gamma$ explizit berechnen.

\begin{lemma}{Erwartungswert der 2--Norm eines N(0,1)--Prozesses}
             {ErwartungswertDer2NormEinesStandardnormalverteiltenProzesses}
Seien $n \in \N$, $\gamma : \Omega \rightarrow \R^n$ unabhängig und
$N(0,1)$--verteilt.\\
Sei $\Gamma:\!\  ]0, \infty[ \rightarrow \R, 
             t \mapsto \int_0^\infty x^{t-1}\exp(-x)dx$.\\
Dann gilt:
\[\E\left(\norm\cdot_{2, n} \circ \gamma \right)
= \sqrt 2\frac{\Gamma\left(\frac{n+1}2\right)}{\Gamma\left(\frac n 2 \right)}
\text.\]
\end{lemma}
\begin{proof}
Seien $Z \coloneqq  \norm\cdot_{2, n}\circ \gamma$ und $Y \coloneqq  Z^2$.\\
Nach Definition ist $Y$ dann $\chi^2_n$--verteilt.\\
Für alle $\nu, \lambda \in \!\ ]0, \infty[$ ist 
\[ f_{\nu, \lambda}: \R \rightarrow \R, t \mapsto
	\begin{cases}
\frac 1 {\Gamma(\nu)}\lambda^\nu t^{\nu -1}\exp\left(-\lambda t\right) &:t> 0\\
0 & :\text{ sonst}
	\end{cases}
\]
die Dichte der Gammaverteilung mit den Parametern $\nu$ und $\lambda$
($f_{\nu, \lambda}$ ist nach~\refclaim{ElementareEigenschaftenDerGammafunktion}
{ElementareEigenschaftenDerGammafunktion1} wohldefiniert). (Vgl.~\cite{Irle},
dort 20.8 und 20.10).\\
Weiter ist eine $\chi_n^2$--verteilte Zufallsgröße auch gammaverteilt mit den
Parametern $\nu = \frac n 2$ und $\lambda = \frac 1 2$.\\
Da $Y, Z \ge 0$, ist $Z = \sqrt\cdot \circ Y$ und
weil $f_{\frac n 2, \frac 1 2}, f_{\frac {n+1}2, \frac 1 2}$ stetige
Dichten sind, gilt:
\begin{align*}
\E(\norm\cdot_{2, n} \circ \gamma) &= \E(Z) = \E(\sqrt\cdot \circ Y)
 = \int_{0}^\infty\sqrt t \Prim^Y(dt) 
 = \int_0^\infty t^{\frac 1 2}f_{\frac n 2, \frac 1 2}(t) dt\\
 &= \frac 1 {\Gamma\left(\frac n 2 \right)} \left(\frac 1 2 \right)^{\frac n 2}
 \int_0^\infty t^{\frac 1 2}t^{\frac n 2 - 1}\exp\left(-\frac 1 2 t \right)dt\\
 &= \frac {\left(\frac 1 2 \right)^{\frac n 2}} {\Gamma\left(\frac n 2\right)} 
\frac{\Gamma\left(\frac{n+1}2 \right)}{\left(\frac 12\right)^{\frac{n+1}{2}}}
\int_0^\infty \frac 1 {\Gamma\left(\frac{n+1}2 \right)}
                      \left(\frac 12\right)^{\frac{n+1}{2}}
                      t^{\frac {n+1} 2 - 1}\exp\left(-\frac 1 2 t \right)dt\\
&=\frac {\left(\frac 1 2 \right)^{\frac n 2}} {\Gamma\left(\frac n 2\right)} 
\frac{\Gamma\left(\frac{n+1}2 \right)}{\left(\frac 12\right)^{\frac{n+1}{2}}}
\int_0^\infty f_{\frac{n+1}2, \frac 1 2}(t) dt
\overset{\substack{f_{\frac{n+1}2, \frac 1 2}\\\text{Dichte}}}= \frac {\left(\frac 1 2 \right)^{\frac n 2}} {\Gamma\left(\frac n 2\right)} 
\frac{\Gamma\left(\frac{n+1}2 \right)}{\left(\frac 12\right)^{\frac{n+1}{2}}}\\
&= \frac 1{\left(\frac 1 2 \right)^{\frac 1 2}} 
   \frac{\Gamma\left(\frac{n+1}2 \right)}
        {\Gamma\left(\frac n 2\right)}
= \sqrt 2\frac{\Gamma\left(\frac{n+1}2\right)}{\Gamma\left(\frac n 2 \right)}
\text. 
\end{align*}
\end{proof}
\parag
Der soeben bestimmte Erwartungswert lässt sich trotz seiner schlecht explizit
berechenbaren Form überraschend gut abschätzen.
Diese Abschätzung ist Inhalt des folgenden Lemmas.

\begin{lemma}{Absch�tzung eines Gammafunktionsquotienten}
             {AbschaetzungEinesGammafunktionsquotienten}
Seien $\Gamma: ]0, \infty[ \rightarrow \R, 
             t \mapsto \int_0^\infty x^{t-1}\exp(-x)dx$ und
$a : \N \rightarrow \R, 
  n \mapsto \sqrt 2 \frac{\Gamma\left(\frac{n+1}2 \right)}
{\Gamma\left(\frac n 2 \right)}$
(nach \refclaim{ElementareEigenschaftenDerGammafunktion}
{ElementareEigenschaftenDerGammafunktion1} ist $a$ wohldefiniert).\\
Dann gilt:
\begin{enumerate}[(1)]
\item\label{AbschaetzungEinesGammafunktionsquotienten1} 
 $\forall \, n \in \N: a_n\cdot a_{n+1} = n$
\item\label{AbschaetzungEinesGammafunktionsquotienten2}
 $\forall \, n \in \N: a_n \le \sqrt n$
\item\label{AbschaetzungEinesGammafunktionsquotienten3}
 $\forall \, n \in \N: a_n \ge \frac{n}{\sqrt{n+1}}$
 \item\label{AbschaetzungEinesGammafunktionsquotienten4}
 $\forall\, n \in \N: \sqrt{\frac 2 \pi} \sqrt n \le a_n$
\end{enumerate}
\end{lemma}
\begin{proof}
Nach \refclaim{ElementareEigenschaftenDerGammafunktion}
{ElementareEigenschaftenDerGammafunktion1} ist f�r alle $n \in \N$: 
$\Gamma\left(\frac n 2 \right) > 0$ und $\Gamma(\frac{n+1}2)> 0$.\\
Daher ist f�r alle $n \in \N$: $a_n > 0$.\\
(1) Sei also $n \in \N$. Dann gilt:
\begin{align*}
a_n a_{n+1} &= 
\sqrt 2 \frac{\Gamma\left(\frac{n+1}2 \right)} {\Gamma\left(\frac n 2 \right)}
\cdot
\sqrt 2\frac{\Gamma\left(\frac{n+2}2 \right)}{\Gamma\left(\frac{n+1}2\right)}
= 2 \frac{\Gamma\left(\frac{n+2}2 \right)}{\Gamma\left(\frac n 2 \right)}
= 2 \frac{\Gamma\left(\frac n 2 + 1\right)}{\Gamma\left(\frac n 2 \right)}
\overset{\refclaim{ElementareEigenschaftenDerGammafunktion}
{ElementareEigenschaftenDerGammafunktion2}}= 
2\frac n2 \frac{\Gamma\left(\frac n2\right)}{\Gamma\left(\frac n2\right)}\\
&=n\text.
\end{align*}
(2)/(4) Sei $n \in \N$.\\
Dann gibt es nach  \refclaim{ExistenzFVerteilterZufallsgr�ssen}
{ExistenzFVerteilterZufallsgr�ssen2}  einen Wahrscheinlichkeitsraum 
$(\Omega', \mfa', \Prim')$ und ein
$\gamma: \Omega' \rightarrow \R^n$ unabh�ngig und $N(0, 1)$--verteilt.\\
Nach \ref{ErwartungswertDer2NormEinesStandardnormalverteiltenProzesses} gilt
dann $\E\left(\norm\cdot_{2, n} \circ \gamma \right) = a(n)$.\\
Da $\gamma$ unabh�ngig und $N(0,1)$--verteilt ist, ist es auch
unkorreliert, zentriert und $L_2-$normiert, sodass mit
\refclaim{ElementareEigenschaftenQuadratintegrierbarerAbbildungen}
{ElementareEigenschaftenQuadratintegrierbarerAbbildungen4} auch
$a(n) = \E\left(\norm\cdot_{2, n} \circ \gamma \right) \le \sqrt n$ gilt.\\
Au�erdem gilt nach \refclaim{AequivalenzkonstantenFuerPNormen}
{AequivalenzkonstantenFuerPNormen1}: $\frac 1 {\sqrt n} \norm\cdot_{1,n}
\le \norm\cdot_{2, n}$.\hfill$(\ast)$\\
Also gilt $0\le\frac 1{\sqrt n}\norm\cdot_{1,n}\circ\gamma\le\norm\cdot_{2,n}
\circ \gamma$ und nach
\ref{ErwartungswertDer2NormEinesStandardnormalverteiltenProzesses} ist
$\norm\cdot_{2, n} \circ \gamma$
integrierbar, also nach \refclaim{Integrierbarkeitskriterien}
{Integrierbarkeitskriterien1} auch $\norm\cdot_{1,n} \circ \gamma$.\\
Ferner ist $\E(\abs{\gamma_1}) = \sqrt{\nicefrac 2 \pi}$ und da die
Komponentenfunktionen von $\gamma$ jeweils
$N(0,1)$--verteilt sind, gilt: $\E(\norm\cdot_{1,n} \circ \gamma) 
= \sqrt{\nicefrac 2 \pi}\cdot n$.\\
Somit gilt nach \refclaim{Integrierbarkeitskriterien}
{Integrierbarkeitskriterien1}
\begin{align*}\E\left(\norm\cdot_{2,n} \circ \gamma \right)
 \overset{(\ast)}\ge \E \left(\frac 1 {\sqrt n} \norm\cdot_{1,n}
\circ \gamma \right) = \frac 1 {\sqrt n} \E \left(\norm\cdot_{1,n}\circ \gamma
\right)
= \frac 1 {\sqrt n} \sqrt{\frac 2 \pi} n = \sqrt{\frac 2 \pi} \sqrt n\text.
\end{align*}
(3) Sei $n \in \N$. 
Dann gilt: 
$a_n \overset{\eqref{AbschaetzungEinesGammafunktionsquotienten1}}= 
\frac n {a_{n+1}}
\overset{\eqref{AbschaetzungEinesGammafunktionsquotienten2}}\ge 
\frac n {\sqrt{n+1}}$.\\
\end{proof}
\parag
Wir formulieren nun einige Sätze, welche uns später dabei nützlich sein werden,
den Gauß'schen Min--Max--Satz auf bestimmte unendliche Familien zu
verallgemeinern.
\begin{lemma}{Erwartungswert des Supremums integrierbarer Funktionen}
            {EndlicheDarstellungDesSupremumsIntegrierbarerFunktionen}
Seien $I$ eine nichtleere, abzählbare Menge und 
$\folge X i I \in L_1(\Omega, \mfa, \Prim, \R)^I$.\\
Es gelte:
\begin{enumerate}[(a)]
\item für alle 
     $\omega \in \Omega$ ist $\menge{X_i(\omega)}{i \in I}$
     nach oben beschränkt,
\item $\Sup{i \in I} X_i \in L_1(\Omega, \mfa, \Prim, \R)$
(wohldefiniert wegen (a) und $I\neq \emptyset$).
\end{enumerate}
Dann gilt:
\begin{enumerate}[(1)]
\item\label{EndlicheDarstellungDesSupremumsIntegrierbarerFunktionen1}
 $\menge{\E\left( \Sup{i \in E} X_i \right)}{E \in \Pot_{ne}(I)} $
    ist nichtleer und nach oben beschränkt.
\item\label{EndlicheDarstellungDesSupremumsIntegrierbarerFunktionen2}
 $\sup\menge{\E\left( \Sup{i \in E} X_i\right)}{E \in \Pot_{ne}(I)} 
    = \E\left(\Sup{i \in I} X_i \right)$.
\end{enumerate}
\end{lemma}
\begin{proof}
Da $I$ nichtleer ist, gibt es ein $i_0 \in I$. Daher ist
$\meng{i_0}\in\Pot_{ne}(I)$, sodass die in (1) angegebene Menge nichtleer
ist.\\
Sei $A \in \Pot_{ne}(I)$. Dann ist auch $A \neq \emptyset$, sodass ein
$a \in A$ existiert.\\
Wegen
\[X_a \le \sup_{i \in A}X_i \le \sup_{i \in I}X_i =: f,\]
der Integrierbarkeit von $X_a$ und nach (b) auch der von $f$
ist nach \refclaim{Integrierbarkeitskriterien}{Integrierbarkeitskriterien2}
\begin{align}\sup_{i \in A}X_i \text{ integrierbar mit }
\E\left(\sup_{i \in A}X_i \right) \le  \E\left(f \right)\text. \tag{$\ast$}
\end{align}
Somit ist die in (1) angegebene Menge auch nach oben beschränkt.\\
Falls nun $I$ endlich ist, so ist wegen $I \in \Pot_{ne}(I)$ für (2) nichts zu
zeigen.\\
Sei also $I$ abzählbar unendlich.
Dann existiert $\iota: \N \rightarrow I$ bijektiv.\\
Setze
\[I_n \coloneqq  \iota(\haken n) \text{ und } f_n\coloneqq  \sup_{i \in I_n}X_i
\text{ für alle } n \in \N\text.\]
Dann ist $\folge I n \N$ eine aufsteigende Mengenfolge mit
$\bigcup\menge{I_n}{n \in \N} = I$ und daher $\folge f n \N$ eine monoton
wachsende Folge nach $(\ast)$ integrierbarer Funktionen, die punktweise gegen
$f$ konvergiert.\\
Somit sind die Voraussetzungen von \ref{MonotoneKonvergenzImIntegrierbarenFall}
erfüllt und es gilt:
\begin{align*}
\E(f)& \overset{\ref{MonotoneKonvergenzImIntegrierbarenFall}}= 
\Limes n \infty \E(f_n) = \sup_{n \in \N}\E(f_n)
= \sup\menge{\E\left(\sup_{i \in I_n}X_i \right)}{n \in \N}\\
&\le\sup\menge{\E\left(\sup_{i \in E}X_i \right)}{E \in \Pot_{ne}(I)}
\overset{(\ast)} \le \E(f)\text.
\end{align*}
\end{proof}

\parag
Wegen der Eigenschaft $\inf(M) = -\sup (-M)$ für alle nichtleeren nach unten
beschränkten Mengen $M \subseteq \R$, kann man aus dem obigen Satz ein
elementares Korollar folgern.

\begin{korollar}{Erwartungswert des Infimums integrierbarer Funktionen}
                {EndlicheDarstellungDesInfimumsIntegrierbarerFunktionen}
Seien $I$ eine nichtleere, abzählbare Menge und 
$\folge X i I \in L_1(\Omega, \mfa,\Prim,\R)^I$.\\
Es gelte:
\begin{enumerate}[(a)]
 \item für alle $\omega \in \Omega$ ist $\menge{X_i(\omega)}{i \in I}$
       nach unten beschränkt,
 \item $\Inf{i\in I} X_i \in L_1(\Omega, \mfa,\Prim,\R)$ (wohldefiniert
       wegen (a) und $I \neq \emptyset$).
\end{enumerate}
Dann gilt:
\begin{enumerate}[(1)]
 \item\label{EndlicheDarstellungDesInfimumsIntegrierbarerFunktionen1}
      $\menge{\E\left(\Inf{i\in E}X_i\right)}{E\in\Pot_{ne}(I)}$ ist nichtleer
      und nach unten beschränkt.
 \item\label{EndlicheDarstellungDesInfimumsIntegrierbarerFunktionen2}
      $\inf\menge{\E\left( \Inf{i\in E} X_i\right)}{E\in\Pot_{ne}(I)} 
      = \E\left(\Inf{i\in I}X_i \right)$.
\end{enumerate}
\end{korollar}
\begin{proof}
Für alle $i \in I$ sei $Y_i \coloneqq  -X_i$.\\
Für alle $i \in I$ ist dann $Y_i$ als Vielfaches einer integrierbaren Funktion
auch integrierbar.\\
Sei $\omega \in \Omega$. 
Dann ist
$ \menge{Y_i(\omega)}{i \in I} = \menge{-X_i(\omega)}{i \in I}
 = -\menge{X_i(\omega)}{i \in I}$.
Daher ist $\menge{Y_i(\omega)}{i \in I}$ wegen $(a)$ nach oben beschränkt.\\
Sei $A\in\Pot_{ne}(I)$. Wegen $A \neq \emptyset$, gibt es ein $a \in A$, sodass
gilt:
\[X_a \ge \inf_{i \in A}X_i \ge \inf_{i \in I}X_i\]
und wegen (b) bedeutet das nach \refclaim{Integrierbarkeitskriterien}
{Integrierbarkeitskriterien2}
\[ \inf_{i \in A}X_i \text{ ist integrierbar, also auch }
  -\inf_{i \in A}X_i = \sup_{i \in A}Y_i\text.\]
Somit gilt:
$\forall\,E\in\Pot_{ne}(I)\cup\meng{I}:\Sup{i \in A}Y_i = -\Inf{i\in A}X_i
\text{ integrierbar}$. \hfill $(\ast)$\\
Also sind die Voraussetzungen von
\ref{EndlicheDarstellungDesSupremumsIntegrierbarerFunktionen} erfüllt, sodass
gilt:
\[\menge{\E\left(\inf_{i\in E}X_i\right)}{E\in\Pot_{ne}(I)}
 \overset{(\ast)}= -\menge{\E\left(\sup_{i\in E}Y_i\right)}{E\in\Pot_{ne}(I)}\]
ist nach \refclaim{EndlicheDarstellungDesSupremumsIntegrierbarerFunktionen}
{EndlicheDarstellungDesSupremumsIntegrierbarerFunktionen1} nach unten
beschränkt und
\begin{align*}
&\, \E\left( \inf_{i \in I} X_i\right)  \overset{(\ast)}{=}
 - \E\left( \sup_{i \in I} Y_i \right)
 \overset{\refclaim{EndlicheDarstellungDesSupremumsIntegrierbarerFunktionen}
{EndlicheDarstellungDesSupremumsIntegrierbarerFunktionen2}}=
- \sup\menge{\E\left(\sup_{i \in E}Y_i\right)}{E\in \Pot_{ne}(I)}\\
\overset{(\ast)}= &\,-\sup\menge{-\E\left(\inf_{i\in E}X_i\right)}{E\in\Pot_{ne}(I)}
 = \inf\menge{\E\left(\inf_{i\in E}X_i\right)}{E\in\Pot_{ne}(I)}\text.
\end{align*}
\end{proof}

\begin{satz}{Fortsetzung einer endlichen Erwartungswertungleichung (I)}
            {AbzaehlbareDarstellungUndInfimumVonSuprema}
Seien $I, J$ nichtleere, abz�hlbare Mengen,
$X, Y: I\times J \rightarrow \mcm(\Omega, \mfa, \R)$, sowie
$\mci := \Pot_{ne}(I)$ und $\mcj := \Pot_{ne}(J)$.\\
Es gelte:
\begin{enumerate}[(a)]
\item $\forall\,\omega\in\Omega\,\forall\, Z \in \meng{X, Y}: 
      \menge{Z(i, j)(\omega)}{(i, j)\in I\times J}$
      beschr�nkt.
\item $\forall\, A \in \mci \cup \meng I \, \forall\, B \in \mcj \cup \meng J:
     \Inf{i \in A}\Sup{j \in B} X(i, j) = \Sup{j \in B}\Inf{i \in A} X(i,j)$.
\item $\forall A \in \mci \cup \meng I  \forall B \in \mcj \cup \meng J:
     \Inf{i \in A}\Sup{j \in B} X(i, j), \Inf{i \in A}\Sup{j \in B} Y(i, j)
     \in L_1(\Omega, \mfa, \Prim, \R)$.   
\item $\forall\, A \in \mci \, \forall\, B \in \mcj:
\E\bigg(\underset{i \in A}\inf\,\underset{j \in B}\sup\, X(i, j) \bigg)
\le \E\bigg(\underset{i \in A}\inf\,\underset{j \in B}\sup\, Y(i, j) \bigg)$.\\
(Alle genannten Abbildungen sind wegen $(a)$ und $I \neq \emptyset \neq J$
 wohldefiniert.)
\end{enumerate}
Dann gilt \[\E\left(\inf_{i \in I}\,\sup_{j \in J}\, X(i, j) \right)
\le \E\left(\inf_{i \in I}\,\sup_{j \in J}\, Y(i, j) \right)\text.\]
\end{satz}
\begin{proof}
Definiere $W_j := \Inf{i\in I} X(i, j)$ f�r alle $j \in J$.\\
Wegen (a) ist $\menge{W_j(\omega)}{j\in J}$ f�r alle $\omega \in \Omega$ nach
oben beschr�nkt.\\
F�r jedes $k\in J$ ist $W_k$ und
\begin{align}
\sup_{j \in J}W_j = \sup_{j \in J}\inf_{i \in I}X(i, j) \overset{(b)}= 
\inf_{i \in I}\sup_{j \in J}X(i, j)\quad \text{ nach (c) integrierbar.}
\label{feeeieq:1}
\end{align}
Damit sind die Voraussetzungen von
\ref{EndlicheDarstellungDesSupremumsIntegrierbarerFunktionen} f�r
$\folge W j J$ erf�llt. Daher ist
\[\menge{\E\left(\inf_{i\in I}\sup_{j \in B}X(i,j) \right)}{B \in \mcj}
\overset{\eqref{feeeieq:1}}= 
 \menge{\E\left(\sup_{j \in B}W_j \right)}{B \in \mcj}\]
nach \refclaim{EndlicheDarstellungDesSupremumsIntegrierbarerFunktionen}
{EndlicheDarstellungDesSupremumsIntegrierbarerFunktionen1} nichtleer und nach
oben beschr�nkt und es gilt:
\begin{align}
& \sup\menge{\E\left(\inf_{i \in I}\sup_{j \in B}X(i, j) \right)}{B \in \mcj}
\notag\\
 =&\,  \sup\menge{\E\left(\sup_{j \in B}W_j \right)}{B \in \mcj}
 \overset{\refclaim{EndlicheDarstellungDesSupremumsIntegrierbarerFunktionen}
 {EndlicheDarstellungDesSupremumsIntegrierbarerFunktionen2}}=
\E\left(\sup_{j \in J}W_j \right)\notag\\ 
\overset{(b)}= &\,\E\left(\inf_{i\in I}\sup_{j\in J}X(i,j)\right)\text.
\label{feeeieq:2}
\end{align}
Seien nun $Z_i := \Sup{j \in J}Y(i, j)$ f�r alle $i \in I$.\\
Wegen (a) ist $\menge{Z_i(\omega)}{i\in I}$ nach unten beschr�nkt f�r alle
$\omega \in \Omega$.\\
F�r alle $l\in I$ ist $Z_l$ und
\[\inf_{i \in I}Z_i = \inf_{i \in I}\sup_{j \in J}Y(i, j) \quad
\text{ nach (c) integrierbar.} \]
Damit sind auch die Vorausssetzungen von
\ref{EndlicheDarstellungDesInfimumsIntegrierbarerFunktionen} erf�llt. Also ist
\[\menge{\E\left(\inf_{i\in A}\sup_{j\in J}Y(i,j)\right)}{A\in\mci}
  =\menge{\E\left(\inf_{i\in A}Z_i\right)}{A\in\mci}\]
nach \refclaim{EndlicheDarstellungDesInfimumsIntegrierbarerFunktionen}
{EndlicheDarstellungDesInfimumsIntegrierbarerFunktionen1} nichtleer und nach
oben beschr�nkt und es gilt:
\begin{align}
& \inf\menge{\E\left(\inf_{i\in A}\sup_{j\in J}Y(i,j)\right)}{A\in\mci}
= \inf\menge{\E\left(\inf_{i \in A}Z_i \right)}{A \in \mci}\notag\\
\overset{\refclaim{EndlicheDarstellungDesInfimumsIntegrierbarerFunktionen}
{EndlicheDarstellungDesInfimumsIntegrierbarerFunktionen2}}=&
\E\left(\inf_{i\in I}Z_i\right)=\E\left(\inf_{i\in I}\sup_{j\in J}Y(i,j)\right)
\text.\label{feeeieq:3}
\end{align}
Seien nun $C \in \mci$ und $D \in \mcj$. Dann gilt
\[ \inf_{i \in I}\sup_{j \in D}X(i, j) \le \inf_{i \in C}\sup_{j \in D}X(i,j)
\quad \text{ und }\quad  
\inf_{i\in C}\sup_{j\in D}Y(i,j)\le \inf_{i \in C}\sup_{j \in J}Y(i, j)\]
und diese vier Abbildungen sind nach (c) integrierbar.\\
Daher gelten diese Ungleichungen auch im Erwartungswert, sodass gilt:
\begin{align*}
\E\left(\inf_{i \in I}\sup_{j \in D}X(i, j) \right) & \le
\E\left(\inf_{i \in C}\sup_{j \in D}X(i, j) \right) \overset{(d)}\le
\E\left(\inf_{i \in C}\sup_{j \in D}Y(i, j) \right)\\
&\le \E\left(\inf_{i \in C}\sup_{j \in J}Y(i, j) \right)\text.
\end{align*}
Somit erh�lt man
\begin{align}
&
\sup\menge{\E\left(\inf_{i\in I}\sup_{j\in B}X(i,j)\right)}{B \in \mcj}\notag\\
\le&\, \inf\menge{\E\left(\inf_{i\in A}\sup_{j\in J}Y(i,j)\right)}{A\in\mci}
 \text. \label{feeeieq:4}\end{align}
Dies liefert schlie�lich
\begin{align*}
& \E\left(\inf_{i\in I}\sup_{j\in J}X(i,j)\right) \overset{\eqref{feeeieq:2}}=
\sup\menge{\E\left(\inf_{i \in I}\sup_{j \in B}X(i, j) \right)}{B \in \mcj}\\
\overset{\eqref{feeeieq:4}}\le &
\inf\menge{\E\left(\inf_{i\in A}\sup_{j\in J}Y(i,j)\right)}{A\in\mci}
\overset{\eqref{feeeieq:3}}= \E\left(\inf_{i\in I}\sup_{j\in J}Y(i,j)\right)
\end{align*}
und damit die Behauptung.
\end{proof}

\begin{satz}{Fortsetzung einer endlichen Erwartungswertungleichung (II)}
            {AbzaehlbareDarstellungUndSupremumVonSuprema}
Seien $I, J$ nichtleere, abzählbare Mengen,
$X, Y: I\times J \rightarrow \mcm(\Omega, \mfa, \R)$, sowie
$\mci := \Pot_{ne}(I)$ und $\mcj := \Pot_{ne}(J)$.\\
Es gelte:
\begin{enumerate}[(a)]
\item $\forall\,\omega\in\Omega\,\forall\, Z\in \meng{X, Y}: 
\menge{Z(i, j)(\omega)}{(i, j)\in I\times J}$
      nach oben beschränkt.
\item $\forall A \in \mci \cup \meng I \forall B \in \mcj \cup \meng J
\forall Z\in \meng{X, Y}:
     \underset{i \in A}\sup\,\underset{j \in B}\sup\, Z(i, j)
     \in L_1(\Omega, \mfa, \Prim, \R)$.
\item $\forall\, A \in \mci \, \forall\, B \in \mcj: 
\E\bigg(\underset{i \in A}\sup\,\underset{j \in B}\sup\, Y(i, j) \bigg)
\le \E\bigg(\underset{i \in A}\sup\,\underset{j \in B}\sup\, X(i, j) \bigg)$.\\
 (Alle genannten Abbildungen sind wegen $(a)$ und
$I \neq \emptyset \neq J$wohldefiniert.)
\end{enumerate}
Dann gilt
\[\E\left(\sup_{i \in I}\sup_{j \in J}Y(i, j) \right) \le
  \E\left(\sup_{i \in I}\sup_{j \in J}X(i, j) \right)\text.\]
\end{satz}
\begin{proof}
Seien $\mcp:= \Pot_{ne}(I \times J)$ und $(i_0, j_0) \in I \times J$.\\
Dann sind $Y(i_0, j_0)$ und
\[\sup_{(i, j) \in I \times J}Y(i, j)= \sup_{i \in I} \sup_{j \in J}Y(i, j)
\text{ nach (b) integrierbar.}\]
Da $I$ und $J$ abzählbar sind, ist auch $I \times J$ abzählbar.
Zusammen mit (a) sind damit die Voraussetzungen von
\ref{EndlicheDarstellungDesSupremumsIntegrierbarerFunktionen} für die Familie
$(Y(i,j))_{(i,j) \in I \times J}$ erfüllt. Also ist
\[M_1:=\menge{\E\left(\sup_{(i, j) \in P}Y(i,j) \right)}{P \in \mcp}\]
nach \refclaim{EndlicheDarstellungDesSupremumsIntegrierbarerFunktionen}
{EndlicheDarstellungDesSupremumsIntegrierbarerFunktionen1} nichtleer und nach
oben beschränkt und es gilt
\begin{align}\sup M_1
\overset{\refclaim{EndlicheDarstellungDesSupremumsIntegrierbarerFunktionen}
{EndlicheDarstellungDesSupremumsIntegrierbarerFunktionen2}}= 
\E\bigg(\sup_{(i, j) \in I \times J}Y(i, j) \bigg)
= \E\left(\sup_{i \in I}\sup_{j \in J}Y(i,j)\right)\text. \label{feeeiieq:1}
\end{align}
Sei nun 
\[M_2 := \menge{\E\left(\sup_{i \in A}\sup_{j \in B}Y(i,j) \right)}
               {A \in \mci, B \in \mcj}\text.\]
Offenbar ist $M_2 \subseteq M_1$ und wegen $I \neq \emptyset\neq J$ auch
nichtleer.
Da $M_1$ nach oben beschränkt ist, ist es auch
$M_2$ und es gilt $\sup M_2 \le \sup M_1$.\\
Sei nun $Q \in \mcp$. Dann gibt es ein
$A_Q \in \mci$ und ein $B_Q \in \mcj$ mit 
$Q \subseteq A_Q \times B_Q$.\\
Da $Q \neq \emptyset$ existiert ein $(i', j') \in Q$, sodass gilt:
\[Y(i',j')\le\sup_{(i,j)\in Q}Y(i,j)\le\sup_{(i,j)\in A_Q\times B_Q}Y(i,j)
=\sup_{i\in A_Q}\sup_{j\in B_Q}Y(i,j)\text.\]
Nach (b) sind die erste und die letzte Abbildung integrierbar, also ist
es nach \refclaim{Integrierbarkeitskriterien}{Integrierbarkeitskriterien2}
auch die zweite Abbildung. Somit gilt:
\[\E\left(\sup_{(i, j) \in Q}Y(i, j) \right) \le 
\E\left(\sup_{(i, j) \in A_Q \times B_Q}Y(i, j) \right)
= \E\left(\sup_{i \in A_Q}\sup_{j \in B_Q}Y(i, j) \right) \in M_2\text.\]
Erhalte: $\forall \, x \in M_1\, \exists \, y \in M_2: x \le y$.\\
Damit ist offensichtlich:
\begin{align}
 \sup M_1\le\sup M_2\text{, also }
\sup M_1=\sup M_2\text. \label{feeeiieq:2}
\end{align}
Seien nun $C \in \mci$ und $D \in \mcj$. Dann gilt
\[\sup_{i\in C}\sup_{j\in D}X(i,j)\le\sup_{i\in I}\sup_{j \in J}X(i,j)\text.\]
Gemäß (b) sind beide Abbildungen integrierbar und daher ist auch
\[\E\left(\sup_{i \in C} \sup_{j \in D}X(i, j) \right)
\le \E\left(\sup_{i \in I} \sup_{j \in J}X(i, j) \right)\text.\]
Daher gilt
\[\E\left(\sup_{i \in C} \sup_{j \in D}Y(i, j) \right)
\overset{(c)}\le\E\left(\sup_{i \in C} \sup_{j \in D}X(i, j) \right)
\le \E\left(\sup_{i \in I} \sup_{j \in J}X(i, j) \right)\]
und damit
\begin{align}
\sup\menge{\E\left(\sup_{i \in A} \sup_{j \in B}Y(i, j) \right)}
            {A \in \mci, B \in \mcj} \le
            \E\left(\sup_{i \in I} \sup_{j \in J}X(i, j) \right)\text.
            \label{feeeiieq:3}
\end{align}
Schließlich gilt also
\begin{align*}
&\,\E\left(\sup_{i \in I}\sup_{j \in J}Y(i,j)\right)
\overset{\eqref{feeeiieq:1}}=
\sup M_1 \overset{\eqref{feeeiieq:2}}= \sup M_2\\
= &\, \sup \menge{\E\left(\sup_{i \in A}\sup_{j \in B}Y(i,j) \right)}
        {A \in \mci, B \in \mcj}
 \overset{\eqref{feeeiieq:3}}\le
 \E\left(\sup_{i\in I}\sup_{j \in J}X(i, j) \right)\text,
\end{align*}
was gerade behauptet wurde.
\end{proof}
\parag
Wir werden an späterer Stelle Ordnungsstatistiken von
standardnormalverteilten Zufallsgrößen benötigen, wobei uns insbesondere
deren Erwartungswerte interessieren werden.

\begin{definition}{Ordnungsstatistik}{Ordnungsstatistik}
Seien $n\in \N$, $k \in \haken n$ und $\gamma: \Omega \rightarrow \R^n$.\\
Dann definieren wir
\[ \gamma^\sigma_k\coloneqq  \mu_{k,n}\circ \gamma \]
und nennen diese Abbildung die \textbf{$k$--te Ordnungsstatistik von $\gamma$}.
\\Ferner definieren wir
\[ \gamma_k^\ast \coloneqq  \mu_{k,n} \circ \pfeil {\abs\cdot} n \circ \gamma\]
und nennen diese Abbildung die \textbf{absolute $k$--te Ordnungsstatistik von
$\gamma$}.\\
Schließlich setzen wir:
\[ \gamma ^\sigma \coloneqq  \left(\gamma^\sigma_i \right)_{i=1}^n\quad
\text{ und }\quad \gamma^\ast \coloneqq  \left(\gamma^\ast_i \right)_{i=1}^n\]
\end{definition}
\parag
Wir notieren einige Eigenschaften von Ordnungsstatistiken, wobei uns
insbesondere deren Verteilung bzw. Dichte interessieren wird.

\begin{lemma}{Eigenschaften von Ordnungsstatistiken}
             {EigenschaftenVonOrdnungsstatistiken}
Seien $n \in \N$, $k \in \haken n$ und $\gamma \in \mcm_n(\Omega,\mfa,\R^n)$.\\
Dann gilt:
\begin{enumerate}[(1)]
	\item\label{EigenschaftenVonOrdnungsstatistiken1}
	$\gamma^\sigma_k, \gamma^\ast_k \in \mcm(\Omega, \mfa, \R)$.
	\item\label{EigenschaftenVonOrdnungsstatistiken2}
	Ist $\Bild \gamma \subseteq [0, \infty[^n$, so ist
	$\gamma_k^\ast = \gamma^\sigma_k$. Ist $\gamma$ unabh�ngig, so ist auch
	\[\gamma_{\abs{\cdot}}: \Omega \rightarrow \R^n, \omega \mapsto
	(\abs{\gamma_i(\omega)})_{i=1}^n\]
	unabh�ngig.
	\item\label{EigenschaftenVonOrdnungsstatistiken3}
	Ist $\gamma$ unabh�ngig und identisch verteilt,
	so gilt:
	\[\forall\, t \in \R: \Prim(\gamma^\sigma_k > t) 
	=\sum_{l=0}^{n-k}\binom n l
	\Prim(\gamma_1\le t)^l(1-\Prim(\gamma_1\le t))^{n-l}\text.\]
	\item\label{EigenschaftenVonOrdnungsstatistiken4}
	Ist $\gamma$ unabh�ngig und identisch verteilt, sodass $\gamma_1$ eine
	stetige Dichte $f$ besitzt, so ist 
	$\Prim(\gamma^\sigma_k \le \cdot)$ stetig differenzierbar und es gilt
	f�r alle $t_0 \in \R$:
	\[\left(\frac d {dt} \Prim(\gamma^\sigma_k\le \cdot) \right)(t_0) 
	= n\binom{n-1}{n-k}
	\Prim(\gamma_1 \le t_0)^{n-k}(1-\Prim(\gamma_1 \le t_0))^{k-1}f(t_0)\text.\]
\end{enumerate}
\end{lemma}
\begin{proof}
(1) Nach
\refclaim{EigenschaftenDesKtenMaximums}{EigenschaftenDesKtenMaximums2c}
ist $\mu_{k,n}$ stetig, also $\mfb_n-\mfb$--messbar und daher ist
$\gamma^\sigma_k=\mu_{k,n} \circ\gamma$ selbst $\mfa-\mfb$--messbar als
Komposition messbarer Funktionen.\\
Die Abbildung $\pfeil {\abs\cdot}n$ ist offenbar stetig und daher ist
$\gamma_k^\ast = \mu_{k,n} \circ\pfeil {\abs\cdot}n\circ \gamma$ als
Komposition stetiger und messbarer Funktionen $\mfa-\mfb$--messbar.\\
(2) Die erste Aussage ist direkt aus der Definition ersichtlich.\\
F�r die zweite seien $(A_i)_{i \in \haken n} \in \mfb^n$. Da die Abbildung
$f:\R \rightarrow \R, x \mapsto -x$ stetig und damit messbar ist, ist
$-A_i = f\inv(A_i) \in \mfb$ f�r alle $i \in \haken n$. Also gilt:
\begin{align*}
&\,\Prim\left(\bigcap_{i \in \haken n} \meng{\abs{\gamma_i} \in A_i} \right)
= \Prim\left(\bigcap_{i \in \haken n} 
\meng{\gamma_i \in A_i \vee \gamma_i \in -A_i} \right)\\
=&\,\Prim\left(\bigcap_{i\in\haken n}\meng{\gamma_i\in(A_i \cup -A_i)} \right)
\overset{\gamma \text{ unabh.}}=
\prod_{i \in \haken n} \Prim\left(\meng{\gamma_i \in (A_i \cup -A_i )}\right)\\
=&\prod_{i \in \haken n} \Prim\left(\meng{\abs{\gamma_i} \in A_i)}\right)\text.
\end{align*}
(3)/(4) Siehe \cite{AB}, dort 2.2.13 f�r die Verteilung und 2.2.2 f�r die
Dichte. Dabei ist zu beachten, dass in dem genannten Buch Ordnungsstatistiken
andersherum definiert sind, d. h. aufsteigend.
\end{proof}

Um konkrete Abschätzungen der Erwartungswerte gewisser absoluter
Ordnungsstatistiken anzugeben, kann man deren Verteilung bzw. Dichte benutzen.
Es wird sich als nützlich herausstellen, die Abbildung
$\Prim(\abs{\gamma} > \cdot)$ statt der Verteilung zu betrachten.
Im Falle
einer $N(0, 1)$--verteilten Zufallsgröße $\gamma$ hat diese Abbildung
besonders gute Eigenschaften, die wir als Nächstes notieren.
Die erste Ungleichung der vierten Aussage beruht dabei maßgeblich auf \cite{RW}, die
Anwendung auf die komplementäre Verteilungsfunktion ist elementar.

\begin{satz}
{Eigenschaften der komplementären \textsl{N}(0,1)--Verteilungsfunktion}
{EigenschaftenDerKomplementaerenVerteilungsfunktionVonN01}
Seien $\gamma: \Omega \rightarrow \R$ $N(0,1)$--verteilt,
$u:\R_{\ge 0} \rightarrow [0,1], t \mapsto \Prim(\abs \gamma > t)$ und
schließlich
$\Phi: \R \rightarrow\!\  ]0,1[$ die Verteilungsfunktion von $N(0,1)$.\\
Dann gilt:
\begin{enumerate}[(1)]
\item \label{EigenschaftenDerKomplementaerenVerteilungsfunktionVonN011}
$u = 2(1 - \restrict{\Phi}{[0, \infty[}) = \sqrt{\frac 2 \pi}\int_\cdot ^\infty
\exp\left(-\frac{s^2}2 \right)ds$, $\Bild u =\!\  ]0,1]$, $u$ ist injektiv
und besitzt daher ein Inverses 
$u\inv:\!\ ]0,1] \rightarrow [0, \infty[$ mit $\Phi \circ u\inv 
= 1- \frac \cdot 2$.
\item \label{EigenschaftenDerKomplementaerenVerteilungsfunktionVonN012}
$u$ ist stetig differenzierbar auf $\R_{>0}$ mit 
\[\forall\,t \in \R_{>0}: u'(t) = 
-\sqrt{\frac 2 \pi}\exp\left(- \frac{t^2}2 \right)\text.\]
\item \label{EigenschaftenDerKomplementaerenVerteilungsfunktionVonN013}
Es gibt $(c_j)_{j \in \N_0}\in [0, \infty[^{\N_0}$ mit
\[\forall \, t \in\!\  ]0,1]: u\inv(t) = \sum_{j=0}^\infty
c_j(1-t)^{2j+1}\text.\]
\item \label{EigenschaftenDerKomplementaerenVerteilungsfunktionVonN014}
$\forall\, t \in \R_{>0}:
\frac{\sqrt{2\pi}}{(\pi-1)t + \sqrt{t^2+2\pi}}\exp\left(-\frac{t^2}2 \right)
\le u(t) \le \sqrt{\frac 2 \pi}\frac 1 t \exp\left(-\frac {t^2}2 \right)$.
\item \label{EigenschaftenDerKomplementaerenVerteilungsfunktionVonN015}
$\forall\,t\in [1,\infty[: u\inv\left(\frac 1 {t+1} \right)\ge \sqrt{\log(t)}$.
\item \label{EigenschaftenDerKomplementaerenVerteilungsfunktionVonN016}
$\forall\, t \in\!\  ]0, \frac 2 {13}]: 
 u \inv(t) \ge \sqrt{\log\left(1+\frac 1 t\right)}$.
\item \label{EigenschaftenDerKomplementaerenVerteilungsfunktionVonN017}
$\forall\, t \in \!\ ]0,\exp\left(- \frac 1 \pi \right)]:
u\inv(t) \le \sqrt{2\log\left(\frac 1 t\right)}$.
\end{enumerate}
\end{satz}
\begin{proof}
(1) Sei $t \in [0, \infty[$. Dann gilt:
\begin{align*}
u(t) &= 1 - \Prim(\abs \gamma \le t) = 1 - \Prim(-t \le \gamma \le t)
= 1 - (\Phi(t) - \Phi(-t))\\
&= 1 - \Phi(t) + 1- \Phi(t) = 2(1- \Phi(t)) = 2\left(\int_t^\infty \frac 1 {\sqrt{2 \pi}} \exp\left(- \frac{s^2}2 \right)ds\right)\\
& = \sqrt{\frac 2 \pi}\int_t ^\infty \exp\left(- \frac{s^2}2 \right)ds\text.
\end{align*}
Wegen $\Bild \Phi = \!\ ]0,1[$ gilt daher $0 < 2(1-\Phi(t)) = u(t)$ und 
wegen $t \ge 0$ und damit $\Phi(t) \ge \frac 1 2$ auch
$u(t) = 2(1-\Phi(t)) \le 1$.\\
Damit ist $\Bild u \subseteq \!\ ]0, 1]$.\\
Weiter ist $\Phi$ invertierbar und zu $y \in\!\  ]0,1]$ ist 
$1 - \frac y 2 \in [\frac 1 2, 1[\, \subseteq \Bild \Phi$, also:
\begin{align}
u\left(\Phi\inv\left(1- \frac y 2 \right) \right) = 2(1-(1-\frac y 2))=y\text.
\label{edkvvneq:1}
\end{align}
Damit ist $\Bild u =\!\  ]0,1]$.\\
Da $\Phi$ streng monoton wachsend ist, ist $u$ streng monoton fallend und somit
insbesondere injektiv, besitzt also ein Inverses auf $]0,1]$\\
Wegen $\eqref{edkvvneq:1}$ ist für alle $z\in\!\ ]0,1]$: 
$u(\Phi\inv(1-\frac z2))=z$, also $u\inv(z) = \Phi\inv(1- \frac z 2)$.\\
(2) Es ist $\Phi$ stetig differenzierbar auf $\R$ mit
\[\forall\,t\in\R:\Phi'(t)=\frac 1{\sqrt{2\pi}}\exp\left(-\frac{t^2}2
 \right)\text.\]
Daher gilt für alle $t \in \R_{>0}:$
\[u'(t)\overset{(1)}= 2(-\Phi'(t))=-\frac 2{\sqrt{2\pi}}\exp\left(-\frac{t^2}2
\right)=-\sqrt{\frac 2 \pi}\exp\left(-\frac{t^2}2 \right)\text.\]
(3) Sei $\erf: \R \rightarrow \R,
x \mapsto \frac 2 {\sqrt \pi}\int_0^x\exp(-t^2)dt$.\\
Dann gilt für alle $x \in \R$:
\begin{align*}
&\, \frac 1 2 + \frac 1 2 \erf\left(\frac x {\sqrt 2} \right)
= \frac 1 2 + \frac 1 {\sqrt \pi}\int_0^{\frac x {\sqrt 2}}\exp(-t^2)dt\\
\overset{\substack{\text{Subst.}\\ s \mapsto \frac s{\sqrt 2}}}=&\, \frac 1 2
+\frac 1{\sqrt\pi}\int_0^x \exp\left(- \frac {t^2}2 \right)\frac 1{\sqrt 2}dt\\
=&\, \spi \int_{-\infty}^0 \exp\left(-\frac {t^2}2 \right)dt 
+ \spi\int_0^x\exp\left(- \frac {t^2}2 \right)dt = \Phi(x)\text.
\end{align*}
Also $\Phi = \frac 1 2\left(1+\erf\left(\frac \cdot{\sqrt 2} \right)\right)$,
d. h. 
\begin{align}
\erf = 2\Phi(\sqrt 2 \cdot) - 1\text. \label{edkvvneq:2}
\end{align}
Für alle $\alpha\in\!\ ]-1,1[$ ist $\beta\coloneqq \frac 12(\alpha+1)\in\!\ ]0,1[$,
also gibt es ein $t \in \R$ mit $\Phi(t) = \beta$ und damit
\[\erf\left(\frac 1 {\sqrt 2}t \right) \overset{\eqref{edkvvneq:2}}= 2\Phi(t)-1
= 2\left(\frac 1 2 (\alpha +1)\right) -1 = \alpha \text.\]
Damit ist $\Bild \erf = \!\ ]-1,1[$.\\
Offensichtlich ist $\erf$ analytisch und es gilt $\erf' = \frac 2{\sqrt \pi}
\exp\left(- (\cdot^2) \right) > 0$.\\
Wegen der Stetigkeit des Integranden ist $\erf$ also streng monoton wachsend
und damit injektiv. Ferner ist damit $\erf\inv:\!\ ]-1,1[ \rightarrow \R$
ebenfalls analytisch.\\
Nach Setzung ist $\erf$ ungerade und damit nach
\refclaim{GeradeUndUngeradeFunktionen}{GeradeUndUngeradeFunktionen4} auch
$\erf\inv$ ungerade.\\
Sei $k \in \N$. Nach
\refclaim{GeradeUndUngeradeFunktionen}{GeradeUndUngeradeFunktionen2} und
\refclaim{GeradeUndUngeradeFunktionen}{GeradeUndUngeradeFunktionen3} ist dann
$(\erf\inv)^{(2k)}$ ungerade und damit nach
\refclaim{GeradeUndUngeradeFunktionen}{GeradeUndUngeradeFunktionen1}:
\begin{align}(\erf\inv)^{(2k)}(0) = 0 \text. \label{edkvvneq:3}\end{align}
Weiter gilt wegen $\id_{]-1,1[} = \erf \circ \erf\inv$ auch
$1 = (\erf\inv)' \cdot \erf' \circ \erf\inv$, also:
\begin{align}
\forall\, p \in \!\ ]-1,1[: (\erf\inv)'(p) = \frac 1 {\erf' ( \erf \inv(p))} 
=\frac {\sqrt \pi}2 \exp(\erf\inv(p)^2)\text. \label{edkvvneq:4}
\end{align}
Damit gilt aber auch für alle $p \in\!\  ]-1,1[$:
\begin{align*}(\erf\inv)''(p) &\overset{\eqref{edkvvneq:4}}=
2\cdot\frac{\sqrt \pi}2 \erf\inv(p)(\erf\inv)'(p)
\exp\left(\erf\inv(p)^2\right)\\
&\overset{\eqref{edkvvneq:4}} = 
2\cdot\erf\inv(p)\cdot ((\erf\inv)'(p))^2\text.\end{align*}
Ferner gilt $(\erf\inv)'(0)= \frac {\sqrt\pi}2 \exp(\erf\inv(0)^2) > 0$, also
gilt nach~\refclaim{GeradeUndUngeradeFunktionen}{GeradeUndUngeradeFunktionen5}:
\begin{align}
(\erf\inv)^{(2k+1)}(0) \ge 0 \text.\label{edkvvneq:5}
\end{align}
Nach~\cite{Dom} (dort Abschnitt 4) ist die Taylorreihe von $\erf\inv$ mit
Entwicklungspunkt $0$ auf $]-1,1[$ konvergent. Wegen der Analytizität von
$\erf\inv$ stellt die besagte Reihe also auf $]-1,1[$ bereits $\erf\inv$ dar.\\
Also liefert die Taylorentwicklung um $0$ für alle $p \in \!\ ]-1,1[$:
\begin{align}
\erf\inv(p) = \sum_{l=0}^\infty \frac{(\erf\inv)^{(l)}(0)}{l!}p^l
\overset{\eqref{edkvvneq:3}}=\sum_{l=0}^\infty
\frac{(\erf\inv)^{(2l+1)}(0)}{(2l+1)!}p^{2l+1}\text.\label{edkvvneq:6}
\end{align}
Wegen~\eqref{edkvvneq:2} gilt weiter: $\Phi\inv=\sqrt 2\erf\inv(2\cdot -1)$.\\
Damit für alle $q \in \!\ ]0,1]$ wegen $1-\frac q 2 \in\!\  ]-1,1[$:
\begin{align*}
u\inv(q) &\overset{(1)}= \Phi\inv\left(1- \frac q 2\right)
= \sqrt 2 \erf\inv\left(2\left(1-\frac q 2 \right)-1 \right)
= \sqrt 2 \erf\inv(1-q)\\
&\overset{\eqref{edkvvneq:6}}= \sum_{l=0}^\infty
 \sqrt 2 \frac{(\erf\inv)^{(2l+1)}(0)}{(2l+1)!}(1-q)^{2l+1}\text.
\end{align*}
Nach~\eqref{edkvvneq:5} sind alle Koeffizienten dieser Reihendarstellung
nichtnegativ, also ist (3) gezeigt.\\
(4) Seien \[V_0 : [0, \infty[ \rightarrow \R, 
x \mapsto 2\exp\left(x^2\right)\int_x^\infty \exp\left(-s^2\right)ds \]
und 
\[g_\pi:[0,\infty[\rightarrow\R,x\mapsto \frac{\pi}{(\pi-1)x+\sqrt{x^2+\pi}}
\text.\]
Nach~\cite{RW} (Notation wie dort) gilt dann:
\begin{align}
g_\pi\le V_0\text.
\label{edkvvneq:7}
\end{align}
Sei $t \in \R_{>0}$. Dann gilt:
\begin{align*}
& \sqrt{\frac 2 \pi}\frac{\pi}{(\pi-1)t + \sqrt{t^2+2\pi}} =
\sqrt{\frac 2 \pi}\cdot \frac{\pi}{(\pi-1)\frac{\sqrt 2}{\sqrt 2}t
+ \sqrt{2\frac{t^2}{2} + 2\pi}}\\
 = & \,
\sqrt{\frac 2 \pi} \cdot \frac \pi {\sqrt 2\left((\pi-1)\frac 1{\sqrt 2}t  +
\sqrt{\left(\frac 1 {\sqrt{2}}t\right)^2 + \pi} \right)} 
= \sqrt{\frac 2 \pi} \frac 1 {\sqrt 2} g_\pi\left(\frac 1 {\sqrt 2} t \right)\\
\overset{\eqref{edkvvneq:7}}\le & \, \sqrt{\frac 2 \pi} \frac 1 {\sqrt 2} 
V_0\left(\frac 1 {\sqrt 2} t \right) = 
\sqrt{\frac 2 \pi}\exp\left(\frac {t^2} 2 \right)
\int_{\frac 1 {\sqrt 2} t}^\infty \exp\left(-s^2 \right)\sqrt 2 ds\\
= & \,\sqrt{\frac 2 \pi}\exp\left(\frac {t^2} 2 \right)
\int_{\frac 1{\sqrt 2}t}^\infty 
\exp\left(-\frac{(\sqrt 2s)^2}{2}\right)\sqrt 2 ds\\
\overset{\substack{\text{Subst.}\\ s \mapsto \sqrt 2 s}}=&\,
\sqrt{\frac 2 \pi}\exp\left(\frac {t^2} 2 \right)
\int_{t}^\infty \exp\left(- \frac{s^2}2 \right)ds
\overset{(1)}= \exp\left( \frac{t^2} 2 \right) u(t)\text,\end{align*}
also nach Umstellung und Kürzung:
\[\frac{\sqrt{2\pi}}{(\pi-1)t+\sqrt{t^2+2\pi}}\exp\left(-\frac{t^2}2\right)
 \le u(t)\text.\]
Weiter gilt auch:
\begin{align*}
t\cdot u(t) &\overset{(1)}= 
\sqrt{\frac 2 \pi} \int_t^\infty t \exp\left(- \frac {s^2}2 \right)ds
\le \sqrt{\frac 2 \pi} \int_t^\infty s \exp\left(- \frac {s^2}2 \right)ds\\
&= \sqrt{\frac 2 \pi}\exp\left(- \frac{t^2}2 \right)\text,
\end{align*}
also
\[u(t)\le\sqrt{\frac 2 \pi} \frac 1 t \exp\left(- \frac{t^2}2 \right)\text.\]
(5) Annahme: es gibt ein $t \in [1, \infty[$ mit 
$u\inv\left(\frac 1 {t+1} \right) < \sqrt{\log(t)}$.\\
Dann gilt:
\begin{align*}
&\frac 1 {t+1} = u\left(u\inv\left(\frac 1{t+1} \right)\right)\\
 \overset{(4)}\ge &\,
\frac{\sqrt{2\pi}}{(\pi-1)u\inv\left(\frac 1{t+1}\right)
+ \sqrt{\left(u\inv\left(\frac 1{t+1}\right)\right)^2+2\pi}}
\exp\left(-\frac 1 2 \left(u\inv\left(\frac 1{t+1}\right)\right)^2\right)\\
 >&\, \frac{\sqrt{2\pi}}{(\pi-1)\sqrt{\log(t)} 
+ \sqrt{\log(t)+2\pi}}\exp\left(-\frac{\log(t)}2\right)\\
=&\, \frac{\sqrt{2\pi}}{(\pi-1)\sqrt{\log(t)} 
+ \sqrt{\log(t)+2\pi}}\frac 1 {\sqrt t}\text.
\end{align*}
Somit auch
\[ \frac{\sqrt t}{t+1} \left((\pi-1)\sqrt{\log(t)} 
+ \sqrt{\log(t)+2\pi} \right) > \sqrt{2\pi}\text.\]
Sei $k: [1, \infty[ \rightarrow \R, 
s \mapsto \frac{\sqrt s}{s+1} \left((\pi-1)\sqrt{\log(s)} 
+ \sqrt{\log(s)+2\pi} \right)$.\\
Dann gilt $k(t) > \sqrt{2\pi}$ und nach
\refclaim{Numeriklemma}{Numeriklemma1} auch $k(t) <\sqrt{2\pi}$. $\lightning$\\
(6) Kontraposition:
Sei $t\in \!\ ]0, 1]$ mit $u\inv(t)<\sqrt{\log\left(1+\frac 1 t\right)}$.\\
Dann gilt:
\begin{align*}
t & = u\left(u\inv(t)\right) \overset{(4)}\ge
\frac{\sqrt{2\pi}}{(\pi-1)u\inv(t) 
+ \sqrt{u\inv(t)^2+2\pi}}\exp\left(-\frac{u\inv(t)^2}2\right)\\
& >  \frac{\sqrt{2\pi}}{(\pi-1) \sqrt{\log\left(1 +\frac 1 t\right)}
+ \sqrt{\log\left(1 +\frac 1 t\right)+2\pi}}
\exp\left(-\frac{\log\left(1 +\frac 1 t\right)}2\right)\\
& = \frac{\sqrt{2\pi}}{(\pi-1) \sqrt{\log\left(1 +\frac 1 t\right)}
+ \sqrt{\log\left(1 +\frac 1 t\right)+2\pi}}
\frac 1 {\sqrt{1+\frac 1 t}}\text.
\end{align*}
Also
\[t\sqrt{1+\frac 1 t} \left((\pi-1) \sqrt{\log\left(1 +\frac 1 t\right)}
+ \sqrt{\log\left(1 +\frac 1 t\right)+2\pi}\right)  > \sqrt{2\pi}\text.\]
Sei $h: [1, \infty[\!\  \rightarrow \R, s \mapsto
\frac{\sqrt{s+1}}s
\bigg( (\pi-1)\sqrt{\log(1+s)} + \sqrt{\log(1+s) + 2\pi} \bigg)$.\\
Dann ist also $\frac 1 t \in [1, \infty[$ und 
$h\left(\frac 1 t \right) > \sqrt{2\pi}$.\\
Nach~\refclaim{Numeriklemma}{Numeriklemma2} gilt daher $\frac 1 t <\frac{13}2$,
also $t > \frac 2 {13}$, d.h. $t \in\!\  ]\frac 2 {13}, 1]$ und damit
$t \notin \!\ ]0, \frac 2 {13}[$.\\
(7) Sei $t\in\!\  ]0,1]$ mit $u\inv(t) > \sqrt{2\log\left(\frac 1 t\right)}$.\\
Dann gilt:
\begin{align*}
t &= u\left(u\inv(t)\right) \overset{(4)}\le 
\sqrt{\frac 2 \pi} \frac 1 {u\inv(t)}
\exp\left(- \frac 1 2 \left(u\inv(t) \right)^2 \right)\\
&< \sqrt{\frac 2 \pi} \frac 1 {\sqrt{2\log\left(\frac 1 t\right)}}
\exp\left(- \frac 1 2 2\log\left(\frac 1 t\right) \right)
= \sqrt{\frac 2 \pi} \frac 1 {\sqrt{2\log\left(\frac 1 t\right)}} t\text.
\end{align*}
Daher also
$\sqrt{2\log\left(\frac 1 t \right)} < \sqrt {\frac 2 \pi}$, also 
 $\log\left(\frac 1 t \right) < \frac 1 \pi$
und somit $t > \exp\left(- \frac 1 \pi \right)$.
\end{proof}
\parag
Als letztes vorbereitendes Hilfsmittel benötigen wir noch zwei im Wesentlichen
elementare Abschätzungen der unvollständigen Betafunktion, die Inhalt des
nachfolgenden Lemmas sind.

\begin{lemma}{Absch�tzung der unvollst�ndigen Beta--Funktion}
             {AbschaetzungDerUnvollstaendigenBetafunktion}
Seien $n \in \N_{\ge 2}$, $k \in \haken{n-1}$ und $u:[0,\infty[ \rightarrow \R,
t \mapsto\sqrt{\frac 2\pi} \int_t^\infty \exp\left(- \frac{s^2}2 \right)ds$.\\
Dann gilt:
\begin{enumerate}[(1)]
	\item \label{AbschaetzungDerUnvollstaendigenBetafunktion1}
	$\forall\, x \in \!\ ]0, \frac k {n-1}[\, \forall\, m \in \haken{n-k-1}_0:
	\int_0^x t^{n-1-m}(1-t)^m dt \le x^{n-m}(1-x)^m$.
	\item \label{AbschaetzungDerUnvollstaendigenBetafunktion2}
	$\forall\, y \in \R_{>0}\, \forall\, l \in \haken {n-1}_0:
	\int_{y}^\infty tu(t)^{n-l}(1-u(t))^ldt \le
	\int_0^{u(y)}t^{n-l-1}(1-t)^ldt$.
\end{enumerate}
\end{lemma}
\begin{proof}
(1) Seien $x \in\!\  ]0, \frac k {n-1}[$ und $m \in \haken{n-k-1}_0$.\\
Falls $m = 0$, so ist
\[\int_0^{x} t^{n-1-m}(1-t)^m dt 
= \int_0^{x} t^{n-1}dt = \frac 1 n x^n \le x^n = x^{n-m}(1-x)^m\]
Falls $m > 0$, so betrachte
$\alpha: [0, x] \rightarrow \R, t \mapsto t^{n-1-m}(1-t)^m$.\\
Dann ist $\alpha$ offenbar stetig differenzierbar auf $]0,x[$ und es gilt
f�r alle $s \in\!\  ]0, x[:$
\begin{align*}
\alpha'(s)& = (n-1-m)s^{n-m-2}(1-s)^m - ms^{n-m-1}(1-s)^{m-1}\\
&= s^{n-m-2}(1-s)^{m-1}\bigg((n-1-m)(1-s)-ms \bigg)\text.
\end{align*}
Ferner gilt f�r alle $s \in \!\ ]0,x[$:
\[\frac{n-m-1}{n-1} \ge \frac{n-(n-k-1)-1}{n-1} = \frac k{n-1} > x > s\text,\]
also
\[\frac{1-s}s = \frac 1 s - 1 > \frac{n-1}{n-m-1}-1 =
\frac{n-1-(n-m-1)}{n-m-1}= \frac m {n-m-1}\]
und damit $(n-1-m)(1-s) > ms$.\\
Somit ist $\alpha'(s)>0$ f�r alle $s \in\!\  ]0,x[$, also
$\restrict{\alpha}{\!]0,x[}$ monoton wachsend.\\
Damit gilt wegen $\alpha(0) = 0$:
\begin{align*}\int_0^{x} t^{n-1-m}(1-t)^m dt
= \int_0^x \alpha(t)dt \le x  \sup_{t \in [0,x]}\alpha(t)
= x  \alpha(x) = x^{n-m}(1-x)^m\text.\end{align*}
(2) Seien $y \in \R_{>0}$ und $l \in \haken {n-1}_0$.\\
Dann ist $n-l \ge 1$ und es gilt:
\begin{align*}
&\int_y^\infty t u(t)^{n-l}(1-u(t))^ldt
= \int_y^\infty u(t)^{n-l-1}(1-u(t))^l t u(t)dt\\
\overset{\refclaim{EigenschaftenDerKomplementaerenVerteilungsfunktionVonN01}
{EigenschaftenDerKomplementaerenVerteilungsfunktionVonN014}}\le
&\, \int_y^\infty  u(t)^{n-l-1}(1-u(t))^l t \cdot 
\sqrt{\frac 2 \pi} \frac 1 t \exp\left(- \frac 1 2 t^2 \right)dt\\
=&\, \int_y^\infty  u(t)^{n-l-1}(1-u(t))^l \cdot \sqrt{\frac 2 \pi}
\exp\left(- \frac 1 2 t^2 \right)dt\\
\overset{\refclaim{EigenschaftenDerKomplementaerenVerteilungsfunktionVonN01}
{EigenschaftenDerKomplementaerenVerteilungsfunktionVonN012}}=&\,
\int_y^\infty  u(t)^{n-l-1}(1-u(t))^l \abs{u'(t)}dt
= \int_{]y, \infty[}u(t)^{n-l-1}(1-u(t))^l \abs{u'(t)}dt\\
\overset{\substack{\text{Transf.}\\\text{mit }u}}=&\,
\int_{]0,u(y)[}t^{n-l-1}(1-t)^ldt = \int_0^{u(y)} t^{n-l-1}(1-t)^ldt
\end{align*}
\end{proof}
\parag
Wir sind nun imstande, Schranken für den Erwartungswert der $k$--ten absoluten
Ordnungsstatistik eines standardnormalverteilten $\gamma$ anzugeben.
Es stellt sich heraus, dass die obere und untere Schranke sich im Wesentlichen nur um
eine multiplikative Konstante unterscheiden und dass vor allem das Verhältnis
der Anzahl zur gewählten Ordnungsstatistik eine Rolle spielt (asymptotisch:
$\E(\gamma_k^\ast) \in \Theta(\nicefrac n k \mapsto
\sqrt{\log\left(\nicefrac n k \right)})$). Wir werden an
an dieser Stelle lediglich auf jene Ordnungsstatistiken eingehen, die an
späterer Stelle von Interesse sein werden. Alle anderen können mit ähnlichen
Abschätzungen erreicht werden, wie sie in
\refclaim{EigenschaftenDerKomplementaerenVerteilungsfunktionVonN01}
{EigenschaftenDerKomplementaerenVerteilungsfunktionVonN015} und
\refclaim{EigenschaftenDerKomplementaerenVerteilungsfunktionVonN01}
{EigenschaftenDerKomplementaerenVerteilungsfunktionVonN016} für gewisse
Spezialfälle durchgeführt wurden.

\begin{satz}{Untere Schranke des Erwartungswerts von abs.
\textsl{N}(0,1)--Ord.st.}
{UntereSchrankeDesErwartungswertsVonAbsolutenN01Ordnungsstatistiken}
Seien $n \in \N$, $k \in \haken n$, $\gamma: \Omega \rightarrow \R^n$
unabhängig und $N(0,1)$--verteilt, sowie $u: [0, \infty[\!\ \rightarrow \R, 
t \mapsto \Prim(\abs{\gamma_1}> t)$
(vgl.~\ref{EigenschaftenDerKomplementaerenVerteilungsfunktionVonN01}).\\
Dann ist $\gamma_k^\ast$ nach
\refclaim{EigenschaftenVonOrdnungsstatistiken}
{EigenschaftenVonOrdnungsstatistiken1} messbar, $u$ nach
\refclaim{EigenschaftenDerKomplementaerenVerteilungsfunktionVonN01}
{EigenschaftenDerKomplementaerenVerteilungsfunktionVonN011} invertierbar und es
gilt:
\begin{enumerate}[(1)]
  \item
   \label{UntereSchrankeDesErwartungswertsVonAbsolutenN01Ordnungsstatistiken1}
	$\E\left(\gamma_k^\ast \right) \ge u\inv\left(\frac k {n+1} \right)$.
	\item
	 \label{UntereSchrankeDesErwartungswertsVonAbsolutenN01Ordnungsstatistiken2}
	 $\E\left(\gamma_1^\ast \right)\ge \sqrt{\log(n)}$.
	 \item
	 \label{UntereSchrankeDesErwartungswertsVonAbsolutenN01Ordnungsstatistiken3}
	 Ist $k \le \frac 2 {13}(n+1)$, so ist
	 $\E\left(\gamma_k^\ast \right) \ge \sqrt{\log\left(1+\frac{n+1}k \right)}$.
\end{enumerate}
\end{satz}
\begin{proof}
(1) Offenbar ist 
$\gamma_k^\ast$ die $k$--te Ordnungsstatistik von 
$\pfeil {\abs\cdot} n \circ \gamma$.\\
Es ist $\Prim(\abs{\gamma_1} \le \cdot) = 1 - \Prim(\abs{\gamma_1} > \cdot)
= 1 - u$, also hat $\abs{\gamma_1}$ die stetige Dichte $-u'$.\\
Es gilt für alle $i \in \N_0$:
\begin{align*}
&\frac{n!}{(n+2i+1)!}\cdot\frac{(n-k+2i+1)!}{(n-k)!}
= \left(\prod_{l=0}^{2i}\frac{1}{n+l+1} \right)
\left(\prod_{r=0}^{2i}(n-k+1+r) \right)\\
=&\, \prod_{l=0}^{2i}\frac{n+l+1-k}{n+l+1}
=\prod_{l=0}^{2i}\left(1- \frac{k}{n+l+1} \right)
\ge \prod_{l=0}^{2i}\left(1- \frac{k}{n+1} \right)\\
=&\, \left(1- \frac{k}{n+1} \right)^{2i+1}\text.\tag {$\ast$}
\end{align*}
Sei $(c_j)_{j\in \N_0} \in [0, \infty[^{\N_0}$ die Folge aus
\refclaim{EigenschaftenDerKomplementaerenVerteilungsfunktionVonN01}
         {EigenschaftenDerKomplementaerenVerteilungsfunktionVonN013} und
$B, \Gamma$ wie in~\ref{ElementareEigenschaftenDerGammafunktion}.\\
Somit gilt wegen \refclaim{EigenschaftenVonOrdnungsstatistiken}
{EigenschaftenVonOrdnungsstatistiken2}:
\begin{align*}
\E(\gamma_k^\ast) &\overset{
\refclaim{EigenschaftenVonOrdnungsstatistiken}
{EigenschaftenVonOrdnungsstatistiken4}}= 
\int_0^\infty t \cdot n \binom{n-1}{n-k}\Prim(\abs{\gamma_1}\le t)^{n-k}
(1 - \Prim(\abs{\gamma_1} \le t))^{k-1}(-u'(t))dt\\
&= n \binom{n-1}{n-k} \int_0^\infty t (1-u(t))^{n-k}u(t)^{k-1}(-u'(t))dt\\
&= n \binom{n-1}{n-k} 
\int_{]0, \infty[} (1-u(t))^{n-k}u(t)^{k-1}u\inv(u(t)) \abs{u'(t)}dt\\
&\overset{\substack{\text{Transf.}\\ \text{mit }u}}=
n \binom{n-1}{n-k} \int_{]0,1[}(1-t)^{n-k}t^{k-1}u\inv(t)dt\\
&\overset{\refclaim{EigenschaftenDerKomplementaerenVerteilungsfunktionVonN01}
         {EigenschaftenDerKomplementaerenVerteilungsfunktionVonN013}}=
n \binom{n-1}{n-k} \int_{]0, 1[} (1-t)^{n-k}t^{k-1}\sum_{j=0}^\infty 
c_j(1-t)^{2j+1} dt\\
&=n \binom{n-1}{n-k}\int_{]0,1[} 
\sum_{j=0}^\infty c_j (1-t)^{n+2j+1-k}t^{k-1}dt\\
&\overset{\substack{\text{Satz v. d.}\\\text{mon. Konv.}}}= 
n \binom{n-1}{n-k} \sum_{j=0}^\infty c_j
\int_{]0,1[}(1-t)^{n+2j+1-k}t^{k-1}dt\\
&=n \binom{n-1}{n-k} \sum_{j=0}^\infty c_j
\int_0^1(1-t)^{n+2j+1-k}t^{k-1}dt\\
&\overset{\text{Def. }B}= n \binom{n-1}{n-k} \sum_{j=0}^\infty c_j
B(n+2j-k+2, k)\\
&\overset{\refclaim{ElementareEigenschaftenDerGammafunktion}
{ElementareEigenschaftenDerGammafunktion4}}=
n \binom{n-1}{n-k} \sum_{j=0}^\infty c_j \frac{\Gamma(n+2j-k+2)\Gamma(k)}
{\Gamma(n+2j+2)}\\
&\overset{\refclaim{ElementareEigenschaftenDerGammafunktion}
{ElementareEigenschaftenDerGammafunktion3}}= 
\sum_{j=0}^\infty c_j \frac{n!}{(k-1)!(n-k)!}
\frac{(k-1)!(n-k+2j+1)!}{(n+2j+1)!}\\
&= \sum_{j=0}^\infty c_j \frac{n!}{(n-k)!}\frac{(n-k+2j+1)!}{(n+2j+1)!}
\overset{(\ast)}\ge \sum_{j=0}^\infty c_j
\left(1-\frac{k}{n+1}\right)^{2j+1}\\
&\overset{\refclaim{EigenschaftenDerKomplementaerenVerteilungsfunktionVonN01}
{EigenschaftenDerKomplementaerenVerteilungsfunktionVonN013}}= 
u\inv\left(\frac k {n+1} \right)\text.
\end{align*}
(2) Es gilt:
$\E(\gamma_1^\ast) \overset{(1)}\ge u\inv\left(\frac 1 {n+1} \right)
\overset{\refclaim{EigenschaftenDerKomplementaerenVerteilungsfunktionVonN01}
{EigenschaftenDerKomplementaerenVerteilungsfunktionVonN015}}\ge
\sqrt{\log(n)}$.\\
(3) Es gelte $k \le \frac 2{13}(n+1)$.\\
Dann ist $\frac{k}{n+1} \le \frac 2 {13}$, also
$\E(\gamma_k^\ast) \overset{(1)}\ge u\inv\left(\frac k {n+1} \right)
\overset{\refclaim{EigenschaftenDerKomplementaerenVerteilungsfunktionVonN01}
{EigenschaftenDerKomplementaerenVerteilungsfunktionVonN016}}\ge
\sqrt{\log\left(1 + \frac{n+1}k \right)}$.
\end{proof}
\parag
Diese untere Abschätzung basiert ausschließlich auf abstrakten Eigenschaften
der verwendeten Funktionen. Die obere Abschätzung gestaltet sich hingegen
anschaulicher, weil die komplementäre Verteilungsfunktion monoton fallend ist,
sodass man bis zu einem gut gewählten Punkt über die konstante 1--Funktion
integrieren und den verbleibenden Rest geeignet abschätzen kann.
Interessanterweise hängt die Wahl dieses Punktes auf natürliche Art und Weise
von dem Verhältnis von $k$ zu $n$ ab.\\
Der Vollständigkeit halber werden wir in diesem Satz auch alle später
benötigten oberen Abschätzungen angeben wobei sich deren Beweise weitgehend
aus der ersten Aussage ergeben.

\begin{satz}
{Obere Schranke des Erwartungswerts von abs.
\textsl{N}(0,1)--Ord.st.}
{ObereAbschaetzungDesErwartungswertsDerKtenAbsolutenN01Ordnungsstatistik}
Seien $k, n \in \N$ mit $k \le \frac n 2$, $\gamma:\Omega\rightarrow \R^n$
unabhängig und $N(0,1)$--verteilt und $u: [0, \infty[ \rightarrow \R,
t \mapsto \Prim(\abs{\gamma_1}> t)$.\\
Dann gilt:
\begin{enumerate}[(1)]
	\item\label
	{ObereAbschaetzungDesErwartungswertsDerKtenAbsolutenN01Ordnungsstatistik1}
	$\forall\, x \in \R_{>0}: u(x) \le \frac k n
	\folgt \E\left(\left(\gamma_k^\ast\right)^2\right) \le x^2 + 4$.
	\item\label
	{ObereAbschaetzungDesErwartungswertsDerKtenAbsolutenN01Ordnungsstatistik2}
	$\E\left(\left(\gamma_k^\ast\right)^2\right) 
	\le 2 \log\left(\frac n k \right) + 4 
	\le \left(2+\frac 4{\log(2)} \right)\log\left(\frac n k \right)$.
	\item\label
	{ObereAbschaetzungDesErwartungswertsDerKtenAbsolutenN01Ordnungsstatistik3}
	$\E\left(\left(\gamma_k^\ast\right)^2 \right)
	\le 2 \log\left(\frac {2n} k \right) + 4 
	\le \left(2+\frac 4{\log(4)} \right)\log\left(\frac {2n} k \right)$.
	\item\label
	{ObereAbschaetzungDesErwartungswertsDerKtenAbsolutenN01Ordnungsstatistik4}
	$\E\left(\gamma_k^\ast\right) \le \sqrt{2+\frac 4{\log(2)}}
	\sqrt{\log\left(\frac n k\right)}$.
	\item\label
	{ObereAbschaetzungDesErwartungswertsDerKtenAbsolutenN01Ordnungsstatistik5}
	$\E\left(\gamma_k^\ast\right) \le \sqrt{2+\frac 4{\log(4)}}
	\sqrt{\log\left(\frac {2n} k\right)}$.
	\item\label
	{ObereAbschaetzungDesErwartungswertsDerKtenAbsolutenN01Ordnungsstatistik6}
	Ist $k \le \frac n {\exp(1)}$, so ist $\E\left(\norm\cdot_{2,k} \circ 
	\left(\gamma_j^\ast\right)_{j\in \haken k}   \right)
	\le \sqrt 2 \sqrt{2+\frac 4{\log(2)}}\sqrt{k\log\left(\frac{n}k\right)}$.
	\item\label
	{ObereAbschaetzungDesErwartungswertsDerKtenAbsolutenN01Ordnungsstatistik7}
	$\E\left(\norm\cdot_{2,k} \circ 
	\left(\gamma_j^\ast\right)_{j\in \haken k}   \right)
	\le \sqrt 2\sqrt{2+\frac 4{\log(4)}}\sqrt{k\log\left(\frac{2n}k\right)}$.
\end{enumerate}
\end{satz}
\begin{proof}
(1) Zunächst gilt offenbar:
\begin{align}
\forall\, t \in [0, \infty[: \left(\gamma_k^\ast\right)^2 > t \iff
\gamma_k^\ast > \sqrt t \text.\label{oadedkanoeq:1}
\end{align}
Sei $x \in \R_{>0}$ mit $u(x) \le\frac k n $.
Wegen $\haken{n-k}_0 \subseteq \haken {n-1}_0$ gilt nach
\refclaim{AbschaetzungDerUnvollstaendigenBetafunktion}
{AbschaetzungDerUnvollstaendigenBetafunktion2}:
\begin{align}\forall\, l \in \haken {n-k}_0:
\int_x^\infty (1-u(t))^l(u(t))^{n-l}tdt
\le \int_0^{u(x)} (1-t)^lt^{n-l-1}dt \label{oadedkanoeq:2}
\end{align}
und wegen $u(x) \le \frac k n \le \frac k {n-1}$ nach
\refclaim{AbschaetzungDerUnvollstaendigenBetafunktion}
{AbschaetzungDerUnvollstaendigenBetafunktion1}:
\begin{align}\forall\, l \in \haken{n-k-1}_0:
\int_0^{u(x)} t^{n-l-1}(1-t)^ldt
\le (u(x))^{n-l}(1-u(x))^l \text.\label{oadedkanoeq:3}
\end{align}
Seien $B, \Gamma$ wie in~\ref{ElementareEigenschaftenDerGammafunktion}.
Wir fassen $\gamma_k^\ast$ als $k$--te Ordnungsstatistik von
$\pfeil {\abs\cdot} n \circ \gamma$ auf.\\
Dann gilt mit \refclaim{EigenschaftenVonOrdnungsstatistiken}
{EigenschaftenVonOrdnungsstatistiken2}:
\begin{align*}
& \int_{x^2}^\infty \Prim\left(\gamma_k^\ast > \sqrt t\right)dt\\
\overset{\refclaim{EigenschaftenVonOrdnungsstatistiken}
{EigenschaftenVonOrdnungsstatistiken3}}= &\,
\int_{x^2}^\infty \sum_{l=0}^{n-k}\binom n l 
\Prim\left(\abs{\gamma_1} \le\sqrt t\right)^l
\left(1-\Prim\left(\abs{\gamma_1}\le \sqrt t\right)\right)^{n-l}dt\\
=&\, \sum_{l=0}^{n-k}\binom n l \int_{x^2}^\infty 
\left(u\left(\sqrt t\right)\right)^{n-l}\left(1-u\left(\sqrt t\right)\right)^l
dt\\
\overset{x>0}=&\, \sum_{l=0}^{n-k}\binom n l \int_{x^2}^\infty 
\left(u\left(\sqrt t\right)\right)^{n-l}\left(1-u\left(\sqrt t\right)\right)^l
2\sqrt t \frac 1 {2\sqrt t}dt\\
\overset{\substack{\text{Subst.}\\t \mapsto \sqrt t}}= &\,
2\sum_{l=0}^{n-k}\binom n l \int_x^\infty (u(t))^{n-l}(1-u(t))^l tdt\\
\overset{\eqref{oadedkanoeq:2}}\le &\,
2\sum_{l=0}^{n-k}\binom n l \int_0^{u(x)} t^{n-l-1}(1-t)^ldt\\
\overset{\eqref{oadedkanoeq:3}}\le &\,
2\binom n {n-k}\int_0^{u(x)}t^{k-1}(1-t)^{n-k}dt +
2\underbrace{
\sum_{l=0}^{n-k-1}\binom n l (u(x))^{n-l}(1-u(x))^l}_
{\overset{\text{Bin. Form.}}\le 1}\\
\le &\, 2 \binom n{n-k}B(k, n-k+1) + 2 
\overset{\refclaim{ElementareEigenschaftenDerGammafunktion}
{ElementareEigenschaftenDerGammafunktion4}}=
2 \frac{n!}{(n-k)!k!}\frac{\Gamma(k)\Gamma(n-k+1)}{\Gamma(n+1)} + 2\\
\overset{\refclaim{ElementareEigenschaftenDerGammafunktion}
{ElementareEigenschaftenDerGammafunktion3}}=&\,
2 \frac{n!}{(n-k)!k!} \frac{(k-1)!(n-k)!}{n!} + 2
= \frac 2 k + 2 \le 4\text.
\end{align*}
Somit gilt wegen $\gamma_k^\ast \ge 0$:
\begin{align*}
\E\left(\left(\gamma_k^\ast\right)^2 \right) 
&= \int_0^\infty \Prim\left(\left(\gamma_k^\ast\right)^2> t\right)dt
\overset{\eqref{oadedkanoeq:1}}=
\int_0^\infty \Prim\left(\gamma_k^\ast> \sqrt t\right)dt\\
&= \int_0^{x^2}\underbrace{\Prim\left(\gamma_k^\ast> \sqrt t\right)}_{\le 1}dt
+ \int_{x^2}^\infty \Prim\left(\gamma_k^\ast> \sqrt t\right)dt
\le x^2 + 4\text.
\end{align*}
(2) Sei $x \coloneqq  u\inv\left(\frac k n \right)$.\\
Dann ist wegen $3 < \pi < 4$ und der Standardabschätzung für $\exp$
\[\exp\left(\frac 1\pi\right)\le\frac 1{1-\frac1\pi}=\frac\pi{\pi-1}< 2\text,\]
also auch $\exp\left(-\frac1\pi \right) \ge \frac 1 2 \ge \frac k n$ und somit
nach \refclaim{EigenschaftenDerKomplementaerenVerteilungsfunktionVonN01}
{EigenschaftenDerKomplementaerenVerteilungsfunktionVonN017}:
$x \le \sqrt{2 \log\left(\frac n k \right)}$.\\
Somit
\begin{align*}
\E\left(\left(\gamma_k^\ast\right)^2 \right) \overset{(1)} \le
2 \log\left(\nicefrac n k \right) + 4 
= \left(2 + \nicefrac 4 {\log\left(\frac n k \right)} \right)
\log\left(\nicefrac n k \right)
 \overset{\frac n k \ge 2}\le
\left(2 + \nicefrac 4 {\log(2)} \right)
\log\left(\nicefrac n k \right)\text.
\end{align*}
(3) Es gilt:
\begin{align*}
\E\left(\left(\gamma_k^\ast\right)^2 \right) &\overset{(2)} \le
2 \log\left(\frac n k \right) + 4
\le 2 \log\left(\frac {2n} k \right) + 4 = \left(2 + \frac 4 {\log\left(\frac {2n} k \right)} \right)
\log\left(\frac {2n} k \right)\\
&\overset{\frac{2n}k \ge 4}\le
\left(2 + \frac 4 {\log(4)} \right)\log\left(\frac {2n} k \right)\text.
\end{align*}
(4)/(5) Gemäß Jensenungleichung ist
$\E\left(\gamma_k^\ast \right)= 
\E\left(\sqrt\cdot \circ \left(\gamma_k^\ast \right)^2 \right)
\le \sqrt{\E\left(\left(\gamma_k^\ast\right)^2 \right)}$, also
folgt (4) aus (2) und (5) aus (3).\\
(6) Sei $k \le \frac n {\exp(1)}$. Dann gilt:
\begin{align*}
& \E\left(\norm\cdot_{2,k} \circ 
	\left(\gamma_j^\ast\right)_{j\in \haken k}   \right)
= \E\left(\sqrt\cdot \circ \sum_{j=1}^k \left(\gamma_j^\ast\right)^2 \right)
\le \sqrt{\sum_{j=1}^k \E\left(\left(\gamma_j^\ast\right)^2 \right)}\\
\overset{(2)}\le&\, \sqrt{\sum_{j=1}^k\left(2+\frac 4{\log(2)} \right)
\log\left(\frac n j\right)}
=\sqrt{2+\frac 4{\log(2)}}\sqrt{\sum_{j=0}^k\log\left(\frac n j\right)}\\
\overset{\ref{SummenVonLogarithmen}}\le&\,
\sqrt{2+\frac 4{\log(2)}}\sqrt{2 k \log\left(\frac n k \right)}
= \sqrt 2 \sqrt{2+\frac 4{\log(2)}}\sqrt{k \log \left(\frac n k \right)}\text.
\end{align*}
(7) Wegen $k \le \frac n 2 = \frac{2n} 4 \le \frac{2n}{\exp(1)}$ folgt diese
Aussage analog zu (6), nur dass (3) statt (2) angewandt wird.
\end{proof}

Die Aussagen $(2), (3), (4), (5)$ des vorigen Satzes lassen sich analog mit
der Voraussetzung $k \le \frac n {\exp(\frac 1 \pi)}$ anstelle von
$k \le \frac n 2$ beweisen, wenn man den Vorfaktor der jeweiligen Abschätzung
entsprechend vergrößert.\parag
Der nachfolgend zitierte Gauß'sche Min--Max--Satz ist Hauptwerkzeug für viele
der späteren Sätze. Im weiteren Verlauf der Arbeit werden wir diesen Satz unter
einer weiteren Voraussetzung auf spezielle unendliche Familien von
Zufallsgrößen übertragen, welche dann mit Hilfe des Erwartungswertlemmas
(\ref{Erwartungswertlemma}) für die Existenzzusicherung gewünschter linearer
Abbildungen verwendet werden.
\newpage
\begin{satz}{Gauß'scher Min--Max--Satz}
            {GaussscherMinMaxSatz}
Seien $r, s \in \N$ und 
$X, Y: \haken r \times \haken s \rightarrow L_2(\Omega, \mfa, \Prim, \R)$
mit 
\[\forall\,i\in\haken r\,\forall\,j\in\haken s\,\exists\,\rho_{i,j}\in\R_{>0}:
X(i, j) \text{ ist }N(0, \rho_{i,j}^2)\text{--verteilt,}\]
\[
\forall\,i\in\haken r\,\forall\,j\in\haken s\,
\exists\,\sigma_{i,j}\in\R_{>0}:
Y(i, j) \text{ ist }N(0, \sigma_{i,j}^2)\text{--verteilt.}\]
Es gelte
\begin{enumerate}[(a)]
	\item $\forall\, i \in \haken r \, \forall\, k,l \in \haken s:
	\norm{X(i,k)-X(i,l)}_2 \le \norm{Y(i,k)-Y(i,l)}_2$,
	\item $\forall\, i, j \in \haken r\,\forall\, k, l \in \haken s:
i \neq j \folgt \norm{X(i,k)-X(j,l)}_2 \ge \norm{Y(i,k)-Y(j,l)}_2$.
\end{enumerate}
Dann gilt
\[\E\left(\min_{i \in \haken r}\max_{j \in \haken s}X(i,j) \right)
\le\E\left(\min_{i \in \haken r}\max_{j \in \haken s}Y(i,j) \right)\text.\]
\end{satz}
\begin{proof}
siehe~\cite{Gor}
\end{proof}

Besonders interessant ist ein Spezialfall des Satzes, in dem die erste
Bedingung stets mit Gleichheit erfüllt ist. Dadurch dualisieren sich in
gewisser Hinsicht die Forderungen an $X, Y$, sodass diese gemäß der erwähnten
Dualisierung vertauscht werden können. Diese
Dualisierung ist Inhalt des folgenden Korollars.

\begin{korollar}
  {Gauß'scher Min--Max--Satz (Gleichheit in der ersten Bed.)}
  {GaussscherMinMaxSatzGleichheitInDerErstenBedingung}
Seien $r, s \in \N$ und 
$X, Y: \haken r \times \haken s \rightarrow L_2(\Omega, \mfa, \Prim, \R)$
mit 
\[\forall\,i\in\haken r\,\forall\,j\in\haken s\,\exists\,\rho_{i,j}\in\R_{>0}:
X(i, j) \text{ ist }N(0, \rho_{i,j}^2)\text{--verteilt,}\]
\[
\forall\,i\in\haken r\,\forall\,j\in\haken s\,
\exists\,\sigma_{i,j}\in\R_{>0}:
Y(i, j) \text{ ist }N(0, \sigma_{i,j}^2)\text{--verteilt.}\]
Es gelte
\begin{enumerate}[(a)]
	\item $\forall\, i \in \haken r \, \forall\, k,l \in \haken s:
	\norm{X(i,k)-X(i,l)}_2 = \norm{Y(i,k)-Y(i,l)}_2$,
	\item $\forall\, i, j \in \haken r\,\forall\, k, l \in \haken s:
i \neq j \folgt \norm{X(i,k)-X(j,l)}_2 \ge \norm{Y(i,k)-Y(j,l)}_2$.
\end{enumerate}
Dann gilt:
\begin{enumerate}[(1)]
	\item\label{GaussscherMinMaxSatzGleichheitInDerErstenBedingung1}
	$\E\left(\underset{{i \in \haken r}}\min\,
	         \underset{{j \in \haken s}}\max \,X_{i,j} \right)
	\le \E\left(\underset{{i \in \haken r}}\min\,
	         \underset{{j \in \haken s}}\max \,Y_{i,j} \right)$.
	\item\label{GaussscherMinMaxSatzGleichheitInDerErstenBedingung2}
	$\E\left(\underset{{i \in \haken r}}\max\,
	         \underset{{j \in \haken s}}\max \,Y_{i,j} \right)
  \le\E\left(\underset{{i \in \haken r}}\max\,
	         \underset{{j \in \haken s}}\max \,X_{i,j} \right)$.
\end{enumerate}
Insbesondere also
\begin{align*}
\E\bigg(\min_{i \in \haken r}\max_{j \in \haken s}X_{i,j} \bigg)
&\le\E\left(\min_{i \in \haken r}\max_{j \in \haken s}Y_{i,j} \right)
\le\E\left(\max_{i \in \haken r}\max_{j \in \haken s}Y_{i,j} \right)\\
&\le\E\left(\max_{i \in \haken r}\max_{j \in \haken s}X_{i,j} \right)\text.
\end{align*}
\end{korollar}
\begin{proof}
(1) Offenbar sind die Voraussetzungen von~\ref{GaussscherMinMaxSatz} erfüllt.\\
Damit ist bereits die erste Aussage gezeigt.\\
(2) Seien zunächst $a \in \haken r$ und $b, c \in \haken s$.\\
Wegen (a) ist dann $\norm{X_{a,b}-X_{a,c}}_2 = \norm{Y_{a,b} -Y_{a,c}}_2$, sodass
also gilt 
\begin{align}
\forall\, i, j \in \haken r\,\forall\, k, l \in \haken s:
\norm{X_{i,k}-X_{j,l}}_2 \ge \norm{Y_{i,k}-Y_{j,l}}_2 \tag{$\ast$}
\end{align}
Wegen $\abs{\haken r \times \haken s}= \abs{\haken{rs}}$ gibt es
$\sigma: \haken {rs} \rightarrow \haken r \times \haken s$ bijektiv.\\
Seien $\rho_1: \haken r \times \haken s \rightarrow \haken r,(x, y) \mapsto x$,
$\rho_2: \haken r \times \haken s \rightarrow \haken s, (x, y) \mapsto y$ und
$\tau_i\coloneqq  \rho_i \circ \sigma$ für alle $i \in \haken 2$.\\
Seien nun $r'\coloneqq 1, s'\coloneqq  rs$ sowie
\begin{align*}
X' & :\haken {r'}\times \haken {s'} \rightarrow L_2(\Omega, \Prim, \mfa, \R),
(t,z) \mapsto Y(\tau_1(z), \tau_2(z))\text,\\
Y' & :\haken {r'}\times \haken {s'} \rightarrow L_2(\Omega, \Prim, \mfa, \R),
(t,z) \mapsto X(\tau_1(z), \tau_2(z))\text.
\end{align*}
Dann sind $X', Y'$ offensichtlich punktweise zentriert und
normalverteilt, weil $X, Y$ es sind.\\
Seien $k, l \in \haken {s'}$. Dann gilt:
\begin{align*}
&\norm{X'(1,k)-X'(1, l)}_2 = \norm{Y(\tau_1(k), \tau_2(k))
- Y(\tau_1(l), \tau_2(l))}_2\\
\overset{(\ast)}\le &\,
\norm{X(\tau_1(k), \tau_2(k))- X(\tau_1(l), \tau_2(l))}_2
= \norm{Y'(1,k)-Y'(1, l)}_2
\end{align*}
Wegen $r' = 1$ gilt weiter trivial
\begin{align*}
\forall i, j \in \haken {r'}\forall m, n \in \haken {s'}:
i \neq j \folgt \norm{X'(i,m)-X'(j,n)}_2 \ge \norm{Y'(i,m)-Y'(j,n)}_2\text.
\end{align*}
Damit sind die Voraussetzungen von~\ref{GaussscherMinMaxSatz} auch für
$X'$ und $Y'$ erfüllt, also gilt:
\begin{align*}
&\E\left(\max_{i \in \haken r}\max_{i \in \haken s} Y(i,j) \right)
= \E\left(\max_{(i, j) \in \haken r \times \haken s} Y(i,j) \right)
= \E\left(\min_{i\in \haken {r'}}\max_{z \in \haken {s'}}
Y(\tau_1(z), \tau_2(z))\right)\\
= &\, \E\left(\min_{i\in \haken {r'}}\max_{z \in \haken {s'}}
X'(i, z)\right)
\overset{\ref{GaussscherMinMaxSatz}}\le
\E\left(\min_{i\in \haken {r'}}\max_{z \in \haken {s'}}
Y'(i, z)\right)\\
=&\,\E\left(\min_{i\in \haken {r'}}\max_{z \in \haken {s'}}
X(\tau_1(z), \tau_2(z)) \right)
=\E\left(\max_{z \in \haken {s'}}X(\tau_1(z), \tau_2(z)) \right)\\
= & \, \E\left(\max_{(i,j) \in \haken r \times \haken s} X(i,j) \right)
= \E\left(\max_{i \in \haken r}\max_{i \in \haken s} X(i,j) \right)\text.
\end{align*}
Somit ist auch (2) gezeigt.
\end{proof}

\begin{korollar}{Satz von Fernique}{SatzVonFernique}
Seien $n \in \N$, $\gamma, \delta : \Omega \rightarrow \R^n$ mit
\begin{align*}
\forall\, i \in \haken n \, \exists\, \rho_i \in \R_{>0}:
\gamma_i \text{ ist } N(0, \rho_i^2)\text{--verteilt,}\\
\forall\, i \in \haken n \, \exists\, \sigma_i \in \R_{>0}:
\delta_i \text{ ist } N(0, \sigma_i^2)\text{--verteilt}
\end{align*}
und es gelte: $\forall\, i, j \in \haken n: \norm{\gamma_i - \gamma_j}_2 \le
\norm{\delta_i - \delta_j}_2$.\\
Dann gilt:
\[\E\left(\max_{i \in \haken n}\gamma_i \right)
\le \E\left(\max_{i \in \haken n}\delta_i \right)\text.\]
\end{korollar}
\begin{proof}
Seien $r\coloneqq  1$, $s \coloneqq  n$ und
\begin{align*}
X: \haken r \times \haken s \rightarrow L_2(\Omega, \mfa, \Prim, \R),
(i, j) \mapsto \gamma_j\text,\\
Y: \haken r \times \haken s \rightarrow L_2(\Omega, \mfa, \Prim, \R),
(i, j) \mapsto \delta_j\text.
\end{align*}
Dann sind $X, Y$ offenbar wohldefiniert.\\
Ferner gilt für alle $i \in \haken r$ und alle $k, l \in \haken s$:
\begin{align*}
\norm{X(i,k)-X(i,l)}_2 &= \norm{\gamma_k - \gamma_l}_2 \overset{\text{Vor.}}\le
\norm{\delta_k - \delta_l}_2 = \norm{Y(i,k)- Y(i,l)}_2\text.
\end{align*}
Außerdem gilt wegen $r = 1$ trivial:
\[\forall j_1, j_2 \in \haken r \forall a, b \in \haken s:
j_1\neq j_2\Rightarrow\norm{X(j_1, a)-X(j_2,b)}_2\ge 
\norm{Y(j_1, a)-Y(j_2,b)}_2 \text.\]
Damit sind die Voraussetzungen von \ref{GaussscherMinMaxSatz} für $X,Y$
erfüllt und es gilt:
\begin{align*}
\E\left(\max_{j \in \haken n}\gamma_j \right) &=
\E\left(\max_{j \in \haken n}X(1, j) \right) \overset{r=1}=
\E\left(\min_{j \in \haken r}\max_{j \in \haken n}X(i, j) \right)\\
&\overset{\ref{GaussscherMinMaxSatz}}\le
\E\left(\min_{j \in \haken r}\max_{j \in \haken n}Y(i, j) \right)
= \E\left(\max_{j \in \haken n}Y(1, j) \right)
= \E\left(\max_{j \in \haken n}\delta_j \right)\text.
\end{align*}
\end{proof}

Obwohl der Gauß'sche Min--Max--Satz keinerlei Aussagen über die Schärfe der
Ungleichung trifft, gibt es durchaus interessante Fälle, in denen diese
optimal ist.
In dem folgenden Korollar kommt noch hinzu, dass die
Beweisstrategie vermuten lässt, man würde grob schätzen, während die erhaltene
Abschätzung allerdings tatsächlich optimal ist.
Das folgende Korollar ist eine
Modifikation einer in \cite{Gor2} vorgenommenen Nebenrechnung.

\begin{korollar}{Maximum zentrierter normalverteilter Zufallsgrößen}
             {MaximumZentrierterNormalverteilterZufallsgroessen}
Seien $n \in \N$, $\sigma \in \left(\R_{>0}\right)^n$,
$\gamma: \Omega \rightarrow \R^n$, $\delta: \Omega \rightarrow \R^n$ mit
$\delta$ unabhängig und
\[\forall\, i \in \haken n: \gamma_i, \delta_i \text{ sind }
N(0,\sigma_i^2)\text{--verteilt.}\]
Dann gilt:
\[\E\left(\max_{i \in \haken n}\gamma_i \right)
\le \sqrt 2 \E\left(\max_{i \in \haken n}\delta_i\right)\]
und diese Ungleichung ist optimal.
\end{korollar}
\begin{proof}
Seien $i, j \in \haken n$.\\
Dann gilt $\sigma_i^2 - 2\sigma_i\sigma_j +\sigma_j^2
= (\sigma_i - \sigma_j)^2 \ge 0$ und damit $2 \sigma_i\sigma_j \le
\sigma_i^2 + \sigma_j^2$, also
\begin{align}
(\sigma_i + \sigma_j)^2 = \sigma_i^2 + 2 \sigma_i\sigma_j + \sigma_j^2
\le \sigma_i^2 + \sigma_i^2 + \sigma_j^2 + \sigma_j^2 
= 2(\sigma_i^2 + \sigma_j^2) \text. \label{mzngeq:1}
\end{align}
Außerdem gilt:
\begin{align*}
\norm{\gamma_i - \gamma_j}_2^2 
&\le \left(\norm{\gamma_i}_2 + \norm{\gamma_j}_2 \right)^2
= (\sigma_i + \sigma_j)^2 \overset{\eqref{mzngeq:1}}\le 
2(\sigma_i^2 + \sigma_j^2)\\
&= 2\left(\norm{\delta_i}_2^2 + \norm{\delta_j}_2^2 \right)
\overset{\delta \text{ unabh.}}=
2 \left(\norm{\delta_i}_2^2 + \norm{\delta_j}_2^2 
- 2 \sprod{\delta_i}{\delta_j}_{L_2} \right)\\
&= 2\norm{\delta_i - \delta_j}_2^2 =
\norm{\sqrt 2 \delta_i - \sqrt 2 \delta_j}_2^2\text.
\end{align*}
Offenbar ist für jedes $k\in \N$ die Abbildung $\sqrt 2\delta_k$
$N(0, 2\sigma_k^2)$--verteilt. Damit sind die Voraussetzungen von
\ref{SatzVonFernique} für $\gamma$ und $\sqrt 2 \delta$ erfüllt, sodass gilt:
\[\E\left(\max_{i \in \haken n}\gamma_i \right)
\overset{\ref{SatzVonFernique}}\le 
\E\left(\max_{i \in \haken n}\sqrt 2\delta_i\right)
\overset{\sqrt 2>0}=\sqrt 2\E\left(\max_{i \in \haken n}\delta_i\right)\text.\]
Für den Beweis der Optimalität seien $Y_1, Y_2$ stochastisch unabhängige
$N(0,1)$--verteilte Zufallsgrößen, sowie $X_1 :=Y_1$ und $X_2:=-Y_1$.\\
Dann sind auch $X_1, X_2$ $N(0,1)$--verteilt, aber nicht stochastisch
unabhängig.\\
Offensichtlich ist $\max\meng{X_1,X_2} = \max\meng{Y_1, -Y_1} = \abs {Y_1}$
und damit
\begin{align}
\E\left(\max_{i \in \haken 2}X_i \right) = \E\left(\abs {Y_1} \right) 
= \sqrt{\frac 2 \pi}\label{mzngeq:2}\text.
\end{align}
Nach \cite{K2} (dort Abschnitt 6.6) gilt:
\begin{align}
\int_{-\infty}^\infty \exp(-t^2)dt = 2 \int_0^\infty \exp(-t^2)dt = \sqrt \pi
\label{mzngeq:3}\text.
\end{align}
Seien $\Phi$ die Verteilung und $\vi$ die Dichte von $N(0,1)$, sowie
$f$ die Dichte von $\max\meng{Y_1, Y_2}$.\\
Nach \refclaim{EigenschaftenVonOrdnungsstatistiken}
{EigenschaftenVonOrdnungsstatistiken4} gilt dann: $f = 2\Phi\vi$,
sodass wir wegen $\forall\, t \in \R: \vi'(t) = -t\vi(t)$ erhalten:
\begin{align}
\E\left(\max_{i \in \haken 2}Y_i \right)&=
\int_{-\infty}^\infty t f(t)dt = 2\int_{-\infty}^\infty \Phi(t)t\vi(t)dt
= - 2\int_{-\infty}^\infty \Phi(t)\vi'(t)dt\notag\\
&= -2\left([\Phi(t)\vi(t)]_{t=-\infty}^\infty - 
\int_{-\infty}^\infty \vi(t)^2dt \right)
= 2 \int_{-\infty}^\infty \frac 1 {2\pi}\exp\left(-t^2\right)dt\notag\\
&= \frac 1 \pi \int_{-\infty}^\infty \exp\left(-t^2\right)dt
\overset{\eqref{mzngeq:3}}= \frac 1 \pi \sqrt \pi = \sqrt{\frac 1 \pi}\text.
\label{mzngeq:4}
\end{align}
Also
\[\E\left(\max_{i \in \haken 2}X_i \right) \overset{\eqref{mzngeq:2}}= 
\sqrt 2 \sqrt{\frac 1 \pi} \overset{\eqref{mzngeq:4}}= 
\sqrt 2 \E\left(\max_{i \in \haken 2}Y_i \right)\text.\]
\end{proof}

\chapter{Gordon'sche Paare und der Gauß'sche Min--Max--Satz}
\label{ch:kapitel3}

In diesem Kapitel führen wir Gordon'sche Paare ein, die nichts
weiter als spezielle unendliche Familien normalverteilter Zufallsgrößen sind.
Neben einigen einfachen Eigenschaften Gordon'scher Paare werden wir vor allem
nachweisen, dass diese Paare ebenfalls die Anwendung des Gauß'schen
Min--Max--Satzes zulassen.
Die Namensgebung geht auf Yehoram Gordon zurück, der
diese Familien in \cite{Gor} (bzw. abgewandelt in \cite{Gor2}) einführte.
\parag
Für dieses Kapitel seien $(\Omega, \mfa, \Prim)$ ein Wahrscheinlichkeitsraum
$m, N \in \N$.\\
Seien $\gamma: \Omega \rightarrow \R^N$, $\delta: \Omega \rightarrow \R^m$
und $\widetilde g: \Omega \rightarrow \R^{mN}$ jeweils $N(0,1)$--verteilt,
$\widetilde g$ sei unabhängig und $\gamma, \delta$ seien vollständig
unabhängig.
Sei weiter $g\coloneqq  \Mat_{m,N} \circ \widetilde g$.\\
Wir werden in diesem Kapitel für $\omega \in \Omega$ die Matrix $g(\omega)$
stets mit der linearen Abbildung $t \mapsto g(\omega)t$ identifizieren.

\begin{lemma}{Integrierbarkeit und Verteilung von Linearkombinationen}
             {IntegrierbarkeitUndVerteilungVonLinearkombinationen}
Seien $t \in \R^m$, $v \in \R^N$, $c \in \R$.\\
Dann gilt:
\begin{enumerate}[(1)]
\item\label{IntegrierbarkeitUndVerteilungVonLinearkombinationen1}
$c\bprod t \delta_m + \bprod v\gamma_N, \, 
\sprod \cdot v_N \circ ev(t, \cdot) \circ g \in L_2(\Omega,\mfa,\Prim,\R)$.
\item\label{IntegrierbarkeitUndVerteilungVonLinearkombinationen2}
$c\bprod t \delta_m + \bprod v\gamma_N = 0 \iff (ct = 0 \wedge v = 0)$.
\item\label{IntegrierbarkeitUndVerteilungVonLinearkombinationen3}
$\sprod \cdot v_N \circ ev(t, \cdot) \circ g = 0 \iff (t = 0 \vee v= 0)$.
\item\label{IntegrierbarkeitUndVerteilungVonLinearkombinationen4}
Ist $ct \neq 0$ oder $v \neq 0$, so gilt:
\[c\bprod t \delta_m + \bprod v\gamma_N \quad{ist}\quad
N\left(0, c^2\norm t_{2,m}^2 +\norm v_{2, N}^2\right)\text{--verteilt.}\]
\item\label{IntegrierbarkeitUndVerteilungVonLinearkombinationen5}
Ist $t \neq 0\neq v $, so gilt:
$\sprod \cdot v_N \circ ev(t, \cdot) \circ g$ ist
$N\left(0, \norm t _{2,m}^2\norm v _{2, N}^2\right)$--verteilt.
\end{enumerate}
\end{lemma}
\begin{proof}
(1) $c\bprod t \delta_m + \bprod v\gamma_N$ ist als Linearkombination
quadratintegrierbarer Funktionen selbst quadratintegrierbar.\\
Ferner gilt f�r alle $\omega \in \Omega$:
\[\bigg(\sprod \cdot v_N \circ ev(t, \cdot) \circ g\bigg)(\omega) =
\sprod {g(\omega)t}v_N = \sum_{i=1}^N\sum_{j=1}^mt_jv_i(g_{ij}(\omega))\text.\]
Damit ist\[
\sprod\cdot v_N\circ ev(t,\cdot)\circ g=\sum_{i=1}^N\sum_{j=1}^mt_jv_ig_{ij}
\text,\]
also ebenfalls als Linearkombination quadratintegrierbarer Abbildungen
quadratintegrierbar.\\
(2) Da $\gamma, \delta$ vollst�ndig unabh�ngig und zentriert sind, gilt
\[\forall\,i\in\haken m\,\forall\,j\in\haken N:\E(\delta_i\gamma_j)=0\text.\]
Da $\E$ das Skalarprodukt auf $L_2(\Omega, \mfa, \Prim, \R)$ erzeugt, sind also
die Komponentenfunktionen von $\gamma$ und $\delta$ paarweise orthogonal und
damit linear unabh�ngig. Damit gilt:
\begin{align*}
c\bprod t \delta _m + \bprod v \gamma_N = 0 
& \iff \left(\sum_{i=1}^mct_i\delta_i \right) 
     + \left(\sum_{j=1}^Nv_j\gamma_j \right) = 0\\
& \iff(\forall\,i\in \haken m: ct_i = 0) \wedge (\forall\,j\in\haken N:v_j=0)\\
& \iff ct = 0 \wedge v = 0\text.
\end{align*}
(3) Weil $\widetilde g$ unabh�ngig und zentriert ist, sind die
Komponentenfunktionen von $\widetilde g$ und damit auch jene von $g$ paarweise
orthogonal, also linear unabh�ngig. Somit folgt wie im Beweis von (1):
\begin{align*}
 \sprod \cdot v_N \circ ev(t, \cdot) \circ g 
 \overset{\substack{\text{Bew.}\\ \text{von }(1)}}=
  \sum_{i=1}^N\sum_{j=1}^mt_jv_ig_{ij} = 0
&\iff \forall\, i \in \haken N \forall\, j \in \haken m: t_jv_i = 0\\
& \iff  t= 0 \vee v = 0\text.
\end{align*}
(4)/(5) Sind $X_1,X_2$ stochastisch unabh�ngige Zufallsgr��en, $a_1,a_2\in \R$
und $\sigma_1,\sigma_2\in\R_{>0}$ derart, dass $X_i$
$N(a_i,\sigma_i^2)$--verteilt sind f�r alle $i \in \haken 2$, so ist nach \cite{Irle} (dort 10.15) $X_1 + X_2$
$N(a_1+a_2, \sigma_1^2 + \sigma_2^2)$--verteilt.\\
Seien $c\in \R\setminus\meng{0}$ und $t \in \R$. Dann gilt falls $c > 0$:
\begin{align*}
\Prim(cX_1 \le t) &= \Prim(X_1 \le \nicefrac t c) 
= \frac 1{\sqrt{2\pi\sigma_1^2}}\int_{-\infty}^{\nicefrac t c}
\exp\left(- \frac 1 2 \left(\frac{s-a_1}{\sigma_1}\right)^2 \right)ds\\
&\overset{\substack{\text{Subst.}\\s \mapsto \frac s c}}=
\frac 1{\sqrt{2\pi\sigma_1^2}}\int_{-\infty}^t 
\exp\left(-\frac 1 2 \left(\frac{\frac s c - a_1}{\sigma_1} \right)^2
 \right)\frac 1 c ds\\
 &=\frac 1 {\sqrt{2 \pi (c\sigma_1)^2}}\int_{-\infty}^t 
 \exp\left(-\frac 12\left(\frac{s-(ca_1)}{(c\sigma_1)}\right)^2 \right)ds
\end{align*}
und falls $c < 0$ wegen der Stetigkeit von $N(a_1, \sigma_1^2)$:
\begin{align*}
&\,\Prim(cX_1 \le t) = \Prim(X_1 \ge \nicefrac tc) 
= 1 - \Prim(X_1 \le \nicefrac t c)\\
=&\,
\frac 1 {\sqrt{2\pi\sigma_1^2}}\int_{\frac t c}^\infty
\exp\left(- \frac 1 2 \left(\frac{s-a_1}{\sigma_1}\right)^2 \right)ds\\
\overset{\substack{\text{Subst.}\\s \mapsto \frac s c}}=&\,
\frac 1 {\sqrt{2\pi\sigma_1^2}}\int_{ t }^{-\infty}
\exp\left(-\frac 12\left(\frac{\frac sc-a_1}{\sigma_1}\right)^2\right)
\frac 1 c ds\\
=&\, \frac 1{\sqrt{2\pi\sigma_1^2}}\frac 1{\abs c}
\int_{-\infty}^t \exp\left(-\frac 12\left(\frac{ s-(ca_1)}{c\sigma_1}\right)^2\right)ds\text.
\end{align*}
Also ist $cX_1$ auch $N(ca_1, (c\sigma_1)^2)$--verteilt.\\
Damit folgen die Aussagen (4) und (5) durch einfache Induktionen.
\end{proof}

Wir halten fest, dass für alle $x \in E^N$ und $\omega \in \Omega$ die
Abbildungen
\[m_{x,\omega} : \R^m \rightarrow E, t \mapsto \bprod{g(\omega)t}{x}_N\]
durch die Darstellung $m_{x,\omega} = \bprod\cdot x_N \circ g(\omega)$ als
Komposition stetiger linearer Abbildungen selbst stetig und linear sind
(vgl. \ref{StetigkeitDesLinearkombinators}).\\
Dies motiviert die folgende Definition (basierend auf \cite{Gor},
\cite{Gor2}).

\begin{definition}{Gaufl'scher Operator}{GaussscherOperator}
Seien $x \in E^N$. Dann bezeichnen wir die Abbildung
\[G_x : \Omega \rightarrow L(\R^m, E), \quad
        \omega \mapsto \bprod\cdot x_N \circ g(\omega)\]
als \textbf{Gaufl'schen Operator zum Tupel $(x,g)$}.
\end{definition}

Der Gauß'sche Operator wird später verwendet, um die Existenz gewisser linearer
Abbildungen zuzusichern, weil er sich als operatorwertige Zufallsvariable
auffassen lässt.\parag
Wir beobachten, dass für alle $t \in \R^m$, $v \in \R^N$ und $c \in \R$ die
Abbildungen $c\bprod t \delta _m + \bprod v \gamma_N$ und
$\sprod \cdot v _N \circ ev(t, \cdot) \circ g$ nach
\refclaim{IntegrierbarkeitUndVerteilungVonLinearkombinationen}
{IntegrierbarkeitUndVerteilungVonLinearkombinationen1} quadratintegrierbar
sind.
Damit ist die folgende Setzung wohldefiniert.

\begin{definition}{Gordon'sches Paar}{GordonschesPaar}
Sei $c \in \R$. Definiere
\begin{align*}
X&: \R^m \times \R^N \rightarrow L_2(\Omega, \mfa, \Prim, \R),
    (t, v) \mapsto c\bprod t \delta _m + \bprod v \gamma_N\text,\\
Y&: \R^m \times \R^N \rightarrow L_2(\Omega, \mfa, \Prim, \R),
    (t, v) \mapsto \sprod \cdot v_N \circ ev(t, \cdot)\circ g\text.
\end{align*}
Dann heißt $(X, Y)$ \textbf{das Gordon'sche Paar zum Tupel}
$(\gamma, \delta, g, c)$.\\
Eine unmittelbare Folgerung aus dieser Definition ist, dass für alle
$t \in \R^m$ und alle $v \in \R^N$ gilt:
$\Bild X(t,v)\subseteq \R$ und $\Bild Y(t, v)\subseteq \R$.\\
Ebenso evident ist
$\forall\, (t, v) \in \R^m\times \R^N\, \forall\, \omega \in \Omega:
Y(t, v)(\omega) = \sprod{g(\omega)t}v_N$.
\end{definition}
\parag
Als nächstes notieren einige elementare Eigenschaften Gordon'scher Paare.

\begin{lemma}{Eigenschaften Gordon'scher Paare}
             {EigenschaftenGordonscherPaare}
Seien $c \in \R$ und $(X, Y)$ das Gordon'sche Paar zu $(\gamma, \delta,g,c)$.\\
Der Raum $\R^m \times \R^N$ sei mit der Produktnorm ausgestattet.\\
Dann gilt:
\begin{enumerate}[(1)]
	\item \label{EigenschaftenGordonscherPaare1}
	$\forall\, t \in \R^m\, \forall\, v \in \R^N:
	\E(X(t,v)) = 0 = \E(Y(t,v))$.
	\item \label{EigenschaftenGordonscherPaare2}
	$\forall\, t \in \R^m\, \forall\, v \in \R^N: \norm{X(t,v)}_2^2 
	= c^2\norm t _{2,m}^2 + \norm v _{2, N}^2$.
	\item \label{EigenschaftenGordonscherPaare3}
	$\forall\, t \in \R^m\, \forall\, v\in \R^N:
	\norm{Y(t, v)}_2^2 = \norm t_{2,m}^2\norm v_{2, N}^2$.
	\item \label{EigenschaftenGordonscherPaare4}
	$X$ ist linear und stetig.
	\item\label{EigenschaftenGordonscherPaare5}
	$Y$ ist bilinear und stetig.
\end{enumerate}
\end{lemma}
\begin{proof}
(1) Seien $t \in \R^m$ und $v \in \R^N$.\\
Falls $ct\neq 0$ oder $v \neq 0$, ist $X(t,v)$ nach
\refclaim{IntegrierbarkeitUndVerteilungVonLinearkombinationen}
         {IntegrierbarkeitUndVerteilungVonLinearkombinationen4}
$N(0, c^2\norm t_{2, m}^2 + \norm v_{2,N}^2)-$verteilt, hat also Erwartungswert
0 und Varianz $c^2\norm t_{2, m}^2 + \norm v_{2,N}^2$ 
(siehe 8.8, 9.2 in~\cite{Irle}).\\
Ist $ct = 0$ und $v = 0$, so ist $X(t, v) = 0$ und hat daher trivialerweise den
Erwartungswert 0 und auch Varianz 
$0 = c^2\norm t_{2, m}^2 + \norm v_{2,N}^2$.\\
Somit ist $\E(X(t, v)) = 0$ und daher
$\norm{X(t,v)}_2^2 = c^2\norm t_{2, m}^2 + \norm v_{2,N}^2$ (vgl. Bemerkung vor~\ref{ElementareEigenschaftenQuadratintegrierbarerAbbildungen}).\\
Falls $t \neq 0$ und $v \neq 0$, so ist $Y(t,v)$ nach~\refclaim{IntegrierbarkeitUndVerteilungVonLinearkombinationen}
         {IntegrierbarkeitUndVerteilungVonLinearkombinationen5}
$N(0, \norm t _{2,m}^2 \norm v_{2, N}^2)-$verteilt, hat also Erwartungswert 0
und Varianz $\norm t_{2,m}^2\norm v_{2,N}^2$.\\
Falls $t = 0$ oder $v = 0$, so ist $Y(t, v) = 0$ und hat wiederum trivial
Erwartungswert 0 und Varianz $0=\norm t _{2,m}^2 \norm v_{2, N}^2$.\\
Also ist $E(Y(t, v))= 0$ und 
somit $\norm{Y(t, v)}_2^2 = \norm t _{2,m}^2 \norm v_{2, N}^2$ (mit dem
gleichen Argument, wie bei $X(t, v)$).\\
Damit sind (1), (2) und (3) gezeigt.\\
(4) Nach~\refclaim{EigenschaftenDesLinearkombinators}
{EigenschaftenDesLinearkombinators1}
sind \[c\bprod\cdot \delta_m:\R^m \rightarrow L_2(\Omega,\mfa, \Prim, \R)
\quad \text{und}\quad 
\bprod \cdot \gamma_N:\R^N \rightarrow L_2(\Omega,\mfa, \Prim,\R)\] linear.
Daher ist deren Summe linear im Tupel der beiden Komponenten.\\
Als lineare Abbildung mit endlichdimensionalem Definitionsraum ist $X$
stetig.\\
(5) Seien $t,s\in \R^m$, $v, w \in \R^N$, $a \in \R$ und $\omega \in \Omega$.\\
Dann gilt wegen der Linearität von $g(\omega)$:
\begin{align*}
Y(t+as, v)(\omega)& = \sprod{g(\omega)(t+as)}v_N 
= \sprod{g(\omega)t + ag(\omega)s}v_N\\
&= \sprod{g(\omega)t}v_N + a\sprod{g(\omega)s}v_N =
Y(t, v)(\omega)+ a Y(s, v)(\omega)\\
&= (Y(t, v) + a Y (s, v))(\omega)\text.
\end{align*}
Außerdem ist
\begin{align*}
Y(t, v+aw)(\omega) &= \sprod{g(\omega)t}{v+aw}_N
= \sprod{g(\omega)t}v_N+a\sprod{g(\omega)t}w_N\\
&= Y(t, v)(\omega) + a Y(t, w)(\omega)
= (Y(t, v) + a Y(t, w))(\omega)\text.
\end{align*}
Somit ist $Y$ bilinear.\\
Wegen (3) ist $\norm{Y(t, v)}_2 = \norm t_{2,m}\norm v_{2,N}$ und da $Y$
bilinear ist, ist es nach der bekannten Charakterisierung der Stetigkeit
bilinearer Abbildungen stetig. Auch hier könnte man alternativ die endliche
Dimension des Definitionsraums nutzen, um die Stetigkeit von $Y$ einzusehen:
$\R^m, \R^N$ und $L_2(\Omega, \mfa, \Prim, \R)$ sind vollständig, $Y$ ist
komponentenweise als lineare Abbildung mit endlichdimensionalem
Definitionsraum stetig, also ist nach dem Satz von Banach--Steinhaus auch
$Y$ stetig.
\end{proof}
\parag
Für die Prüfung der Voraussetzungen im Gauß'schen Min--Max--Satz sind die
$L_2$--Abstände innerhalb Gordon'scher Paare von besonderem Interesse.
Mit dem eben bewiesenen Lemma können diese einfach bestimmt werden.
Diese Abstände wurden bereits von Gordon in \cite{Gor}, \cite{Gor2} in modifizierter Weise bestimmt.
\newpage
\begin{lemma}{Abstände in Gordon'schen Paaren}
             {AbstaendeInGordonschenPaaren}
Seien $c \in \R$ und $(X, Y)$ das Gordon'sche Paar zum Tupel 
$(\gamma, \delta, g, c)$.\\
Dann gilt:
\begin{enumerate}[(1)]
	\item \label{AbstaendeInGordonschenPaaren1}
	$\forall\, t, s \in \R^m\, \forall\, v, w \in \R^N: 
	\norm{X(t, v) - X(s, w)}_2^2 = c^2\norm{t-s}_{2,m}^2 + \norm{v-w}_{2,N}^2$
	\item \label{AbstaendeInGordonschenPaaren2}
	$\forall\, t, s \in \R^m\, \forall\, v, w \in \R^N:$
	\[\norm{Y(t, v) - Y(s, w)}_2^2 
	= \norm t _{2,m}^2\cdot\norm v_{2,N}^2 + \norm s_{2,m}^2\cdot\norm w_{2,N}^2
	  - 2\sprod t s _m\sprod v w_N\]
	\item \label{AbstaendeInGordonschenPaaren3}
	$\forall\, t, s \in S^{m-1}\, \forall\, v, w \in \R^N:$
	\[\norm{X(t, v) - X(s, w)}_2^2 - \norm{Y(t, v) - Y(s, w)}_2^2
	= 2(1-\sprod t s _m)(c^2-\sprod v w_N)\]
	\item \label{AbstaendeInGordonschenPaaren4}
	$\forall\, t \in S^{m-1}\, \forall\, v, w \in \R^N: 
	\norm{X(t, v) - X(t, w)}_2^2 = \norm{Y(t, v) - Y(t, w)}_2^2$
	\item \label{AbstaendeInGordonschenPaaren5}
	$\forall\, t, s \in S^{m-1}\, \forall\, v, w \in \kabg {\ell_2^N}0{\abs c}:$
	\[\norm{Y(t, v) - Y(s, w)}_2^2 \le \norm{X(t, v) - X(s, w)}_2^2\]
\end{enumerate}
\end{lemma}
\begin{proof}
(1) Seien $t, s \in \R^m$ und $v, w \in \R^N$.\\
Dann gilt:
\begin{align*}\norm{X(t, v) - X(s, w)}_2^2
&\overset{\refclaim{EigenschaftenGordonscherPaare}
                  {EigenschaftenGordonscherPaare4}}=
\norm{X(t-s, v-w)}_2^2 \\
&\overset{\refclaim{EigenschaftenGordonscherPaare}
                  {EigenschaftenGordonscherPaare2}}=
         c^2\norm{t - s}_{2,m}^2 + \norm{v-w}_{2, N}^2\text.
\end{align*}
(2) Seien $t, s \in \R^m$ und $v, w \in \R^N$.\\
Da $\widetilde g$ unabhängig und $N(0,1)-$verteilt ist, gilt
\begin{align}
\forall\, i, k \in \haken N\, \forall\, j, l \in \haken m:
\E(g_{ij}g_{kl}) \neq 0 \iff (i,j)=(k,l) \label{aigpeq:1}
\end{align}
und
\begin{align}
\forall\, i \in \haken N\, \forall\, j \in \haken m:
\E(g_{ij}^2) = 1 \text.\label{aigpeq:2}
\end{align}
Somit gilt mit der ausgeschriebenen Darstellung von $Y$ (vgl.
\ref{IntegrierbarkeitUndVerteilungVonLinearkombinationen}):
\begin{align}
\E(Y(t,v)Y(s,w)) & = 
\E\left(\left(\sum_{i=1}^N\sum_{j=1}^m t_jv_ig_{ij} \right)
\left(\sum_{k=1}^N\sum_{l=1}^m s_lw_kg_{kl} \right) \right)\notag\\
&= \sum_{i=1}^N\sum_{j=1}^m\sum_{k=1}^N\sum_{l=1}^mt_jv_is_lw_k
\E(g_{ij}g_{kl})
\overset{\substack{\eqref{aigpeq:1},\\\eqref{aigpeq:2}}}=
\sum_{i=1}^N\sum_{j=1}^mt_jv_is_jw_i\notag\\
&= \sum_{i=1}^N v_iw_i\sum_{j=1}^mt_js_j = \sprod t s _m \sprod v w_N\text.
\label{aigpeq:3}
\end{align}
Daher erhält man
\begin{align*}
\norm{Y(t, v) - Y(s, w)}_2^2 & = \norm{Y(t,v)}^2_2 + \norm{Y(s, w)}_2^2 -
2\E(Y(t,v)Y(s,w))\\
&\overset{\eqref{aigpeq:3}}= \norm{Y(t,v)}^2_2 + \norm{Y(s, w)}_2^2 -
2\sprod t s _m \sprod v w _N\\
&\overset{\refclaim{EigenschaftenGordonscherPaare}
{EigenschaftenGordonscherPaare3}}=
\norm t _{2,m}^2\norm v_{2, N}^2 + \norm s _{2, m}^2 \norm w_{2, N}^2
- 2 \sprod t s _m \sprod v w _N\text.
\end{align*}
(3) Seien $t, s \in S^{m-1}$ und $v, w \in \R^N$.\\
Es gilt:
\begin{align}
\norm{v-w}_{2, N}^2-\norm v_{2, N}^2 - \norm w_{2, N}^2
= - 2 \sprod v w_N\text.\label{aigpeq:4}
\end{align}
Damit erhält man:
\begin{align*}
&\norm{X(t, v) - X(s, w)}_2^2 - \norm{Y(t, v) - Y(s, w)}_2^2\\
\overset{(2)}=&
\,\norm{X(t, v) - X(s, w)}_2^2- \norm t _{2,m}^2\norm v_{2,N}^2
 - \norm s_{2,m}^2\norm w_{2,N}^2
	  + 2\sprod t s _m\sprod v w_N\\
=&\, \norm{X(t, v) - X(s, w)}_2^2 - \norm v_{2,N}^2 -
 \norm w_{2,N}^2+ 2\sprod t s _m\sprod v w_N\\
\overset{(1)}=&\,c^2\norm{t-s}_{2,m}^2 + \norm{v-w}_{2,N}^2 - \norm v_{2,N}^2 -
 \norm w_{2,N}^2+ 2\sprod t s _m\sprod v w_N\\
\overset{\eqref{aigpeq:4}}=& \,c^2\norm{t-s}_{2,m}^2 - 2 \sprod v w_N
 + 2\sprod t s _m\sprod v w_N\\
=&\, c^2(\norm t _{2,m}^2 + \norm s_{2, m}^2 - 2\sprod t s_m)- 2 \sprod v w_N
+ 2\sprod t s _m\sprod v w_N\\
=&\, c^2(2 -2\sprod t s_m) - 2 \sprod v w_N + 2\sprod t s _m\sprod v w_N\\
=&\, 2c^2(1-\sprod t s _m) - 2 \sprod v w_N(1-\sprod t s_m)
= 2(1-\sprod t s_m)(c^2-\sprod v w_N)\text.
\end{align*}
(4) Seien $t \in S^{m-1}$ und $v, w \in \R^N$. Dann gilt:
\begin{align*}
&\norm{X(t, v) - X(t, w)}_2^2 - \norm{Y(t, v) - Y(t, w)}_2^2
\overset{(3)}= 2(1-\sprod t t_m)(c^2- \sprod v w_N)\\
=&\, 2(1-\norm t_{2,m}^2)(c^2-\sprod v w_N)= 0\text.
\end{align*}
Also $\norm{X(t, v) - X(s, w)}_2^2 = \norm{Y(t, v) - Y(s, w)}_2^2$.\\
(5) Seien $t, s \in S^{m-1}$ und $v, w \in \kabg{\ell_2^N}0 {\abs c}$.\\
Wegen Cauchy--Schwarz ist dann
$\sprod t s _m \le \norm t _{2,m} \norm s_{2,m} = 1$ und analog dazu
$\sprod v w_m \le \norm v_{2,N}\norm w_{2, N} \le \abs c^2 = c^2$.\\
Also ist $(1-\sprod t s_m) \ge 0$ und $(c^2- \sprod v w_N) \ge 0$.\\
Somit
\[\norm{X(t, v) - X(s, w)}_2^2 - \norm{Y(t, v) - Y(s, w)}_2^2
\overset{(3)}= 2\underbrace{(1-\sprod t s_m)}_{\ge 0}
\underbrace{(c^2 \sprod v w_N)}_{\ge 0} \ge 0 \]
und daher $\norm{Y(t, v) - Y(s, w)}_2^2 \le \norm{X(t, v) - X(s, w)}_2^2$.
\end{proof}
\parag
Da das \enquote{Min--Max} aus dem Gauß'schen Min--Max--Satz für Gordon'sche Paare
sinnvollerweise zu einem \enquote{Inf--Sup} umgewandelt werden sollte, brauchen wir
die Beschränktheit gewisser Mengen, welche wir in dem nachfolgenden Lemma
notieren.

\begin{lemma}{Beschränktheit Gordon'scher Paare}
             {BeschraenktheitGordonscherPaare}
Seien $c \in \R$ und $(X, Y)$ das Gordon'sche Paar zu $(\gamma,\delta,g,c)$.\\
Seien $\emptyset \neq T \subseteq \R^m$ beschränkt,
$\emptyset \neq V \subseteq \R^N$ beschränkt und $\omega \in \Omega$.\\
Dann gilt:
\begin{enumerate}[(1)]
	\item \label{BeschraenktheitGordonscherPaare1}
	$\forall\, s \in \R^m\, \forall\, u \in \R^N: 
	\abs{X(s,u)(\omega)} \le \abs{c}\norm s_{2, m}\norm{\delta(\omega)}_{2,m}
	+ \norm u_{2, N} \norm{\gamma(\omega)}_{2, N}$\absatz
	$\forall\, s \in \R^m\, \forall\, u \in \R^N: \abs{Y(s,u)(\omega)}
	\le \opnorm{g(\omega)}{\R^m}{\R^N}\norm s_{2, m} \norm u_{2, N}$\absatz
	Insbesondere sind 
	$\menge{X(t, v)(\omega)}{t \in T, v \in V}$ und 
	$\menge{Y(t, v)(\omega)}{t \in T, v \in V}$ beschränkt.
  \item \label{BeschraenktheitGordonscherPaare2}
  Die folgenden Abbildungen sind wohldefiniert und $\R$--wertig:
  \[\inf_{t \in T}\sup_{v \in V}X(t, v), \  
    \sup_{t \in T}\sup_{v \in V}X(t, v), \ 
    \inf_{t \in T}\sup_{v \in V}Y(t, v), \ 
    \sup_{t \in T}\sup_{v \in V}Y(t, v)\text.\]
  \item \label{BeschraenktheitGordonscherPaare3}
  \begin{align*}
  \inf_{t \in T}\sup_{v \in V}X(t, v)(\omega)
  &= \left(\inf_{t \in T}(c\sprod t {\delta(\omega)} _m) \right)
  + \left(\sup_{v \in V}\sprod v {\gamma(\omega)}_N \right)\\
  &= \sup_{v \in V}\inf_{t \in T}X(t, v)(\omega)
  \end{align*}
  \item \label{BeschraenktheitGordonscherPaare4}
  \[\sup_{t \in T}\sup_{v \in V}X(t, v) = 
  \left(\sup_{t \in T}(c\sprod t {\delta(\omega)} _m) \right)
  + \left(\sup_{v \in V}\sprod v {\gamma(\omega)}_N \right)\]
\end{enumerate}
\end{lemma}
\begin{proof} Da $T$ beschränkt ist, gibt es ein $a \in \R_{>0}$ mit
$\forall\, t \in T: \norm t_{2, m} \le a$ und da auch $V$ beschränkt ist, gibt
es ein $b \in \R_{>0}$ mit $\forall\, v \in V: \norm v_{2, N} \le b$.\\
(1) Seien $t \in T$ und $v \in V$. Dann gilt:
\begin{align*}
\abs{X(t,v)(\omega)}&= \abs{c \sprod t {\delta(\omega)}_m 
+ \sprod v {\gamma(\omega)}_N}
\le \abs c\abs{\sprod t{\delta(\omega)}_m}+\abs{\sprod v {\gamma(\omega)}_N}\\
&\le \abs c \norm t_{2, m}\norm{\delta(\omega)}_{2,m}
 + \norm v_{2, N}\norm{\gamma(\omega)}_{2, N}\\
&\le \abs c a \norm{\delta(\omega)}_{2,m} + b\norm{\gamma(\omega)}_{2, N}
\end{align*}
und da $g(\omega)$ wegen der endlichen Dimension von $\R^m$ und $\R^N$ stetig
ist auch
\begin{align*}
\abs{Y(t, v)(\omega)} &= 
\abs{\sprod {g(\omega)t} v_N} \le \norm{g(\omega)t}_{2,N}\norm v_{2, N}\\
&\le \opnorm{g(\omega)} {\R^m}{\R^N} \norm t_{2, m}\norm v_{2, N}
\le \opnorm{g(\omega)} {\R^m}{\R^N} a b\text.
\end{align*}
Die Ungleichungen aus $(1)$ ergeben sich durch die jeweils vorletzte
Abschätzung.\\
Da $\norm{\gamma(\omega)}_{2, N}$, $\norm{\delta(\omega)}_{2,m}$ und
$\opnorm{g(\omega)} {\R^m}{\R^N}$ nach Voraussetzung endlich sind, sind die
beiden genannten Mengen beschränkt.\\
(2) Folgt unmittelbar aus der Beschränktheit der in (1) angegebenen Mengen.\\
(3)/(4) Sei $s \in T$. Dann ist
\begin{align*}\sup_{v \in V}X(s, v)(\omega) 
= \sup_{v\in V}\left(c\sprod s {\delta(\omega)}_m 
  + \sprod v{\gamma(\omega)}_N\right)
  = c\sprod s {\delta(\omega)}_m 
  + \sup_{v\in V}\sprod v{\gamma(\omega)}_N\text.\end{align*}
Somit ist
\begin{align*}
\inf_{t \in T}\sup_{v \in V}X(t, v)(\omega) &=
\inf_{t\in T}\left(c\sprod t {\delta(\omega)}_m 
  + \sup_{v\in V}\sprod v{\gamma(\omega)}_N \right)\\
  &=\left(\inf_{t \in T}(c\sprod t {\delta(\omega)} _m) \right)
  + \left(\sup_{v \in V}\sprod v {\gamma(\omega)}_N \right) \tag {$\ast$}\\
\sup_{t \in T}\sup_{v \in V}X(t, v)(\omega) 
 &=\sup_{t\in T}\left(c\sprod t {\delta(\omega)}_m 
  + \sup_{v\in V}\sprod v{\gamma(\omega)}_N \right)\\
&=\left(\sup_{t \in T}(c\sprod t {\delta(\omega)} _m) \right)
  + \left(\sup_{v \in V}\sprod v {\gamma(\omega)}_N \right)\text.
\end{align*}
Schließlich gilt auch
\begin{align*}
\sup_{v \in V}\inf_{t\in T}X(t, v)(\omega) &=
\sup_{v \in V}\inf_{t \in T}(c\sprod t {\delta(\omega)}_m 
  + \sprod v{\gamma(\omega)}_N)\\
  &=\sup_{v \in V}\left(\sprod v{\gamma(\omega)}_N + 
  \left(\inf_{t \in T}(c\sprod t {\delta(\omega)}_m)\right) \right)\\
  &= \left(\sup_{v \in V}\sprod v{\gamma(\omega)}_N \right) + 
  \left(\inf_{t \in T}(c\sprod t {\delta(\omega)}_m)\right)\\
  &\overset{(\ast)}= \inf_{t \in T}\sup_{v \in V}X(t, v)(\omega)\text.
\end{align*}
\end{proof}
\parag
Die in Kapitel 2 getroffenen Vorbereitungen für die Verallgemeinerung gewisser
Erwartungswertungleichungen wurde nur für abzählbare Mengen getroffen.
Allerdings ist $\R^n$ für alle $n \in \N$ mit beliebiger Norm separabel, sodass
es interessant zu wissen ist, wie sich Gordon'sche Paare auf dicht liegenden
Mengen verhalten.
Die Antwort ist nicht sehr überraschend -- die betrachteten
Infima bzw. Suprema stimmen darauf mit denen der ganzen Menge überein.

\begin{lemma}{Gordon'sche Paare auf dichten Teilmengen}
             {GordonschePaareAufDichtenTeilmengen}
Seien $c\in\R$ und $(X,Y)$ das Gordon'sche Paar zum Tupel
$(\gamma,\delta,g,c)$.\\
Seien $\emptyset \neq T \subseteq \R^m$ beschr�nkt und
$\emptyset \neq V \subseteq \R^N$ beschr�nkt, sowie $T' \subseteq T$
dicht in $T$ und $V' \subseteq V$ dicht in V.\\
Dann sind $T'\neq \emptyset$, $V' \neq \emptyset$ und es gilt:
\begin{enumerate}[(1)]
	\item \label{GordonschePaareAufDichtenTeilmengen1}
	$ \underset{v \in V}\sup\, \bprod v \gamma _N 
	 =\underset{v \in V'}\sup\,\bprod v \gamma_N$.
	\item \label{GordonschePaareAufDichtenTeilmengen2}
	 $\underset{t \in T}\inf\, c\bprod t \delta_m 
	 =\underset{t \in T'}\inf\, c\bprod t \delta_m\,$ und $\, 
	 \underset{t \in T} \sup\, c\bprod t \delta_m 
	  =\underset{t \in T'}\sup\, c\bprod t \delta_m$.\\
	  Insbesondere gilt
 \[\inf_{t \in T}\sup_{v \in V}X(t, v) = \inf_{t\in T'}\sup_{v \in V'}X(t,v)\, \text{ und }
 \,\sup_{t\in T}\sup_{v\in V}X(t,v)=\sup_{t\in T'}\sup_{v\in V'}X(t,v)\text.\]
	\item \label{GordonschePaareAufDichtenTeilmengen3}
	$\Inf{t \in T}\Sup{v \in V}Y(t, v) = \Inf{t\in T'}\Sup{v \in V}Y(t, v)\,$ und
	$\,\Sup{t \in T}\Sup{v \in V}Y(t, v) = \Sup{t\in T'}\Sup{v \in V}Y(t, v)$.
	\item \label{GordonschePaareAufDichtenTeilmengen4}
	$\Inf{t \in T}\Sup{v \in V}Y(t, v) = \Inf{t\in T'}\Sup{v \in V'}Y(t,v)\,$ und
	$\,\Sup{t \in T}\Sup{v \in V}Y(t, v) = \Sup{t\in T'}\Sup{v \in V'}Y(t, v)$.
\end{enumerate}
\end{lemma}
\begin{proof}
Da die leere Menge stets abgeschlossen ist und 
$\overline {T'}^T = T \neq \emptyset$, sowie 
$\overline{V'}^V= V \neq \emptyset$ gilt, sind $T'$ und $V'$ nichtleer.\\
Da $V$ beschr�nkt ist, gibt es $d \in \R_{>0}$ mit 
$\forall\,v \in V: \norm v_{2, N} \le d$.\\
Sei $\omega \in \Omega$.\\
(1) Die Abbildung $\sprod \cdot {\gamma(\omega)}_N$ ist lipschitzstetig und da
lipschitzstetige Bilder beschr�nkter Mengen beschr�nkt sind, ist die Menge
$\sprod \cdot {\gamma(\omega)}_N(V)$ beschr�nkt.\\
Wegen $V \neq \emptyset$ und $V'$ dicht in $V$, gilt:
\begin{align*}
\left(\sup_{v \in V}\bprod v \gamma_N\right)(\omega)
&= \sup_{v \in V}\sprod v {\gamma(\omega)}_N 
= \sup \left(\sprod \cdot {\gamma(\omega)}_N \right)(V) \\
&\overset{\refclaim{StetigeFunktionenUndDichteTeilmengen}
{StetigeFunktionenUndDichteTeilmengen1}}=
 \sup \left(\sprod \cdot {\gamma(\omega)}_N \right)(V')
 = \left(\sup_{v \in V'}\bprod v \gamma_N\right)(\omega)
\end{align*}
(2) Die Abbildung $c \sprod \cdot {\delta(\omega)}_m$ ist ebenfalls
lipschitzstetig und daher auch die Menge 
$(c \sprod \cdot {\delta(\omega)}_m)(T)$ nach oben und unten beschr�nkt.\\
Da $T \neq \emptyset$ und $T'$ dicht in $T$, folgt wie in (1):
\begin{align*}
\left(\inf_{t \in T}\bprod t \delta_m\right) (\omega)
&\overset{\refclaim{StetigeFunktionenUndDichteTeilmengen}
{StetigeFunktionenUndDichteTeilmengen2}}=
\left(\inf_{t \in T'}\bprod t \delta_m\right) (\omega)\text,\\
\left(\sup_{t \in T}\bprod t \delta_m\right) (\omega)
&\overset{\refclaim{StetigeFunktionenUndDichteTeilmengen}
{StetigeFunktionenUndDichteTeilmengen1}}=
= \left(\sup_{t \in T'}\bprod t \delta_m\right) (\omega)\text.
\end{align*}
Weiter gilt wegen der Beschr�nktheit von $T, T', V, V'$:
\begin{align*}
\inf_{t \in T'}\sup_{v \in V'}X(t, v) &
\overset{\refclaim{BeschraenktheitGordonscherPaare}
{BeschraenktheitGordonscherPaare3}}=
\inf_{t \in T'}c\bprod t \delta_m + \sup_{v \in V'}\bprod v \gamma_N\\
&\overset{(1),(2)}= 
\inf_{t \in T}c\bprod t \delta_m + \sup_{v \in V}\bprod v \gamma_N
\overset{\refclaim{BeschraenktheitGordonscherPaare}
{BeschraenktheitGordonscherPaare3}}=\inf_{t \in T}\sup_{v \in V}X(t, v)\text,\\
\sup_{t \in T'}\sup_{v \in V'}X(t, v) &
\overset{\refclaim{BeschraenktheitGordonscherPaare}
{BeschraenktheitGordonscherPaare4}}=
\sup_{t \in T'}c\bprod t \delta_m + \sup_{v \in V'}\bprod v \gamma_N\\
&\overset{(1),(2)}= 
\sup_{t \in T}c\bprod t \delta_m + \sup_{v \in V}\bprod v \gamma_N
\overset{\refclaim{BeschraenktheitGordonscherPaare}
{BeschraenktheitGordonscherPaare4}}=\sup_{t \in T}\sup_{v \in V}X(t, v)\text.
\end{align*}
(3)/(4) F�r jedes $w\in V$ ist die Abbildung $\sprod \cdot{g(\omega)^\top w}_m$
lipschitzstetig und linear und es gilt f�r alle $x \in \R^m$:
\begin{align}
\abs{\sprod x {g(\omega)^\top w}_m} &\le
\norm x_{2, m}\norm{g(\omega)^\top w}_{2, m} 
\le \norm x_{2, m} \opnorm{g(\omega)^\top} {\R^N}{\R^m} \norm w _{2, N}\notag\\
&\le \norm x_{2, m} \opnorm{g(\omega)^\top} {\R^N}{\R^m} d\text.
\label{gpadteq:1}
\end{align}
Somit ist $\menge{\sprod x {g(\omega)^\top v}}{v \in V}$ beschr�nkt
und $\left(\sprod \cdot {g(\omega)^\top v}_m \right)_{v \in V}$ eine
Familie gleichm��ig lipschitzstetiger Abbildungen.\\
Somit ist nach \ref{GleichmaessigeLipschitzstetigkeit} die Abbildung
\[f: \R^m \rightarrow \R, t \mapsto 
 \sup\menge{\sprod t {g(\omega)^\top v}_m}{v \in V}\]
wohldefiniert und lipschitzstetig.\\
Da $T$ beschr�nkt ist, ist also auch $f(T)$ beschr�nkt, sodass gilt
\begin{align*}
\left(\inf_{t \in T}\sup_{v \in V}Y(t, v) \right)(\omega)&=
\inf_{t \in T} \sup_{v \in V}\sprod {g(\omega)t} v_N =
\inf_{t \in T} \sup_{v \in V}\sprod t {g(\omega)^\top v}_N\\
&= \inf_{t \in T} f(t)
\overset{\refclaim{StetigeFunktionenUndDichteTeilmengen}
{StetigeFunktionenUndDichteTeilmengen2}}=
\inf_{t \in T'} f(t)
 = \left(\inf_{t \in T'}\sup_{v \in V}Y(t, v) \right)(\omega)\text.
\end{align*}
Es ist $T \times V'$ dicht in $T \times V$, wenn man $\R^m \times \R^N$ mit der
Produktnorm ausstattet.\\
Die Abbildung $\sprod{g(\omega)\cdot}{\cdot\cdot}_N$ ist
als Komposition stetiger Funktionen stetig.\\
Daher gilt
\begin{align*}
\left(\sup_{t \in T}\sup_{v \in V}Y(t, v) \right)(\omega)&= 
\sup_{(t,v)\in T\times V} Y(t, v)(\omega) 
= \sup_{(t,v)\in T\times V} \sprod{g(\omega)t}v_N\\
&\overset{\refclaim{StetigeFunktionenUndDichteTeilmengen}
{StetigeFunktionenUndDichteTeilmengen1}}=
\sup_{(t, v) \in T \times V'}\sprod{g(\omega)t}v_N
= \left(\sup_{t \in T}\sup_{v \in V'}Y(t, v) \right)(\omega)
\end{align*}
Damit ist zun�chst (3) gezeigt.\\
F�r jedes $s \in T$ ist die Abbildung $\sprod {g(\omega)s}\cdot_N$
lipschitzstetig, sodass gilt:
\begin{align}
\sup_{v \in V}\sprod{g(\omega)s}v_N
\overset{\refclaim{StetigeFunktionenUndDichteTeilmengen}
{StetigeFunktionenUndDichteTeilmengen2}}= 
\sup_{v \in V'}\sprod{g(\omega)s}v_N\text. \label{gpadteq:2}
\end{align}
Wegen $\eqref{gpadteq:1}$ ist (analog zu $f$) auch 
\[h : \R^m \rightarrow \R, 
t \mapsto \sup\menge{\sprod t {g(\omega)^\top v}_m}{v \in V'}\]
nach \ref{GleichmaessigeLipschitzstetigkeit} wohldefiniert und
lipschitzstetig.\\
Erhalte daher
\begin{align*}
\left(\inf_{t \in T}\sup_{v \in V}Y(t, v) \right)(\omega)&=
\inf_{t \in T} \sup_{v \in V}\sprod {g(\omega)t} v_N
\overset{\eqref{gpadteq:2}}=
\inf_{t \in T} \sup_{v \in V'}\sprod {g(\omega)t} v_N\\
&= \inf_{t \in T} \sup_{v \in V'}\sprod t {g(\omega)^\top v}_N
= \inf_{t \in T} h(t)\overset{\refclaim{StetigeFunktionenUndDichteTeilmengen}
{StetigeFunktionenUndDichteTeilmengen2} }= \inf_{t \in T'} h(t)\\
&= \inf_{t \in T'} \sup_{v \in V'}\sprod t {g(\omega)^\top v}_N
= \left(\inf_{t \in T'}\sup_{v \in V'}Y(t, v) \right)(\omega)\text.
\end{align*}
Es ist auch $T' \times V'$ dicht in $T \times V$, wenn man $\R^m \times \R^N$
mit der Produktnorm ausstattet. Daher gilt auch:
\begin{align*}
\left(\sup_{t \in T}\sup_{v \in V}Y(t, v) \right)(\omega) &=
\sup_{(t, v) \in T\times V}Y(t, v)(\omega) = 
\sup_{(t, v) \in T\times V}\sprod {g(\omega)t} v_N\\
&\overset{\refclaim{StetigeFunktionenUndDichteTeilmengen}
{StetigeFunktionenUndDichteTeilmengen1}}=
\sup_{(t, v) \in T'\times V'}\sprod {g(\omega)t} v_N
= \left(\sup_{t \in T'}\sup_{v \in V'}Y(t, v) \right)(\omega)\text.
\end{align*}
\end{proof}
\parag
Das eben bewiesene Lemma ermöglicht es uns bei gegebenen $T\subseteq \R^m$
und $V \subseteq \R^N$ die im Gauß'schen Min--Max--Satz vorkommenden
Infima und Suprema zunächst als abzählbare Extrema darzustellen.
Daher können wir nun einfach die Messbarkeit dieser Extrema im Falle beschränkter
$T,V$ einsehen.

\begin{satz}{Messbarkeit von Extrema in Gordon'schen Paaren}
            {MessbarkeitVonExtremaInGordonschenPaaren}
Seien $c \in \R$ und $(X,Y)$ das Gordon'sche Paar zum Tupel
$(\gamma,\delta,g,c)$.\\
Seien $\emptyset \neq T \subseteq \R^m$ beschr�nkt und
$\emptyset \neq V \subseteq \R^N$ beschr�nkt.\\
Dann gilt:
\begin{enumerate}[(1)]
	\item\label{MessbarkeitVonTeilextremaInGordonschenPaaren1}
	 $\underset{t \in T} \inf\, \underset{v \in V}\sup\, X(t, v), \, 
	\underset{t \in T} \sup\, \underset{v \in V}\sup\, X(t, v) \in 
	\mcm(\Omega, \mfa, \R)$.
	\item\label{MessbarkeitVonTeilextremaInGordonschenPaaren2}
	$\underset{t \in T} \inf\, \underset{v \in V}\sup\, Y(t, v), \, 
	\underset{t \in T} \sup\, \underset{v \in T}\sup\, Y(t, v) \in 
	\mcm(\Omega, \mfa, \R)$.
\end{enumerate}
\end{satz}
\begin{proof}
Da $\R^m$ und $\R^N$ unabh�ngig von der Normwahl separabel sind, sind es auch
$T$ und $V$ (als metrische R�ume).
Daher existieren $T' \subseteq T$ abz�hlbar und dicht in $T$, sowie
$V'\subseteq V$ abz�hlbar und dicht in $V$.
Als Teilmengen beschr�nkter Mengen sind $T'$ und $V'$ beschr�nkt und wegen der
jeweiligen Dichtheit in den nichtleeren Mengen $T$ und $V$ auch nichtleer.\\
(1)
Sei $s \in \R^m$.\\
Dann ist $X(s, w)$ f�r jedes $w \in \R^N$ nach 
\refclaim{IntegrierbarkeitUndVerteilungVonLinearkombinationen}
{IntegrierbarkeitUndVerteilungVonLinearkombinationen1} quadratintegrierbar
und da $\Bild X(s, w) \subseteq \R$ insbesondere $\mfa-\mfb-$messbar.\\
Da $V'$ dicht in $V$ und $\meng s$ nichtleer und beschr�nkt sind, sowie 
$\meng s$ dicht in $\meng s$ ist, gilt
\begin{align*}
 \sup_{v \in V}X(s, v) = \inf_{t \in \meng s}\sup_{v \in V}X(t, v)
\overset{\refclaim{GordonschePaareAufDichtenTeilmengen}
{GordonschePaareAufDichtenTeilmengen2}}= 
 \inf_{t \in \meng s}\sup_{v \in V'}X(t, v)
= \sup_{v \in V'}X(s, v)\text. \tag {$\ast$}
\end{align*}
Da $V'$ abz�hlbar ist, ist die letztgenannte Abbildung als abz�hlbares Supremum
$\mfa-\mfb$--messbarer Abbildungen selbst $\mfa-\overline\mfb$--messbar, wobei
$\overline \mfb$ die Borel'sche $\sigma-$Algebra auf $\overline\R$ bezeichne.\\
Also gilt:
\[\forall\, t \in \R^m: \sup_{v \in V'}X(t, v) \text{ ist } \mfa-\overline\mfb
\text{--messbar.}\]
Da $T'$ abz�hlbar ist, sind Abbildungen \[
\inf_{t \in T'}\sup_{v \in V'}X(t,v)
 \quad \text{ und }\quad
\sup_{t \in T'}\sup_{v \in V'}X(t,v) \]
abz�hlbare Infima, bzw. Suprema, 
$\mfa-\overline\mfb$--messbarer Abbildungen und als solche selbst
$\mfa-\overline\mfb$--messbar. Da $T'$ und $V'$ nichtleer und beschr�nkt sind,
sind die obigen Abbildungen nach
\refclaim{BeschraenktheitGordonscherPaare}{BeschraenktheitGordonscherPaare2}
aber $\R$--wertig und damit bereits $\mfa-\mfb$--messbar.\\
Da $T'$ dicht in $T$ und $V'$ dicht in $V$ sind, gilt
\begin{align*}
\inf_{t \in T}\sup_{v \in V}X(t, v) 
&\overset{\refclaim{GordonschePaareAufDichtenTeilmengen}
{GordonschePaareAufDichtenTeilmengen2}}= 
\inf_{t \in T'} \sup_{v \in V'}X(t,v) \in \mcm(\Omega, \mfa, \R)\\
\sup_{t \in T}\sup_{v \in V}X(t, v) 
&\overset{\refclaim{GordonschePaareAufDichtenTeilmengen}
{GordonschePaareAufDichtenTeilmengen2}}=
 \sup_{t \in T'}\sup_{v \in V'}X(t,v) \in \mcm(\Omega, \mfa, \R)\text.
\end{align*}
(2) Sei $s \in \R^m$.\\
Dann ist $Y(s, w)$ f�r jedes $w \in \R^N$ nach
\refclaim{IntegrierbarkeitUndVerteilungVonLinearkombinationen}
{IntegrierbarkeitUndVerteilungVonLinearkombinationen1} quadratintegrierbar und
mit $\Bild Y(s, w) \subseteq \R$ insbesondere $\mfa-\mfb$--messbar.\\
Die Menge $\meng s$ ist dicht in $\meng s$ und $V'$ ist dicht in $V$, sodass
gilt:
\[\sup_{v \in V} Y (s, v) = \inf_{t \in \meng s}\sup_{v \in V}Y(t,v)
\overset{\refclaim{GordonschePaareAufDichtenTeilmengen}
{GordonschePaareAufDichtenTeilmengen4}}=
\inf_{t \in \meng s}\sup_{v \in V'}Y(t,v) = \sup_{v \in V'}Y(s,v)\text.\]
Wegen der Abz�hlbarkeit von $V'$ ist diese Abbildung als abz�hlbares Supremum
$\mfa-\mfb$--messbarer Abbildungen selbst $\mfa-\overline\mfb$--messbar.\\
Also gilt:\[
\forall\,t\in\R^m:\sup_{v\in V'}Y(t,v)\text{ ist } \mfa-\overline\mfb
\text{--messbar.}\]
Damit sind wegen der Abz�hlbarkeit von $T'$ auch 
\[\inf_{t \in T'}\sup_{v\in V'}Y(t, v)
 \quad \text{ und }\quad 
  \sup_{t \in T'}\sup_{v\in V'}Y(t, v)\]
als abz�hlbare Infima, bzw. Suprema, $\mfa-\overline\mfb$--messbarer
Abbildungen selbst $\mfa-\overline\mfb$--messbar.
Da $T'$ und $V'$ nichtleer und beschr�nkt sind, sind die obigen Abbildungen
nach
\refclaim{BeschraenktheitGordonscherPaare}{BeschraenktheitGordonscherPaare2}
 allerdings $\R$--wertig und damit bereits
$\mfa-\mfb$--messbar.\\
Schlie�lich gilt wegen $T'\subseteq T$ dicht in $T$ und $V'\subseteq V$ dicht
in $V$:
\begin{align*}
\inf_{t \in T}\sup_{v \in V}Y(t, v)
&\overset{\refclaim{GordonschePaareAufDichtenTeilmengen}
{GordonschePaareAufDichtenTeilmengen4}}=
\inf_{t \in T'}\sup_{v\in V'}Y(t, v) \in \mcm(\Omega, \mfa, \R)\text,\\
\sup_{t \in T}\sup_{v \in V}Y(t, v)
&\overset{\refclaim{GordonschePaareAufDichtenTeilmengen}
{GordonschePaareAufDichtenTeilmengen4}}=
\sup_{t \in T'}\sup_{v\in V'}Y(t, v) \in \mcm(\Omega, \mfa, \R)\text.
\end{align*}
\end{proof}
    
\parag
Da uns an späterer Stelle Erwartungswertungleichungen von Extrema Gordon'scher
Paare interessieren werden, ist es wichtig, auch die Integrierbarkeit selbiger
einzusehen.

\begin{lemma}{Integrierbarkeit von Extrema in Gordon'schen Paaren}
             {IntegrierbarkeitVonExtremaInGordonschenPaaren}
Seien $c \in \R$, $(X, Y)$ das Gordon'sche Paar zum Tupel
$(\gamma, \delta, g, c)$.\\
Seien $\emptyset \neq T \subseteq \R^m$ beschränkt
und $\emptyset \neq V \subseteq \R^N$ beschränkt.\\
Dann gilt:
\begin{enumerate}[(1)]
	\item \label{IntegrierbarkeitVonExtremaInGordonschenPaaren1}
	$\Inf{t \in T}\Sup{v \in V}X(t, v)$ und $\Sup{t \in T}\Sup{v \in V}X(t, v)$
	sind integrierbar.\absatz
	Insbesondere sind $\Inf{t \in T}\bprod t \delta_m$
	und $\Sup{t \in T}\bprod t \delta_m$ integrierbar.\\
	Außerdem ist für jeden
	normierten $\R$--Vektorraum $(E, \norm\cdot_E)$ und $x \in E^N$
die Abbildung $\norm\cdot_E \circ \bprod \cdot x_n \circ \gamma$ integrierbar.
	\item \label{IntegrierbarkeitVonExtremaInGordonschenPaaren2}
	$\Inf{t \in T}\Sup{v \in V}Y(t, v)$ und $\Sup{t \in T}\Sup{v \in V}Y(t, v)$
	sind integrierbar.\\
	Sind $(E, \norm\cdot_E)$ ein normierter $\R$--Vektorraum, 
	$x \in E^N$, $G_x$ der Gauß'sche Operator zum Tupel $(x,g)$, so sind
	\[\inf_{t \in T}\, \norm\cdot_E \circ ev(t,\cdot)\circ G_x \quad\text{ und } 
	  \quad \sup_{t \in T}\, \norm\cdot_E \circ ev(t,\cdot)\circ G_x \quad
	\text{integrierbar.}\]
\end{enumerate}
\end{lemma}
\begin{proof}
Wir zeigen, dass für alle $Z \in \meng{X, Y}$ ein
 $f_Z: \Omega \rightarrow \R$ integrierbar existiert mit
\[(\ast)\quad
\forall\, t \in T \, \forall\, v \in V \, \forall\, \omega \in \Omega:
 \abs{Z(t, v)(\omega)} \le f_Z(\omega)\text.\]
Ist nämlich $Z \in \meng{X, Y}$ und $f_Z :\Omega \rightarrow \R$ integrierbar, 
sodass $(\ast)$ erfüllt ist, so gilt für alle $\omega \in \Omega$:
$\forall\,t \in T \, \forall\, v \in V:
-f_Z(\omega) \le Z(t, v) (\omega) \le f_Z(\omega)$\\
und damit sogleich
\[-f_Z(\omega) \le \inf_{t \in T}\inf_{v \in V} Z(t, v)(\omega)
\le \inf_{t \in T}\sup_{v \in V} Z (t, v) (\omega)
\le \sup_{t \in T}\sup_{v\in V} Z(t, v)(\omega) \le f_Z(\omega)\text.\]
Somit gilt insbesondere
\[-f_Z \le \inf_{t \in T}\sup_{v \in V} Z (t, v)
\le \sup_{t \in T}\sup_{v\in V} Z(t, v) \le f_Z\text.\]
Da $f_Z$ und $-f_Z$ integrierbar und nach
\ref{MessbarkeitVonExtremaInGordonschenPaaren} die Abbildungen
\[\inf_{t\in T}\sup_{v\in V}Z(t, v),\quad \sup_{t \in T}\sup_{v\in V} Z(t, v)\]
messbar sind, sind die letztgenannten Abbildungen nach
\refclaim{Integrierbarkeitskriterien}{Integrierbarkeitskriterien2}
integrierbar.\\
Da $T$ beschränkt ist, gibt es ein $a \in \R_{>0}$ mit 
$\forall\, t \in T: \norm t_{2, m} \le a$.\\
Da auch $V$ beschränkt ist, gibt
es ein $b \in \R_{>0}$ mit $\forall\, v \in V: \norm v_{2, N} \le b$.\\
(1) Sei 
$f_X:=a\abs c(\norm\cdot_{2,m}\circ\delta) + \norm\cdot_{2, N} \circ\gamma$.\\
Nach \refclaim{ElementareEigenschaftenQuadratintegrierbarerAbbildungen}
{ElementareEigenschaftenQuadratintegrierbarerAbbildungen4} sind
$\norm\cdot_{2,m} \circ \delta$ und $\norm\cdot_{2, N} \circ\gamma$ 
$\mfa-\mfb$--messbar und nach
\ref{ErwartungswertDer2NormEinesStandardnormalverteiltenProzesses} sind diese
beiden Abbildungen auch integrierbar.\\
Damit ist es auch $f_X$ als Linearkombination integrierbarer Abbildungen.\\
Seien nun $t \in T$, $v \in V$ und $\omega \in \Omega$.\\
Dann gilt:
\begin{align*}
\abs{X(t, v)(\omega)}&\overset{\refclaim{BeschraenktheitGordonscherPaare}
{BeschraenktheitGordonscherPaare1}}\le
\abs c \norm t _{2, m}\norm {\delta(\omega)}_{2, m} 
+ \norm v _{2, N}\norm{\gamma(\omega)}_{2, N}\\
&\le\abs ca\norm{\delta(\omega)}_{2,m}+b\norm{\gamma(\omega)}_{2,N}=f_X(\omega)
\text.\end{align*}
Damit ist (1) (ohne den Zusatz) nach obiger Begründung gezeigt.\\
Sei $(X', Y')$ das Gordon'sche Paar zum Tupel $(\gamma, \delta, g, 1)$ und
$V' := \meng 0$.\\
Dann ist nach (1) die Abbildung
\[\sup_{t \in T}\bprod t \delta_m =
\sup_{t \in T}\bprod t \delta_m + \sup_{v \in V'} \bprod v \gamma_N
\overset{\refclaim{BeschraenktheitGordonscherPaare}
{BeschraenktheitGordonscherPaare4}}=
\sup_{t \in T}\sup_{v\in V'}X'(t, v) \text{ integrierbar.}\]
Analog ist
\[\inf_{t \in T}\bprod t \delta _m =\inf_{t \in T}\bprod t \delta_m + \sup_{v \in V'} \bprod v \gamma_N
\overset{\refclaim{BeschraenktheitGordonscherPaare}
{BeschraenktheitGordonscherPaare4}}=
\inf_{t \in T}\sup_{v\in V'}X'(t, v) \text{ integrierbar.}\]
Seien nun $(E, \norm\cdot_E)$ ein normierter $\R$-Vektorraum, $x \in E^N$,
$T':= \meng 0$, sowie $V'':= \menge{\pfeil \vi N (x)}{\vi \in \kabg{E'}01}$.\\
Wegen
\[\forall\, \vi \in \kabg{E'}01: \abs{\pfeil\vi N (x)}
\overset{\ref{EigenschaftenDerErweiterungImNormiertenFall}}\le
\opnorm \vi E \R \norm x _{1, N, E} \le \norm x _{1, N, E}\]
ist $V''$ beschränkt und offensichtlich nichtleer.\\
Also gilt für alle $\omega \in \Omega$
\begin{align*}
&\,\left(\norm\cdot_E \circ \bprod \cdot x \circ \gamma \right)(\omega)
=\norm{\bprod {\gamma(\omega)} x_N}_E
 \overset{\ref{NormUndDualeNorm}}= 
\sup_{\vi \in \kabg{E'}01}\vi\left(\bprod{\gamma(\omega)}x_N \right)\\
\overset{\refclaim{EigenschaftenDesLinearkombinators}
{EigenschaftenDesLinearkombinators2}}= &\,
\sup_{\vi \in \kabg{E'}01}\sprod{\gamma(\omega)}{\pfeil \vi N (x)}_N
= \sup_{v \in V''}\sprod{\gamma(\omega)}v_N
=  \sup_{v \in V''}\sprod v{\gamma(\omega)}_N\\
=&\, \left(\sup_{v \in V''}\bprod v \gamma_N\right)(\omega)
\end{align*}
und daher
\begin{align*}
\norm\cdot_E \circ \bprod \cdot x \circ \gamma &=
\sup_{v \in V''}\bprod v \gamma_N 
= \sup_{t \in T'}\bprod t \delta_m + \sup_{v \in V''}\bprod v \gamma_N\\
&\overset{\refclaim{BeschraenktheitGordonscherPaare}
{BeschraenktheitGordonscherPaare4}}= \sup_{t \in T'}\sup_{v \in V''}X'(t,v)
\ \text{ integrierbar nach (1).}
\end{align*}
(2) Die Abbildung $\Mat_{m, N}$ (vgl. Bezeichnungen.(\ref{Bez12})) ist linear
und bijektiv.\\
Da die $\R$--Vektorräume $\R^{mN}$ und $\R^{m\times N}$ endlichdimensional
sind, ist $\Mat_{m,N}$ als Abbildung von $(\R^{mN}, \norm\cdot_{2, mN})$ nach
$(\R^{m \times N}, \opnorm \cdot {\R^m}{\R^N})$ stetig.\\
Wegen der Injektivität von $\Mat_{m,N}$ ist nach
\refclaim{EigenschaftenDerBildhalbnorm}{EigenschaftenDerBildhalbnorm2} 
\[\norm\cdot_{\Mat_{m, N}}: \R^{mN} \rightarrow [0, \infty[,
x \mapsto \opnorm {\Mat_{m, N}(x)}{\R^m}{\R^N}\]
eine Norm auf dem $\R^{mN}$ und da $\Mat_{m, N}$ stetig ist, existiert nach
\refclaim{EigenschaftenDerBildhalbnorm}{EigenschaftenDerBildhalbnorm3} 
ein $d \in \R_{>0}$ mit
$\norm\cdot_{\Mat_{m, N}} \le d \norm\cdot_{2, mN}$.\hfill $(\ast\ast)$\\
Sei nun $f_Y := abd \left(\norm\cdot_{2, mN} \circ \widetilde g \right)$.\\
Da $\widetilde g$ nach Voraussetzung $N(0, 1)$--verteilt und unabhängig ist,
ist $\norm\cdot_{2, mN} \circ \widetilde g$ nach
\refclaim{ElementareEigenschaftenQuadratintegrierbarerAbbildungen}
{ElementareEigenschaftenQuadratintegrierbarerAbbildungen4}
$\mfa-\mfb$--messbar und nach
\ref{ErwartungswertDer2NormEinesStandardnormalverteiltenProzesses} 
auch integrierbar, also ist es auch $f_Y$ als Vielfaches einer integrierbaren
Abbildung.\\
Seien $t \in T$, $v \in V$ und $\omega \in \Omega$.\\
Dann gilt:
\begin{align*}
&\abs{Y(t, v)(\omega)} \overset{\refclaim{BeschraenktheitGordonscherPaare}
{BeschraenktheitGordonscherPaare1}}\le 
 \opnorm{g(\omega)}{\R^m}{\R^N}\norm t _{2, m}\norm v_{2, N}
 \le\, ab\opnorm{g(\omega)}{\R^m}{\R^N}\\
  =&\, ab \opnorm{\Mat_{m, N}(\widetilde g (\omega))}{\R^m}{\R^N}
 = ab \norm{\widetilde g (\omega)}_{\Mat_{m, N}} \overset{(\ast\ast)}\le
 abd \norm{\widetilde g (\omega)}_{2, mN} = f_Y(\omega) \text.
\end{align*}
Somit ist auch (2) (ohne den Zusatz) gezeigt.\\
Seien $(E, \norm\cdot_E)$ ein normierter $\R$--Vektorraum, $x \in E^N$ und
$G_x$ der Gauß'sche Operator zum Tupel $(x,g)$.\\
Sei $V'':= \menge{\pfeil \vi N(x)}{\vi \in \kabg{E'}01}$.\\
Wie im Beweis des Zusatzes von (1) gesehen, ist $V''$ beschränkt und
nichtleer.\\
Es gilt für alle $s\in T$ und alle $\omega \in \Omega$:
\begin{align*}
\sup_{v \in V''}Y(s,v) &= \sup_{v\in V''}\sprod{g(\omega)s}{v}_N
=\sup_{\vi \in \kabg{E'}01}\sprod{g(\omega)s}{\pfeil\vi N (x)}_N\\
&\overset{\refclaim{EigenschaftenDesLinearkombinators}
{EigenschaftenDesLinearkombinators2}}=
\sup_{\vi \in \kabg{E'}01}\vi\left(\bprod{g(\omega)s}x_N \right)
\overset{\ref{NormUndDualeNorm}}= \norm{\bprod{g(\omega)s}x_N}_E\\
&= \norm{G_x(\omega)s}_E 
= \left(\norm\cdot_E\circ ev(s,\cdot)\circ G_x\right)(\omega)\text.
\end{align*}
Somit gilt:
\begin{align*}
\inf_{t \in T}\, \norm\cdot_E \circ ev(t,\cdot)\circ G_x
&= \inf_{t \in T}\sup_{v \in V''}Y(t, v)\text,\\
\sup_{t \in T}\, \norm\cdot_E \circ ev(t,\cdot)\circ G_x
&= \sup_{t \in T}\sup_{v \in V''}Y(t, v)\end{align*}
und damit sind diese beiden Abbildungen nach (2) integrierbar.
\end{proof}
\parag
Als Nächstes stellen wir fest, dass endliche Teilfamilien Gordon'scher Paare
den Voraussetzungen des Gauß'schen Min--Max--Satzes genügen.

\begin{lemma}{Gauß'scher Min--Max--Satz für Gordon'sche Paare (endlich)}
             {GordonschePaareUndDerGaussscheMinMaxSatzEndlich}
Seien $\emptyset \neq T \subseteq S^{m-1}$ und
$\emptyset \neq V \subseteq \R^N\setminus\meng 0$ beschränkt.\\
Seien $c \in \R_{>0}$ mit
$\forall\, v \in V: \norm v _{2, N} \le c$
und $(X, Y)$ das Gordon'sche Paar zum Tupel $(\gamma, \delta, g, c)$.
Schließlich seien $\mct := \Pot_{ne}(T)$ und $\mcv := \Pot_{ne}(V)$.\\
Dann gilt:
\begin{enumerate}[(1)]
	\item \label{GordonschePaareUndDerGaussscheMinMaxSatzEndlich1}
	$\forall\, A \in \mct\, \forall\, B \in \mcv:
	\E\left(\Inf{t \in A}\Sup{v \in B}X(t, v) \right)
	\le \E\left(\Inf{t \in A}\Sup{v \in B}Y(t, v) \right)$.
	\item \label{GordonschePaareUndDerGaussscheMinMaxSatzEndlich2}
	$\forall\, A \in \mct\, \forall\, B \in \mcv:
	\E\left(\Sup{t \in A}\Sup{v \in B}Y(t, v) \right)
	\le \E\left(\Sup{t \in A}\Sup{v \in B}X(t, v) \right)$.
\end{enumerate}
\end{lemma}
\begin{proof}
Seien $A \in \mct$ und $B \in \mcv$.\\
Da $A$ und $B$ endlich und nichtleer sind, existieren $r, s \in \N$ und
$\alpha : \haken r \rightarrow A$ bijektiv und $\beta: \haken s \rightarrow B$
bijektiv.\\
Seien $i, j \in \haken r$ und $k, l \in \haken s$.\\
Nach \refclaim{IntegrierbarkeitUndVerteilungVonLinearkombinationen}
{IntegrierbarkeitUndVerteilungVonLinearkombinationen1} sind 
$X(\alpha(i), \beta(j))$ und $Y(\alpha(i), \beta(j))$ quadratintegrierbar und
damit
\begin{align*}
X'& : \haken r \times \haken s \rightarrow L_2(\Omega, \mfa, \Prim, \R),\ 
(a, b) \mapsto X(\alpha(a), \beta(b))\\
Y'& : \haken r \times \haken s \rightarrow L_2(\Omega, \mfa, \Prim, \R),\ 
(a, b) \mapsto Y(\alpha(a), \beta(b))
\end{align*}
wohldefiniert.\\
Da $\alpha$ und $\beta$ bijektiv, sowie $A$ und $B$ endlich sind, gilt
offenbar:
\begin{align}\inf_{t \in A}\sup_{v \in B}X(t, v) =
  \min_{a \in \haken r}\max_{b \in \haken s}X(\alpha(a), \beta(b))
  = \min_{a \in \haken r}\max_{b \in \haken s}X'(a, b) \label{gpugmmseeq:1}
\end{align}
und analog:
\begin{align}
\inf_{t \in A}\sup_{v\in B}Y(t, v) 
= \min_{a \in \haken r}\max_{b \in \haken s}Y'(a, b)\text,\label{gpugmmseeq:2}
\end{align}
sowie
\begin{align}\sup_{t \in A}\sup_{v\in B}X(t, v) 
&=\max_{a\in \haken r}\max_{b \in \haken s}X'(a, b)\text,\label{gpugmmseeq:3}\\
\sup_{t \in A}\sup_{v\in B}Y(t, v) 
&=\max_{a \in \haken r}\max_{b \in \haken s}Y'(a, b)\text.\label{gpugmmseeq:4}
\end{align}
Wegen $\alpha(i) \in S^{m-1}$ ist $\alpha(i) \neq 0$ und wegen
$\beta(i) \in V \subseteq \R^N\setminus\meng 0$ ist $\beta(i) \neq 0$.\\
Somit ist nach
\refclaim{IntegrierbarkeitUndVerteilungVonLinearkombinationen}
{IntegrierbarkeitUndVerteilungVonLinearkombinationen2}: 
$X(\alpha(i), \beta(j))\neq 0$ und nach
\refclaim{IntegrierbarkeitUndVerteilungVonLinearkombinationen}
{IntegrierbarkeitUndVerteilungVonLinearkombinationen3}:
$Y(\alpha(i), \beta(j)) \neq 0$.\\
Mit $\norm{\alpha(i)}_{2,m}=1$ gilt nach
\refclaim{IntegrierbarkeitUndVerteilungVonLinearkombinationen}
{IntegrierbarkeitUndVerteilungVonLinearkombinationen4}
\[X'(i,j)=X(\alpha(i), \beta(j)) \ \text{ ist }\ 
N(0, c^2 + \norm {\beta(j)}_{2, N}^2)\text{--verteilt}\] und 
nach \refclaim{IntegrierbarkeitUndVerteilungVonLinearkombinationen}
{IntegrierbarkeitUndVerteilungVonLinearkombinationen5}
\[Y'(i, j)=Y(\alpha(i), \beta(j)) \ \text{ ist } \ 
N(0, \norm{\beta(i)}_{2, N}^2)\text{--verteilt.}\]
Weiter gilt ebenfalls wegen $\alpha(i) \in S^{m-1}$:
\begin{align*}
&\, \norm{X'(i, k) - X'(i, l)}_2  = \norm{X(\alpha(i), \beta(k)) 
- X(\alpha(i), \beta(l))}_2\\ \overset{\refclaim{AbstaendeInGordonschenPaaren}
{AbstaendeInGordonschenPaaren4}} = &\, \norm{Y(\alpha(i), \beta(k)) 
- Y(\alpha(i), \beta(l))}_2 = \norm{Y'(i, k) - Y'(i, l)}_2\text.
\end{align*}
Es ist auch $\alpha(j) \in S^{m-1}$, sowie 
$\beta(k), \beta(l)\in V$ und damit nach Voraussetzung 
$\norm{\beta(k)}_{2,N} \le c$, $\norm{\beta(l)}_{2,N} \le c$, also gilt:
\begin{align*}
&\, \norm{Y'(i, k) - Y'(j, l)}_2 
= \norm{Y(\alpha(i), \beta(k)) - Y(\alpha(j), \beta(l))}_2\\
\overset{\refclaim{AbstaendeInGordonschenPaaren}
{AbstaendeInGordonschenPaaren5}}\le &\,
\norm{X(\alpha(i), \beta(k)) - X(\alpha(j), \beta(l))}_2
= \norm{X'(i, k) - X'(j, l)}_2\text.
\end{align*}
Damit erfüllen $X', Y'$ alle Voraussetzungen von
\ref{GaussscherMinMaxSatzGleichheitInDerErstenBedingung}, sodass gilt:
\begin{align*}
\E\left(\inf_{t \in A}\sup_{v \in B}X(t, v) \right)
&\overset{\eqref{gpugmmseeq:1}}=
\E\left(\min_{a \in \haken r}\max_{b \in \haken s}X'(a, b)\right)
\overset{\refclaim{GaussscherMinMaxSatzGleichheitInDerErstenBedingung}
{GaussscherMinMaxSatzGleichheitInDerErstenBedingung1}}\le
\E\left(\min_{a \in \haken r}\max_{b \in \haken s}Y'(a, b)\right)\\
& \overset{\eqref{gpugmmseeq:2}}= 
\E\left(\inf_{t \in A}\sup_{v \in B}Y(t, v) \right)
\end{align*}
und
\begin{align*}
\E\left(\sup_{t \in A}\sup_{v \in B}X(t, v) \right)
&\overset{\eqref{gpugmmseeq:3}}=
\E\left(\max_{a \in \haken r}\max_{b \in \haken s}X'(a, b)\right)\\
&\overset{\refclaim{GaussscherMinMaxSatzGleichheitInDerErstenBedingung}
{GaussscherMinMaxSatzGleichheitInDerErstenBedingung2}}\le
\E\left(\max_{a \in \haken r}\max_{b \in \haken s}Y'(a, b)\right)
\overset{\eqref{gpugmmseeq:4}}=
\E\left(\sup_{t \in A}\sup_{v \in B}Y(t, v) \right)\text.
\end{align*}
\end{proof}
\parag
Wir können nun den Gauß'schen Min--Max--Satz für Gordon'sche Paare formulieren.
Im Wesentlichen verwenden wir dazu die Ergebnisse aus dem zweiten Kapitel.
Es ist allerdings auch noch ein gewisser Mehraufwand vonnöten, weil in der
späteren Anwendung die Situation $0 \in V$ durchaus auftreten wird, wodurch
der letzte Satz nicht direkt anwendbar ist.
Ist allerdings $0$ auch ein
Häufungspunkt von $V$, so kann man zu einer dichten Teilmenge übergehen, die
$0$ nicht enthält.
Auf die Voraussetzung an die $0$ als Häufungspunkt kann
nicht verzichtet werden, weil z. B. für $V \coloneqq  \meng 0$, $T\coloneqq \meng{e_1^m}$ und
$c \coloneqq  1$ die anderen Voraussetzungen offensichtlich gegeben sind, aber
die Bilinearität von $Y$ bereits liefert, dass der zweite Erwartungswert in
der ersten Aussage $0$ ist, während der erste Erwartungswert allerdings $1$
ist.

\begin{satz}{Gordon'sche Paare und der Gauß'sche Min--Max-Satz}
            {GordonschePaareUndDerGaussscheMinMaxSatz}
Seien $\emptyset \neq T \subseteq S^{m-1}$ und 
$\emptyset\neq V \subseteq \R^N$ beschränkt.\\
Seien $c \in \R_{>0}$ mit $\forall\, v \in V: \norm v _{2,N} \le c$ und es gelte
\[0 \in V \folgt 0 \text{ Häufungspunkt von }V\]
Schließlich sei $(X,Y)$ das Gordon'sche Paar zum Tupel $(\gamma,\delta,g,c)$.\\
Dann sind die Abbildungen
\[\inf_{t \in T}\sup_{v \in V}X(t, v),\  \sup_{t \in T}\sup_{v \in V}X(t, v),\ 
\inf_{t \in T}\sup_{v \in V}Y(t, v),\  \sup_{t \in T}\sup_{v \in V}Y(t, v)\]
nach~\ref{IntegrierbarkeitVonExtremaInGordonschenPaaren} integrierbar und es
gilt:
\begin{enumerate}[(1)]
	\item \label{GordonschePaareUndDerGaussscheMinMaxSatz1}
	$  \E\left(\Inf{t \in T}\Sup{v\in V}X(t, v)\right)
	\le\E\left(\Inf{t \in T}\Sup{v\in V}Y(t, v)\right)$
	\item \label{GordonschePaareUndDerGaussscheMinMaxSatz2}
	$  \E\left(\Sup{t \in T}\Sup{v\in V}Y(t, v)\right)
	\le\E\left(\Sup{t \in T}\Sup{v\in V}X(t, v)\right)$
\end{enumerate}
Insbesondere gilt
\begin{align*}
\E\bigg(\Inf{t \in T}\Sup{v\in V}X(t, v)\bigg) & \le 
\E\left(\Inf{t \in T}\Sup{v\in V}Y(t, v)\right) \le 
\E\left(\Sup{t \in T}\Sup{v\in V}Y(t, v)\right)\\& \le
\E\left(\Sup{t \in T}\Sup{v\in V}X(t, v)\right)
\end{align*}
\end{satz}
\begin{proof}
Zunächst gibt es wegen der Separabilität von $\R^m$ und $\R^N$ ein
$T' \subseteq T$ abzählbar und dicht in $T$, sowie ein
$\widetilde V \subseteq V$ abzählbar und dicht in $V$.\\
Falls $0 \in V$, so ist nach Voraussetzung $0$ ein Häufungspunkt von $V$, also
gibt es nach 
\refclaim{DichtheitUndHaeufungspunkte}{DichtheitUndHaeufungspunkte2} ein
$V'' \subseteq V$ abzählbar und dicht in $V$ mit $0 \notin V''$.\\
Sei also
\[V' \coloneqq  \begin{cases}\widetilde V &:\text{ falls } 0 \notin V\\
V'' & : \text{ sonst} \end{cases}\text.\]
Dann ist insbesondere $0 \notin V'$.\\
Seien $\mct \coloneqq  \Pot_{ne}(T')$ und $\mcv\coloneqq  \Pot_{ne}(V')$.\\
Nach
\refclaim{BeschraenktheitGordonscherPaare}{BeschraenktheitGordonscherPaare1}
sind für alle $\omega \in \Omega$ die Mengen
\begin{align}\menge{X(t, v)(\omega)}{(t, v) \in T \times V},\,
\menge{Y(t, v)(\omega)}{(t, v) \in T \times V} \text{ beschränkt.}\tag {$a$}
\end{align}
Da jedes $C \in \mct\cup\meng{T'}$ und jedes $D\in\mcv\cup\meng{V'}$
beschränkt und nichtleer ist, gilt nach
\refclaim{BeschraenktheitGordonscherPaare}{BeschraenktheitGordonscherPaare3}:
\begin{align}
\forall\, A \in \mct\cup\meng{T'}\, \forall\, B \in \mcv\cup\meng{V'}:
 \inf_{t \in A}\sup_{v \in V}X(t, v) = \sup_{v \in V}\inf_{t \in A}X(t,v)
 \tag {$b$}
\end{align}
und nach~\ref{IntegrierbarkeitVonExtremaInGordonschenPaaren}:
\begin{align*}
 &\forall\, A \in \mct\cup\meng{T'}\, \forall\, B \in \mcv\cup\meng{V'}:\\
 &\quad \quad
 \inf_{t \in A}\sup_{v \in B}X(t, v), \, \inf_{t \in A}\sup_{v \in B}Y(t, v)
 \overset{\refclaim{IntegrierbarkeitVonExtremaInGordonschenPaaren}
 {IntegrierbarkeitVonExtremaInGordonschenPaaren1}}\in L_1(\Omega,\mfa,\Prim,\R)
 \tag  {$c$}\\
 &\wedge\quad
 \sup_{t \in A}\sup_{v \in B}X(t, v), \, \sup_{t \in A}\sup_{v \in B}Y(t, v)
 \overset{\refclaim{IntegrierbarkeitVonExtremaInGordonschenPaaren}
 {IntegrierbarkeitVonExtremaInGordonschenPaaren2}}\in L_1(\Omega,\mfa,\Prim,\R)
  \text.\tag {$b'$}
\end{align*}
Da für jedes $C \in \mct \cup \meng {T'}$ nach Setzung 
$\emptyset \neq C \subseteq S^{m-1}$ gilt, sowie jedes 
$D \in \mcv \cup \meng{V'}$
nichtleer und beschränkt ist mit $0 \notin D$, gilt nach
\ref{GordonschePaareUndDerGaussscheMinMaxSatzEndlich}
\begin{align*}
&\forall\, A \in \mct\, \forall\, B \in \mcv:\\
& \quad\quad\E\left(\inf_{t \in A}\sup_{v \in B}X(t, v) \right)
 \overset{\refclaim{GordonschePaareUndDerGaussscheMinMaxSatzEndlich}
{GordonschePaareUndDerGaussscheMinMaxSatzEndlich1}}\le
 \E\left(\inf_{t \in A}\sup_{v \in B}Y(t, v)\right) \tag {$d$}\\
& \wedge \quad
 \E\left(\sup_{t \in A}\sup_{v \in B}Y(t, v) \right)
 \overset{\refclaim{GordonschePaareUndDerGaussscheMinMaxSatzEndlich}
{GordonschePaareUndDerGaussscheMinMaxSatzEndlich2}}\le
 \E\left(\sup_{t \in A}\sup_{v \in B}X(t, v)\right) \text.\tag {$c'$}
\end{align*}
Seien $I\coloneqq T'$, $J\coloneqq V'$, $X'\coloneqq  X\vert_{I\times J} = X\vert_{T' \times V'}$ und
$Y'\coloneqq Y\vert_{I \times J}= Y\vert_{T' \times V'}$.\\
Mit $(a)$, $(b)$, $(c)$, $(d)$, sind die Voraussetzungen von 
\ref{AbzaehlbareDarstellungUndInfimumVonSuprema} und
wegen $(a)$, $(b')$, $(c')$ auch jene von
\ref{AbzaehlbareDarstellungUndSupremumVonSuprema} für $X', Y'$
und $I, J$ erfüllt.\\
Es gilt:
\[\inf_{t \in T'}\sup_{v \in V'}X'(t,v) =
  \inf_{t \in T'}\sup_{v \in V'}X(t, v)\text, \ 
  \inf_{t \in T'}\sup_{v \in V'}Y'(t,v) =
  \inf_{t \in T'}\sup_{v \in V'}Y(t, v)\text,\]
\[\sup_{t \in T'}\sup_{v \in V'}X'(t,v) =
  \sup_{t \in T'}\sup_{v \in V'}X(t, v), \ 
  \sup_{t \in T'}\sup_{v \in V'}Y'(t,v) =
  \sup_{t \in T'}\sup_{v \in V'}Y(t, v)\text.\]
Somit folgt:
\begin{align}\E\left(\inf_{t \in T'}\sup_{v \in V'}X(t, v)\right) &=
\E\left(\inf_{t \in T'}\sup_{v \in V'}X'(t, v)\right) 
\overset{\ref{AbzaehlbareDarstellungUndInfimumVonSuprema}}\le
\E\left(\inf_{t \in T'}\sup_{v \in V'}Y'(t, v)\right)\notag\\
&= \E\left(\inf_{t \in T'}\sup_{v \in V'}Y(t, v)\right) \label{gpudgmmseq:1}
\end{align}
und
\begin{align}\E\left(\sup_{t \in T'}\sup_{v \in V'}Y(t, v)\right) &=
\E\left(\sup_{t \in T'}\sup_{v \in V'}Y'(t, v)\right) 
\overset{\ref{AbzaehlbareDarstellungUndSupremumVonSuprema}}\le
\E\left(\sup_{t \in T'}\sup_{v \in V'}X'(t, v)\right)\notag\\
&= \E\left(\sup_{t \in T'}\sup_{v \in V'}X(t, v)\right)\text.
\label{gpudgmmseq:2}
\end{align}
Schließlich ist $T'$ dicht in $T$ und $V'$ dicht in $V$, sodass wegen
\ref{GordonschePaareAufDichtenTeilmengen} gilt:
\begin{align*}
\E\left(\inf_{t \in T}\sup_{v\in V}X(t,v) \right)
&\overset{\refclaim{GordonschePaareAufDichtenTeilmengen}
{GordonschePaareAufDichtenTeilmengen2}}=
\E\left(\inf_{t \in T'}\sup_{v\in V'}X(t,v) \right)\\
& \overset{\eqref{gpudgmmseq:1}}\le
\E\left(\inf_{t \in T'}\sup_{v\in V'}Y(t,v) \right)
\overset{\refclaim{GordonschePaareAufDichtenTeilmengen}
{GordonschePaareAufDichtenTeilmengen4}}=
\E\left(\inf_{t \in T}\sup_{v\in V}Y(t,v) \right)\text,\\
\E\left(\sup_{t \in T}\sup_{v\in V}Y(t,v) \right)
&\overset{\refclaim{GordonschePaareAufDichtenTeilmengen}
{GordonschePaareAufDichtenTeilmengen4}}=
\E\left(\sup_{t \in T'}\sup_{v\in V'}Y(t,v) \right)\\
& \overset{\eqref{gpudgmmseq:2}}\le
\E\left(\sup_{t \in T'}\sup_{v\in V'}X(t,v) \right)
\overset{\refclaim{GordonschePaareAufDichtenTeilmengen}
{GordonschePaareAufDichtenTeilmengen2}}=
\E\left(\sup_{t \in T}\sup_{v\in V}X(t,v) \right)\text.
\end{align*}
Damit sind (1) und (2) gezeigt.\\
Zusätzlich gilt
\[\inf_{t \in T}\sup_{v \in V}Y(t,v) \le \sup_{t \in T}\sup_{v \in V}Y(t,v)\]
und da die beiden genannten Abbildungen nach
\ref{IntegrierbarkeitVonExtremaInGordonschenPaaren} integrierbar sind, gilt
auch
\begin{align} \E\left(\inf_{t \in T}\sup_{v \in V}Y(t,v) \right) 
\le \E\left(\sup_{t \in T}\sup_{v \in V}Y(t,v) \right)
\label{gpudgmmseq:3}
\end{align}
und daher
\begin{align*}
\E\bigg(\Inf{t \in T}\Sup{v\in V}X(t, v)\bigg) & \overset{(1)}\le 
\E\left(\Inf{t \in T}\Sup{v\in V}Y(t, v)\right)
\overset{\eqref{gpudgmmseq:3}}\le 
\E\left(\Sup{t \in T}\Sup{v\in V}Y(t, v)\right)\\& \overset{(2)}\le
\E\left(\Sup{t \in T}\Sup{v\in V}X(t, v)\right)\text.
\end{align*}
\end{proof}
\parag
Die Idee, den Gauß'schen Min--Max--Satz auf Gordon'sche Paare anzuwenden
stammte von Gordon.
Tatsächlich sind die Gordon'schen Paare genau so
konstruiert, dass der Gauß'sche Min--Max--Satz nach einigen technischen
Vorüberlegungen direkt angewandt werden kann.

Es ergeben sich nun zwei Korollare, welche unter zusätzlichen Annahmen an $V$
Aussagen über die in den Zusätzen von
\ref{IntegrierbarkeitVonExtremaInGordonschenPaaren} vorgestellten Abbildungen
treffen.\\
Das nachfolgende Korollar stammt im Wesentlichen aus \cite{Gor}
(und wurde in \cite{Gor2} verwendet).
In der vorliegenden Form wird es in
dieser Arbeit vorrangig für den Beweis des Satzes von Milman und Schechtman
verwendet.
Dort werden wir die Normdarstellung durch ein Supremum auf die
spezielle Norm aus \ref{Sortierungsnorm} anwenden.

\begin{korollar}{Gordon'sche Paare und normiterierende Mengen}
{GordonschePaareUndNormiterierendeMengen}
Seien $(E, \norm\cdot_E)$ ein normierter $\R$--Vektorraum, $x\in E^N$ und
$G_x$ der Gauß'sche Operator zum Tupel $(x,g)$.\\
Seien $\emptyset \neq T \subseteq S^{m-1}$,
$\emptyset\neq V\subseteq \R^N$ und $c \in \R_{> 0}$ mit
$\forall\, v \in V: \norm v_{2, N} \le c$.\\
Es gelte:
\begin{enumerate}[(a)]
	\item $0 \in V \folgt 0$ Häufungspunkt von $V$.
	\item $\forall\, w \in \R^N: \sup\menge{\sprod w v _N}{v \in V} 
	= \norm{\bprod w x_N}_E$.
\end{enumerate}
Dann sind nach \refclaim{IntegrierbarkeitVonExtremaInGordonschenPaaren}
{IntegrierbarkeitVonExtremaInGordonschenPaaren1}
\[\norm\cdot_E \circ \bprod \cdot x _N \circ \gamma, \quad
\inf_{t\in T}\bprod t \delta_m, \quad \sup_{t\in T}\bprod t \delta_m\]
und nach \refclaim{IntegrierbarkeitVonExtremaInGordonschenPaaren}
{IntegrierbarkeitVonExtremaInGordonschenPaaren2}
\[\inf_{t \in T}\, \norm\cdot_E \circ ev(t,\cdot)\circ G_x, \quad
\sup_{t \in T}\, \norm\cdot_E \circ ev(t,\cdot)\circ G_x \]
integrierbar und es gilt:
\begin{enumerate}[(1)]
	\item \label{GordonschePaareUndNormiterierendeMengen1}
	\begin{align*}
	& \E\left( \norm\cdot_E \circ \bprod \cdot x _N \circ \gamma\right)
	+ c \E\left(\inf_{t\in T}\bprod t \delta_m \right)\\
	\le &\,\E\left(\inf_{t \in T} \norm\cdot_E \circ ev(t,\cdot)\circ G_x \right)
	\le \E\left(\sup_{t \in T} \norm\cdot_E \circ ev(t,\cdot)\circ G_x \right)\\
	\le &\,\E\left( \norm\cdot_E \circ \bprod \cdot x _N \circ \gamma\right)
	+ c \E\left(\sup_{t\in T}\bprod t \delta_m \right)\text.
	\end{align*}
	\item \label{GordonschePaareUndNormiterierendeMengen2}
	Ist $T = -T$, so ist
	$\E\left(\Inf{t\in T}\bprod t \delta_m \right)
	= - \E\left(\Sup{t\in T}\bprod t \delta_m \right)$.
\end{enumerate}
\end{korollar}
\begin{proof}
(1) Sei $(X,Y)$ das Gordon'sche Paar zum Tupel $(\gamma,\delta,g,c)$.\\
Dann sind die Voraussetzungen von
\ref{GordonschePaareUndDerGaussscheMinMaxSatz} erfüllt.\\
Ferner gilt für alle $\omega \in \Omega$:
\begin{align}
\left(\sup_{v \in V}\bprod v \gamma_N \right)(\omega)
&= \sup_{v \in V}\sprod v {\gamma(\omega)_N}
=\sup_{v \in V}\sprod{\gamma(\omega)}v_N
\overset{\text{Vor.}}= \norm{\bprod {\gamma(\omega)} x}_N\notag\\
&= \left(\norm\cdot_E \circ \bprod \cdot x_N \circ \gamma\right) (\omega)\text,
\label{gpunmeq:1}
\end{align}
sowie für alle $s \in T$:
\begin{align}
\left(\sup_{v \in V}Y(s,v)\right)(\omega)&= \sup_{v \in V}\sprod{g(\omega)s}v_N
\overset{\text{Vor.}}= \norm{\bprod {g(\omega)s} x_N}_E
= \norm{G_x(\omega)s}_E \notag\\
&=\left(\norm\cdot_E \circ ev(s,\cdot)\circ G_x\right)(\omega)\text.
\label{gpunmeq:2}
\end{align}
Somit gilt wegen der Integrierbarkeit aller vorkommenden Abbildungen:
\begin{align*}
&\E\left( \norm\cdot_E \circ \bprod \cdot x _N \circ \gamma\right)
	+ c \E\left(\inf_{t\in T}\bprod t \delta_m \right)\\
	\overset{c>0}=&\,
	\E\left( \norm\cdot_E \circ \bprod \cdot x _N \circ \gamma\right)
	+  \E\left(\inf_{t\in T}c\bprod t \delta_m \right)\\
\overset{\eqref{gpunmeq:1}}= &\, \E\left(\sup_{v \in V}\bprod v \gamma_N
+ \inf_{t\in T}c\bprod t \delta_m\right)\\
\overset{
\refclaim{BeschraenktheitGordonscherPaare}{BeschraenktheitGordonscherPaare3}}=&
\,\E \left(\inf_{t \in T}\sup_{v \in V}X(t, v) \right)
\overset{\refclaim{GordonschePaareUndDerGaussscheMinMaxSatz}
{GordonschePaareUndDerGaussscheMinMaxSatz1}}
\le \E\left(\inf_{t \in T}\sup_{v \in V}Y(t, v) \right)\\
\overset{\eqref{gpunmeq:2}}= & \,
\E\left(\inf_{t \in T}\norm\cdot_E \circ ev(t,\cdot) \circ G_x \right)
\le
\E\left(\sup_{t \in T}\norm\cdot_E \circ ev(t,\cdot) \circ G_x \right)\\
\overset{\eqref{gpunmeq:1}}= 
&\, \E\left(\sup_{t \in T}\sup_{v \in V}Y(t, v) \right)
\overset{\refclaim{GordonschePaareUndDerGaussscheMinMaxSatz}
{GordonschePaareUndDerGaussscheMinMaxSatz2}}\le
\E\left(\sup_{t \in T}\sup_{v \in V}X(t, v) \right)\\
\overset{\refclaim{BeschraenktheitGordonscherPaare}
{BeschraenktheitGordonscherPaare3}} = &\,
\E\left(\sup_{v \in V}\bprod v \gamma_N
+ \sup_{t\in T}c\bprod t \delta_m\right)\\
\overset{\eqref{gpunmeq:1}}=&\,
\E\left( \norm\cdot_E \circ \bprod \cdot x _N \circ \gamma\right)
	+  \E\left(\sup_{t\in T}c\bprod t \delta_m \right)
\end{align*}
(2) Es gelte $T=-T$ und sei $\omega \in \Omega$. Dann gilt:
\begin{align*}
\left(\inf_{t\in T}\bprod t \delta_m\right)(\omega)
&= \inf_{t \in T}\sprod t {\delta(\omega)}_m
= - \sup_{t\in T}\left(-\sprod t {\delta(\omega)}\right)
= - \sup_{t\in T}\left(\sprod {-t} {\delta(\omega)}\right)\\
&=- \sup_{t\in -T}\left(\sprod t {\delta(\omega)}\right)
=- \sup_{t\in T}\left(\sprod t {\delta(\omega)}\right)
= \left(- \sup_{t\in T}\bprod t \delta_m \right)(\omega)\text.
\end{align*}
Da die genannten Abbildungen integrierbar sind, gilt somit (2).
\end{proof}
Das nächste Korollar ist in der vorgestellten Form eine Nebenrechnung aus
\cite{Gor2} (verkürzt auch in \cite{Gor} verwendet).

\begin{korollar}{Gordon'sche Paare und subeuklidische Bildhalbnormen}
{GordonschePaareUndSubeuklidischeBildhalbnormen}
Seien $(E, \norm\cdot_E)$ ein normierter $\R$--Vektorraum, $0\neq x \in E^N$,
$G_x$ der Gauß'sche Operator zum Tupel $(x,g)$ und $\emptyset\neq T
\subseteq S^{m-1}$.\\
Es gelte $\norm\cdot_E \circ \bprod \cdot x_N \le \norm\cdot_{2, N}$.\\
Dann sind nach~\refclaim{IntegrierbarkeitVonExtremaInGordonschenPaaren}
{IntegrierbarkeitVonExtremaInGordonschenPaaren1}
\[\norm\cdot_E \circ \bprod \cdot x _N \circ \gamma, \quad
\inf_{t\in T}\bprod t \delta_m, \quad \sup_{t\in T}\bprod t \delta_m\]
und nach nach~\refclaim{IntegrierbarkeitVonExtremaInGordonschenPaaren}
{IntegrierbarkeitVonExtremaInGordonschenPaaren2}
\[\inf_{t \in T}\, \norm\cdot_E \circ ev(t,\cdot)\circ G_x, \quad
\sup_{t \in T}\, \norm\cdot_E \circ ev(t,\cdot)\circ G_x \]
integrierbar und es gilt:
\begin{enumerate}[(1)]
	\item \label {GordonschePaareUndSubeuklidischeBildhalbnormen1}
\begin{align*}
& \E\left( \norm\cdot_E \circ \bprod \cdot x _N \circ \gamma\right)
	+ \E\left(\inf_{t\in T}\bprod t \delta_m \right)\\
	\le &\,\E\left(\inf_{t \in T} \norm\cdot_E \circ ev(t,\cdot)\circ G_x \right)
	\le \E\left(\sup_{t \in T} \norm\cdot_E \circ ev(t,\cdot)\circ G_x \right)\\
	\le &\,\E\left( \norm\cdot_E \circ \bprod \cdot x _N \circ \gamma\right)
	+ \E\left(\sup_{t\in T}\bprod t \delta_m \right)\text.
\end{align*}
	\item \label {GordonschePaareUndSubeuklidischeBildhalbnormen2}
	Ist $T = - T$, so gilt
	$\E\left(\Inf{t\in T}\bprod t \delta_m \right)
	= - \E\left(\Sup{t\in T}\bprod t \delta_m \right)$.
\end{enumerate}
\end{korollar}
\begin{proof}
(1) Sei $V\coloneqq  \menge{\pfeil \varphi N (x)}{\varphi \in \kabg{E'}01}$.\\
Dann ist $V$ nichtleer und, wie bereits im Beweis von
\refclaim{IntegrierbarkeitVonExtremaInGordonschenPaaren}
{IntegrierbarkeitVonExtremaInGordonschenPaaren1} gesehen, beschränkt.\\
Sei $v \in V$. Dann gibt es ein $\varphi \in \kabg{E'}01$ mit $v=\pfeil\varphi N(x)$.\\
Ferner gilt:
\begin{align*}
\norm v_{2, N}^2 &= \sprod v v_N = \sprod v {\pfeil \varphi N(x)}_N
\overset{\refclaim{EigenschaftenDesLinearkombinators}
{EigenschaftenDesLinearkombinators2}}=
\varphi\left(\bprod v x_N \right) \le \abs{\varphi\left(\bprod v x_N \right)}\\
&\le \opnorm \varphi E \R\norm{\bprod v x_N}_E\overset{\varphi \in \kabg{E'}01}\le
\norm{\bprod v x_N}_E \overset{\text{Vor.}}\le \norm v _{2, N}
\end{align*}
Damit ist $\norm v _{2, N} \le 1$.\\
Daher ist $c\coloneqq 1$ eine obere Schranke von $\menge{\norm w _{2, N}}{w \in V}$.
\hfill $(\ast)$\\
Offensichtlich ist $0 \in V$. Wegen $x \neq 0$ existiert ein $i \in \haken N$
mit $x_i \neq 0$.\\ Nach~\ref{NormUndDualeNorm} existiert ein 
$\alpha \in \kran{E'}01$ mit $\alpha(x_i) = \norm {x_i}_E$.\\
Definiere $v_n \coloneqq  \frac 1 n \pfeil \alpha N (x)$ für alle $n \in \N$.
Dann ist $(v_n)_{n \in \N}$ eine Nullfolge.\\
Für alle $n \in \N$ gilt wegen $\opnorm \alpha E \R= 1$: $v_n \neq 0$ und
nach Definition von $V$ auch $v_n \in V$, also $v_n \in V\setminus \meng 0$
und damit $0$ ein Häufungspunkt von $V$.\\
Schließlich gilt für alle $y\in\R^N$:
\begin{align*}
\sup_{w\in V} \sprod y w _N 
&= \sup_{\tau \in \kabg{E'}01}\sprod y {\pfeil \tau N (x)}_N
\overset{\refclaim{EigenschaftenDesLinearkombinators}
{EigenschaftenDesLinearkombinators2}}= \sup_{\tau \in \kabg{E'}01} 
\tau\left(\bprod y x_N \right)\\& \overset{\ref{NormUndDualeNorm}}=
\norma{\bprod y x_N}_E\text.
\end{align*}
Damit sind wegen $(\ast)$ die Voraussetzungen von~\ref{GordonschePaareUndNormiterierendeMengen} mit $c=1$ erfüllt und die Behauptung ist dann genau die Aussage
\refclaim{GordonschePaareUndNormiterierendeMengen}
{GordonschePaareUndNormiterierendeMengen1}.\\
(2) wurde bereits in~\refclaim{GordonschePaareUndNormiterierendeMengen}
{GordonschePaareUndNormiterierendeMengen2} bewiesen.
\end{proof}

\chapter{Banachraumtheoretische Resultate}
\label{ch:kapitel4}

Für dieses Kapitel seien
$\Gamma: ]0, \infty[ \rightarrow \R,t \mapsto \int_0^\infty x^{t-1}\exp(-x)dx$
(vgl. \ref{ElementareEigenschaftenDerGammafunktion}) und
$a: \N \rightarrow \R,
n\mapsto\sqrt 2\frac{\Gamma\left(\frac{n+1}2\right)}
                    {\Gamma\left(\frac n2 \right)}$
(vgl. \ref{AbschaetzungEinesGammafunktionsquotienten}).\absatz
In diesem Kapitel werden uns zwei bestimmte Zufallsgrößen interessieren, die
wir nun im Zusammenhang mit einem abkürzenden Begriff einführen wollen.

\begin{defsatz}{Auswertungspaar}{Auswertungspaar}
Seien $m,n \in \N$, $((\Omega, \mfa, \Prim), (\gamma,\delta), (n,m))$
ein unabhängiges $(N(0,1), 2)$--Tupel (zu $(2, (n,m))$)
(vgl.~\ref{NormalNTupel}), $(E, \norm\cdot_E)$ ein normierter
$\R$--Vektorraum, $x \in E^n$ und schließlich 
$\emptyset \neq T \subseteq S^{m-1}$.\\
Definiere
\[Z_{\gamma, x}\coloneqq  \norm\cdot_E \circ \bprod \cdot x_n \circ \gamma,\quad
  M_{\delta, T}\coloneqq  \sup_{t \in T}\bprod t \delta_m\]
Dann sind $Z_{\gamma, x}$ und $M_{\delta, T}$ nach
\ref{IntegrierbarkeitVonExtremaInGordonschenPaaren} wohldefiniert und
integrierbar.\\
Das Paar $(Z_{\gamma, x}, M_{\delta, T})$ nennen wir
\textbf{Auswertungspaar bezüglich $(\gamma, \delta, x, T)$}.
\end{defsatz}

Wie wir in \ref{GordonschePaareUndNormiterierendeMengen} und
\ref{GordonschePaareUndSubeuklidischeBildhalbnormen} gesehen haben, stehen
Auswertungspaare bei geeigneter Wahl des $x$ in engem Zusammenhang zum
Gauß'schen Min--Max--Satz für Gordon'sche Paare.

Das Vektortupel $x$ der Auswertungspaare wird in diesem Kapitel meist aus einem
klassischen Satz der Banachraumtheorie stammen und dazu verwendet werden, mit
wahrscheinlichkeitstheoretischen Mitteln die Existenz bestimmter linearer
Abbildungen zuzusichern.
Die Menge $T$ wird die Eigenschaften dieser Abbildungen regulieren.
So wird zum Beispiel $T = S^{m-1}$ die Abbildungen injektiv machen.

Wir sind nun in der Lage, einen wichtigen Satz der Banachraumtheorie zu
beweisen.
Dieser sichert die Existenz gewisser linearer Abbildungen zu,
welche sich besonders gut dazu eignen, den Banach--Mazur--Abstand zweier
Räume sehr einfach abzuschätzen.

\begin{satz}{Existenz eingeschränkter Fast--Isometrien}
            {ExistenzEingeschraenkterFastIsometrien}
Seien $m, N \in \N$,
$((\Omega, \mfa, \Prim), (\gamma, \delta, \widetilde g), (N, m, mN))$
ein unabhängiges $(N(0,1),3)$--Tupel (zu $(3, (N, m, mN))$) und
$g\coloneqq \Mat_{m,N} \circ \widetilde g$.\\
Seien $(E, \norm\cdot_E)$ ein normierter $\R$--Vektorraum,
$0\neq x \in E^n$ und $\emptyset \neq T \subseteq S^{m-1}$.\\
Seien $G_x$ der Gauß'sche Operator zum Tupel $(x,g)$ 
(vgl.~\ref{GaussscherOperator}) und $(Z_{\gamma,x}, M_{\delta, T})$ das
Auswertungspaar zu $(\gamma, \delta, x, T)$ (vgl.~\ref{Auswertungspaar}).\\
Es gelte:
\begin{enumerate}[(a)]
	\item $\norm\cdot_E \circ \bprod \cdot x _N \le \norm\cdot_{2, N}$,
	\item $T = -T$,
	\item $\E\left(M_{\delta, T}\right) < \E\left(Z_{\gamma, x} \right)$.
\end{enumerate}
Dann gilt $\E(Z_{\gamma, x}) > 0$, $\E(M_{\delta, T}) \ge 0$ und es gibt 
$\mca \subseteq L(\ell_2^m, (E, \norm\cdot_E))$ mit:
\begin{enumerate}[(1)]
  \item\label{ExistenzEingeschraenkterFastIsometrien1}
  $\mca \subseteq \Bild(G_x)$, es gibt $Q \in \mfa$ mit $\Prim(Q)>0$ und
  $Q \subseteq G_x\inv(\mca)$, insbesondere also $\mca \neq \emptyset$.
	\item \label{ExistenzEingeschraenkterFastIsometrien2}
	$\forall\, A \in \mca: \inf\menge{\norm{At}_E}{t \in T} > 0$\\
	Falls $T = S^{m-1}$, so ist jedes $A \in \mca$ injektiv.
	\item \label{ExistenzEingeschraenkterFastIsometrien3}
	$\forall\, A \in \mca:
	\dfrac{\sup\menge{\norm{At}_E}{t \in T}}{\inf\menge{\norm{At}_E}{t \in T}}
	\le \dfrac{\E\left(Z_{\gamma, x} \right) + \E\left(M_{\delta, T}\right)}
	{\E\left(Z_{\gamma, x} \right) - \E\left(M_{\delta, T}\right)}
	=\dfrac{1 + \frac{\E(M_{\delta, T})}{\E(Z_{\gamma, x})}}
	       {1 - \frac{\E(M_{\delta, T})}{\E(Z_{\gamma, x})}}$\\
	Ist $T = S^{m-1}$, so gilt für jedes $A \in \mca$: 
	$A\inv\in \Inv(A(\R^m),\R^m)$ und
	\[\opnorm A {\R^m} E \opnorm{A\inv}{A(\R^m)}{\R^m}
	\le \dfrac{1 + \frac{\E(M_{\delta, T})}{\E(Z_{\gamma, x})}}
	       {1 - \frac{\E(M_{\delta, T})}{\E(Z_{\gamma, x})}}\text.\]
\end{enumerate}
Ist zusätzlich $\eps \in\!\  ]0, 1[$ mit
$\E\left(M_{\delta, T}\right) \le \eps\E\left(Z_{\gamma, x} \right)$, so gilt
auch
\begin{enumerate}
\item[(3')] \label{ExistenzEingeschraenkterFastIsometrien3strich}
$\forall\, A \in \mca:
	\dfrac{\sup\menge{\norm{At}_E}{t \in T}}{\inf\menge{\norm{At}_E}{t \in T}}
	\le \dfrac{1 + \eps}	{1 - \eps}$\\
	Gilt $T = S^{m-1}$, so ist für jedes $A \in \mca$:
	\[\opnorm A {\R^m} E \opnorm{A\inv}{A(\R^m)}{\R^m}
	\le \dfrac{1 + \eps}	{1 - \eps}\text.\]
\end{enumerate}
\end{satz}
\begin{proof}
Zunächst sind offenbar die Voraussetzungen von~\ref{GordonschePaareUndSubeuklidischeBildhalbnormen} erfüllt.\\
Seien
\[f\coloneqq  \inf_{t \in T} \, \norm \cdot_E \circ ev(t, \cdot) \circ G_x
\quad \text{ und }\quad
  h\coloneqq  \sup_{t \in T} \, \norm \cdot_E \circ ev(t, \cdot) \circ G_x.\]
Nach~\refclaim{IntegrierbarkeitVonExtremaInGordonschenPaaren}
{IntegrierbarkeitVonExtremaInGordonschenPaaren2} sind $f, h$ integrierbar.\\
Wegen (b) und~\refclaim{GordonschePaareUndSubeuklidischeBildhalbnormen}
{GordonschePaareUndSubeuklidischeBildhalbnormen2} gilt weiter:
\begin{align} \E\left(\inf_{t\in T}\bprod t \delta_m \right)
	= - \E\left(\sup_{t\in T}\bprod t \delta_m \right) = 
	- \E\left(M_{\delta,T}\right)\text.\tag {$\ast$}
\end{align}
Somit gilt wegen der Integrierbarkeit von $Z_{\gamma, x}$ und $M_{\delta, T}$
(siehe~\ref{Auswertungspaar}):
\begin{align*}
0 & \overset{(c)}< \E\left(Z_{\gamma, x}\right) - \E\left(M_{\delta, T} \right)
\overset{(\ast)}= \E\left(Z_{\gamma, x}\right) 
+ \E\left(\inf_{t\in T}\bprod t \delta_m \right)\\
&\overset{\refclaim{GordonschePaareUndSubeuklidischeBildhalbnormen}
{GordonschePaareUndSubeuklidischeBildhalbnormen1}}\le \E(f) \le \E(h)
\overset{\refclaim{GordonschePaareUndSubeuklidischeBildhalbnormen}
{GordonschePaareUndSubeuklidischeBildhalbnormen1}}\le
\E\left(Z_{\gamma, x}\right) + \E\left(M_{\delta, T} \right)\text.
\end{align*}
Somit $\E(Z_{\gamma,x}) - \E(M_{\delta, T}) \le
\E(Z_{\gamma,x}) + \E(M_{\delta, T})$ und daher $\E(M_{\delta, T}) \ge 0$, also
gilt wegen $(c)$ auch $\E(Z_{\gamma,x})>\E(M_{\delta, T}) \ge 0$.\\
Seien nun $a\coloneqq \E(Z_{\gamma,x}) - \E(M_{\delta, T})$ und
$b\coloneqq \E(Z_{\gamma,x}) + \E(M_{\delta, T})$.\\
Wegen $f \ge 0$ und $h \ge 0$ und nach obiger Rechnung gilt:
\[ 0 < a \le \E(f) \le E(h) \le b\text.\]
Mit der Integrierbarkeit von $f$ und $h$ sind somit die Voraussetzungen von~\ref{Erwartungswertlemma} für $f, h, a, b$ erfüllt und für
$Q \coloneqq \menge{\omega \in \Omega}{f(\omega)> 0 \wedge ah(\omega) \le bf(\omega)}$
gilt daher (gemäß~\ref{Erwartungswertlemma}): $Q\in \mfa$ und $\Prim(Q) > 0$.\\
Sei nun $\mca\coloneqq G_x(Q)$.
Dann ist $G_x\inv(\mca) \supseteq Q$.\\
Sei $A \in \mca$. Dann existiert ein $\omega \in Q$ mit $A = G_x(\omega)$.\\
Wegen $\omega \in Q$ ist also
\[0 < f(\omega) = \inf_{t \in T}\norm{G_x(\omega)t}_E 
    = \inf_{t \in T}\norm{At}_E \]
und außerdem $ah(\omega) \le bf(\omega)$, was wegen $f(\omega)> 0$ und $a >0$
auch $\frac{h(\omega)}{f(\omega)} \le \frac b a$ bedeutet, sodass mit
$\E(Z_{\gamma, x})>0$ gilt:
\begin{align*}
&\,\frac{\sup\menge{\norm{At}_E}{t \in T}}
       {\inf\menge{\norm{At}_E}{t \in T}}
 =\frac{\sup\menge{\norm{G_x(\omega)t}_E}{t \in T}}
       {\inf\menge{\norm{G_x(\omega)t}_E}{t \in T}}
 = \frac{h(\omega)}{f(\omega)}\\ \le & \,\frac b a = 
 \frac{\E(Z_{\gamma,x})+\E(M_{\delta, T})}{\E(Z_{\gamma,x})-\E(M_{\delta, T})}
 = \frac{\E(Z_{\gamma,x})(1 + \frac{\E(M_{\delta, T})}{\E(Z_{\gamma,x})})}
        {\E(Z_{\gamma,x})(1 - \frac{\E(M_{\delta, T})}{\E(Z_{\gamma,x})})}
 = \frac{1 + \frac{\E(M_{\delta, T})}{\E(Z_{\gamma,x})}}
        {1 - \frac{\E(M_{\delta, T})}{\E(Z_{\gamma,x})}}\text.
 \end{align*}
Ist $T = S^{m-1}$, so folgt mit $\inf\menge{\norm{At}_E}{t \in S^{m-1}} > 0$
und $S^{m-1} = \kran{\ell^2_m}01$ nach~\ref{InjektivitätLinearerAbbildungenZwischenNormiertenRäumen} die Injektivität
von $A$.\\
Weiter ist der Raum 
$(A(\R^m), \restrict{\norm\cdot_E}{A(\R^m)})$ ein endlichdimensionaler
normierter Raum und damit ein Banachraum.\\
Fasse nun $A$ als Abbildung $A:\R^m \rightarrow A(\R^m)$ auf.\\
Dann ist $A$ bijektiv, linear und stetig, sodass nach~\ref{OperatornormDerInversen} gilt:
\[\opnorm{A\inv}{A(\R^m)}{\R^m} 
= \frac 1 {\inf\menge{\norm{At}_E}{t \in S^{m-1}}}\text.\]
Wegen $\opnorm A {\R^m} E = \sup\menge{\norm{At}_E}{t \in S^{m-1}}$ gilt also
\[\opnorm A {\R^m} E \opnorm{A\inv}{A(\R^m)}{\R^m}
= \frac{\sup\menge{\norm{At}_E}{t \in S^{m-1}}}
       {\inf\menge{\norm{At}_E}{t \in S^{m-1}}}\text.\]
Damit sind $(1)$, $(2)$, $(3)$ gezeigt.\\
Sei nun $\eps\in\!\ ]0, 1[$ mit $\E(M_{\delta, T})\le \eps \E(Z_{\gamma,x})$.\\
Wegen $\E(Z_{\gamma,x}) > 0$ ist also 
$\frac{\E(M_{\delta, T})}{\E(Z_{\gamma,x})} \le \eps$ und daher
\[\frac{1 + \frac{\E(M_{\delta, T})}{\E(Z_{\gamma,x})}}
       {1 - \frac{\E(M_{\delta, T})}{\E(Z_{\gamma,x})}}
\le \frac{1 + \eps}{1 - \frac{\E(M_{\delta, T})}{\E(Z_{\gamma,x})}}
\le \frac{1 + \eps}{1 - \eps} \text.\]
Wegen $(3)$ ist damit auch $(3')$ gezeigt.
\end{proof}
\parag
In den späteren Anwendungen werden die ersten beiden Voraussetzungen dieses
Satzes ($(a),(b)$) sehr einfach zu prüfen sein und die Arbeit wird im Verifizieren von $(c)$ stecken.

Interessiert man sich nur für injektive Abbildungen, so kann man direkt
$T \coloneqq  S^{m-1}$ setzen.
Es stellt sich allerdings heraus, dass die
Voraussetzungen an das $x$ und das Auswertungspaar mit Hilfe des
Banach--Mazur--Abstandes implizit ausgedrückt werden können.
Die dadurch entstehende Bedingung ist einerseits handlicher, weil es eine einzige
Ungleichung ist, andererseits weniger einfach, weil in der Anwendung der
Banach--Mazur--Abstand bestimmt werden müsste.
Wir fassen dies in dem folgenden Korollar zusammen.

\begin{korollar}{Existenz fast isometrischer Abbildungen}
                {ExistenzFastIsometrischerAbbildungen}
Seien $n \in \N$ und $(E, \norm\cdot_E)$ ein normierter, $n$--dimensionaler $\R$--Vektorraum.\\
Seien weiter $\eps\in\!\ ]0,1[$, und $m \in \N$ mit
$m \le \frac 2 \pi \cdot \left(\frac{\eps}
{d_{BM}((E, \norm\cdot_E), \ell_2^n)} \right)^2 n$.\\
Dann existiert ein injektives $A \in L(\ell_2^m, E)$ mit:
	\[\opnorm A{\ell_2^m}{E}\opnorm{A\inv}{A\left(\ell_2^m\right)}{\ell_2^m}
  \le \frac {1+\eps}{1-\eps}\text.\]
Insbesondere existiert $F \le E$ mit $\dim F = m$ und 
\[d_{BM}\left(\left(F,\restrict{\norm\cdot_E}F\right),\ell_2^m\right)
\le\frac{1+\eps}{1 - \eps}\text.\]
\end{korollar}
\begin{proof}
Sei $d \coloneqq  d_{BM}((E, \norm\cdot_E), \ell_2^n)$.\\
Nach~\ref{NormäquivalenzbasisImEndlichdimensionalen} existiert eine Basis 
$x \in E^n$ mit
$\frac 1 d \norm\cdot_{2, n} \le \norm\cdot_E \circ \bprod \cdot x_n
\le \norm\cdot_{2, n}$.\hfill $(\ast)$\\
Seien $T\coloneqq S^{m-1}$,
\[\bigg((\Omega, \mfa, \Prim), (\gamma,\delta,\widetilde g),(n, m, mn)\bigg)
\text{ ein }(N(0,1), 3)\text{--Tupel zu }(3,(n,m,mn))\] (existiert nach
\ref{ExistenzNormalNTupel}) und $(Z_{\gamma, x}, M_{\delta,T})$ das
Auswertungspaar zu $(\gamma, \delta, x, T)$.\\
Wegen $(\ast)$ gilt insbesondere $x \neq 0$ und
$\norm\cdot_E \circ \bprod\cdot x_n \le \norm\cdot_{2, n}$.\hfill $(a)$\\
Es gilt: $T \neq \emptyset$, $T = - T$ und $T \subseteq S^{m-1}$.\hfill $(b)$\\
Schließlich gilt:
\begin{align*}
\E\left(M_{\delta,T}\right)
&=\E\left(\sup_{t\in S^{m-1}}\bprod t\delta_m \right)
\overset{\ref{IterierteEuklidischeNorm}}=
\E\left(\norm\cdot_{2,m}\circ \delta \right)
\overset{\ref{ErwartungswertDer2NormEinesStandardnormalverteiltenProzesses}}
=a_m
\overset{\refclaim{AbschaetzungEinesGammafunktionsquotienten}
{AbschaetzungEinesGammafunktionsquotienten2}}\le \sqrt m
\overset{\text{Vor.}}\le \frac \eps d \sqrt{\frac 2 \pi n}\\
&\overset{\refclaim{AbschaetzungEinesGammafunktionsquotienten}
{AbschaetzungEinesGammafunktionsquotienten4}}\le  \frac \eps d a_n
\overset{\ref{ErwartungswertDer2NormEinesStandardnormalverteiltenProzesses}}=
\frac \eps d\E\left(\norm\cdot_{2,n} \circ \gamma \right)
\overset{(\ast)}\le 
\eps \E\left(\norm\cdot_E \circ \bprod \cdot x_n \circ \gamma \right)\\
&= \eps\E\left(Z_{\gamma,x} \right)\text.\tag {$c$}
\end{align*}
Wegen $(a), (b), (c)$ sind die Voraussetzungen von
\ref{ExistenzEingeschraenkterFastIsometrien} erfüllt und wegen $T = S^{m-1}$
existiert nach~\ref{ExistenzEingeschraenkterFastIsometrien} ein
$\emptyset \neq \mca \subseteq L(\ell_2^m, E)$, sodass jedes $A \in \mca$
injektiv ist und gilt:
\[\opnorm A {\ell_2^m}E \opnorm{A\inv}{A(\R^m)}{\ell_2^m} \le
\frac {1+\eps}{1-\eps}\text{.}\]
Wegen $\mca \neq \emptyset$ ist damit die Aussage (ohne den Zusatz) gezeigt.\\
Seien $A_0 \in \mca$ und $F\coloneqq  \Bild A_0$. Da $A_0$ injektiv ist, ist
$\dim F = m$ und damit ist $A_0 \in \Inv(\ell_2^m, F)$, sodass gilt:
\begin{align*}
d_{BM}\left(\ell_2^m, \left(F, \restrict{\norm\cdot_E}F\right) \right)
&= \inf\menge{\opnorm S {\ell_2^m}{F}\opnorm {S\inv}F{\ell_2^m}}
{S \in \Inv(\ell_2^m, F)}\\
& \le \opnorm {A_0} {\ell_2^m}E \opnorm{{A_0}\inv}{F}{\ell_2^m}
\le \frac {1+\eps}{1-\eps}\text.
\end{align*}
%Für alle $\omega \in \Omega$ gilt dann:
%5\begin{align*}
%(\norm\cdot_E \circ \bprod \cdot x_N\circ \gamma )(\omega)
%&= \norm{\bprod{\gamma(\omega)}x_N}_E \\
%&\overset{(\ast)}\ge 
%\frac 1 d \norm{\gamma(\omega)}_{2, N} = 
%\left(\frac 1 d \norm\cdot_{2, N} \circ \gamma \right)(\omega) \quad (\ast\ast)
%\end{align*}
%und
%\begin{align*}\sup_{t \in S^{m-1}}\bprod t \delta_m
% \overset{IterierteEuklidischeNorm}= \norm\cdot_{2, m}\circ \delta\quad 
%(\ast\ast\ast)\end{align*}
%Die jeweils ersten Abbildungen in den (Un)Gleichungen sind nach
%\refclaim{IntegrierbarkeitVonExtremaInGordonschenPaaren}
%{IntegrierbarkeitVonExtremaInGordonschenPaaren1} integrierbar, die jeweils
%letztgenannten sind es nach
%\ref{ErwartungswertDer2NormEinesStandardnormalverteiltenProzesses}.\\
%Wegen $(\ast)$ ist insbesondere $\norm\cdot_E\circ \bprod \cdot x_N \le
%\norm\cdot_{2, N}$.\\
%Ferner gilt $S^{m-1}= -S^{m-1}$ und
%\begin{align*}
%\E\left(\sup_{t \in S^{m-1}}\bprod t \delta_m \right)
%&%\overset{
%\substack{\ref{ErwartungswertDer2NormEinesStandardnormalverteiltenProzesses},
%\\(\ast\ast\ast)}}=
%a_m \overset{\refclaim{AbschaetzungEinesGammafunktionsquotienten}
%{AbschaetzungEinesGammafunktionsquotienten2}}\le \sqrt m \le
%\sqrt {\frac 1 2}\sqrt N \frac{\eps}d 
% \overset{\refclaim{AbschaetzungEinesGammafunktionsquotienten}
%{AbschaetzungEinesGammafunktionsquotienten4}}\le \eps \frac 1 d a_N\\
%&
%\overset{\ref{ErwartungswertDer2NormEinesStandardnormalverteiltenProzesses}}= 
%\E\left(\eps \frac 1 d \norm\cdot_{2, N}\circ \gamma  \right)
%\overset{(\ast\ast)}\le \eps\E(\norm\cdot_E\circ \bprod \cdot x_N \circ \gamma)
%\end{align*}
%Damit sind die Voraussetzungen von~\ref{ExistenzBanachMazurnaherAbbildungen}
%erfüllt und es gibt ein $A \in L(\R^m, E)$ injektiv mit
%$\opnorm A {\R^m}E\opnorm{A\inv}{A(\R^m)} {\R^m} \le \frac {1+\eps}{1-\eps}$.\\
%Sei $F\coloneqq \Bild A$. Dann ist $F \le E$ und da $A$ injektiv ist, ist $\dim F=m$.\\
%Ferner gilt wegen der Surjektivität von $A$ auf $F$:
%\begin{align*}
%d_{BM}((F, \norm\cdot_E\vert_{F}), l_2^m) &= 
%\inf\menge{ \opnorm{B}{\R^m}F \opnorm {B \inv}F {\R^m}}{B \in \Inv(F, \R^m)}\\
%&\le \opnorm A {\R^m}E\opnorm{A\inv}{A(\R^m)} {\R^m} \le \frac {1+\eps}{1-\eps}
%\end{align*}
\end{proof}

Mit den letzten beiden Sätzen können wir nun beispielsweise einfach und
elegant die Existenz gewisser Unterräume von $\ell_p^n$ zusichern.

\begin{korollar}{Abstände bezüglich der p--Norm}{AbstaendeBezueglichDerPnorm}
Seien $n \in \N$, $p \in [1, \infty]$, $\eps \in \,\! ]0,1[$ und $m \in \N$.\\
Dann gilt:
\begin{enumerate}[(1)]
\item\label{AbstaendeBezueglichDerPnorm1} Ist $p <\infty$ und $m\le\frac 2\pi\eps^2 n^{1-\abs{\frac2p-1}}$,
so gibt es ein injektives $A\in L(\ell_2^m, \ell_p^n)$, sodass gilt
\[\opnorm A
 {\ell_2^m}{\ell_p^n}\opnorm{A\inv}{A\left(\ell_2^m\right)}{\ell_2^m}
\le \frac {1+\eps}{1-\eps}\text.\]
\item\label{AbstaendeBezueglichDerPnorm2} Ist $p = \infty$ und $m \le \eps^2 \log(n)$, 
so gibt es ein injektives $A\in L(\ell_2^m, \ell_\infty^n)$, sodass gilt
\[\opnorm A
 {\ell_2^m}{\ell_\infty^n}\opnorm{A\inv}{A\left(\ell_2^m\right)}{\ell_2^m}
\le \frac {1+\eps}{1-\eps}\text.\]
\end{enumerate}
\end{korollar}
\begin{proof}
(1) Es gelte $p < \infty$ und
$m\le\frac 2\pi\eps^2 n^{1-\abs{\frac2p-1}}$.\\
Ist $p < 2$, so gilt nach \refclaim{AequivalenzkonstantenFuerPNormen}
{AequivalenzkonstantenFuerPNormen3}:
$\norm\cdot_{2,n}\le\norm\cdot_{p,n}\le n^{\frac 1p-\frac 12}\norm\cdot_{2,n}$,
also nach \refclaim{EigenschaftenDesBanachMazurAbstands}
{EigenschaftenDesBanachMazurAbstands3}
\begin{align}
d_{BM}\left(\ell_p^n, \ell_2^n \right) \le n ^{\frac 1 p - \frac 1 2}
= n ^{\abs{\frac 1 p - \frac 1 2}}\text.\label{abdpneq:1}
\end{align}
Ist $p \ge 2$, so gilt nach \refclaim{AequivalenzkonstantenFuerPNormen}
{AequivalenzkonstantenFuerPNormen3}:
$\norm\cdot_{p,n}\le\norm\cdot_{2,n}\le n^{\frac 12-\frac 1p}\norm\cdot_{p,n}$,
also nach \refclaim{EigenschaftenDesBanachMazurAbstands}
{EigenschaftenDesBanachMazurAbstands3}
\begin{align}
d_{BM}\left(\ell_2^n, \ell_p^n\right) \le n ^{\frac 1 2 - \frac 1 p}
= n ^{\abs{\frac 1 2 - \frac 1 p}} = n ^{\abs{\frac 1 p -\frac 12}}\text.
\label{abdpneq:2}
\end{align}
Somit gilt
\[m \le \frac 2 \pi \eps^2 n^{1- \abs{\frac 2 p - 1}}
= \frac 2 \pi \eps^2 \frac n {\left(n^{\abs{\frac 1 p - \frac 1 2}}\right)^2}
\overset{\substack{\eqref{abdpneq:1},\\\eqref{abdpneq:2}}}
\le\frac 2 \pi \eps^2 \frac n {d_{BM}(\ell_p^n, \ell_2^n)^2}\]
und damit sind die Voraussetzungen von
\ref{ExistenzFastIsometrischerAbbildungen} erfüllt.\\
Die Aussage (1) ist genau die Aussage von
\ref{ExistenzFastIsometrischerAbbildungen}.\\
(2) Es gelte $p=\infty$ und $m \le \eps^2 \log(n)$.\\
Seien $(E,\norm\cdot_E):=\ell_\infty^n$, $T:=S^{m-1}$, $x:= (e_i^n)_{i=1}^n$,
\[\bigg((\Omega, \mfa, \Prim), (\gamma,\delta,\widetilde g),(n, m, mn)\bigg)
\text{ ein }(N(0,1), 3)\text{--Tupel zu }(3,(n,m,mn))\] (existiert nach
\ref{ExistenzNormalNTupel}) und $(Z_{\gamma, x}, M_{\delta_T})$ das
Auswertungspaar zu $(\gamma, \delta, x, T)$.\\
Es gilt $x \neq 0$ und wegen $\bprod \cdot x_n = \id_{\R^n}$ gilt
\begin{align}\norm\cdot_E \circ\bprod \cdot x _n
= \norm\cdot_E = \norm\cdot_{\infty, n}
\overset{\refclaim{AequivalenzkonstantenFuerPNormen}
{AequivalenzkonstantenFuerPNormen2}}\le \norm\cdot_{2, n}\text. \tag {$a$}
\end{align}
Ferner ist $\emptyset \neq T$, $T=-T$ und $T \subseteq S^{m-1}$.\hfill $(b)$\\
Schließlich gilt:
\begin{align*}
\E\left(M_{\delta, T} \right) 
&= \E\left(\sup_{t \in S^{m-1}} \bprod t \delta_m \right)
\overset{\ref{IterierteEuklidischeNorm}}= 
\E\left(\norm\cdot_{2,m} \circ \delta \right)
\overset{\ref{ErwartungswertDer2NormEinesStandardnormalverteiltenProzesses}}=
a_m \overset{\refclaim{AbschaetzungEinesGammafunktionsquotienten}
{AbschaetzungEinesGammafunktionsquotienten2}}\le \sqrt m\\
&\overset{\text{Vor.}}\le \eps \sqrt{\log(n)}
\overset{
\refclaim{UntereSchrankeDesErwartungswertsVonAbsolutenN01Ordnungsstatistiken}
{UntereSchrankeDesErwartungswertsVonAbsolutenN01Ordnungsstatistiken2}}\le
\eps\E\left(\gamma_1^\ast \right) \overset{\ref{Ordnungsstatistik}}=
\eps \E\left(\norm\cdot_{\infty, n}\circ\gamma\right)
=\eps\E\left(Z_{\gamma,x}\right)\text. \tag {$c$}
\end{align*}
Wegen $(a), (b), (c)$ sind die Voraussetzungen von
\ref{ExistenzEingeschraenkterFastIsometrien} erfüllt und wegen 
$T = S^{m-1}$ gibt es 
$\emptyset \neq \mca \subseteq L(\ell_2^m, \ell_\infty^n)$ mit
\[\forall\, A \in \mca: A \text{ injektiv } \wedge \, 
\opnorm A
{\ell_2^m}{\ell_\infty^n}\opnorm{A\inv}{A\left(\ell_2^m\right)}{\ell_2^m}
\le \frac {1+\eps}{1-\eps}\text.\]
Wegen $\mca \neq \emptyset$ ist damit insbesondere (2) gezeigt.
\end{proof}
\parag
Bereits in Kapitel 1 hatten wir in \ref{AequivalenzkonstantenFuerPNormen}
verschiedene $p$--Normen verglichen und später das Dvoretzky--Rogers--Lemma
als (in gewisser Hinsicht) Verschärfung dieser Vergleiche kennengelernt.
Der auf dem Dvoretzky--Rogers--Lemma beruhende Satz von Dvoretzky ermöglicht es nun,
unter einer sehr ähnlichen Bedingung wie in \refclaim{AbstaendeBezueglichDerPnorm}{AbstaendeBezueglichDerPnorm2} den
vorher verwendeten Raum $\ell_\infty^n$ durch einen beliebigen normierten Raum
zu ersetzen und das gleiche Resultat zu erhalten.

\begin{satz}{Satz von Dvoretzky}{SatzVonDvoretzky}
Seien $\eps \in\!\  ]0,1[$, $n \in \N_{\ge 4}$, $m \in \haken n$ mit
$m\le\frac 14\eps^2\log\left(\floor{\nicefrac{\sqrt n}2}\right)$ und
schließlich
$(E, \norm\cdot_E)$ ein $n$--dimensionaler, normierter $\R$--Vektorraum.\\
Dann existiert ein $A \in L(\ell_2^m, E)$ injektiv mit
\[\opnorm A {\ell_2^m}E\opnorm{A\inv}{A(\ell_2^m)}{\ell_2^m}  \le
\frac {1+\eps}{1-\eps}\text.\]
Insbesondere existiert ein $F \le E$ mit $\dim F = m$ und
$d_{BM}(F, \ell_2^m) \le \dfrac{1+\eps}{1-\eps}$.
\end{satz}
\begin{proof}
Sei $N\coloneqq \floor{\nicefrac{\sqrt n}2}$. Wegen $n \ge 4$ ist $N \ge 1$ und es gilt
$4N^2 \le n$.\\
Nach \ref{DvoretzkyRogersLemma} gibt es $x \in E^N$ mit
\begin{align}
\frac 1 2 \norm\cdot_{\infty, N} \le \norm\cdot_E \circ \bprod \cdot x_N
\le \norm\cdot_{2, N}\text. \label{svdeq:1}
\end{align}
Seien $T \coloneqq  S^{m-1}$,
\[\bigg((\Omega, \mfa, \Prim), (\gamma,\delta,\widetilde g),(N, m, mN)\bigg)
\text{ ein }(N(0,1), 3)\text{--Tupel zu }(3,(N,m,mN))\] (existiert nach
\ref{ExistenzNormalNTupel}) und $(Z_{\gamma, x}, M_{\delta,T})$ das
Auswertungspaar zu $(\gamma, \delta, x, T)$.\\
Wegen $\eqref{svdeq:1}$ gilt $x \neq 0$ (erste Ungleichung) und
$\norm\cdot_E \circ \bprod\cdot x_N \le \norm\cdot_{2, N}$. \hfill $(a)$\\
Es ist $T \neq \emptyset$, $T = -T$ und $T \subseteq S^{m-1}$. \hfill $(b)$\\
Außerdem gilt:
\begin{align*}
\eps\E\left(Z_{\gamma, x} \right) 
& = \eps\E\left(\norm\cdot_E\circ \bprod \cdot x _N \circ \gamma \right)
\overset{\eqref{svdeq:1}}\ge
\frac 12\eps\E\left(\norm\cdot_{\infty, N}\circ\gamma \right)\\
&\overset{\ref{Ordnungsstatistik}}= \frac 1 2\eps\E(\gamma_1^\ast) \overset{
\refclaim{UntereSchrankeDesErwartungswertsVonAbsolutenN01Ordnungsstatistiken}
{UntereSchrankeDesErwartungswertsVonAbsolutenN01Ordnungsstatistiken2}}\ge
\frac 12\eps\sqrt{\log(N)}
=\frac 12\eps\sqrt{\log\left(\floor{\frac{\sqrt n}2}\right)}
\overset{\text{Vor.}}\ge \sqrt m\\
&\overset{\refclaim{AbschaetzungEinesGammafunktionsquotienten}
{AbschaetzungEinesGammafunktionsquotienten2}}\ge a_m
\overset{\ref{ErwartungswertDer2NormEinesStandardnormalverteiltenProzesses}}=
\E\left(\norm\cdot_{2, m} \circ \delta \right)
\overset{\ref{IterierteEuklidischeNorm}}= 
\E\left(\sup_{t \in T}\bprod t \delta_m \right) = \E(M_{\delta, T})\text.
\tag{$c$}
\end{align*}
Mit $(a),(b),(c)$ sind die Voraussetzungen von
\ref{ExistenzEingeschraenkterFastIsometrien} erfüllt, also existiert wegen
$T=S^{m-1}$ und \refclaim{ExistenzEingeschraenkterFastIsometrien}
{ExistenzEingeschraenkterFastIsometrien2} ein $A\in L(\ell_2^m,E)$ injektiv mit
\begin{align}\opnorm A {\ell_2^m}E\opnorm{A\inv}{A(\ell_2^m)}{\ell_2^m} \overset{
\ref{ExistenzEingeschraenkterFastIsometrien}.(3')} \le
\frac {1+\eps}{1-\eps}\text. \label{svdeq:2}
\end{align}
Sei nun $F \coloneqq  A(\R^m)$.\\
Da $A$ injektiv ist, ist $\dim F = m$, $F \le E$ und
da $F = \Bild(A)$ gilt auch
\[d_{BM}\left(\left(F, \restrict{\norm\cdot_E}F\right), \ell_2^m \right)
\le\opnorm A {\ell_2^m}F\opnorm{A\inv}{F}{\ell_2^m} \overset{\eqref{svdeq:2}}\le
\frac {1+\eps}{1-\eps}\text.\]
\end{proof}
\parag
Der wesentliche Aspekt des Beweises ist die Existenz eines $x \in E^n$ mit
der verwendeten Abschätzung nach oben und unten.
In der Dvoretzky--Rogers--Variante kann diese Abschätzung nicht weiter verbessert werden,
weil die Voraussetzung an das $n$ bzw. $N$ wesentlich in die Abschätzung eingeht.
Allerdings gibt es abstraktere Verbesserungen, wie die
sogenannte proportionale Dvoretzky--Rogers--Faktorisierung (vgl. \cite{BS}).
Die dort verwendete Konstante für die Abschätzung nach unten
($\nicefrac 12$ bei Dvoretzky--Rogers) ist allerdings schlecht explizit bestimmbar, weil sie
sich aus axiomatisch existenten Faktoren zusammensetzt.
Nichtsdestoweniger
werden wir auf genau diese proportionale Faktorisierung später im Satz von
Milman--Schechtman zurückkommen.\\
Wir beweisen nun das Lemma von Johnson und Lindenstrauss, wobei wir im Beweis
zum ersten Mal ein anderes $T$ als $S^{m-1}$ verwenden werden.

\begin{satz}{Johnson--Lindenstrauss--Lemma}{JohnsonLindenstraussLemma}
Seien $\eps \in\!\  ]0,1[$, $m, n, N \in \N$ mit $n \ge 2$ und
$\log(n) \le \eps^2 \frac 1{4\left(2+\nicefrac 4 {\log(2)} \right)}\cdot
\frac{N^2}{N+1}\text{,}$ sowie $v : \haken n \rightarrow \R^m$ injektiv.\\
Dann existiert ein $A\in L(\ell_2^m, \ell_2^N)$ mit
\[\forall\, i, j \in \haken n: i \neq j \folgt
1 - \eps \le \frac{\norm{A(v(i)) - A(v(j))}_{2,N}}{\norm{v(i)-v(j)}_{2,m}}
\le 1 + \eps\]
\end{satz}
\begin{proof}
Seien $(E, \norm\cdot_E)\coloneqq  \ell_2^N$, $x \coloneqq  (e_i^N)_{i=1}^N$ und
\[T\coloneqq  \mengea{\frac{v(i)-v(j)}{\norm{v(i)-v(j)}_{2,m}}}{i, j \in \haken n,
\, i \neq j}\text{.}\]
Dann ist $\bprod \cdot x_N = \id_{\R^N}$ und damit $x \neq 0$ und
$\norm\cdot_E\circ \bprod \cdot x _N = \norm\cdot_{2, N}$. \hfill $(a)$\\
Wegen $n\ge 2$ ist $T\neq \emptyset$, $T \subseteq S^{m-1}$ und $T=-T$. 
\hfill$(b)$\\
Seien $k \coloneqq  \abs T$, $\alpha \coloneqq  (N, m, mN, k)$,
\begin{align*}
\bigg((\Omega,\mfa,\Prim),(\gamma,\delta,\widetilde g,\eta),\alpha\bigg)
\text{ ein }(N(0,1), 4)\text{--Tupel zu }(4,\alpha)\end{align*}
 (existiert nach
\ref{ExistenzNormalNTupel}), sowie $(Z_{\gamma, x}, M_{\delta, T})$ das
Auswertungspaar zu $(\gamma, \delta, x, T)$.\\
F¸r jedes $s \in T$ ist die Abbildung $\bprod s \delta_m$ nach
\refclaim{IntegrierbarkeitUndVerteilungVonLinearkombinationen}
{IntegrierbarkeitUndVerteilungVonLinearkombinationen4} wegen $s \in S^{m-1}$
ebenfalls $N(0,1)$--verteilt.\\
Wegen~\ref{MaximumZentrierterNormalverteilterZufallsgroessen} (mit
$\sigma_i\coloneqq 1$ f¸r alle $i \in \haken k$) gilt daher
\begin{align}\E\left(\sup_{t \in T} \bprod t \delta_m \right) \le
\sqrt 2 \E\left(\max_{i \in \haken k}\eta_i \right)\text. \label{jlleq:1}
\end{align}
Außerdem gilt
\begin{align}
\max_{i \in \haken k}\eta_i \le \max_{i \in \haken k}\abs{\eta_i}
= \norm\cdot_{\infty, k} \circ \eta \overset{\ref{Ordnungsstatistik}}
= \eta_1^\ast \text.\label{jlleq:2}
\end{align}
Es ist:
\begin{align}
k = \abs{T} \le n^2\text{ und da } n \ge 2\text{ und }
v\text{ injektiv auch } k \ge 2\text.\label{jlleq:3}
\end{align}
Somit gilt:
\begin{align*}
\E(M_{\delta, T}) & = \E\left(\sup_{t \in T}\bprod t \delta_m \right)
\overset{\eqref{jlleq:1}}\le 
\sqrt 2 \E\left(\max_{i \in \haken k}\eta_i \right)
\overset{\eqref{jlleq:2}}\le \sqrt 2 \E\left(\eta_1^\ast \right)\\
&\overset{\refclaim
{ObereAbschaetzungDesErwartungswertsDerKtenAbsolutenN01Ordnungsstatistik}{ObereAbschaetzungDesErwartungswertsDerKtenAbsolutenN01Ordnungsstatistik4}}
\le \sqrt 2 \sqrt{2+\frac 4{\log(2)}}
	\sqrt{\log\left(k \right)}\overset{\eqref{jlleq:3}}\le
	\sqrt 2 \sqrt{2+\frac 4{\log(2)}}
	\sqrt{\log\left(n^2 \right)}\\
&=
	\sqrt 2 \sqrt{2+\frac 4{\log(2)}}
	\sqrt{2\log\left(n \right)}
	= 2 \sqrt{2+\frac 4{\log(2)}}\sqrt{\log(n)}\\
	&\overset{\substack{\text{Vor.}\\ \text{an }n}}\le
	2 \sqrt{2+\frac 4{\log(2)}} \eps \frac 1 {2 \sqrt{2+\frac 4{\log(2)}}}\cdot
	\frac N {\sqrt{N+1}} = \eps \frac N {\sqrt{N+1}}\\
	&\overset{\refclaim{AbschaetzungEinesGammafunktionsquotienten}
	{AbschaetzungEinesGammafunktionsquotienten3}}\le \eps a_N
	\overset{\ref{ErwartungswertDer2NormEinesStandardnormalverteiltenProzesses}}=
	\eps\E\left(\norm\cdot_{2, N} \circ \gamma \right) \overset{(a)}=
	\eps \E\left( \norm\cdot_E \circ \bprod\cdot x_N \circ \gamma\right)\\
	&= \eps \E(Z_{\gamma, x})\text. \tag{$c$}
\end{align*}
Wegen $(a), (b), (c)$ sind die Voraussetzungen von~\ref{ExistenzEingeschraenkterFastIsometrien} erf¸llt und daher existiert ein
$\emptyset \neq \mcb \subseteq L(\ell_2^m, \ell_2^N)$ mit:
\begin{align}
\forall\, B \in \mcb: \inf_{t \in T} \norm{Bt}_{2,N} > 0 \quad \wedge \quad
\frac{\sup\menge{\norm{Bt}_{2,N}}{t\in T}}{\inf\menge{\norm{Bt}_{2,N}}{t\in T}}
\le \frac{1+\eps}{1-\eps}\text. \label{jlleq:4}
\end{align}
Sei nun
$\mca \coloneqq  \menge{\frac{1-\eps}{\inf\menge{\norm{Bt}_{2,N}}{t \in T}}B}
{B \in \mcb}$.\\
Offensichtlich ist auch $\emptyset \neq \mca \subseteq L(\ell_2^m,\ell_2^N)$.\\
Sei $A \in \mca$. Dann gibt es ein $B \in \mcb$ mit
$A = (\nicefrac{(1-\eps)}{\inf\menge{\norm{Bt}_{2,N}}{t \in T}})B$.\\
Seien $\alpha\coloneqq  \inf\menge{\norm{Bt}_{2,N}}{t \in T}$ und $i, j \in \haken n$
mit $i \neq j$.\\
Dann ist $t_{ij}\coloneqq  \frac{v(i)-v(j)}{\norm{v(i)-v(j)}_{2,m}} \in T$ und es gilt:
\begin{align*}
&\,\frac{\norm{A(v(i)) - A(v(j))}_{2,N}}{\norm{v(i)-v(j)}_{2,m}}=
\norma{A\left(\frac{v(i)-v(j)}{\norm{v(i)-v(j)}_{2,m}} \right)}_{2, N}
= \norm{A(t_{ij})}_{2, N}\\
=&\, \frac {1-\eps}\alpha \norm{B(t_{ij})}_{2,N}
\le \frac {1-\eps}\alpha \sup_{t \in T}\,\norm{Bt}_{2,N}
\overset{\text{Def. }\alpha}=(1-\eps)
\frac{\sup\menge{\norm{Bt}_{2,N}}{t \in T}}
           {\inf\menge{\norm{Bt}_{2,N}}{t \in T}}\\
\overset{\eqref{jlleq:4}}\le&\, 
(1 - \eps)\frac{1+\eps}{1-\eps} = (1+ \eps)\text.
\end{align*}
Außerdem gilt vollkommen analog:
\begin{align*}
&\,\frac{\norm{A(v(i)) - A(v(j))}_{2,N}}{\norm{v(i)-v(j)}_{2,m}}=
\frac {1-\eps}\alpha \norm{B(t_{ij})}_{2,N}
\ge \frac {1-\eps}\alpha \inf_{t \in T}\,\norm{Bt}_{2,N}
= 1-\eps\text.
\end{align*}
Da $\mca$ nichtleer ist, ist damit die Behauptung gezeigt.
\end{proof}
\parag
Wir können also eine endliche Menge von normierten Vektorendifferenzen so
abbilden, dass alle Bilder \enquote{fast} normiert sind.
Unter etwas schärferen Voraussetzungen können aber auch bestimmte unendliche Mengen (nicht
notwendigerweise spezieller Differenzen) ähnlich abgebildet werden.

\begin{satz}{Restricted Isometry Property}{RIP}
Seien $\eps \in\!\  ]0,1[$, $m \in \N$, $N \in \haken {m-1}$ und
$\eta \in \!\ ]0, \frac \eps{2+\eps}[$.\\
Sei $k \in \N$ mit $k \le \frac m {\exp(1)}$ und
$k\log\left(\frac m k \right) \le \frac 1 {4 \left(1+ \frac 2{\log(2)}
 \right)}\eta^2 \frac {N^2}{N+1}$.\\
Sei schließlich
\[\Sigma_k^m \coloneqq  \menge{x \in\R^m}{\abs{\menge{i\in\haken m}{x_i\neq 0}}\le k}\]
Dann gibt es $A \in L(\ell_2^m, \ell_2^N)$ mit
\begin{enumerate}[(1)]
	\item \label{RIP1}
	$\forall\, x \in \Sigma_k^m \cap S^{m-1}: Ax \neq 0$
	\item \label{RIP2}
	$\forall\, x, y \in \Sigma_k^m\cap S^{m-1}:
	1-\eps < \frac{\norm{Ax}_{2, N}}{\norm{Ay}_{2, N}} < 1 + \eps$
\end{enumerate}
\end{satz}
\begin{proof}
Seien $(E, \norm\cdot_E)\coloneqq  \ell_2^N$, $T\coloneqq  \Sigma_k^m\cap S^{m-1}$ und
$x \coloneqq  \left(e_i^N\right)_{i=1}^N$.\\
Dann ist $\bprod\cdot x_N = \id_{\R^N}$ und damit $x \neq 0$ und
$\norm\cdot_E\circ\bprod \cdot x_N  = \norm\cdot_{2, N}$.
\hfill $(a)$\\
Ferner ist $T \neq \emptyset$ (da $e_1^m \in T$), $T \subseteq S^{m-1}$ und
offenbar $T = -T$. \hfill $(b)$\\
Sei
\[\bigg((\Omega, \mfa, \Prim), (\gamma,\delta,\widetilde g),(N, m, mN)\bigg)
\text{ ein }(N(0,1), 3)\text{--Tupel zu }(3,(N,m,mN))\] (existiert nach~\ref{ExistenzNormalNTupel}) und $(Z_{\gamma, x}, M_{\delta,T})$ das
Auswertungspaar zu $(\gamma, \delta, x, T)$.\\
Seien $s \in T$ und $\omega \in \Omega$.\\
Da $\abs{\menge{i \in \haken m}{s_i\neq 0}} \le k$, gilt ($\mu_{j,n}$ gemäß~\ref{ktesMaximum} für alle $j \in \haken n$):
\begin{align*}
&\,\bprod s \delta_m (\omega) = \sprod s {\delta(\omega)}_m
\le \sum_{j=1}^m \abs{s_j}\abs{\delta_j(\omega)}\\
\le&\, \sum_{j=1}^k \left(\mu_{j,n} \circ \pfeil {\abs\cdot} m\right)(s)
     \left(\mu_{j,n} \circ\pfeil {\abs\cdot} m\right)(\delta(\omega))
\overset{\ref{Ordnungsstatistik}}= \sum_{j=1}^k\left(\mu_{j,n}\circ\pfeil
{\abs\cdot} m \right)(s)\delta_j^\ast(\omega)\\
\le &\, \norma{\left(\left(\mu_{j,n}\circ\pfeil
{\abs\cdot} m \right)(s)\right) _{j=1}^k}_{2, k}
\norma{\left(\delta^\ast_j(\omega) \right)_{j=1}^k}_{2, k}
\overset{s \in T}= \norma{\left(\delta^\ast_j(\omega) \right)_{j=1}^k}_{2, k}\\
=&\, \left(\norm\cdot_{2,k}\circ(\delta_j^\ast)_{j\in \haken k} \right)(\omega)
\text.\tag {$\ast$}
\end{align*}
Somit gilt wegen $k \le \frac m {\exp(1)}$:
\begin{align*}
\E\left(M_{\delta, T} \right)&= \E\left(\sup_{t \in T}\bprod t \delta_m \right)
\overset{(\ast)}\le
\E\bigg(\norm\cdot_{2,k} \circ (\delta_j^\ast)_{j \in \haken k} \bigg)\\
&\overset{\refclaim
{ObereAbschaetzungDesErwartungswertsDerKtenAbsolutenN01Ordnungsstatistik}{ObereAbschaetzungDesErwartungswertsDerKtenAbsolutenN01Ordnungsstatistik6}}
\le \sqrt{2}\sqrt{2 + \frac 4 {\log(2)}}\sqrt{k \log\left(\frac m k \right)}
= 2 \sqrt{1 + \frac 2 {\log(2)}}\sqrt{k \log\left(\frac m k \right)}\\
&\overset{\substack{\text{Vor.}\\ \text{an }k}}\le \eta \frac N {\sqrt{N+1}}
\overset{\refclaim{AbschaetzungEinesGammafunktionsquotienten}
{AbschaetzungEinesGammafunktionsquotienten3}}\le \eta a_N
\overset{\ref{ErwartungswertDer2NormEinesStandardnormalverteiltenProzesses}}=
\eta \E \left(\norm\cdot_{2, N} \circ \gamma \right)
= \eta\E(Z_{\gamma, x})\text. \tag{$c$}
\end{align*}
Offenbar ist $\eta \in \!\  ]0,1[$ und mit $(a), (b), (c)$ sind die Voraussetzungen von~\ref{ExistenzEingeschraenkterFastIsometrien} erfüllt,
sodass ein $\emptyset \neq\mca\subseteq L(\ell_2^m, \ell_2^N)$ existiert mit
\[\forall\, A \in \mca:\,\inf\menge{\norm{At}_{2,N}}{t \in T} > 0\,\wedge
\,\frac{\sup\menge{\norm{At}_{2,N}}{t\in T}}
       {\inf\menge{\norm{At}_{2,N}}{t\in T}}\le\frac {1+ \eta}{1-\eta}\text.\]
Seien $A\in \mca$, $u_A \coloneqq  \inf\menge{\norm{At}_{2,N}}{t \in T}$ und
$o_A\coloneqq \sup\menge{\norm{At}_{2,N}}{t \in T}$.\\
Seien $x, y \in T = \Sigma_k^m\cap S^{m-1}$.\\
Da $u_A>0$ ist insbesondere $\norm{Ay}_{2,N}>0$, also $Ay\neq 0$, womit (1)
gezeigt ist.\\
Weiter gilt:
\[\frac{\norm{Ax}_{2, N}}{\norm{Ay}_{2,N}} 
\le \frac {o_A} {u_A} \le \frac {1+\eta}{1-\eta}
< \frac{1 + \frac \eps{2+\eps}}{1-\frac \eps{2+\eps}}
= \frac{\frac{2+2\eps}{2+\eps}}{\frac{2+\eps-\eps}{2+\eps}} = 1 + \eps\text,\]
sowie wegen $\eps \in\!\ ]0, 1[$ und damit $(1+\eps)(1-\eps) = 1 - \eps^2 < 1$:
\[\frac{\norm{Ax}_{2, N}}{\norm{Ay}_{2,N}} \ge\frac {u_A} {o_A} 
\ge \frac {1- \eta}{1+\eta}
> \frac{1-\frac \eps{2+\eps}}{1 + \frac \eps{2+\eps}}
= \frac{\frac{2+\eps-\eps}{2+\eps}}{\frac{2+\eps+\eps}{2+\eps}}= 
\frac 1 {1+\eps} > 1-\eps\text.\]
Also insgesamt
$1 - \eps < \nicefrac{\norm{Ax}_{2, N}}{\norm{Ay}_{2,N}} < 1+\eps$,
sodass auch (2) gilt.
\end{proof}
\parag
Die \enquote{Ausdünnung} der Vektoren in diesem Satz lässt höchstens ein Drittel der
Komponenten von $0$ verschieden.
Auch wenn dies zunächst sehr schwach erscheint
(tatsächlich wird die Anzahl der von $0$ verschiedenen Komponenten noch stärker
eingeschränkt), ist die Angabe einer linearen Abbildung mit den besagten
Eigenschaften mit klassischen Methoden nicht ohne weiteres möglich.\parag
In den Sätzen \ref{ExistenzFastIsometrischerAbbildungen},
\ref{AbstaendeBezueglichDerPnorm}, \ref{JohnsonLindenstraussLemma} und
\ref{RIP} wurde im Wesentlichen stets
\ref{ExistenzEingeschraenkterFastIsometrien} verwendet, in den Formulierungen
jedoch stets ein einzelner Operator und nicht wie in
\ref{ExistenzEingeschraenkterFastIsometrien} eine (in gewisser Hinsicht große)
Menge solcher Operatoren.
Der Grund für diese Abschwächung ist der, dass die \enquote{Größe} der Menge sich nicht ohne den Gauß'schen Operator auf
einem bestimmten Wahrscheinlichkeitsraum beschreiben lässt, während die
Aussagen der genannten Sätze für ihre Formulierung keinen solchen Raum
benötigen.\parag
Wir kommen nun zu einem weiteren Resultat, welches sich unter Zuhilfenahme des
Gauß'schen Min--Max--Satzes überraschend einfach beweisen lässt.
Dieses ist der Satz von Gluskin, welcher besagt, dass es in $\ell_1^m$ zu $N \in \haken m$
stets $N$--codimensionale Unterräume gibt, deren Abstand zum $\ell_2^{m-N}$ von
der Größenordnung $\sqrt{\frac m N \log\left(\frac{2m}N\right)}$ ist.
Eine konkrete (wenn auch recht grobe) Abschätzung ist mit unseren Methoden ebenfalls möglich.

\begin{satz}{Satz von Gluskin}{SatzVonGluskin}
Seien $m \in \N$, $\eps \in\!\ ]0, 1[$, $N\coloneqq  \ceil{(1-\eps)m}$ und
\[c \coloneqq  2\sqrt \pi (\sqrt 2 +1) \sqrt{2 + \frac 4{\log(4)}}
\sqrt{\frac{\log\left(\pi(\sqrt 2+1)^2\left(1+\frac 2{\log(4)}\right) \right)}
{\log(2)} + 1}\text.\]
Dann existiert $F \le \ell_1^m$ mit $\dim F = m - N$ und
\[d_{BM}(F,\ell_2^{m-N})<
c\sqrt{\frac 1{1-\eps}\log\left(\frac 2{1-\eps}\right)}\text.\]
\end{satz}
\begin{proof}
Seien $\rho \coloneqq  (\sqrt 2 +1)^2\left(2+\frac 4{\log(4)} \right)$,
$\eta \coloneqq  \frac 2 \pi \frac 1 \rho (1-\eps)$ und
$t_\eta \coloneqq  \frac 4 \eta \log\left(\frac 2 \eta \right)$.\\
Dann ist $c^2=4\pi\rho \left(\frac{\log\left(\frac{\pi\rho}2\right)}{\log(2)}+1\right)$ und damit:
\begin{align}
&2t_\eta = \frac{4\pi\rho}{1-\eps}\log\left(\frac{\pi\rho}{1-\eps} \right)
= \frac{4\pi\rho}{1-\eps}\log\left(\frac{\pi\rho}{2}\frac{2}{1-\eps} \right)
\notag\\
=& \frac{4\pi\rho}{1-\eps}\left(\log\left(\frac{\pi\rho}2 \right) 
 + \log\left( \frac 2 {1-\eps}\right) \right)
 \overset{\eps>-1}=
 \frac{4\pi\rho}{1-\eps}\left(\frac{\log\left(\frac{\pi\rho}2
 \right)}{\log\left( \frac 2 {1-\eps}\right)}+1 \right)\log\left(\frac 2
 {1-\eps} \right)\notag\\
 =&4\pi\rho\left(\frac{\log\left(\frac{\pi\rho}2 \right)}
 {\log\left( \frac 2 {1-\eps}\right)}+1 \right)
 \frac 1{1-\eps}\log\left(\frac 2 {1-\eps} \right)
 \overset{\frac 2{1-\eps}>2}<
 c^2
 \frac 1{1-\eps}\log\bigg(\frac 2 {1-\eps} \bigg). \label{svgeq:1}
\end{align}
Falls $N=m$ gilt, ist die Aussage trivial, sodass wir $N<m$ annehmen können.\\
\underline{1. Fall:} $m \le 2 t_\eta$.\\
Sei $F \coloneqq  \left\langle \menge{e_i^m}{i \in \haken {m-N}}\right\rangle$.\\
Offenbar ist dann $(F,\restrict{\norm\cdot_1}{F})$ isometrisch isomorph zu
$\ell_1^{m-N}$ und damit nach~\refclaim{EigenschaftenDesBanachMazurAbstands}
{EigenschaftenDesBanachMazurAbstands2}:
$d_{BM}((F,\restrict{\norm\cdot_1}{F}), \ell_1^{m-N}) = 1$.\\
Weiter gilt nach
\refclaim{AequivalenzkonstantenFuerPNormen}{AequivalenzkonstantenFuerPNormen1}
$\norm\cdot_{2, m-N} \le\norm\cdot_{1, m-N} \le \sqrt{m-N}\norm\cdot_{2, m-N}$
und somit nach~\refclaim{EigenschaftenDesBanachMazurAbstands}
{EigenschaftenDesBanachMazurAbstands3}: $d_{BM}(\ell_1^{m-N}, \ell_2^{m-N}) \le
\sqrt{m-N}$.\\
Somit gilt
\begin{align*}
d_{BM}\left((F,\restrict{\norm\cdot_1}F), \ell_2^{m-N}\right)
&\overset{\refclaim{EigenschaftenDesBanachMazurAbstands}
{EigenschaftenDesBanachMazurAbstands1}}\le
d_{BM}((F,\restrict{\norm\cdot_1}{F}), \ell_1^{m-N})d_{BM}
(\ell_1^{m-N}, \ell_2^{m-N})\\
& \le \sqrt{m-N} \le \sqrt{m - (1-\eps)m} = \sqrt{\eps m} < \sqrt{m}
\le \sqrt{2t_\eta}\\
&\overset{\eqref{svgeq:1}}\le
c \sqrt{\frac 1 {1-\eps}\log\left(\frac 2{1-\eps} \right)}
\text.
\end{align*}
Damit ist also $F$ ein Unterraum wie behauptet.\\
\underline{2. Fall:} $m > 2t_\eta$.\\
Dann ist $1<m\frac1{2t_\eta}=m\left(\frac1{t_\eta}-\frac1{2t_\eta}\right)$ und
damit $\N \cap \left[\frac 1 2 \frac m {t_\eta}, \frac m {t_\eta}\right]
\neq \emptyset$.\\
Sei $k\in\N\cap\left[\frac12\frac m{t_\eta},\frac m{t_\eta}\right]$.\\
Offenbar ist $\eta < 1$ und damit $t_\eta > 4 \log(2) = \log(16) > 2$.\\
Also ist $k \le \frac m {t_\eta} < \frac m 2$.\\
Seien $(E,\norm\cdot_E)\coloneqq \ell_2^N$, $T\coloneqq S^{m-1}\cap\kabg{\ell_1^m} 0{\sqrt k}$,
$x \coloneqq  (e_i)_{i = 1}^N$,
\[\bigg((\Omega, \mfa, \Prim), (\gamma,\delta,\widetilde g),(N, m, mN)\bigg)
\text{ ein }(N(0,1), 3)\text{--Tupel zu }(3,(N,m,mN))\] (existiert nach~\ref{ExistenzNormalNTupel}) und $(Z_{\gamma, x}, M_{\delta,T})$ das
Auswertungspaar zu $(\gamma, \delta, x, T)$.\\
Dann gilt $\bprod \cdot x_N = \id_{\R^N}$ und damit $x \neq 0$ und
$\norm\cdot_E\circ\bprod\cdot x _N=\norm\cdot_{2,N}$. \hfill $(a)$\\
Offenbar gilt $T \neq \emptyset$ (wegen $k \ge 1$ ist $e_1^m \in T$),
$T \subseteq S^{m-1}$, $T=-T$. \hfill $(b)$\\
Ferner ist $T$ als Schnitt $1$--symmetrischer Mengen selbst $1$--symmetrisch
und damit gilt ($\mu_{i,n}$ für alle $i \in \haken m$ gemäß~\ref{ktesMaximum})
\begin{align}
\forall\ a \in \R^m : \sup_{t \in T} \sprod t a_m = \sup_{t \in T} \sprod
t {((\mu_{i,n}\circ \pfeil {\abs\cdot} m)(a))_{i=1}^m}_m \text.\label{svgeq:2}
\end{align}
Sei $\omega \in \Omega$. Dann gilt für alle $s \in T$:
\begin{align}
 \sum_{i=k+1}^m s_i\delta_i^\ast(\omega) 
&\le \sum_{i=k+1}^m \abs{s_i}\delta_i^\ast(\omega) 
\overset{\refclaim{EigenschaftenDesKtenMaximums}
{EigenschaftenDesKtenMaximums1}}\le
\delta_{k+1}^\ast(\omega) \sum_{i=k+1}^m \abs{s_i} \notag\\
&\overset{\refclaim{EigenschaftenDesKtenMaximums}
{EigenschaftenDesKtenMaximums1}}\le \delta_k^\ast(\omega) \norm s _{1,m}
\le \sqrt k \cdot \delta_k^\ast(\omega) \text.\label{svgeq:3}
\end{align}
Weiter gilt wegen $k \le \frac m {t_\eta}$ auch $\frac m k \ge t_\eta$.\\
Da $\eta \in\!\  ]0,1[$ gilt nach Setzung von
$t_\eta$ gemäß~\refclaim{Numeriklemma}{Numeriklemma4}:
\begin{align}
\log\left(\frac{2m}k \right) < \eta \frac m k\text. \label{svgeq:4}
\end{align}
Somit gilt:
\begin{align*}
&\,\E\left(M_{\delta,T}\right)=\E\left(\sup_{t \in T} \bprod t \delta_m \right)
\overset{\eqref{svgeq:2}}=
\E\left(\sup_{t\in T}\bprod t{\delta^\ast}_m \right)\\
= &\, \E\left(\sup_{t \in T}\left(\left(\sum_{i=1}^k t_i\delta_i^\ast\right) 
+ \left(\sum_{i=k+1}^m t_i\delta_i^\ast \right) \right) \right)\\
\le &\, \E\left(\sup_{t \in T}\sum_{i=1}^k t_i\delta_i^\ast \right) 
+ \E\left(\sup_{t \in T} \sum_{i=k+1}^m t_i\delta_i^\ast\right) \\
\overset{\eqref{svgeq:3}}\le &\,
\E\left(\sup_{t \in T}\sum_{i=1}^k t_i\delta_i^\ast \right)
+ \sqrt k \cdot \E\left(\delta_k^\ast \right)\\
\le &\, \E\left(\sup_{t \in T}\norm{t}_{2,m}\left(
\norm\cdot_{2,k} \circ \left(\delta_i^\ast \right)_{i=1}^k \right) \right)
+ \sqrt k \cdot \E\left(\delta_k^\ast \right)\\
\overset{T \subseteq S^{m-1}}=
&\, \E\left(\norm\cdot_{2,k} \circ \left(\delta_i^\ast \right)_{i=1}^k \right)
+ \sqrt k \cdot \E\left(\delta_k^\ast \right)\\
\overset{\substack{\refclaim
{ObereAbschaetzungDesErwartungswertsDerKtenAbsolutenN01Ordnungsstatistik}
{ObereAbschaetzungDesErwartungswertsDerKtenAbsolutenN01Ordnungsstatistik5},\\
\refclaim
{ObereAbschaetzungDesErwartungswertsDerKtenAbsolutenN01Ordnungsstatistik}
{ObereAbschaetzungDesErwartungswertsDerKtenAbsolutenN01Ordnungsstatistik7}}}
\le&\,\sqrt 2\sqrt{2+ \frac 4 {\log(4)}} \sqrt{k \log\left(\frac {2m}k\right)}
+\sqrt k \cdot\sqrt{2 + \frac 4 {\log(4)}}\sqrt{\log\left(\frac {2m}k\right)}\\
=&\, \left(\sqrt{2} + 1\right)\sqrt{2+ \frac 4 {\log(4)}}
\sqrt{k \log\left(\frac {2m}k\right)}\\
\overset{\eqref{svgeq:4}}<
&\, \left(\sqrt{2} + 1\right)\sqrt{2+ \frac 4 {\log(4)}}
\sqrt{k \eta \frac mk}\\
=&\, \left(\sqrt{2} + 1\right)\sqrt{2+ \frac 4 {\log(4)}}
\sqrt{\eta m} = 
\sqrt{\frac 2 \pi} \sqrt{(1-\eps)m} \le \sqrt{\frac 2 \pi}
\sqrt N \overset{\refclaim{AbschaetzungEinesGammafunktionsquotienten}
{AbschaetzungEinesGammafunktionsquotienten4}}\le a_N\\
\overset{~\ref{ErwartungswertDer2NormEinesStandardnormalverteiltenProzesses}}
=&\, \E\left(\norm\cdot_{2, N} \circ \gamma \right)
= \E\left(\norm\cdot_{2, N}\circ \bprod \cdot x_N \circ \gamma \right)
= \E\left(\norm\cdot_E \circ \bprod \cdot x_N \circ \gamma \right) \tag{$c$}
\end{align*}
Also existiert nach~\ref{ExistenzEingeschraenkterFastIsometrien} ein
$A \in L(\ell_2^m, E)$ mit $\inf\menge{\norm{At}_{2, N}}{t \in T} >0$.\\
Wegen $\dim \Kern A = m - \dim \Bild A \ge m - N$ existiert ein $F \le \Kern A$
mit $\dim F = m-N$.\\
Sei $w \in F\setminus \meng 0$.\\
Dann ist $w \in \R^m$ und damit nach
\refclaim{AequivalenzkonstantenFuerPNormen}{AequivalenzkonstantenFuerPNormen1}:
$\norm w _{1,m} \le \sqrt m \norm w _{2,m}$.\\
Wegen $w \neq 0$ ist $\frac 1 {\norm w_{2,m}}w\in S^{m-1}$, also wegen $Aw = 0$
bereits $w \in \R^m \setminus T$.\\
Somit gilt nach Setzung von $T$:
\[\norma{\frac 1 {\norm w_{2,m}}w}_{1,m}>\sqrt k \text{, }\,\text{ also }\,
 \norm w _{1, m} > \sqrt k \norm w _{2,m}\text.\]
Daher gilt: $\forall\, v \in F: \sqrt k \norm v_{2,m} \le \norm v_{1,m}
\le \sqrt m \norm v _{2, m}$.\\
Also gilt
\begin{align}
d_{BM}\left((F,\restrict{\norm\cdot_{1,m}}F), 
(F, \restrict{\norm\cdot_{2,m}}F) \right)
\overset{\refclaim{EigenschaftenDesBanachMazurAbstands}
{EigenschaftenDesBanachMazurAbstands3}} \le \sqrt{\frac m k} \text.
\label{svgeq:5}
\end{align}
Da $F \le \R^m$, ist $\left(F, 
\restrict{\sprod\cdot{\cdot\cdot}_m}{F \times F}\right)$ ein
Hilbertraum und damit sind $\left(F,\restrict{\norm\cdot_{2,m}}F\right)$
und $\ell_2^{m-N}$ wegen gleicher Dimension isometrisch isomorph.\\
Wegen $k \ge \frac 1 2 \frac m {t_\eta}$ gilt schließlich
\begin{align*}
&d_{BM}\left(\left(F,\restrict{\norm\cdot_{1,m}}F\right),\ell_2^{m-N} \right)\\
\overset{\refclaim{EigenschaftenDesBanachMazurAbstands}
{EigenschaftenDesBanachMazurAbstands1}}
\le &\, d_{BM}\left(\left(F,\restrict{\norm\cdot_{1,m}}F\right),
                    \left(F,\restrict{\norm\cdot_{2,m}}F\right) \right)
    d_{BM}\left(\left(F,\restrict{\norm\cdot_{2,m}}F\right),
                \ell_2^{m-N} \right)\\
\overset{\refclaim{EigenschaftenDesBanachMazurAbstands}
{EigenschaftenDesBanachMazurAbstands2}}= &\,
d_{BM}\left(\left(F,\restrict{\norm\cdot_{1,m}}F\right),
                    \left(F,\restrict{\norm\cdot_{2,m}}F\right) \right)
 \overset{\eqref{svgeq:5}}\le\sqrt{\frac m k}\le\sqrt{2 t_\eta}\\
\overset{\eqref{svgeq:1}}\le
 &\,c \sqrt{\frac 1 {1-\eps}\log\left(\frac 2 {1-\eps} \right)}
\text.
\end{align*}
\end{proof}

Die Konstante $c$ in dem Satz ist in etwa $48.2$.
Beachtet man, dass die Abschätzungen für die Ordnungsstatistiken eher kanonisch als filigran waren,
so ist es naheliegend anzunehmen, dass diese Konstante etwas verbessert werden kann.
Ferner ist auch die Abschätzung des $t_\eta$ aus dem Numeriklemma (vgl.
\ref{Numeriklemma}) nur im Hinblick auf die Größenordnung genau -- mit weiteren
numerischen Methoden könnte auch diese noch verbessert werden.

Trotz der augenscheinlichen Größe von $c$ ist auch hier allerdings eher die
Wachstumsordnung des Abstandes interessant, welche, wie schon zu Beginn erwähnt, $d_{BM}(F, \ell_2^{m-N}) \in \mco(\nicefrac m N \mapsto \sqrt{(\nicefrac m N) \cdot \log(\nicefrac {2m}N)})$ erfüllt.

Interessant an diesem Beweis ist auch, dass im Gegensatz zu allen bisher
mithilfe von \ref{ExistenzEingeschraenkterFastIsometrien} bewiesenen Sätzen
nur \refclaim{ExistenzEingeschraenkterFastIsometrien}
{ExistenzEingeschraenkterFastIsometrien1} verwendet wird und nicht etwa die
sonst nützliche Abschätzung aus
\refclaim{ExistenzEingeschraenkterFastIsometrien}
{ExistenzEingeschraenkterFastIsometrien2}.
Hierbei bestand die Aufgabe darin,
einen Unterraum zu finden, welcher das gegebene $T$ nicht schneidet (was im
zwei-- und dreidimensionalen Fall sehr anschaulich ist), woraus die für den
Beweis interessante Abschätzung folgt.\parag
Zum Abschluss beweisen wir den Satz von Milman und Schechtman.
Dieser stellt eine interessante Ergänzung zum Satz von Dvoretzky dar, auf die im Anschluss
an den Beweis eingegangen wird.

Der vorgestellte Beweis beruht auf jenen von Gordon in \cite{Gor2}
und Guédon in \cite{Gued}.
Die Idee, die in \ref{Sortierungsnorm} vorgestellte
Sortierungsnorm für eine wesentliche Abschätzung zu benutzen, stammt von
Gordon und wurde von Guédon in \cite{Gued} verwendet.


\begin{satz}{Satz von Milman--Schechtman}{SatzVonMilmanSchechtman}
Es gibt ein $C_{BS} \in [1, \infty[$ derart, dass für alle $m, n \in \N$
mit $\frac 1 {10} \log(n) \le m \le n$ und jeden $n$--dimensionalen normierten
$\R$--Vektorraum $(E, \norm\cdot_E)$ ein $F \le E$ existiert mit $\dim F = m$
und
\[d_{BM}\left((F, \restrict{\norm\cdot_E}F), \ell_2^m \right)
\le 3 \frac{\sqrt{2 \log\left(\frac {15}2 \right)}}
{\sqrt{\log\left(\frac {15}2 \right)} - 1}C_{BS}
\sqrt{\frac m {\log\left(1+ \frac n m \right)}}\text.\]
\end{satz}
\begin{proof}
Nach \cite{BS} (dort: Theorem 2) gibt es ein $C_{BS} \in \R_{>0}$, sodass für
alle normierten Räume $(X, \norm\cdot_X)$ ein $M \in \N$ mit
$M > \frac 1 2 \dim X$ und $y \in X^M$ existieren mit
$\norm\cdot_{\infty,M}\le\norm\cdot_X\circ\bprod\cdot y_M
\le C_{BS}\norm\cdot_{2,M}$.\\
Wegen $\norm{e_1^M}_{\infty, M} = \norm{e_1^M}_{2, N}$ ist $C_{BS} \ge 1$.\\
Seien also $m, n \in \N$ mit $\frac 1 {10} \log(n) \le m \le n$ und
$(E, \norm\cdot_E)$ ein $n$--dimensionaler normierter $\R$--Vektorraum.\\
Nach dem oben zitierten Satz gibt es also ein $\widetilde M \in \N$ mit
$\widetilde M > \frac n 2$ und ein $\widetilde y \in E^{\widetilde M}$ mit
$\norm\cdot_{\infty,\widetilde M}
\le\norm\cdot_E\circ\bprod\cdot {\widetilde y}_{\widetilde M}
\le C_{BS}\norm\cdot_{2,\widetilde M}$.\\
Sei $N \coloneqq  \floor{\frac n 2}$. Dann ist $N \le \frac n 2 < \widetilde M$ und
daher $x\coloneqq \left(\frac 1{C_{BS}}\widetilde y_i \right)_{i=1}^N$ wohldefiniert
und es gilt (durch Auffüllung mit Nullen auf den Indizes ab $N+1$):
\begin{align}\frac 1{C_{BS}}
\norm\cdot_{\infty, N} \le \norm\cdot_E \circ \bprod\cdot x_N \le
\norm\cdot_{2,N} \text.\label{svmseq:1}\end{align}
Damit ist $x \neq 0$, sogar $x_i \neq 0$, da
$\frac 1 {C_{BS}} = \frac 1 {C_{BS}} \norm{e_i^N}_{\infty,N} \le \norm{x_i}_E$
für alle $i \in \haken N$.\\
Seien nun $T\coloneqq S^{m-1}$,
\[\bigg((\Omega, \mfa, \Prim), (\gamma,\delta,\widetilde g),(N, m, mN)\bigg)
\text{ ein }(N(0,1), 3)\text{--Tupel zu }(3,(N,m,mN))\] (existiert nach
\ref{ExistenzNormalNTupel}), $G_x \coloneqq  \bprod \cdot x_N \circ g$ und $(Z_{\gamma, x}, M_{\delta,T})$ das
Auswertungspaar zu $(\gamma, \delta, x, T)$.\\
Seien weiter
\[f: = \inf_{t \in T} \, \norm\cdot_E \circ \bprod \cdot x _N \circ G_x,
\quad
h\coloneqq  \sup_{t \in T} \, \norm\cdot_E \circ \bprod \cdot x _N \circ G_x\text.\]
Nach \ref{IntegrierbarkeitVonExtremaInGordonschenPaaren} sind
\begin{align}
f, h\text{ integrierbar und es ist } f \le h\text.\label{svmseq:2}
\end{align}
\underline{1. Fall:} $m \le \frac 2 {13}(N+1)$\absatz
\underline{1.1. Fall:} $\E(Z_{\gamma, x}) < 2a_m$\absatz
Seien $\mu_{j, N}$ für alle $j \in \haken N$ gemäß \ref{ktesMaximum},
\[\norm\cdot_{(m)}:\R^N \rightarrow [0, \infty[, v \mapsto
\frac 1 m \sum_{i=1}^m \mu_{i, N}\left(\pfeil {\abs\cdot}N (v)\right)\]
und
\[ V\coloneqq  \mengea{z \in \R^N\setminus\meng 0}
    {\left(\forall\, i \in \haken N: \abs{z_i} \in\meng{0,\nicefrac 1m} \right)
    \wedge \left(\abs{\menge{i \in \haken N}{\abs{z_i} \neq 0}}\le m \right)}
    \text.\]
Dann ist $V \subseteq \R^N$, nichtleer,
beschränkt (da endlich) und es ist $0 \notin V$. \hfill $(a)$\\
Nach \refclaim{Sortierungsnorm}{Sortierungsnorm3} ist
\begin{align}
\norm\cdot_{(m)}\text{ eine
Norm auf dem }\R^N\text{ mit } 
\norm\cdot_{(m)} \le \norm\cdot_{\infty, N}\text.\label{svmseq:3}
\end{align}
Ferner ist nach \refclaim{Sortierungsnorm}{Sortierungsnorm1}
$\sup\menge{\norm v_{2, N}}{v \in V} = \frac 1{\sqrt m}$ \hfill $(b)$\\
und nach \refclaim{Sortierungsnorm}{Sortierungsnorm2}
\begin{align}
\forall\, w \in \R^N: \sup_{v \in V}\sprod w v_N = \norm w_{(m)}\text.\tag{$c$}
\end{align}
Für alle $s \in \R^m$ und alle $\omega \in \Omega$ gilt:
\begin{align}
&\,(\bprod \cdot x \circ ev(s, \cdot) \circ g)(\omega)
= \bprod {g(\omega)s}x_N = (\bprod \cdot x_N \circ g(\omega))(s)\notag\\
=&\, ev(s, \cdot) \bigg(\bprod \cdot x _N \circ g(\omega) \bigg)
= (ev(s, \cdot) \circ \bprod \cdot x _N \circ g)(\omega) \text.\label{svmseq:4}
\end{align}
Wegen $N = \floor {\frac n 2}$ ist $N+1\ge \frac n 2$, also auch
$\frac {N+1}m \ge \frac n {2m}$ und damit
\begin{align}\left(1+\frac{N+1}m \right)^2
\ge 1 + 2 \frac{N+1}m \ge 1 + \frac n m\, , \text{ also } 1 + \frac{N+1}m
\ge \sqrt{1+ \frac n m} \text.\label{svmseq:5}\end{align}
Wegen $m \le \frac 2 {13}(N+1)$ und damit
$1+\frac{N+1}m \ge 1+\frac{13}2 = \frac{15}2$ gilt weiter:
\begin{align}
&\,\E\left(\norm \cdot_{(m)} \circ \gamma \right) - \frac {a_m}{\sqrt m}
\ge \E\left(\gamma_m^\ast \right) - \frac {a_m}{\sqrt m}
\overset{\refclaim{AbschaetzungEinesGammafunktionsquotienten}
{AbschaetzungEinesGammafunktionsquotienten2}}\ge 
\E\left(\gamma_m^\ast \right)-1\notag\\
\overset{\refclaim
{UntereSchrankeDesErwartungswertsVonAbsolutenN01Ordnungsstatistiken}
{UntereSchrankeDesErwartungswertsVonAbsolutenN01Ordnungsstatistiken3}}\ge
&\,
\sqrt{\log\left(1 + \frac {N+1}m \right)}-1\notag\\
=&\, \sqrt{\log\left(1 + \frac {N+1}m \right)} - 
\left(1 - \frac{\sqrt{\log\left(\frac {15}2 \right)} -1}
{\sqrt{\log\left(\frac {15}2 \right)}} \right)
\sqrt{\log\left(\frac {15}2 \right)}\notag\\
\ge &\,\sqrt{\log\left(1 + \frac {N+1}m \right)} - 
\left(1 - \frac{\sqrt{\log\left(\frac {15}2 \right)} -1}
{\sqrt{\log\left(\frac {15}2 \right)}} \right)
\sqrt{\log\left(1 + \frac {N+1}m \right)}\notag\\
=&\, \left(\frac{\sqrt{\log\left(\frac {15}2 \right)} -1}
{\sqrt{\log\left(\frac {15}2 \right)}} \right)
\sqrt{\log\left(1 + \frac {N+1}m \right)}\notag\\
\overset{\eqref{svmseq:5}}\ge 
&\, \left(\frac{\sqrt{\log\left(\frac {15}2 \right)} -1}
{\sqrt{\log\left(\frac {15}2 \right)}} \right)
\sqrt{\frac 1 2\log\left(1 + \frac {n}m \right)}\notag\\
=& \underbrace{\left(\frac{\sqrt{\log\left(\frac {15}2 \right)} -1}
{\sqrt{2\log\left(\frac {15}2 \right)}} \right)}_{>0, \text{ da }
\frac{15}2 > \exp(1)}
\sqrt{\log\left(1 + \frac {n}m \right)} \eqqcolon \alpha \text.\label{svmseq:6}
\end{align}
Seien $E_1\coloneqq  \R^N$, $\norm\cdot_{E_1}\coloneqq  \norm\cdot_{(m)}$,
$x_1 \coloneqq  (e_i^N)_{i=1}^N$ und $G_{x_1}\coloneqq  \bprod \cdot {x_1}_N \circ g = g$.\\
Wegen $(a), (b), (c)$ sind die Voraussetzungen von
\ref{GordonschePaareUndNormiterierendeMengen} für $(E_1, \norm\cdot_{E_1})$,
$x_1$, $T$, $V$ und $c\coloneqq \frac 1 {\sqrt{m}}$ erfüllt und damit gilt:
\begin{align}
\E\left(\norm\cdot_{(m)} \circ \gamma \right) - \frac {a_m}{\sqrt m}
& \overset{\ref{ErwartungswertDer2NormEinesStandardnormalverteiltenProzesses}}
= \E\left(\norm\cdot_{(m)} \circ \bprod \cdot {x_1}_N \circ\gamma \right)
- \frac 1 {\sqrt m} \E\left(\norm\cdot_{2, m} \circ \delta \right)\notag\\
&\overset{\ref{IterierteEuklidischeNorm}}=
\E\left(\norm\cdot_{(m)} \circ \bprod \cdot {x_1}_N \circ\gamma \right)
+ \frac 1 {\sqrt m}\E\left(\inf_{t \in T} \bprod t \delta_m \right)\notag\\
&\overset{\refclaim{GordonschePaareUndNormiterierendeMengen}
{GordonschePaareUndNormiterierendeMengen1}}\le
\E\left(\inf_{t \in T}\, \norm\cdot_{(m)}\circ ev(t,\cdot)\circ G_{x_1} \right)
\text.\label{svmseq:7}
\end{align}
Außerdem sind wegen $x \neq 0$ und \eqref{svmseq:1} auch die Voraussetzungen
von \ref{GordonschePaareUndSubeuklidischeBildhalbnormen} für
$(E, \norm\cdot_E)$, $x$ und $T$ erfüllt, sodass gilt:
\begin{align}
\E\left(\sup_{t \in T}\, \norm\cdot_{E}\circ ev(t,\cdot)\circ G_{x} \right)
&\overset{
\refclaim{GordonschePaareUndSubeuklidischeBildhalbnormen}
{GordonschePaareUndSubeuklidischeBildhalbnormen1}}
\le \E\left(\norm\cdot_E \circ \bprod \cdot x_N \circ \gamma \right)
+ \E\left(\sup_{t \in T} \bprod t \delta_m \right)\notag\\
&\overset{\ref{IterierteEuklidischeNorm}}=
\E\left(\norm\cdot_E \circ \bprod \cdot x_N \circ \gamma \right)
+ \E\left(\norm\cdot_{2, m}\circ \delta \right)\notag\\
&\overset{\ref{ErwartungswertDer2NormEinesStandardnormalverteiltenProzesses}}=
\E\left(\norm\cdot_E \circ \bprod \cdot x_N \circ \gamma \right) + a_m\notag\\
&\overset{\text{1.1 Fall}}< 3a_m \text.\label{svmseq:8}
\end{align}
Wir erhalten also insgesamt:
\begin{align*}
0 & \overset{\eqref{svmseq:6}}<
\frac 1 {C_{BS}}\left(\frac{\sqrt{\log\left(\frac {15}2 \right)} -1}
{\sqrt{2\log\left(\frac {15}2 \right)}} \right)
\sqrt{\log\left(1 + \frac {n}m \right)}\\
&\overset{\eqref{svmseq:6}}\le 
\frac 1 {C_{BS}}\left(\E\left(\norm \cdot_{(m)} \circ \gamma \right) 
- \frac {a_m}{\sqrt m}\right)\\
&\overset{\eqref{svmseq:7}}\le
\E\left(\inf_{t \in T}\, 
\frac 1 {C_{BS}}\norm\cdot_{(m)}\circ ev(t,\cdot)\circ G_{x_1} \right)\\
&=
\E\left(\inf_{t \in T}\, 
\frac 1 {C_{BS}}\norm\cdot_{(m)}\circ ev(t,\cdot)\circ g \right)\\
&\overset{\eqref{svmseq:3}}\le
\E\left(\inf_{t \in T}\, 
\frac 1 {C_{BS}}\norm\cdot_{\infty, N}\circ ev(t,\cdot)\circ g \right)\\
&\overset{\eqref{svmseq:1}}\le
\E\left(\inf_{t \in T}\, 
\norm\cdot_{E}\circ \bprod \cdot x_N \circ ev(t,\cdot)\circ g \right)\\
&\overset{\eqref{svmseq:4}}=
\E\left(\inf_{t \in T}\, 
\norm\cdot_{E}\circ ev(t,\cdot)\circ \bprod \cdot x_N \circ g \right)
=\E\left(\inf_{t \in T}\, \norm\cdot_{E}\circ ev(t,\cdot)\circ G_x \right)\\
&= \E(f) \overset{\eqref{svmseq:2}}\le \E(h)
= \E \left(\sup_{t \in T}\, \norm\cdot_{E}\circ ev(t,\cdot)\circ G_x \right)
\overset{\eqref{svmseq:8}}< 3 a_m\text.
\end{align*}
Da $f \ge 0, h \ge 0$ und beide integrierbar sind, gibt es nach
\ref{Erwartungswertlemma} ein $\omega_0 \in \Omega$ mit
$f(\omega_0)>0$ und $\frac{h(\omega_0)}{f(\omega_0)} \le \frac{3a_m}\alpha$.\\
Sei $A\coloneqq  G_x(\omega_0)$. Dann ist $A \in L(\ell_2^m, E)$ und es gilt:
\[\inf_{t \in S^{m-1}}\norm{At}_E =
  \inf_{t \in T} \norm{G_x(\omega_0)t}_E =
  \left(\inf_{t \in T} \norm\cdot_E \circ ev(t, \cdot) \circ G_x\right)
  (\omega_0) = f(\omega_0) > 0\]
und damit ist $A$ nach
\ref{InjektivitätLinearerAbbildungenZwischenNormiertenRäumen} injektiv.\\
Sei $F\coloneqq  \Bild A$. Da $A$ injektiv ist, ist $\dim F = m$ und da jeder
endlichdimensionale normierte Raum vollständig ist, gilt:
\begin{align*}
&\,d_{BM}\bigg((F, \norm\cdot_E\vert_F), \ell_2^m \bigg)
\le \opnorm A {\ell_2^m}{F} \opnorm{A\inv}{F}{\ell_2^m}
\overset{\ref{OperatornormDerInversen}}=
\frac{\Sup{t \in T} \norm{At}_E}{\Inf{t \in T}\norm{At}_E}\\
=&\, \frac{\Sup{t \in T}\norm{G_x(\omega_0)t}_E}
       {\Inf{t \in T}\norm{G_x(\omega_0)t}_E}
= \frac{\left(\Sup{t \in T} \norm\cdot_E \circ ev(t, \cdot) \circ G_x\right)
  (\omega_0)}
  {\left(\Inf{t \in T} \norm\cdot_E \circ ev(t, \cdot) \circ G_x\right)
  (\omega_0)} = \frac{h(\omega_0)}{f(\omega_0)} \le \frac{3a_m}\alpha\\
\overset{C_{BS} \ge 1}\le&\, \frac{3a_m}
{\frac 1 {C_{BS}}\left(\frac{\sqrt{\log\left(\frac {15}2 \right)} -1}
{\sqrt{2\log\left(\frac {15}2 \right)}} \right)
\sqrt{\log\left(1 + \frac {n}m \right)}}
= 3 \frac{\sqrt{2 \log\left(\frac {15}2 \right)}}
{\sqrt{\log\left(\frac {15}2 \right)} - 1}C_{BS}
\frac{a_m}{\sqrt{\log\left(1 + \frac {n}m \right)}}\\
\overset{\refclaim{AbschaetzungEinesGammafunktionsquotienten}
{AbschaetzungEinesGammafunktionsquotienten2}}\le &\,
3 \frac{\sqrt{2 \log\left(\frac {15}2 \right)}}
{\sqrt{\log\left(\frac {15}2 \right)} - 1}C_{BS}
\sqrt{\frac m {\log\left(1 + \frac {n}m \right)}}\text.
\end{align*}
Damit ist der Fall dieser Fall abgeschlossen.\absatz
\underline{1.2. Fall:} $\E(Z_{\gamma, x}) \ge 2 a_m$\absatz
Sei dann zusätzlich $\eps \coloneqq  \frac 1 2$.\\
Wegen \eqref{svmseq:1} ist $x \neq 0$ und 
$\norm\cdot_E\circ \bprod \cdot x_N \le \norm\cdot_{2,N}$. \hfill $(a')$\\
Ferner ist $T\neq \emptyset$, $T \subseteq S^{m-1}$ und $T=-T$. \hfill $(b')$\\
Schließlich gilt:
\begin{align}\E(M_{\delta, T}) \overset{\ref{IterierteEuklidischeNorm}}=
 \E\left(\norm\cdot_{2,m} \circ \delta \right)
 \overset{\ref{ErwartungswertDer2NormEinesStandardnormalverteiltenProzesses}}=
 a_m = \frac 1 2 2a_m \le \eps\E(Z_{\gamma,x})\text. \tag{$c'$}
 \end{align}
Wegen $(a'), (b'), (c')$ sind die Voraussetzungen von
\ref{ExistenzEingeschraenkterFastIsometrien} erfüllt und es gibt ein 
$A \in L(\ell_2^m, E)$ injektiv mit
\begin{align}
\opnorm A {\ell_2^m} E \opnorm{A\inv}{A(E)}{\ell_2^m} \le
\frac{1+\eps}{1-\eps} = \frac{\nicefrac 3 2}{\nicefrac 1 2} = 3\text.
\label{svmseq:9}
\end{align}
Sei nun $\vi:\!\ ]1,\infty[\rightarrow\R,s\mapsto\frac{\log(s)}{\log(1+s)}$.\\
Dann ist $\vi$ offenbar stetig differenzierbar und es gilt für alle
$z\in\!\ ]1,\infty[$:
\[\vi'(z)=\frac{\frac{\log(1+z)}{z}-\frac{\log(z)}{1+z}}{\log(1+z)^2}>0\text,\]
also $\vi$ monoton wachsend, sodass wegen $n \ge 2$ und
$m \ge \frac 1{10}\log(n)$ gilt:
\begin{align}\frac m{\log\left(1+\frac n m \right)}
\ge \frac 1 {10}\frac{\log(n)}{\log(n+1)}
= \frac 1 {10}\vi(n) \ge \frac 1 {10}\vi(2) 
=\frac 1 {10}\frac{\log(2)}{\log(3)}\text. \label{svmseq:10}\end{align}
Ferner gilt:
\begin{align}\frac{\sqrt{2 \log\left(\frac {15}2 \right)}}
{\sqrt{\log\left(\frac {15}2 \right)} - 1}
\sqrt{\frac 1 {10}\frac{\log(2)}{\log(3)}} > 1\text. \label{svmseq:11}
\end{align}
Sei nun $F\coloneqq  \Bild A$. Dann gilt:
\begin{align*}
d_{BM}\bigg((F, \norm\cdot_E\vert_F), \ell_2^m \bigg)
&\le \opnorm A {\ell_2^m} F \opnorm{A\inv}{F}{\ell_2^m} 
\overset{\eqref{svmseq:9}}\le 3
\overset{C_{BS}\ge 1}\le 3C_{BS}\\
&\overset{\eqref{svmseq:11}}< 3\frac{\sqrt{2 \log\left(\frac {15}2 \right)}}
{\sqrt{\log\left(\frac {15}2 \right)} - 1}
C_{BS}\sqrt{\frac 1 {10}\frac{\log(2)}{\log(3)}}\\
&\overset{\eqref{svmseq:10}}\le
3\frac{\sqrt{2 \log\left(\frac {15}2 \right)}}
{\sqrt{\log\left(\frac {15}2 \right)} - 1}
C_{BS}\sqrt{\frac m{\log\left(1+\frac n m \right)}}\text.
\end{align*}
Damit ist auch dieser Fall abgeschlossen.\absatz
\underline{2. Fall:} $m > \frac 2 {13}(N+1)$\absatz
Wegen $N = \floor{\frac n 2} \ge \frac n 2 - 1$ ist 
$m > \frac 2 {13}(N+1)\ge \frac 2 {13}\cdot\frac n 2 = \frac n {13}$ und
damit $\log\left(1+\frac n m \right)< \log(14)$.\\
Ferner gilt
\[3\frac{\sqrt{2 \log\left(\frac {15}2 \right)}}
{\sqrt{\log\left(\frac {15}2 \right)} - 1} \frac 1 {\sqrt{\log(14)}}> 1\text.\]
Sei nun $F \le E$ mit $\dim F = m$.\\
Dann ist nach dem Satz von John
\begin{align*}
d_{BM}\left((F, \restrict{\norm\cdot_E}F), \ell_2^m \right)
&\le \sqrt m \overset{C_{BS}\ge 1}\le C_{BS}\sqrt{m}
< 3\frac{\sqrt{2 \log\left(\frac {15}2 \right)}}
{\sqrt{\log\left(\frac{15}2 \right)}-1}C_{BS}\frac{\sqrt m}{\sqrt{\log(14)}}\\
& \le  3\frac{\sqrt{2 \log\left(\frac {15}2 \right)}}
{\sqrt{\log\left(\frac {15}2 \right)} - 1}
C_{BS}\sqrt{\frac m{\log\left(1+\frac n m \right)}}\text.
\end{align*}
\end{proof}

Die Konstante $C_{BS}$ lässt sich zwar nicht gut bestimmen, der zweite Teil der
angegebenen Konstante allerdings kann mit elementaren Methoden angenähert
werden und beträgt etwa $14.4$ (ist also im Vergleich zur Konstante bei
Gluskin recht klein).
Die Aufteilung des Beweises in die verschiedenen Fälle
ist nicht willkürlich -- im Falle von $m \le \frac 2{13}(n+1)$ ist nämlich eine
sehr wohlwollende Abschätzung nach unten für $\E(\gamma_m^\ast)$ bekannt (vgl.
\ref{UntereSchrankeDesErwartungswertsVonAbsolutenN01Ordnungsstatistiken}),
welche nicht nur die gesuchte Wachstumsordnung widerspiegelt, sondern dies auch
noch mit Konstante $1$ erfüllt.
Abstrahiert man sich von diesem wohlwollenden
Fall so braucht man nur jene $m \le n$ zu betrachten, für welche ein
$\alpha \in \R_{>0}$ existiert, sodass einerseits $\E(\gamma_m^\ast) > 1$ gilt
und $\E(\gamma_m^\ast)-1\ge\alpha \sqrt{\log\left(1+\frac n m \right)}$ ist.
Die vorgenommene Abschätzung, um von $\frac{N+1}m$ auf $\frac n m$ zu kommen
ist zwar etwas grob, sichert allerdings ohne größere Mühe eine multiplikative
Konstante vor dem logarithmischen Term.
Zu beachten gilt aber auch, dass die
oben erwähnte Abschätzung mit $\alpha > 0$ sich im Nenner des gesamten
Ausdrucks wiederfindet, sodass ein möglicherweise sehr kleines $\alpha$ eine
schlechtere Abschätzung liefert, weil der berechenbare Teil der Konstante dann
den Faktor $\frac 1 \alpha$ enthält.\parag
Wie vor dem Satz erwähnt, stellt der Satz von Milman--Schechtman eine Ergänzung
zum Satz von Dvoretzky dar.
Um diesen Sachverhalt genau notieren zu können,
beobachten wir zunächst, dass es für alle $c \in \!\ ]0, \nicefrac 1 2[$ ein
$n_0 \in \N$ gibt, sodass für alle $n \in \N_{\ge n_0}$ gilt:
\begin{align*}
n&\ge\max\meng{\left(\frac 2{1-2c}\right)^2,\frac 1 {c^{10}}}\text.\tag{$\ast$}
\end{align*}
Seien nun $c\in\!\ ]0,\nicefrac 12[$ und $n\in\N$.
Dann gilt:
\begin{align*}
\frac{\sqrt n} 2 - 1 \ge c\sqrt n & \iff\left(\frac 1 2-c\right)\sqrt n \ge 1\\
&\iff n\ge\left(\frac1{\frac 12-c}\right)^2=\left(\frac2{1-2c}\right)^2\tag 1\\
\frac 1 {10}\log(n) \le \frac 1 4 \log(c \sqrt n)&\iff
\log\left(n^{\nicefrac 2 5} \right) \le \log(c\sqrt n)\\
&\iff n^{\nicefrac 25}\le cn^{\nicefrac 1 2} \iff n \ge \frac 1 {c^{10}}
\text.\tag 2
\end{align*}
Erfüllt $n$ also die Eigenschaft $(\ast)$, so gilt
\begin{align*}
\frac 1 {10} \log(n) &\overset{(2)}\le \frac 1 4 \log\left(c\sqrt n \right)
\overset{(1)}\le \frac 1 4 \log\left(\frac{\sqrt n}2-1 \right)
< \frac 1 4 \log\left(\floor{\frac{\sqrt n}2}\right)\text. \tag 3
\end{align*}
Ist also $(E,\norm\cdot_E)$ ein $n$--dimensionaler $\R$--Vektorraum,
$m \in\haken n$ und erfüllt $n$ die Eigenschaft $(\ast)$, so ist der Satz von
Dvoretzky oder der Satz von Milman--Schechtman anwendbar, sodass es in beiden
Fällen ein $F \le E$ mit $\dim F = m$ gibt, dessen
Banach--Mazur--Abstand zu $\ell_2^m$ gut abgeschätzt werden kann.

Die Bedingungen $(1),(2)$ können zu konkreten Berechnungen benutzt werden.
Seien dazu $f:\!\ ]0, \nicefrac 1 2[ \rightarrow \R,
x \mapsto \left(\frac x{1-2x}\right)^2$ und
$g:\!\ ]0, \nicefrac 1 2[ \rightarrow \R, x \mapsto \frac 1 {x^{10}}$.\\
Dann ist $f$ streng monoton wachsend und $g$ streng monoton fallend.\\
Die Abbildung $h : \R \rightarrow \R, x \mapsto x^5 + x - \nicefrac 12$
besitzt nach dem Zwischenwertsatz eine Nullstelle $c_0 \in \R$ (und auch
genau eine, weil $h$ streng monoton wachsend ist).
Offensichtlich ist
$c_0 < \nicefrac 1 2$, denn $h(\nicefrac 1 2) > 0$ und $c_0 > 0$, denn
$h(0) < 0$.\\
Zudem gilt $f(c_0) = g(c_0)$. Das bedeutet, dass die Kenntnis von $c_0$ ein
bezüglich obiger Berechnungen minimales $n \in \N$ liefert.\\
Um eine Größenordnung des $n$ zu erhalten, betrachten wir $c_1\coloneqq \frac 9{20}$.\\
Es gilt $h(c_1) = \frac{59049}{3200000}+ \frac 9 {20}- \frac 1 2 =
- \frac{100951}{3200000}$ (also $c_1$ \enquote{nahe bei} $c_0$).\\
Sei $N \in \N$ minimal mit der Eigenschaft $(\ast)$ bezüglich $c_1$.\\
Durch $(1)$ hat man $N \ge 400$, während $(2)$ liefert
\[N \ge \left(\frac {20}9 \right)^{10} = \frac{10240000000000}{3486784401}
\in\!\  ]2936, 2937[\text.\]
Somit ergänzen sich die Sätze bereits mindestens auf $\N_{\ge 2937}$ (mit
$c_2\coloneqq  \frac{476}{1000}$ erhält man sogar $\N_{\ge 1737}$).

\bibliographystyle{plain}
\bibliography{Literatur}

\chapter*{Danksagung}
Ich danke Herrn Dr. Schütt für die Auswahl dieses interessanten Themas und die
gute Betreuung dieser Arbeit.\\
Vielen Dank an Christian Gießen, Paul Krühner, Sebastian Preugschat und
Insa Stucke, welche mir fachlich und freundschaftlich zur
Seite gestanden haben.\\
Ich danke auch Ole Griebenow und Arne Tüchsen für ihr lebhaftes
Interesse.\\
Mein besonderer Dank gilt Anja Hildebrandt für ihre Geduld und ihren Beistand.

\chapter*{Erklärung}
Ich versichere, diese Arbeit selbstständig und nur mit den angegebenen Mitteln
und Quellen angefertigt zu haben.
\parag\parag\parag\parag
Kiel, den \today \hfill Nikita Danilenko
\end{document}
